\xchapter{The Order for Daily Morning Prayer}

\textbf{The Book of Common Prayer (1928)}

THE
ORDER FOR

DAILY
MORNING PRAYER.

\emph{¶ The Minister shall begin the Morning Prayer
by reading one or more of the following Sentences of Scripture.}

\emph{¶
On any day, save a Day of Fasting or Abstinence, or on any day when the
Litany or Holy Communion is immediately to follow, the Minister may, at
his discretion, pass at once from the Sentences to the Lord's Prayer,
first pronouncing,}The Lord be with you\emph{.
Answer.} And with thy spirit. \emph{Minister.}Let us pray.

\emph{¶
And Note, that when the Confession and Absolution are omitted, the Minister
may, after the Sentences, pass to the Versicles,}O
Lord open thou our lips,\emph{etc., in which case the Lord's Prayer shall
be said with the other prayers, immediately after}The Lord be with
you, \emph{etc., and before the Versicles and Responses which follow, or,
in the Litany, as there appointed.}

THE
LORD is in his holy temple: let all the earth keep
silence before him. \emph{Hab}. ii. 20.

I was glad when they said unto me, We will go into
the house of the LORD. \emph{Psalm} cxxii.
1.

Let the words of my mouth, and the meditation of my
heart, be alway acceptable in thy sight, O LORD,
my strength and my redeemer. \emph{Psalm} xix. 14.

O send out thy light and thy truth, that they may lead
me, and bring me unto thy holy hill, and to thy dwelling. \emph{Psalm}
xliii. 3.

Thus saith the high and lofty One that inhabiteth eternity,
whose name is Holy; I dwell in the high and holy place, with him also
that is of a contrite and humble spirit, to revive the spirit of the humble,
and to revive the heart of the contrite ones. \emph{Isaiah} lvii.
15.

The hour cometh, and now is, when the true worshippers
shall worship the Father in spirit and in truth: for the Father seeketh
such to worship him. \emph{St. John} iv. 23.

Grace be unto you, and peace, from God our Father,
and from the Lord Jesus Christ. \emph{Phil}. i. 2. 

Repent ye, for the Kingdom of heaven is at hand. St. Matt. iii.
2.

Prepare ye the way of the LORD,
make straight in the desert a highway for our God. \emph{Isaiah}
xl. 3.

Behold, I bring you good tidings of great joy, which shall be to all people.
For unto you is born this day in the city of David a Saviour, which is
Christ the Lord. \emph{St. Luke}ii. 10, 11.

From the rising of the sun even unto the going down of the same my Name
shall be great among the Gentiles; and in every place incense shall be
offered unto my Name, and a pure offering: for my Name shall be great
among the heathen, saith the LORD of hosts. \emph{Mal.}i. 11.

Awake, awake; put on thy strength, O Zion; put on thy
beautiful garments, O Jerusalem. \emph{Isaiah} lii. 1.

Rend your heart, and not your garments, and turn unto the LORD
your God: for he is gracious and merciful, slow to anger, and of great
kindness, and repenteth him of the evil. \emph{Joel} ii. 13.

The sacrifices of God are a broken spirit: a broken
and a contrite heart, O God, thou wilt not despise. \emph{Psalm}
li. 17.

I will arise and go to my father, and will say unto
him, Father, I have sinned against heaven, and before thee, and am no
more worthy to be called thy son. \emph{St. Luke} xv. 18, 19.

Is it nothing to you, all ye that pass by? behold, and see if there be
any sorrow like unto my sorrow which is done unto me, wherewith the LORD
hath afflicted me. Lam. i. 12.

In whom we have redemption through his blood, the forgiveness
of sins, according to the riches of his grace. \emph{Eph}. i.
7.

He is risen. The Lord is risen indeed. \emph{St. Mark} xvi. 6;
\emph{St. Luke} xxiv. 34.

This is the day which the LORD
hath made; we will rejoice and be glad in it. \emph{Psalm} cxviii.
24.

Seeing that we have a great High Priest, that is passed into the heavens,
Jesus the Son of God, let us come boldly unto the throne of grace, that
we may obtain mercy, and find grace to help in time of need. \emph{Heb}.
iv. 14, 16.

Ye shall receive power, after that the Holy Ghost is come upon you: and
ye shall be witnesses unto me both in Jerusalem, and in all Judæa,
and in Samaria, and unto the uttermost part of the earth. \emph{Acts}
i. 8.

Because ye are sons, God hath sent forth the Spirit
of his Son into your hearts, crying, Abba, Father. \emph{Gal}.
iv. 6.

Holy, holy, holy, Lord God Almighty, which was, and is, and is to come.
\emph{Rev}. iv. 8.

Honour the LORD with thy substance, and with the
first-fruits of all thine increase: so shall thy barns be filled with
plenty, and thy presses shall burst out with new wine. \emph{Prov}.
iii. 9, 10.

The
LORD by wisdom hath founded the earth; by under-standing
hath he established the heavens. By his knowledge the depths are broken
up, and the clouds drop down the dew. \emph{Prov}. iii. 19, 20.

\begin{quote}

\begin{quote}

\emph{¶
Then the Minister shall say,}

\end{quote}

\end{quote}

DEARLY
beloved brethren, the Scripture moveth us, in sundry places, to acknowledge
and confess our manifold sins and wickedness; and that we should not dissemble
nor cloak them before the face of Almighty God our heavenly Father; but
confess them with an humble, lowly, penitent, and obedient heart; to the
end that we may obtain forgiveness of the same, by his infinite goodness
and mercy. And although we ought, at all times, humbly to acknowledge
our sins before God; yet ought we chiefly so to do, when we assemble and
meet together to render thanks for the great benefits that we have received
at his hands, to set forth his most worthy praise, to hear his most holy
Word, and to ask those things which are requisite and necessary, as well
for the body as the soul. Wherefore I pray and beseech you, as many as
are here present, to accompany me with a pure heart, and humble voice,
unto the throne of the heavenly grace, saying—

\emph{¶
Or he shall say,}

LET us humbly confess
our sins unto Almighty God.

\begin{quote}

\begin{quote}

\emph{A
General Confession.}

\end{quote}

\end{quote}

\emph{¶
To be said by the whole Congregation. after the Minister, all kneeling.}

ALMIGHTY
and most merciful Father; We have erred, and strayed from thy ways like
lost sheep. We have followed too much the devices and desires of our own
hearts. We have offended against thy holy laws. We have left undone those
things which we ought to have done; And we have done those things which
we ought not to have done; And there is no health in us. But thou, O Lord,
have mercy upon us, miserable offenders. Spare thou those, O God, who
confess their faults. Restore thou those who are penitent; According to
thy promises declared unto mankind in Christ Jesus our Lord. And grant,
O most merciful Father, for his sake; That we may hereafter live a godly,
righteous, and sober life, To the glory of thy holy Name. Amen.

\emph{The Declaration of Absolution, or Remission of Sins.}

\emph{¶
To be made by the Priest alone, standing; the People still kneeling.}

\emph{¶
But}NOTE\emph{, That
the Priest, at his discretion, may use, instead of what follows, the Absolution
from the Order for the Holy Communion.}\emph{}

ALMIGHTY
God, the Father of our Lord Jesus Christ, who desireth not the death of
a sinner, but rather that he may turn from his wickedness and live, hath
given power, and commandment, to his Ministers, to declare and pronounce
to his people, being penitent, the Absolution and Remission of their sins.
He pardoneth and absolveth all those who truly repent, and unfeignedly
believe his holy Gospel.

Wherefore let us beseech him to grant us true repentance,
and his Holy Spirit, that those things may please him which we do at this
present; and that the rest of our life hereafter may be pure and holy;
so that at the last we may come to his eternal joy; through Jesus Christ
our Lord. \emph{Amen}.

\emph{¶
Then the Minister shall kneel, and say the Lord’s Prayer; the People
still kneeling, and repeating it with him, both here, and wheresoever
else it is used in Divine Service.}

OUR
Father, who art in heaven, Hallowed be thy Name. Thy kingdom come. Thy
will be done on earth, As it is in heaven. Give us this day our daily
bread. And forgive us our trespasses, As we forgive those who trespass
against us. And lead us not into temptation; But deliver us from evil:
For thine is the kingdom, and the power, and the glory, for ever and ever.
Amen.

\begin{quote}

\begin{quote}

\emph{¶
Then likewise he shall say,}

\end{quote}

\end{quote}

O
Lord, open thou our lips.\emph{}

\emph{Answer}. And our mouth
shall show forth thy praise.

\begin{quote}

\begin{quote}
\emph{¶
Here, all standing up, the Minister shall say,}
\end{quote}

\end{quote}

Glory be to the Father, and to the Son, and to the Holy Ghost;

\emph{Answer}. As it was in the beginning, is now,
and ever shall be, world without end. Amen.

\emph{Minister}. Praise ye the Lord.

\emph{Answer}. The Lord’s Name be praised.

\emph{¶
Then shall be said or sung the following Canticle; except on those days
for which other Canticles are appointed; and except also, that Psalm
95 may be used in this place.}

\emph{¶
But Note, That on Ash Wednesday and Good Friday the Venite may be omitted.}

\emph{}\emph{}\emph{} \emph{}

¶
\emph{On the days hereafter named, here may be sung or said: On the Sundays in Advent.}— Our King and
Saviour draweth nigh : O come, let us adore him.

\emph{On Christmas Day and until the Epiphany}. —
Alleluia. Unto us a child is born : O come, let us adore him.

\emph{On the Epiphany and seven days after, and on the Feast
of the Transfiguration.} — The Lord hath manifested forth his
glory : O come, let us adore him.

\emph{On Monday in Easter Week, and until Ascension Day.}
— Alleluia. The Lord is risen indeed : O come let us adore him.
Alleluia.

\emph{On Ascension Day and until Whitsunday.} —
Alleluia. Christ the Lord ascended into heaven : O come, let us adore
him. Alleluia.

\emph{On Whitsunday and six days after.} —
Alleluia. The Spirit of the Lord filleth the world : O come, let us adore
him. Alleluia.

\emph{On Trinity Sunday.} — Father, Son, and
Holy Ghost, One God : O come, let us adore him.

\emph{On the Purification and the Annunciation.}—The
Word was made flesh : O come, let us adore him.

\emph{On other Festivals for which a proper Epistle and
Gospel are ordered.} — The Lord is glorious in his saints :
O come, let us adore him.

\emph{}\emph{}\emph{} \emph{}

\begin{quote}

\begin{quote}

\emph{Venite,
exultemus Domino.}

\end{quote}

\end{quote}

O
COME, let us sing unto the LORD; * let us heartily
rejoice in the strength of our salvation.

Let us come before his presence with thanksgiving;
* and show ourselves glad in him with psalms.

For the LORD is a great God;
* and a great King above all gods.

In his hand are all the corners of the earth; * and
the strength of the hills is his also.

The sea is his, and he made it; * and his hands prepared
the dry land.

O come, let us worship and fall down, * and kneel before
the LORD our Maker.

For he is the Lord our God; * and we are the people
of his pasture, and the sheep of his hand.

O worship the LORD in the beauty
of holiness; * let the whole earth stand in awe of him.

For he cometh, for he cometh to judge the earth; *
and with righteousness to judge the world, and the people with his truth.

\emph{¶
Then shall follow a Portion of the}PSALMS,\emph{according to the Use of this Church. And at the end of every Psalm, and
likewise at the end of the}Venite, Benedictus es, Benedictus, Jubilate,\emph{may be, and at the end of the whole Portion, or Selection from the Psalter,
shall be, sung or said the}Gloria Patri:

GLORY
be to the Father, and to the Son, * and to the Holy Ghost;

As it was in the beginning, is now, and ever shall
be, * world without end. Amen.

\emph{¶
Then shall be read the First Lesson, according to the Table or Calendar.
And}NOTE,\emph{That before every Lesson, the
Minister shall say,}Here beginneth such a Chapter [\emph{or}Verse
of such a Chapter]. of such a Book; \emph{and after every Lesson,}Here
endeth the First [\emph{or}the Second] Lesson.

\emph{¶
Here shall be said or sung the following Hymn.}

\emph{¶
But,}Note,
\emph{That on any day when the}Holy Communion\emph{is immediately to
follow the Minister at his discretion, after any one of the following
Canticles of Morning Prayer has been said or sung, may pass at once to
the Communion Service.}

\begin{quote}

\begin{quote}

\emph{Te
Deum laudamus.}

\end{quote}

\end{quote}

WE
praise thee, O God; we acknowledge thee to be the Lord.

All the earth doth worship thee, the Father everlasting.

To thee all Angels cry aloud; the Heavens, and all
the Powers therein;

To thee Cherubim and Seraphim continually do cry,

Holy, Holy, Holy, Lord God of Sabaoth;

Heaven and earth are full of the Majesty of thy glory.

The glorious company of the Apostles praise thee.

The goodly fellowship of the Prophets praise thee.

The noble army of Martyrs praise thee.

The holy Church throughout all the world doth acknowledge
thee;

The Father of an infinite Majesty;

Thine adorable, true and only Son;

Also the Holy Ghost the Comforter.

THOU
art the King of Glory, O Christ.

Thou art the everlasting Son of the Father.

When thou tookest upon thee to deliver man, thou didst
humble thyself to be born of a Virgin.

When thou hadst overcome the sharpness of death, thou
didst open the Kingdom of Heaven to all believers.

Thou sittest at the right hand of God, in the glory
of the Father.

We believe that thou shalt come to be our Judge.

We therefore pray thee, help thy servants, whom thou
hast redeemed with thy precious blood.

Make them to be numbered with thy Saints, in glory
everlasting.

O
LORD, save thy people, and bless thine heritage. 

Govern them and lift them up for ever.

Day by day we magnify thee;

And we worship thy Name ever, world without end.

Vouchsafe, O Lord, to keep us this day without sin.

O Lord, have mercy upon us, have mercy upon us.

O Lord, let thy mercy be upon us, as our trust is in
thee.

O Lord, in thee have I trusted; let me never be confounded.

\begin{quote}

\begin{quote}

\emph{¶
Or this Canticle.}

\emph{Benedictus
es Domine.}

\end{quote}

\end{quote}

BLESSED
art thou, O Lord God of our fathers; * praised and exalted above all for
ever.

Blessed art thou for the Name of thy Majesty; * praised
and exalted above all for ever.

Blessed art thou in the temple of thy holiness; * Praised
and exalted above all for ever.

Blessed art thou that beholdest the depths, and dwellest
between the Cherubim: * praised and exalted above all for ever.

Blessed art thou on the glorious throne of thy Kingdom:
* praised and exalted above all for ever.

Blessed art thou in the firmament of heaven: * praised
and exalted above all for ever.

\begin{quote}

\begin{quote}

\emph{¶
Or this,}

\emph{Benedicite,
omnia opera Domini.}

\end{quote}

\end{quote}

O
ALL ye Works of the Lord, bless ye the Lord: * praise him, and magnify
him for ever.

O ye Angels of the Lord, bless ye the Lord: * praise
him, and magnify him for ever.

O
YE Heavens, bless ye the Lord: * praise him, and magnify him for ever.

O ye Waters that be above the firmament, bless ye the
Lord: * praise him, and magnify him for ever.

O all ye Powers of the Lord, bless ye the Lord: * praise
him, and magnify him for ever.

O ye Sun and Moon, bless ye the Lord: * praise him,
and magnify him for ever.

O ye Stars of heaven, bless ye the Lord: * praise him,
and magnify him for ever.

O ye Showers and Dew, bless ye the Lord: * praise him,
and magnify him for ever.

O ye winds of God, bless ye the Lord: * praise him,
and magnify him for ever.

O ye Fire and Heat, bless ye the Lord * praise him,
and magnify him for ever.

O ye Winter and Summer, bless ye the Lord: * praise
him, and magnify him for ever.

O ye Dews and Frosts, bless ye the Lord: * praise him,
and magnify him for ever.

O ye Frost and Cold, bless ye the Lord: * praise him,
and magnify him for ever.

O ye Ice and Snow, bless ye the Lord * praise him,
and magnify him for ever.

O ye Nights and Days, bless ye the Lord: * praise him,
and magnify him for ever.

O ye Light and Darkness, bless ye the Lord: * praise
him, and magnify him for ever.

O ye Lightnings and Clouds, bless ye the Lord * praise
him, and magnify him for ever.

O
LET the Earth bless the Lord: * yea, let it praise him, and magnify him
for ever.

O ye Mountains and Hills, bless ye the Lord: * praise
him, and magnify him for ever.

O all ye Green Things upon the earth, bless ye the
Lord: * praise him, and magnify him for ever.

O ye Wells, bless ye the Lord: * praise him, and magnify
him for ever.

O ye Seas and Floods, bless ye the Lord: * praise him,
and magnify him for ever.

O ye Whales, and all that move in the waters, bless
ye the Lord: * praise him, and magnify him for ever.

O all ye Fowls of the air, bless ye the Lord: * praise
him, and magnify him for ever.

O all ye Beasts and Cattle, bless ye the Lord: * praise
him, and magnify him for ever.

O ye Children of Men, bless ye the Lord: * praise him,
and magnify him for ever.

O
LET Israel bless the Lord: * praise him, and magnify him for ever.

O ye Priests of the Lord, bless ye the Lord: * praise
him, and magnify him for ever.

O ye Servants of the Lord, bless ye the Lord: * praise
him, and magnify him for ever.

O ye Spirits and Souls of the Righteous, bless ye the
Lord: * praise him, and magnify him for ever.

O ye holy and humble Men of heart, bless ye the Lord:
* praise him, and magnify him for ever.

LET
us bless the Father, and the Son, and the Holy Ghost: * praise him, and
magnify him for ever.

\emph{¶
Then shall be read, in like manner, the Second Lesson, taken out of the New Testament, according
to the Table or Calendar.}

\emph{¶
And after that shall be sung or said the Hymn following.}

\emph{¶
But}Note,\emph{That, save on the Sundays in Advent, the latter portion thereof may be
omitted.}

\begin{quote}

\begin{quote}

\emph{Benedictus}.
St. Luke i. 68.

\end{quote}

\end{quote}

BLESSED
be the Lord God of Israel; * for he hath visited and redeemed his people;

And hath raised up a mighty salvation for us, * in
the house of his servant David;

As he spake by the mouth of his holy Prophets, * which
have been since the world began;

That we should be saved from our enemies, * and from
the hand of all that hate us.

To perform the mercy promised to our forefathers, * and to remember his
holy covenant;

To perform the oath which he sware to our forefather
Abraham, * that he would give us;

That we being delivered out of the hand of our enemies
* might serve him without fear;

In holiness and righteousness before him, * all the
days of our life.

And thou, child, shalt be called the prophet of the
Highest: * for thou shalt go. before the face of the Lord to prepare his
ways;

To give knowledge of salvation unto his people * for
the remission of their sins,

Through the tender mercy of our God; * whereby the
day-spring from on high hath visited us;

To give light to them that sit in darkness, and in
the shadow of death, * and to guide our feet into the way of peace.

\begin{quote}

\begin{quote}

\emph{¶ Or this Psalm.}

\emph{Jubilate
Deo}. Psalm c.

\end{quote}

\end{quote}

O
BE joyful in the LORD, all ye lands: * serve the
LORD with gladness, and come before his presence
with a song.

Be ye sure that the LORD he
is God; it is he that hath made us, and not we ourselves; * we are his
people, and the sheep of his pasture.

O go your way into his gates with thanksgiving, and
into his courts with praise; * be thankful unto him, and speak good of
his Name.

For the LORD is gracious, his
mercy is everlasting; * and his truth endureth from generation to generation.

\emph{¶
Then shall be said the Apostles’ Creed by the Minister and the People,
standing. And any Churches may, instead of the words,}He descended
into hell,\emph{use the words,}He went into the place of departed
spirits,\emph{which are considered as words of the same meaning in the
Creed.}

I
BELIEVE in God the Father Almighty, Maker of heaven and earth:

And in Jesus Christ his only Son our Lord: Who was
conceived by the Holy Ghost, Born of the Virgin Mary: Suffered under Pontius
Pilate, Was crucified, dead, and buried: He descended into hell; The third
day he rose again from the dead: He ascended into heaven, And sitteth
on the right hand of God the Father Almighty: From thence he shall come
to judge the quick and the dead.

I believe in the Holy Ghost: The holy Catholic Church;
The Communion of Saints: The Forgiveness of sins: The Resurrection of
the body: And the Life everlasting. Amen.

\begin{quote}

\begin{quote}

\emph{¶ Or the Creed commonly called the Nicene.}

\end{quote}

\end{quote}

I
BELIEVE in one God the Father Almighty, Maker of heaven and earth, And
of all things visible and invisible:

And in one Lord Jesus Christ, the only-begotten Son
of God; Begotten of his Father before all worlds, God of God, Light of
Light, Very God of very God; Begotten, not made; Being of one substance
with the Father; By whom all things were made: Who for us men and for
our salvation came down from heaven, And was incarnate by the Holy Ghost
of the Virgin Mary, And was made man: And was .crucified also for us under
Pontius Pilate; He suffered and was buried: And the third day he rose
again according to the Scriptures: And ascended into heaven, And sitteth
on the right hand of the Father: And he shall come again, with glory,
to judge both the quick and the dead; Whose kingdom shall have no end.

And I believe in the Holy Ghost, The Lord, and Giver
of Life, Who proceedeth from the Father and the Son; Who with the Father
and the Son together is worshipped and glorified; Who spake by the Prophets:
And I believe one Catholic and Apostolic Church: I acknOwledge one Baptism
for the remission of sins: And I look for the Resurrection of the dead:
And the Life of the world to come. Amen.

\emph{¶
And after that, these Prayers following, the people devoutly kneeling;
the Minister first pronouncing,}

\emph{Answer. Minister.} The
Lord be with you.

And with thy spirit. 

Let us pray.

\emph{¶
Here, if it hath not already been said, shall follow the Lord's Prayer.}

\emph{Minister.}\\
O Lord,
show thy mercy upon us.

\emph{Answer.} And grant us thy
salvation.

\emph{Minister.} O God, make clean
our hearts within us.

\emph{Answer.} And take not thy
Holy Spirit from us.

\emph{¶
Then shall follow the Collect for the Day, except when the Communion Service is read; and then the Collect for the
day shall be omitted here.}

\begin{quote}

\begin{quote}

\emph{A
Collect for Peace.}

\end{quote}

\end{quote}

O
GOD, who art the author of peace and lover of concord, in knowledge of
whom standeth our eternal life, whose service is perfect freedom; Defend
us thy humble servants in all assaults of our enemies; that we, surely
trusting in thy defence, may not fear the power of any adversaries, through
the might of Jesus Christ our Lord. \emph{Amen.}

\begin{quote}

\begin{quote}

\emph{A
Collect for Grace.}

\end{quote}

\end{quote}

O
LORD, our heavenly Father, Almighty and everlasting God, who hast safely
brought us to the beginning of this day; Defend us in the same with thy
mighty power; and grant that this day we fall into no sin, neither run
into any kind of danger; but that all our doings, being ordered by thy
governance, may be righteous in thy sight; through Jesus Christ our Lord.
\emph{Amen}.

\emph{¶
The following Prayers shall be omitted here when the Litany is said,
and may be omitted when the Holy Communion is to follow.}

\emph{¶
And}NOTE,\emph{That the Minister may here end the Morning Prayer with such general intercessions
taken out of this Book, as he shall think fit, or with the Grace.}

\emph{A
Prayer for The President of the United States, and all in Civil Authority.}

O
LORD, our heavenly Father, the high and mighty Ruler of the universe,
who dost from thy throne behold all the dwellers upon earth; Most heartily
we beseech thee, with thy favour to behold and bless thy servant THE
PRESIDENT OF THE UNITED
STATES, and all others in authority; and so replenish
them with the grace of thy Holy Spirit, that they may always incline to
thy will, and walk in thy way. Endue them plenteously with heavenly gifts;
grant them in health and prosperity long to live; and finally, after this
life, to attain everlasting joy and felicity; through Jesus Christ our
Lord. \emph{Amen}.

\begin{quote}

\begin{quote}

\emph{¶ Or this.}

\end{quote}

\end{quote}

O
LORD our Governor, whose glory is in all the world; We commend this nation
to thy merciful care, that being guided by thy Providence, we may dwell
secure in thy peace. Grant to THE PRESIDENT
OF THE UNITED STATES,
and to all in Authority, wisdom and strength to know and to do thy will.
Fill them with the love of truth and righteousness; and make them ever
mindful of their calling to serve this people in thy fear; through Jesus
Christ our Lord, who liveth and reigneth with thee and the Holy Ghost,
one God, world without end. \emph{Amen}.

\begin{quote}

\begin{quote}

\emph{A
Prayer for the Clergy and People.}

\end{quote}

\end{quote}

ALMIGHTY
and everlasting God, from whom cometh every good and perfect gift; Send
down upon our Bishops, and other Clergy, and upon the Congregations committed
to their charge, the healthful Spirit of thy grace; and, that they may
truly please thee, pour upon them the continual dew of thy blessing. Grant
this, O Lord, for the honour of our Advocate arid Mediator, Jesus Christ.
\emph{Amen}.

\begin{quote}

\begin{quote}

\emph{A
Prayer for all Conditions of Men.}

\end{quote}

\end{quote}

O
GOD, the Creator and Preserver of all mankind, we humbly beseech thee
for all sorts and conditions of men; that thou wouldest be pleased to
make thy ways known unto them, thy saving health unto all nations. More
especially we pray for thy holy Church universal; that it may be so guided
and governed by thy good Spirit, that all who profess and call themselves
Christians may be led into the way of truth, and hold the faith in unity
of spirit, in the bond of peace, and in righteousness of life. Finally,
we commend to thy fatherly goodness all those who are any ways afflicted,
or distressed, in mind, body, or estate; [*\emph{especially
those for whom our prayers are desired};] that it may please thee
to comfort and relieve them, according to. their several necessities;
giving them patience under their sufferings, and a happy issue out of
all their afflictions. And this we beg for Jesus Christ’s sake.
\emph{Amen.}

\begin{quote}

\begin{quote}

\emph{A
General Thanksgiving.}

\end{quote}

\end{quote}

ALMIGHTY
God, Father of all mercies, we, thine unworthy servants, do give thee
most humble and hearty thanks for all thy goodness and loving-kindness
to us and to all men; [*\emph{particularly to
those who desire now to offer up their praises and thanksgivings for thy
late mercies vouchsafed unto them.}] We bless thee for our creation,
preservation, and all the blessings of this life; but above all, for thine
inestimable love in the redemption of the world by our Lord Jesus Christ;
for the means of grace, and for the hope of glory. And, we beseech thee,
give us that due sense of all thy mercies, that our hearts may he unfeignedly
thankful; and that we show forth thy praise, not only with our lips, but
in our lives, by giving up our selves to thy service, and by walking before
thee in holiness and righteousness all our days; through Jesus Christ
our Lord, to whom, with thee and the Holy Ghost, be all honour and glory,
world without end. \emph{Amen}.

\emph{¶
Note, That the General Thanksgiving may be said by the Congregation with
the Minister.}

\begin{quote}

\begin{quote}

\emph{A
Prayer of St. Chrysostom.}

\end{quote}

\end{quote}

ALMIGHTY
God, who hast given us grace at this time with one accord to make our
common supplications unto thee; and dost promise that when two or three
are gathered together in thy Name thou wilt grant their requests; Fulfil
now, O Lord, the desires and petitions of thy servants, as may be most
expedient for them; granting us in this world knowledge of thy truth,
and in the world to come life everlasting. \emph{Amen.}

\begin{quote}

\begin{quote}

2
\emph{Cor}. xiii. 14.

\end{quote}

\end{quote}

THE
grace of our Lord Jesus Christ, and the love of God, and the fellowship
of the Holy Ghost, be with us all evermore. \emph{Amen}.

\begin{quote}

\begin{quote}

\emph{Here
endeth the Order of Morning Prayer.}

\end{quote}

\end{quote}

\begin{quote}

\emph{Return to the U.
S. Book of Common Prayer (1928)}

\end{quote}

\xchapter{The Order for Daily Evening Prayer}

\textbf{The Book of Common Prayer (1928)}

THE ORDER FOR

DAILY
EVENING PRAYER

\emph{¶
The Minister shall begin the Evening Prayer by reading one or more of
the following Sentences of Scripture; and then he shall say that which
is written after them. But he may, at his discretion, pass at once from
the Sentences to the Lord’s Prayer.}

\emph{¶
And}NOTE,\emph{that when the Confession and
Absolution are omitted, the Minister may, after the Sentences, pass to
the Versicles,}O Lord open thou our lips,\emph{etc., in which case
the Lord's Prayer shall be said with the other prayers, immediately after}The Lord be with you,\emph{etc., and before the Versicles and Responses
which follow.}

THE
LORD is in his holy temple: let all the earth keep
silence before him. \emph{Hab}. ii. 20.

LORD, I have loved the habitation
of thy house, and the place where thine honour dwelleth. \emph{Psalm}
xxvi. 8.

Let my prayer be set forth in thy sight as the incense;
and let the lifting up of my hands be an evening sacrifice. \emph{Psalm}
cxli. 2.

O worship the LORD in the beauty
of holiness; let the whole earth stand in awe of him. \emph{Psalm}
xcvi. 9.

Let the words of my mouth, and the meditation of my
heart, be alway acceptable in thy sight, O LORD,
my strength and my redeemer. \emph{Psalm} xix. 14, 15.

Watch
ye, for ye know not when the master of the house cometh, at even, or at
midnight, or at the cock-crowing, or in the morning: lest coming suddenly
he find you sleeping. \emph{St. Mark} xiii, 35, 36. 

Behold,
the tabernacle of God is with men, and he will dwell with them, and they
shall be his people, and God himself shall be with them, and be their
God. \emph{Rev}. xxi. 3.

And the Gentiles shall come to thy light, and kings to the brightness
of thy rising. \emph{Isaiah} lx. 3.

I acknowledge my transgressions: and my sin is ever before me. \emph{Psalm}
li. 3.

To the Lord our God belong mercies and forgivenesses,
though we have rebelled against him; neither have we obeyed the voice
of the LORD our God, to walk in his laws which
he set before us. \emph{Dan}. ix. 9, 10.

If we say that we have no sin, we deceive ourselves,
and the truth is not in us; but if we confess our sins, God is faithful
and just to forgive us our sins, and to cleanse us from all unrighteousness.
1 \emph{St. John} i. 8, 9.

All
we like sheep have gone astray; we have turned every one to his own way;
and the LORD hath laid on him the iniquity of us
all. \emph{Isaiah} liii. 6.

Thanks be to God, which giveth us the victory through our Lord Jesus Christ.
1 \emph{Cor}. xv. 57.

If ye then be risen with Christ, seek those things
which are above, where Christ sitteth on the right hand of God. \emph{Col}.
iii. 1. 

Christ is not entered into the holy places made with hands, which are
the figures of the true; but into heaven itself, now to appear in the
presence of God for us. \emph{Heb}. ix. 24.

There is a river, the streams whereof shall make glad the city of God,
the holy place of the tabernacles of the Most High. \emph{Psalm}
xlvi. 4. 

The
Spirit and the bride say, Come. And let him that heareth say, Come. And
let him that is athirst come. And whosoever will, let him take the water
of life freely. Rev. xxii. 17.

Holy, holy, holy, is the LORD of hosts: the whole
earth is full of his glory. \emph{Isaiah} vi. 3.

LET
us humbly confess our sins unto Almighty God.

\begin{quote}

\begin{quote}

\emph{¶
Or else he shall say as followeth.}

\end{quote}

\end{quote}

DEARLY
beloved brethren, the Scripture moveth us, in sundry places, to acknowledge
and confess our manifold sins and wickedness; and that we should not dissemble
nor cloak them before the face of Almighty God our heavenly Father; but
confess them with an humble, lowly, penitent, and obedient heart; to the
end that we may obtain forgiveness of the same, by his infinite goodness
and mercy. And although we ought, at all times, humbly to acknowledge
our sins before God; yet ought we chiefly so to do, when we assemble and
meet together to render thanks for the great benefits that we have received
at his hands, to set forth his most worthy praise, to hear his most holy
Word, and to ask those things which are requisite and necessary, as well
for the body as the soul. Wherefore I pray and beseech you, as many as
are here present, to accompany me with a pure heart, and humble voice,
unto the throne of the heavenly grace, saying—

\begin{quote}

\begin{quote}

\emph{A
General Confession.}

\end{quote}

\end{quote}

\emph{¶
To be said by the whole Congregation, after the Minister, all kneeling.}

ALMIGHTY
and most merciful Father; We have erred, and strayed from thy ways like
lost sheep. We have followed too much the devices and desires of our own
hearts. We have offended against thy holy laws. We have left undone those
things which we ought to have done; And we have done those things which
we ought not to have done; And there is no health in us. But thou, O Lord,
have mercy upon us, miserable offenders. Spare thou those, O God, who
confess their faults. Restore thou those who are penitent; According to
thy promises declared unto mankind in Christ Jesus our Lord. And grant,
O most merciful Father, for his sake; That we may hereafter live a godly,
righteous, and sober life, To the glory of thy holy Name. Amen.

\emph{The
Declaration of Absolution, or Remission of Sins.}

\emph{¶
To be made by the Priest alone, standing; the People still kneeling,}

ALMIGHTY
God, the Father of our Lord Jesus Christ, who desireth not the death of
a sinner, but rather that he may turn from his wickedness and live, hath
given power, and commandment, to his Ministers, to declare and pronounce
to his people, being penitent, the Absolution and Remission of their sins.
He pardoneth and absolveth all those who truly repent, and unfeignedly
believe his holy Gospel.

Wherefore let us beseech him to grant us true repentance,
and his Holy Spirit, that those things may please him which we do at this
present; and that the rest of our life hereafter may be pure and holy;
so that at the last we may come to his eternal joy; through Jesus Christ
our Lord. \emph{Amen}.

\begin{quote}

\begin{quote}

\emph{¶
Or this.}

\end{quote}

\end{quote}

THE
Almighty and Merciful God grant you Absolution and Remission of all your
sins, true repentance, amendment of life, and the grace and consolation
of the Holy Spirit. \emph{Amen}. 

\emph{¶
Then the Minister shall kneel, and say the Lord’s Prayer; the People
still kneeling, and repeating it with him.}

OUR
Father, who art in heaven, Hallowed be thy Name. Thy kingdom come. Thy
will be done, On earth as it is in heaven. Give us this day our daily
bread. And forgive us our trespasses, As we forgive those who trespass
against us. And lead us not into temptation; But deliver us from evil:
For thine is the kingdom, and the power, and the glory, for ever and ever.
Amen.

\begin{quote}

\begin{quote}

\emph{¶ Then likewise he shall say,}

\end{quote}

\end{quote}

O Lord,
open thou our lips.\emph{}

\emph{Answer}. And our mouth shall
show forth thy praise.

\begin{quote}

\begin{quote}

\emph{¶ Here, all standing up, the Minister shall say,}

\end{quote}

\end{quote}

Glory
be to the Father, and to the Son, and to the Holy Ghost;

\emph{Answer}. As it was in the beginning, is now,
and ever shall be, world without end. Amen.

\emph{Minister.} Praise ye the Lord.

\emph{Answer}. The Lord’s Name be praised.

\emph{¶
Then shall follow a Portion of the}PSALMS,\emph{according to the Use of this Church. And at the end of every Psalm,
and likewise at the end of the}Magnificat, Cantate Domino, Bonum
est confiteri, Nunc dimittis, Deus misereatur, Benedic anima mea,\emph{may be sung or said the}Gloria Patri;\emph{and at the end of the
whole portion or Selection of Psalms for the day, shall be sung or said
the}Gloria Patri,\emph{or else the Gloria in excelsis, as followeth.}

\begin{quote}

\begin{quote}

\emph{Gloria
in excelsis.}

\end{quote}

\end{quote}

GLORY
be to God on high, and on earth peace, good will towards men. We praise
thee, we bless thee, we worship thee, we glorify thee, we give thanks
to thee for thy great glory, O Lord God, heavenly King, God the Father
Almighty.

O Lord, the only-begotten Son, Jesus Christ; O Lord
God, Lamb of God, Son of the Father, that takest away the sins of the
world, have mercy upon us. Thou that takest away the sins of the world,
receive our prayer. Thou that sittest at the right hand of God the Father,
have mercy upon us.

For thou only art holy; thou only art the Lord; thou
only, O Christ, with the Holy Ghost, art most high in the glory of God
the Father. Amen.

\emph{¶
Then shall be read the First Lesson, according to the Table or Calendar.}

\emph{¶
After which shall be sung or said the Hymn called}Magnificat\emph{,
as followeth.}

\emph{¶
But,}NOTE,\emph{That the Minister, at his discretion, may omit one of the Lessons in
Evening Prayer, the Lesson read being followed by one of the Evening
Canticles.}\emph{}

\begin{quote}

\begin{quote}

\emph{Magnificat}.
St. Luke i. 46.

\end{quote}

\end{quote}

MY
soul doth magnify the Lord, * and my spirit hath rejoiced in God my Saviour.

For he hath regarded * the lowliness of his handmaiden.

For behold, from henceforth * all generations shall
call me blessed.

For he that is mighty hath magnified me; * and holy
is his Name.

And his mercy is on them that fear him * throughout
all generations.

He hath showed strength with his arm; * he hath scattered
the proud in the imagination of their hearts.

He hath put down the mighty from their seat, * and
hath exalted the humble and meek.

He hath filled the hungry with good things; * and the
rich he hath sent empty away.

He remembering his mercy hath holpen his servant Israel;
* as he promised to our forefathers, Abraham and his seed, for ever.

\begin{quote}

\begin{quote}

\emph{¶ Or this Psalm.}

\emph{Cantate
Domino.} Psalm xcviii.

\end{quote}

\end{quote}

O
SING unto the LORD a new song; * for he hath done
marvellous things.

With his own right hand, and with his holy arm, * hath
he gotten himself the victory.

The LORD declared his salvation;
* his righteousness hath he openly showed in the sight of the heathen.

He hath remembered his mercy and truth toward the house
of Israel; * and all the ends of the world have seen the salvation of
our God.

Show yourselves joyful unto the LORD,
all ye lands; * sing, rejoice, and give thanks.

Praise the LORD upon the harp;
* sing to the harp with a psalm of thanksgiving.

With trumpets also and shawms, * O show yourselves
joyful before the LORD, the King.

Let the sea make a noise, and all that therein is *
the round world, and they that dwell therein.

Let the floods clap their hands, and let the hills
be joyful together before the LORD * for he cometh to judge
the earth.

With righteousness shall he judge the world * and the
people with equity.

\begin{quote}

\begin{quote}

\emph{¶
Or this.}

\emph{Bonum
est confiteri}. Psalm xcii.

\end{quote}

\end{quote}

IT
is a good thing to give thanks unto the LORD, *
and to sing praises unto thy Name, O Most Highest;

To tell of thy loving-kindness early in the morning,
* and of thy truth in the night season;

Upon an instrument of ten strings, and upon the lute;
* upon a loud instrument, and upon the harp.

For thou, LORD, hast made me
glad through thy works; * and I will rejoice in giving praise for the
operations of thy hands.

¶
\emph{Then a Lesson of the New Testament, as it is appointed.}

¶\emph{And after that shall be sung or said the Hymn called Nunc dimittis, as
followeth.}

\begin{quote}

\begin{quote}

\emph{Nunc
dimittis}. St. Luke ii. 29.

\end{quote}

\end{quote}

LORD,
now lettest thou thy servant depart in peace, * according to thy word.

For mine eyes have seen * thy salvation, 

Which thou hast prepared * before the face of all people;

To be a light to lighten the Gentiles, * and to be
the glory of thy people Israel.

\begin{quote}

\begin{quote}

\emph{¶
Or else this Psalm.}

\emph{Deus
misereatur}. Psalm lxvii.

\end{quote}

\end{quote}

GOD
be merciful unto us, and bless us, * and show us tile light of his countenance,
and be merciful unto us;

That thy way may be known upon earth, * thy saving
health among all nations.

Let the people praise thee, O God; * yea, let all the
people praise thee.

O let the nations rejoice and be glad; * for thou shalt
judge the folk righteously, and govern the nations upon earth.

Let the people praise thee, O God; * yea, let all the
people praise thee.

Then shall the earth bring forth her increase; * and
God, even our own God, shall give us his blessing.

God shall bless us; * and all the ends of the world
shall fear him.

\begin{quote}

\begin{quote}

\emph{¶
Or this.}

\emph{Benedic,
anima mea}. Psalm ciii.

\end{quote}

\end{quote}

PRAISE
the LORD, O my soul; * and all that is within me,
praise his holy Name.

Praise the LORD, O my soul,
* and forget not all his benefits:

Who forgiveth all thy sin, * and healeth all thine
infirmities;

Who saveth thy life from destruction, * and crowneth
thee with mercy and loving-kindness.

O praise the LORD, ye angels
of his, ye that excel in strength; * ye that fulfill his commandment,
and hearken unto the voice of his word.

O praise the LORD, all ye his
hosts; * ye servants of his that do his pleasure.

O speak good of the LORD, all
ye works of his, in all places of his dominion: * praise thou the LORD,
O my soul.

\emph{¶
Then shall be said the Apostles’ Creed by the Minister and the People,
standing. And any Churches may, instead of the words,}He descended
into hell,\emph{use the words,}He went into the place of departed
spirits, \emph{which are considered as words of the same meaning in the
Creed.}

I
BELIEVE in God the Father Almighty, Maker of heaven and earth:

And in Jesus Christ his only Son our Lord: Who was
conceived by the Holy Ghost, Born of the Virgin Mary: Suffered under Pontius
Pilate, Was crucified, dead, and buried: He descended into hell; The third
day he rose again from the dead: He ascended into heaven, And sitteth
on the right hand of God the Father Almighty: From thence he shall come
to judge the quick and the dead. 

I believe in the Holy Ghost: The holy Catholic Church;
The Communion of Saints: The Forgiveness of sins: The Resurrection of
the body: And the Life everlasting. Amen.

\begin{quote}

\begin{quote}

\emph{¶
Or the Creed commonly called the Nicene.}

\end{quote}

\end{quote}

I
BELIEVE
in one God the Father Almighty, Maker of heaven and earth, And of all
things visible and invisible:

And in one Lord Jesus Christ, the only-begotten Son
of God; Begotten of his Father before all worlds, God of God, Light of
Light, Very God of very God; Begotten, not made; Being of one substance
with the Father; By whom all things were made: Who for us men and for
our salvation came down from heaven, And was incarnate by the Holy Ghost
of the Virgin Mary, And was made man: And was crucified also for us under
Pontius Pilate; He suffered and was buried: And the third day lie rose
again according to the Scriptures: And ascended into heaven, And sitteth
on the right hand of the Father: And he shall come again, with glory,
to judge both the quick and the dead; Whose kingdom, shall have no end.

And I believe in the Holy Ghost, The Lord, and Giver
of Life, Who proceedeth from the Father and the Son; Who with the Father
and the Son together is worshipped and glorified; Who spake by the Prophets:
And I believe one Catholic and Apostolic Church: I acknowledge one Baptism
for the remission of sins: And I look for the Resurrection of the dead:
And the Life of the world to come. Amen.

\emph{¶
And after that, these Prayers following, the people devoutly kneeling;
the minister first pronouncing,}

\emph{Answer. Minister.} The
Lord be with you.

And with thy spirit. 

Let us pray.

\emph{¶ Here, if it hath not already been said, shall follow the Lord's
Prayer.}

\emph{Minister}.
O Lord, show thy
mercy upon us.

\emph{Answer.} And
grant us thy salvation.

\emph{Minister.} O Lord, save the
State.

\emph{Answer.} And
mercifully hear us when we call upon thee.

\emph{Minister.} Endue thy Ministers
with righteousness.

\emph{Answer.} And make thy chosen
people joyful.

\emph{Minister.} O Lord, save thy
people.

\emph{Answer.} And bless thine
inheritance.

\emph{Minister.} Give peace in our
time, O Lord.

\emph{Answer.} For
it is thou, Lord, only, that makest us dwell in safety.

\emph{Minister.} O God, make clean
our hearts within us.

\emph{Answer.} And take not thy
Holy Spirit from us.

\emph{¶
Then shall be said the Collect of the Day, and after that the Collects
and Prayers following.}

\begin{quote}

\emph{A
Collect for Peace.}

\end{quote}

O
GOD, from whom all holy desires, all good counsels, and all just works
do proceed; Give unto thy servants that peace which the world cannot give;
that our hearts may be set to obey thy commandments, and also that by
thee, we, being defended from the fear of our enemies, may pass our time
in rest and quietness; through the merits of Jesus Christ our Saviour.
\emph{Amen}.

\begin{quote}

\emph{A
Collect for Aid against Perils.}

\end{quote}

LIGHTEN
our darkness, we beseech thee, O Lord; and by thy great mercy defend us
from all perils and dangers of this night; for the love of thy only Son,
our Saviour, Jesus Christ. \emph{Amen}.

\emph{¶
In places where it may be convenient, here followeth the Anthem.}

\emph{¶
The Minister may here end the Evening Prayer with such Prayer, or Prayers,
taken out of this Book, as he shall think fit.}

\emph{A
Prayer for The President of the United States, and all in Civil Authority.}

ALMIGHTY
God, whose kingdom is everlasting and power infinite; Have mercy upon
this whole land; and so rule the hearts of thy servants THE
PRESIDENT OF THE UNITED
STATES, \emph{The Governor of this State}, and
all others in authority, that they, knowing whose ministers they are,
may above all things seek thy honour and glory; and that we and all the
People, duly considering whose authority they bear, may faithfully and
obediently honour them, according to thy blessed Word and ordinance; through
Jesus Christ our Lord, who with thee and the Holy Ghost liveth and reigneth
ever, one God, world without end. \emph{Amen}.

\begin{quote}

\emph{A
Prayer for the Clergy and People.}

\end{quote}

ALMIGHTY
and everlasting God, from whom cometh every good and perfect gift; Send
down upon our Bishops, and other Clergy, and upon the Congregations committed
to their charge, the healthful Spirit of thy grace; and, that they may
truly please thee, pour upon them the continual dew of thy blessing. Grant
this, O Lord, for the honour of our Advocate and Mediator, Jesus Christ.
\emph{Amen}.

\begin{quote}

\emph{A
Prayer for all Conditions of Men.}

\end{quote}

O
GOD, the Creator and Preserver of all mankind, we humbly beseech thee
for all sorts and conditions of men; that thou wouldest be pleased to
make thy ways known unto them, thy saving health unto all nations. More
especially we pray for thy holy Church universal; that it may be so guided
and governed by thy good Spirit, that all who profess and call themselves
Christians may be led into the way of truth, and hold the faith in unity
of spirit, in the bond of peace, and in righteousness of life. Finally,
we commend to thy fatherly goodness all those who are any ways afflicted,
or distressed, in mind, body, or estate; [*\emph{especially those
for whom our prayers are desired};] that it may please thee to comfort
and relieve them, according to their several necessities; giving them
patience under their sufferings, and a happy issue out of all their afflictions.
And this we beg for Jesus Christ’s sake. \emph{Amen}.

\begin{quote}

\emph{A
General Thanksgiving.}

\end{quote}

ALMIGHTY
God, Father of all mercies, we, thine unworthy servants, do give thee
most humble and hearty thanks for all thy goodness and loving-kindness
to us, and to all men; [*\emph{particularly to
those who desire now to offer up their praises and thanksgivings for thy
late mercies vouchsafed unto them}.]. We bless thee for our creation,
preservation, and all the blessings of this life; but above all, for thine
inestimable love in the redemption of the world by our Lord Jesus Christ;
for the means of grace, and for the hope of glory. And, we beseech thee,
give us that due sense of all thy mercies, that our hearts may he unfeignedly
thankful: and that we show forth thy praise, not only with our lips, but
in our lives, by giving up our selves to thy service, and by walking before
thee in holiness and righteousness all our days; through Jesus Christ
our Lord, to whom, with thee and the Holy Ghost, be all honour and glory,
world without end. \emph{Amen}.

\emph{¶}Note,\emph{That the General Thanksgiving may be said by the Congregation
with the Minister.}

\begin{quote}

\emph{A
Prayer of St. Chrysostom.}

\end{quote}

ALMIGHTY
God, who hast given us grace at this time with one accord to make our
common supplications unto thee; and dost promise that when two or three
are gathered together in thy Name thou wilt grant their requests; Fulfil
now, O Lord, the desires and petitions of thy servants, as may be most
expedient for them; granting us in this world knowledge of thy truth,
and in the world to come life everlasting. Amen.

\begin{quote}

2 \emph{Cor}.
xiii. 14.

\end{quote}

THE
grace of our Lord Jesus Christ, and the love of God, and the fellowship
of the Holy Ghost, be with us all evermore. \emph{Amen}.

\begin{quote}

\emph{Here
endeth the Order of Evening Prayer.}

\end{quote}

\begin{quote}

\emph{Return to the U.
S. Book of Common Prayer (1928)}

\end{quote}

\xchapter{The Litany}

\textbf{The Book of Common Prayer (1928)}

The
Litany

or General
Supplication

¶
\emph{To be used after the Third Collect at Morning or Evening Prayer; or
before the Holy Communion; or separately.}

O
GOD the Father, Creator of heaven and earth; 

\emph{Have mercy upon us.}

O God the Son, Redeemer of the world; 

\emph{Have mercy upon us.}

O God the Holy Ghost, Sanctifier of the faithful; 

\emph{Have mercy upon us.}

O holy, blessed, and glorious Trinity, one God; 

\emph{Have mercy upon us.}

REMEMBER
not, Lord, our offences, nor the offences of our forefathers; neither
take thou vengeance of our sins: Spare us, good Lord, spare thy people,
whom thou hast redeemed with thy most precious blood, and be not angry
with us for ever. 

\emph{Spare us, good Lord.}

FROM
all evil and mischief; from sin; from the crafts and assaults of the devil;
from thy wrath, and from everlasting damnation,

\emph{Good Lord, deliver us.}

From all blindness of heart; from pride, vainglory,
and hypocrisy; from envy, hatred, and malice, and all uncharitableness,

\emph{Good Lord, deliver us.}

From all inordinate and sinful affections; and from
all the deceits of the world, the flesh, and the devil, 

\emph{Good Lord, deliver us.}

From lightning and tempest; from earthquake, fire,
and flood; from plague, pestilence, and famine; from battle and murder,
and from sudden death, 

\emph{Good Lord, deliver us.}

From all sedition, privy conspiracy, and rebellion;
from all false doctrine, heresy, and schism; from hardness of heart, and
contempt of thy Word and Commandment, 

\emph{Good Lord, deliver us}.

By the mystery of thy holy Incarnation; by thy holy
Nativity and Circumcision; by thy Baptism, Fasting, and Temptation, 

\emph{Good Lord, deliver us.}

By thine Agony and Bloody Sweat; by thy Cross and Passion;
by thy precious Death and Burial; by thy glorious Resurrection and Ascension,
and by the Coming of the Holy Ghost, 

\emph{Good Lord, deliver us.}

In all time of our tribulation; in all time of our
prosperity; in the hour of death, and in the day of judgment, 

\emph{Good Lord, deliver us.}

WE
sinners do beseech thee to hear us, O Lord God; and that it may please
thee to rule and govern thy holy Church universal in the right way; 

\emph{We beseech thee to hear us, good Lord.}

That it may please thee so to rule the heart of thy
servant, The President of the United States, that he may above all things
seek thy honour and glory; 

\emph{We beseech thee to hear us, good Lord.}

That it may please thee to bless and preserve all Christian
Rulers and Magistrates, giving them grace to execute justice, and to maintain
truth; 

\emph{We beseech thee to hear us, good Lord.}

That it may please thee to illuminate all Bishops,
Priests, and Deacons, with true knowledge and understanding of thy Word;
and that both by their preaching and living they may set it forth, and
show it accordingly; 

\emph{We beseech thee to hear us, good Lord.}

That it may please thee to send forth labourers into
thy harvest; 

\emph{We beseech thee to hear us, good Lord.}

That it may please thee to bless and keep all thy people; 

\emph{We beseech thee to hear us, good Lord.}

That it may please thee to give to all nations unity, peace, and concord;

\emph{We beseech thee to hear us, good Lord.}

That it may please thee to give us an heart to love and fear thee, and
diligently to live after thy commandments; 

\emph{We beseech thee to hear us, good Lord.}

That it may please thee to give to all thy people increase
of grace to hear meekly thy Word, and to receive it with pure affection,
and to bring forth the fruits of the Spirit; 

\emph{We beseech thee to hear us, good Lord.}

That it may please thee to bring into the way of truth
all such as have erred, and are deceived; 

\emph{We beseech thee to hear us, good Lord.}

That it may please thee to strengthen such as do stand;
and to comfort and help the weak-hearted; and to raise up those who fall;
and finally to beat down Satan under our feet; 

\emph{We beseech thee to hear us, good Lord.}

That it may please thee to succour, help, and comfort,
all who are in danger, necessity, and tribulation; 

\emph{We beseech thee to hear us, good Lord.}

That it may please thee to preserve all who travel
by land, by water, or by air, all women in child-birth, all sick persons,
and young children; and to show thy pity upon all prisoners and captives;

\emph{We beseech thee to hear us, good Lord.}

That it may please thee to defend, and provide for,
the fatherless children, and widows, and all who are desolate and oppressed;

\emph{We beseech thee to hear us, good Lord.}

That it may please thee to have mercy upon all men;

\emph{We beseech thee to hear us, good Lord.}

That it may please thee to forgive our enemies, persecutors,
and slanderers, and to turn their hearts; 

\emph{We beseech thee to hear us, good Lord.}

That it may please thee to give and preserve to our
use the kindly fruits of the earth, so that in due time we may enjoy them;

\emph{We beseech thee to hear us, good Lord.}

That it may please thee to give us true repentance;
to forgive us all our sins, negligences, and ignorances; and to endue
us with the grace of thy Holy Spirit to amend our lives according to thy
holy Word; 

\emph{We beseech thee to hear us, good Lord.}

Son of God, we beseech thee to hear us. 

\emph{Son of God, we beseech thee to hear us.}

O Lamb of God, who takest away the sins of the world;

\emph{Grant us thy peace.}

O Lamb of God, who takest away the sins of the world;

\emph{Have mercy upon us.}

O Christ, hear us. 

\emph{O Christ, hear us.}

Lord, have mercy upon us. 

\emph{Lord, have mercy upon us.}

Christ, have mercy upon us. 

\emph{Christ, have mercy upon us.}

Lord, have mercy upon us. 

\emph{Lord, have mercy upon us.}

¶ \emph{Then shall the Minister, and
the People with him, say the Lord's Prayer.}

OUR Father, who art in
heaven, Hallowed be thy Name. Thy kingdom come. Thy will be done, On earth
as it is in heaven. Give us this day our daily bread. And forgive us our
trespasses, As we forgive those who trespass against us. And lead us not
into temptation, But deliver us from evil. Amen.

¶ \emph{The Minister
may, at his discretion, omit all that followeth, to the Prayer, We humbly
beseech thee, O Father, etc.}

\emph{Minister.}

O LORD, deal not with us
according to our sins. 

\emph{Neither reward us according to our iniquities.}

Let us pray.

O GOD, merciful Father,
who despisest not the sighing of a contrite heart, nor the desire of such
as are sor-rowful; Mercifully assist our prayers which we make before
thee in all our troubles and adversities, whensoever they oppress us;
and graciously hear us, that those evils which the craft and subtilty
of the devil or man worketh against us, may, by thy good providence, be
brought to nought; that we thy servants, being hurt by no persecutions,
may evermore give thanks unto thee in thy holy Church; through Jesus Christ
our Lord. \emph{Amen.}

¶ \emph{Minister and People.}

O Lord, arise, help us, and deliver
us for thy Name's sake.

\emph{Minister.}

O GOD, we have heard with
our ears, and our fathers have declared unto us, the noble works that
thou didst in their days, and in the old time before them.

¶ \emph{Minister and People.}

O Lord, arise, help us, and deliver
us for thine honour.

\emph{Minister.}

Glory be to the Father, and to the
Son, and to the Holy Ghost; 

\emph{As it was in the beginning, is now, and ever shall
be, world without end. Amen.}

From our enemies defend us, O Christ. 

\emph{Graciously look upon our afflictions.}

With pity behold the sorrows of our hearts. 

\emph{Mercifully forgive the sins of thy people.}

Favourably with mercy hear our prayers. 

\emph{O Son of David, have mercy upon us.}

Both now and ever vouchsafe to hear us, O Christ. 

\emph{Graciously hear us, O Christ; graciously hear us,
O Lord Christ.}

O Lord, let thy mercy be showed upon
us; 

\emph{As we do put our trust in thee.}

Let us pray.

WE humbly beseech thee,
O Father, mercifully to look upon our infirmities; and, for the glory
of thy Name, turn from us all those evils that we most justly have deserved;
and grant, that in all our troubles we may put our whole trust and confidence
in thy mercy, and evermore serve thee in holiness and pureness of living,
to thy honour and glory; through our only Mediator and Advocate, Jesus
Christ our Lord. \emph{Amen.}

¶ \emph{The Minister
may end the Litany here, or at his discretion add other prayers from this
Book.}

A Penitential Office

For Ash Wednesday.

¶ \emph{On the First
Day of Lent, the Office ensuing may be read immediately after the Prayer,
We humbly beseech thee, O Father, in the Litany; or it may be used with
Morning Prayer, or Evening Prayer, or as a separate Office.}

¶ \emph{The same Office
may be read at other times, at the discretion of the Minister.}

¶ \emph{The Minister
and the People kneeling, then shall be said by them this Psalm following.}

\emph{Miserere mei, Deus}. Psalm li.

HAVE mercy upon me, O God,
after thy great goodness; * according to the multitude of thy mercies
do away mine offences.

Wash me throughly from my wickedness, * and cleanse
me from my sin.

For I acknowledge my faults, * and my sin is ever before
me.

Against thee only have I sinned, and done this evil
in thy sight; * that thou mightest be justified in thy saying, and clear
when thou art judged.

Behold, I was shapen in wickedness, * and in sin hath
my mother conceived me.

But lo, thou requirest truth in the inward parts, *
and shalt make me to understand wisdom secretly.

Thou shalt purge me with hyssop, and I shall be clean;
* thou shalt wash me, and I shall be whiter than snow.

Thou shalt make me hear of joy and gladness, * that
the bones which thou hast broken may rejoice.

Turn thy face from my sins, * and put out all my misdeeds.

Make me a clean heart, O God, * and renew a right spirit
within me.

Cast me not away from thy presence, * and take not
thy holy Spirit from me.O give me the comfort of thy help again, * and
stablish me with thy free Spirit.

Then shall I teach thy ways unto the wicked, * and
sinners shall be converted unto thee.

Deliver me from blood-guiltiness, O God, thou that
art the God of my health; * and my tongue shall sing of thy righteousness.

Thou shalt open my lips, O Lord, * and my mouth shall
show thy praise.

For thou desirest no sacrifice, else would I give it
thee; * but thou delightest not in burnt-offerings.

The sacrifice of God is a troubled spirit: * a broken
and contrite heart, O God, shalt thou not despise.

Glory be to the Father, and to the Son, * and to the
Holy Ghost;

As it was in the beginning, is now, and ever shall
be, * world without end. Amen.

¶ \emph{If the Litany
hath been already said. the Minister may pass at once to O Lord, save
thy servants; etc.}

\begin{quote}

\begin{quote}

Lord, have mercy
upon us.

\emph{Christ, have mercy upon us.}

Lord, have mercy upon us.

\end{quote}

\end{quote}

OUR
Father, who art in heaven, Hallowed be thy Name. Thy kingdom come. Thy
will be done, On earth as it is in heaven. Give us this day our daily
bread. And forgive us our trespasses, As we forgive those who trespass
against us. And lead us not into temptation, But deliver us from evil.
Amen.

\begin{quote}

\begin{quote}

O Lord, save thy servants; 

\emph{That put their trust in thee.}

Send unto them help from above. 

\emph{And evermore mightily defend them.}

\end{quote}

\end{quote}

Help us, O God our Saviour. 

\emph{And for the glory of thy Name deliver us; be
merciful to us sinners, for thy Name's sake.}

O Lord, hear our prayer. 

\emph{And let our cry come unto thee.}

Let us pray.

O LORD, we beseech thee,
mercifully hear our prayers, and spare all those who confess their sins
unto thee; that they, whose consciences by sin are accused, by thy merciful
pardon may be absolved; through Christ our Lord. \emph{Amen}.

O MOST mighty God, and
merciful Father, who hast compassion upon all men, and who wouldest not
the death of a sinner, but rather that he should turn from his sin, and
be saved; Mercifully forgive us our trespasses; re-ceive and comfort us,
who are grieved and wearied with the burden of our sins. Thy property
is always to have mercy; to thee only it appertaineth to forgive sins.
Spare us there-fore, good Lord, spare thy people, whom thou hast re-deemed;
enter not into judgment with thy servants; but so turn thine anger from
us, who meekly acknowledge our transgressions, and truly repent us of
our faults, and so make haste to help us in this world, that we may ever
live with thee in the world to come; through Jesus Christ our Lord. \emph{Amen}.

¶ \emph{Then shall the
People say this that followeth, after the Minister.}

TURN thou us, O good Lord,
and so shall we be turned. Be favourable, O Lord, Be favourable to thy
people, Who turn to thee in weeping, fasting, and praying. For thou art
a merciful God, Full of compassion, Long-suffering, and of great pity.
Thou sparest when we deserve punishment, And in thy wrath thinkest upon
mercy. Spare thy people, good Lord, spare them, And let not thine heritage
be brought to confusion. Hear us, O Lord, for thy mercy is great, And
after the multitude of thy mercies look upon us; Through the merits and
mediation of thy blessed Son, Jesus Christ our Lord. Amen.

¶ \emph{Then the Minister shall say,}

O GOD, whose nature and
property is ever to have mercy and to forgive; Receive our humble petitions;
and though we be tied and bound with the chain of our sins, yet let the
pitifulness of thy great mercy loose us; for the honour of Jesus Christ,
our Mediator and Advocate. \emph{Amen.}

THE Lord bless us, and
keep us. The Lord make his face to shine upon us, and be gracious unto
us. The Lord lift up his countenance upon us, and give us peace, both
now and evermore. \emph{Amen.}

\begin{quote}

\emph{Return to the U.
S. Book of Common Prayer (1928)}

\end{quote}

\xchapter{Prayers and Thanksgivings}

\textbf{The Book of Common Prayer (1928)}

Prayers
and Thanksgivings

PRAYERS.

\emph{¶
To be used before the Prayer for all Conditions of Men, or, when that
is not said, before the final Prayer of Thanksgiving or of Blessing, or
before the Grace.}

\begin{quote}

\begin{quote}

\emph{A
Prayer for Congress.}

\emph{¶
To be used during their Session.}

\end{quote}

\end{quote}

MOST
gracious God, we humbly beseech thee, as for the people of these United
States in general, so especially for their Senate and Representatives
in Congress assembled; that thou wouldest be pleased to direct and prosper
all their consultations, to the advancement of thy glory, the good of
thy Church, the safety, honour, and welfare of thy people; that all things
may be so ordered and settled by their endeavours, upon the best and surest
foundations, that peace and happiness, truth and justice, religion and
piety, may be establish among us for all generations. These and all other
necessaries, for them, for us, and thy whole Church, we humbly beg in
the Name and mediation of Jesus Christ, our most blessed Lord and Saviour.
\emph{Amen}.

\begin{quote}

\begin{quote}

\emph{For
a State Legislature.}

\end{quote}

\end{quote}

O
GOD, the fountain of wisdom, whose statutes are good and gracious and
whose law is truth; We beseech thee so to guide and bless the Legislature
of this State, that it may ordain for our governance only such things
as please thee, to the glory of thy Name and the welfare of the people;
through Jesus Christ, thy Son, our Lord. \emph{Amen}.

\begin{quote}

\begin{quote}

\emph{For
Courts of Justice.}

\end{quote}

\end{quote}

ALMIGHTY
God, who sittest in the throne judging right; We humbly beseech thee to
bless the courts of justice and the magistrates in all this land; and
give unto them the spirit of wisdom and understanding, that they may
discern the truth and impartially administer the law in the fear of thee
alone; through him who shall come to be our judge, thy Son, our Saviour,
Jesus Christ. \emph{Amen}.

\begin{quote}

\begin{quote}

\emph{For
Our Country.}

\end{quote}

\end{quote}

ALMIGHTY
God, who hast given us this good land for our heritage; We humbly beseech
thee that we may always prove ourselves a people mindful of thy favour
and glad to do thy will. Bless our land with honourable industry, sound
learning, and pure manners. Save us from violence, discord, and confusion;
from pride and arrogancy, and from every evil way. Defend our liberties,
and fashion into one united people the multitudes brought hither out of
many kindreds and tongues. Endue with the spirit of wisdom those to whom
in thy Name we entrust the authority of government, that there may be
justice and peace at home, and that, through obedience to thy law, we
may show forth thy praise among the nations of the earth. In the time
of prosperity, fill our hearts with thankfulness, and in the day of trouble,
suffer not our trust in thee to fail; all which we ask through Jesus Christ
our Lord. \emph{Amen.}

\emph{A
Prayer to be used at the Meetings of Convention.}

ALMIGHTY
and everlasting God, who by thy Holy Spirit didst preside in the Council
of the blessed Apostles, and hast promised, through thy Son Jesus Christ,
to be with thy Church to the end of the world; We beseech thee to be with
the Council of thy Church \emph{here} assembled in thy Name and Presence.
Save \emph{us} from all error, ignorance, pride, and prejudice; and
of thy great mercy vouchsafe, we beseech thee, so to direct, sanctify,
and govern us in our work, by the mighty power of the Holy Ghost, that
the comfortable Gospel of Christ may be truly preached, truly received,
and truly followed, in all places, to the breaking down the kingdom of
sin, Satan, and death; till at length the whole of thy dispersed sheep,
being gathered into one fold, shall become partakers of everlasting life;
through the merits and death of Jesus Christ our Saviour. \emph{Amen}.

¶
\emph{During, or before, the session of any General or Diocesan Convention,
the above Prayer may be used by all Congregations of this Church. or of
the Diocese concerned; the clause,}here assembled in thy Name,\emph{being changed to}now assembled [or about to assemble] in thy Name
and Presence;\emph{and the clause,}govern us in our work,\emph{to}govern
them in their work\emph{.}

\begin{quote}

\begin{quote}

\emph{For
the Church}.

\end{quote}

\end{quote}

O
GRACIOUS Father, we humbly beseech thee for thy holy Catholic Church;
that thou wouldst be pleased to fill it with all truth, in all peace.
Where it is corrupt, purify it; where it is in error, direct it; where
in anything it is amiss, reform it. Where it is right, establish it; where
it is in want, provide for it; where it is divided, reunite it; for the
sake of him who died and rose again, and ever liveth to make intercession
for us, Jesus Christ, thy Son, our Lord. \emph{Amen}.

\begin{quote}

\begin{quote}

\emph{For
the Unity of God’s People.}

\end{quote}

\end{quote}

O
GOD, the Father of our Lord Jesus Christ, our only Saviour, the Prince
of Peace; Give us grace seriously to lay to heart the great dangers we
are in by our unhappy divisions. Take away all hatred and prejudice, and
whatsoever else may hinder us from godly, union and concord: that as there
is but one Body and one Spirit, and one hope of our calling, one Lord,
one Faith, one Baptism, one God and Father of us all, so we may be all
of one heart and of one soul, united in one holy bond of truth and peace,
of faith and charity, and may with one mind and one mouth glorify thee;
through Jesus Christ our Lord. \emph{Amen}.

\emph{For Missions.}

O GOD, who hast made of
one blood all nations of men for to dwell on the face of the whole earth,
and didst send thy blessed Son to preach peace to them that are far off
and to them that are nigh; Grant that all men everywhere may seek after
thee and find thee. Bring the nations into thy fold, pour out thy Spirit
upon all flesh, and hasten thy kingdom; through the same thy Son Jesus
Christ our Lord. \emph{Amen}.

\begin{quote}

\begin{quote}

\emph{¶
Or this}

\end{quote}

\end{quote}

ALMIGHTY
God, whose compassions fail not, and whose loving-kindness reacheth unto
the world’s end; We give thee humble thanks for opening heathen
lands to the light of thy truth; for making paths in the deep waters and
highways in the desert; and for planting thy Church in all the earth.
Grant, we beseech thee, unto us thy servants, that with lively faith we
may labour abundantly to make known to all men thy blessed gift of eternal
life; through Jesus Christ our Lord. \emph{Amen.}

\emph{For
those who are to be admitted into Holy Orders.}

¶
\emph{To be used in the Weeks preceding the stated Times of Ordination.}

ALMIGHTY
God, our heavenly Father, who hast purchased to thyself an universal Church
by the precious blood of thy dear Son; Mercifully look upon the same,
and at this time so guide and govern the minds of thy servants the Bishops
and Pastors of thy flock, that they may lay hands suddenly on no man,
but faithfully and wisely make choice of fit persons, to serve in the
sacred Ministry of thy Church, And to those who shall be ordained to any
holy function, give thy grace and heavenly benediction; that both by their
life and doctrine they may show forth thy glory, and set forward the salvation
of all men; through Jesus Christ our Lord. \emph{Amen}.

\begin{quote}

\begin{quote}

¶
\emph{Or this.}

\end{quote}

\end{quote}

ALMIGHTY
God, the giver of all good gifts, who of thy divine providence hast appointed
divers Orders in thy Church; Give thy grace, we humbly beseech thee, to
all those who are to be called to any office and administration in the
same; and so replenish them with the truth of thy doctrine, and endue
them with innocency of life, that they may faithfully serve before thee,
to the glory of thy great Name, and the benefit of thy holy Church; through
Jesus Christ our Lord. Amen.

\begin{quote}

\begin{quote}

\emph{For
the Increase of the Ministry}.

\end{quote}

\end{quote}

O
ALMIGHTY God, look mercifully upon the world which thou hast redeemed
by the blood of thy dear Son, and incline the hearts of many to dedicate
themselves to the sacred ministry of thy Church; through the same thy
Son Jesus Christ our Lord. \emph{Amen}.

\begin{quote}

\begin{quote}

\emph{For
Fruitful Seasons.}

\end{quote}

\end{quote}

\emph{¶
To be used on Rogation-Sunday and the Rogation-days.}

ALMIGHTY
God, who hast blessed the earth that it should be fruitful and bring forth
whatsoever is needful for the life of man, and hast commanded us to work
with quietness, and eat our own bread; Bless the labours of the husbandman,
and grant such seasonable weather that we may gather in the fruits of
the earth, and ever rejoice in thy goodness, to the praise of thy holy
Name; through Jesus Christ our Lord. \emph{Amen}.

\begin{quote}

\begin{quote}

\emph{¶
Or this.}

\end{quote}

\end{quote}

O
GRACIOUS Father, who openest thine hand and fillest all things living
with plenteousness; We beseech thee of thine infinite goodness to hear
us, who now make our prayers and supplications unto thee. Remember not
our sins, but thy promises of mercy. Vouchsafe to bless the lands and
multiply the harvests of the world. Let thy breath go forth that it may
renew the face of the earth. Show thy loving-kindness, that our land may
give her increase; and so fill us with good things that the poor and needy
may give thanks unto thy Name; through Christ our Lord. \emph{Amen}.

\begin{quote}

\begin{quote}

\emph{For
Rain.}

\end{quote}

\end{quote}

O
GOD, heavenly Father, who by thy Son Jesus Christ hast promised to all
those who seek thy kingdom, and the righteousness thereof, all things
necessary to their bodily sustenance; Send us, we beseech thee, in this
our necessity, such moderate rain and showers, that we may receive the
fruits of the earth to our comfort, and to thy honour; through Jesus Christ
our Lord. \emph{Amen}.

\begin{quote}

\begin{quote}

\emph{For
Fair Weather.}

\end{quote}

\end{quote}

ALMIGHTY
and most merciful Father, we humbly beseech thee, of thy great goodness,
to restrain those immoderate rains, wherewith thou hast afflicted us.
And we pray thee to send us such seasonable weather, that the earth may,
in due time, yield her increase for our use and benefit; through Jesus
Christ our Lord. \emph{Amen}.

\begin{quote}

\begin{quote}

\emph{In
Time of Dearth and Famine.}

\end{quote}

\end{quote}

O
GOD, heavenly Father, whose gift it is that the rain doth fall, and the
earth bring forth her increase; Behold, we beseech thee, the afflictions
of thy people; increase the fruits of the earth by thy heavenly benediction;
and grant that the scarcity and dearth, which we now most justly suffer
for our sins, may, through thy goodness, be mercifully turned into plenty;
for the love of Jesus Christ our Lord, to whom, with thee and the Holy
Ghost, be all honour and glory, now and for ever. \emph{Amen}.

\begin{quote}

\begin{quote}

\emph{In
Time of War and Tumults.}

\end{quote}

\end{quote}

ALMIGHTY
God, the supreme Governor of all things, whose power no creature is able
to resist, to whom it belongeth justly to punish sinners, and to be merciful
to those who truly repent; Save and deliver us, we humbly beseech thee,
from the hands of our enemies; that we, being armed with thy defence,
may be preserved evermore from all perils, to glorify thee, who art the
only giver of all victory; through the merits of thy Son, Jesus Christ
our Lord. \emph{Amen}.

\begin{quote}

\begin{quote}

\emph{In
Time of Calamity.}

\end{quote}

\end{quote}

O
GOD, merciful and compassionate, who art ever ready to hear the prayers
of those who put their trust in thee; Graciously hearken to us who call
upon thee, and grant us thy help in this our need; through Jesus Christ
our Lord. \emph{Amen}.

\begin{quote}

\begin{quote}

\emph{For
the Army.}

\end{quote}

\end{quote}

O
LORD God of Hosts, stretch forth, we pray thee, thine almighty arm to
strengthen and protect the soldiers of our country; Support them in the
day of battle, and in the time of peace keep them safe from all evil;
endue them with courage and loyalty; and grant that in all things they
may serve without reproach; through Jesus Christ our Lord. \emph{Amen}.

\begin{quote}

\begin{quote}

\emph{For
the Navy.}

\end{quote}

\end{quote}

O
ETERNAL Lord God, who alone spreadest out the heavens, and rulest the
raging of the sea; Vouchsafe to take into thy almighty and most gracious
protection our country’s Navy, and all who serve therein. Preserve
them from the dangers of the sea, and from the violence of the enemy;
that they may be a safeguard unto the United States of America, and a
security for such as pass on the seas upon their lawful occasions; that
the inhabitants of our land may in peace and quietness serve thee our
God, to the glory of thy Name; through Jesus Christ our Lord. \emph{Amen}.

\emph{Memorial
Days.}

ALMIGHTY
God, our heavenly Father, in whose hands are the living and the dead;
We give thee thanks for all those thy servants who have laid down their
lives in the service of our country. Grant to them thy mercy and the light
of thy presence, that the good work which thou hast begun in them may
be perfected; through Jesus Christ thy Son our Lord. \emph{Amen}.

\emph{For
Schools, Colleges. and Universities}.

ALMIGHTY
God, we beseech thee, with thy gracious favour, to behold our universities,
colleges, and schools, that knowledge may be increased among us, and all
good learning flourish and abound; bless all who teach and all who learn;
and grant that in humility of heart they may ever look unto thee, who
art the fountain of all wisdom; through Jesus Christ our Lord. \emph{Amen}.

\begin{quote}

\begin{quote}

\emph{For
Religious Education.}

\end{quote}

\end{quote}

ALMIGHTY
God, our Heavenly Father, who hast committed to thy Holy Church the care
and nurture of thy children; Enlighten with thy wisdom those who teach
and those who learn; that, rejoicing in the knowledge of thy truth, they
may worship thee and serve thee from generation to generation; through
Jesus Christ our Lord. \emph{Amen}.

\begin{quote}

\begin{quote}

\emph{For
Children.}

\end{quote}

\end{quote}

O
LORD Jesus Christ, who dost embrace children with the arms of thy mercy,
and dost make them living members of thy Church; Give them grace, we pray
thee, to stand fast in thy faith, to obey thy word, and to abide in thy
love; that being made strong by thy Holy Spirit they may resist temptation
and overcome evil; and may rejoice in the life that now is, and dwell
with thee in the life that is to come; through thy merits, O merciful
Saviour, who with the Father and the Holy Ghost livest and reignest one
God, world without end. \emph{Amen}.

\begin{quote}

\begin{quote}

\emph{For
Those About to be Confirmed}.

\end{quote}

\end{quote}

O
GOD, who through the teaching of thy Son Jesus Christ didst prepare the
disciples for the coming of the Comforter; Make ready, we beseech thee,
the hearts and minds of thy servants who at this time are seeking to be
strengthened by the gift of the Holy Spirit through the laying on of hands,
that, drawing near with penitent and faithful hearts, they may evermore
be filled with the power of his divine indwelling; through the same Jesus
Christ our Lord. \emph{Amen}.

\begin{quote}

\begin{quote}

\emph{For
Christian Service}.

\end{quote}

\end{quote}

O
LORD our heavenly Father, whose blessed Son came not to be ministered
unto, but to minister; We beseech thee to bless all who, following in
his steps, give themselves to the service of their fellow men. Endue them
with wisdom, patience, and courage, that they may strengthen the weak
and raise up those who fall; and, being inspired by thy love, may worthily
minister in thy Name to the suffering, the friendless, and the needy;
for the sake of him who laid down his life for us, the same thy Son, our
Saviour, Jesus Christ. \emph{Amen}.

\begin{quote}

\begin{quote}

\emph{For
Social Justice}.

\end{quote}

\end{quote}

ALMIGHTY
God, who hast created man in thine own image; Grant us grace fearlessly
to contend against evil, and to make no peace with oppression; and, that
we may reverently use our freedom, help us to employ it in the maintenance
of justice among men and nations, to the glory of thy holy Name; through
Jesus Christ our Lord. Amen.

\begin{quote}

\begin{quote}

\emph{For
Every Man in his Work}.

\end{quote}

\end{quote}

ALMIGHTY
God, our heavenly Father, who declarest thy glory and showest forth thy
handiwork in the heavens and in the earth; Deliver us, we beseech thee,
in our several callings, from the service of mammon, that we may do the
work which thou givest us to do, in truth, in beauty, and in righteousness,
with singleness of heart as thy servants, and to the benefit of our fellow
men; for the sake of him who came among us as one that serveth, thy Son,
Jesus Christ our Lord. \emph{Amen}.

\emph{For
the Family of Nations.}

ALMIGHTY
God, our heavenly Father, guide, we beseech thee, the Nations of the world
into the way of justice and truth, and establish among them that peace
which is the fruit of righteousness, that they may become the Kingdom
of our Lord and Saviour Jesus Christ. \emph{Amen.}

\emph{In Time of Great Sickness and Mortality.}

O
MOST mighty and merciful God, in this time of grievous sickness, we flee
unto thee for succour. Deliver us, we beseech thee, from our peril; give
strength and skill to all those who minister to the sick; prosper the
means made use of for their cure; and grant that, perceiving how frail
and uncertain our life is, we may apply our hearts unto that heavenly
wisdom which leadeth to eternal life; through Jesus Christ our Lord. \emph{Amen}.
\emph{For
a Sick Person}.

O
FATHER of mercies and God of all comfort, our only help in time of need;
We humbly beseech thee to behold, visit, and relieve thy sick \emph{servant}
[\emph{N}.] for whom our prayers are desired. Look upon \emph{him}
with the eyes of thy mercy; comfort \emph{him} with a sense of thy goodness;
preserve \emph{him} from the temptations of the enemy; and give \emph{him}
patience under \emph{his} affliction. In thy good time, restore \emph{him}
to health, and enable \emph{him} to lead the residue of \emph{his}
life in thy fear, and to thy glory; and grant that finally \emph{he}
may dwell with thee in life everlasting; through Jesus Christ our Lord.
\emph{Amen}.

\emph{For a Sick Child.}

O
HEAVENLY Father, watch with us, we pray thee, over the sick \emph{child}
for whom our prayers are offered, and grant that \emph{he} may be restored
to that perfect health which it is thine alone to give; through Jesus
Christ our Lord. \emph{Amen}.

\emph{For
a Person under Affliction.}

O
MERCIFUL God, and heavenly Father, who hast taught us in thy holy Word
that thou dost not willingly afflict or grieve the children of men; Look
with pity, we beseech thee, upon the sorrows of thy \emph{servant} for
whom our prayers are offered. Remember \emph{him}, O Lord, in mercy;
endue \emph{his soul} with patience; comfort \emph{him} with a sense
of thy goodness; lift up thy countenance upon \emph{him}, and give \emph{him}
peace; through Jesus Christ our Lord. \emph{Amen}.

\emph{For
a Person, or Persons, going to Sea.}

O
ETERNAL God, who alone spreadest out the heavens, and rulest the raging
of the sea; We commend to thy almighty protection, thy \emph{servant},
for whose preservation on the great deep our prayers are desired. Guard
\emph{him}, we beseech thee, from the dangers of the sea, from sickness,
from the violence of enemies, and from every evil to which \emph{he}
may be exposed. Conduct him in safety to the haven where he would be,
with a grateful sense of thy mercies; through Jesus Christ our Lord. \emph{Amen}.

\emph{For
Prisoners.}

O
GOD, who sparest when we deserve punishment, and in thy wrath rememberest
mercy; We humbly beseech thee, of thy goodness, to comfort and succour
all prisoners [\emph{especially those who are condemned to die}]. Give
them a right understanding of themselves, and of thy promises; that, trusting
wholly in thy mercy, they may not place their confidence anywhere but
in thee. Relieve the distressed, protect the innocent, awaken the guilty;
and forasmuch as thou alone bringest light out of darkness, and good out
of evil, grant to these thy servants, that by the power of thy Holy Spirit
they may be set free from the chains of sin, and may be brought to newness
of life; through Jesus Christ our Lord. \emph{Amen}.

\emph{A
Bidding Prayer}

\emph{¶ To be used before Sermons, or on Special
Occasions.}

\emph{¶
And}NOTE\emph{,
That the Minister, in his discretion, may omit any of the clauses in this
Prayer, or may add others, as occasion may require.}

GOOD
Christian People, I bid your prayers for Christ’s holy Catholic
Church, the blessed company of all faithful people; that it may please
God to confirm and strengthen it in purity of faith, in holiness of life,
and in perfectness of love, and to restore to it the witness of visible
unity; and more especially for that branch of the same planted by God
in this land, whereof we are members; that in all things it may work according
to God’s will, serve him faithfully, and worship him acceptably.

Ye shall pray for the President of these United States,
and for the Governor of this State, and for all that are in authority;
that all, and every one of them, may serve truly in their several callings
to the glory of God, and the edifying and well-governing of the people,
remembering the account they shall be called upon to give at the last
great day.

Ye shall also pray for the ministers of God’s
Holy Word and Sacraments; for Bishops (and herein more especially for
the Bishop of this Diocese), that they may minister faithfully and wisely
the discipline of Christ; likewise for all Priests and Deacons (and herein
more especially for the Clergy here residing), that they may shine as
lights in the world, and in all things may adorn the doctrine of God our
Saviour.

And ye shall pray for a due supply of persons fitted
to serve God in the Ministry and in the State; and to that end, as well
as for the good education of all the youth of this land, ye shall pray
for all schools, colleges, and seminaries of sound and godly learning,
and for all whose hands are open for their maintenance; that whatsoever
tends to the advancement of true religion and useful learning may for
ever flourish and abound.

Ye shall pray for all the people of these United States,
that they may live in the true faith and fear of God, and in brotherly
charity one towards another.

Ye shall pray also for all who travel by land or sea;
for all prisoners and captives; for all who are in sickness or in sorrow;
for all who have fallen into grievous sin; for all who, through temptation,
ignorance, helplessness, grief, trouble, dread, or the near approach of
death, especially need our prayers.

Ye shall also praise God for rain and sunshine; for
the fruits of the earth; for the products of all honest industry; and
for all his good gifts, temporal and spiritual, to us and to all men.

Finally, ye shall yield unto God most high praise and
hearty thanks for the wonderful grace and virtue declared in all his saints,
who have been the choice vessels of his grace and the lights of the world
in their several generations; and pray unto God, that we may have grace
to direct our lives after their good examples; that, this life ended,
we may be made partakers with them of the glorious resurrection, and the
life everlasting.

And now, brethren, summing up all our petitions, and
all our thanksgivings, in the words which Christ hath taught us, we make
bold to say,

OUR
Father, who art in heaven, Hallowed be thy Name. Thy kingdom come. Thy
will be done on earth, As it is in heaven. Give us this day our daily
bread. And forgive us our trespasses, As we forgive those who trespass
against us. And lead us not into temptation; But deliver us from evil:
For thine is the kingdom, and the power, and the glory, for ever and ever.
Amen.

COLLECTS.
\emph{}

\emph{¶
To be used after the Collects of Morning or Evening Prayer, or Communion,
at the discretion of the Minister.}

O
LORD Jesus Christ, who saidst unto thine Apostles, Peace I leave with
you, my peace I give unto you; Regard not our sins, but the faith of thy
Church; and grant to it that peace and unity which is according to thy
will, who livest and reignest with the Father and the Holy Ghost, one
God, world without end. \emph{Amen}.

ASSIST
us mercifully, O Lord, in these our supplications and prayers, and dispose
the way of thy servants towards the attainment of everlasting salvation;
that, among all the changes and chances of this mortal life, they may
ever be defended by thy most gracious and ready help; through Jesus Christ
our Lord. \emph{Amen}.

GRANT,
we beseech thee, Almighty God, that the words which we have heard this
day with our outward ears, may, through thy grace, be so grafted inwardly
in our hearts, that they may bring forth in us the fruit of good living,
to the honour and praise of thy Name; through Jesus Christ our Lord. \emph{Amen}.

DIRECT
us, O Lord, in all our doings, with thy most gracious favour, and further
us with thy continual help; that in all our works begun, continued, and
ended in thee, we may glorify thy holy Name, and finally, by thy mercy,
obtain everlasting life; through Jesus Christ our Lord. \emph{Amen}.

ALMIGHTY
God, the fountain of all wisdom, who knowest our necessities before we
ask, and our ignorance in asking; We beseech thee to have compassion upon
our infirmities; and those things which for our unworthiness we dare not,
and for our blindness we cannot ask, vouchsafe to give us, for the worthiness
of thy Son Jesus Christ our Lord. \emph{Amen}.

ALMIGHTY
God, who hast promised to hear the petitions of those who ask in thy Son’s
Name; We beseech thee mercifully to incline thine ears to us who have
now made our prayers and supplications unto thee; and grant that those
things which we have faithfully asked according to thy will, may effectually
be obtained, to the relief of our necessity, and to the setting forth
of thy glory; through Jesus Christ our Lord. \emph{Amen}.

THANKSGIVINGS.

\emph{¶
To be used after the General Thanksgiving, or, when that is not said,
before the final Prayer of Blessing or the Benediction.}

\emph{A Thanksgiving
to Almighty God for the Fruits of the Earth and all the other Blessings
of his merciful Providence.}

MOST
gracious God, by whose knowledge the depths are broken up, and the clouds
drop down the dew; We yield thee unfeigned thanks and praise for the return
of seed-time and harvest, for the increase of the ground and the gathering
in of the fruits thereof, and for all the other blessings of thy merciful
providence bestowed upon this nation and people. And, we beseech thee,
give us a just sense of these great mercies; such as may appear in our
lives by an humble, holy, and obedient walking before thee all our days;
through Jesus Christ our Lord, to whom, with thee and the Holy Ghost,
be all glory and honour, world without end. \emph{Amen}.

\emph{The
Thanksgiving of Women after Child-birth}

\emph{¶
To be said when any, Woman, being present in Church, shall have desired
to return thanks to Almighty God for her safe deliverance.}

O
ALMIGHTY God, we give thee humble thanks for that thou hast been graciously
pleased to preserve, through the great pain and peril of childbirth, \emph{this
woman}, thy \emph{servant}, who desireth now to offer \emph{her}
praises and thanksgivings unto thee. Grant, we beseech thee, most merciful
Father, that \emph{she}, through thy help, may both faithfully live
and walk according to thy will in this life present, and also may be \emph{partaker}
of everlasting glory in the life to come; through Jesus Christ our Lord.
\emph{Amen}.

\emph{For Rain.}

O
GOD, our heavenly Father, by whose gracious providence the former and
the latter rain descend upon the earth, that it may bring forth fruit
for the use of man; We give thee humble thanks that it hath pleased thee
to send us rain to our great comfort, and to the glory of thy holy Name;
through Jesus Christ our Lord. \emph{Amen}.

\begin{quote}

\begin{quote}

\emph{For
Fair Weather.}

\end{quote}

\end{quote}

O
LORD God, who hast justly humbled us by thy late visitation of us with
immoderate rain and waters, and in thy mercy hast relieved and comforted
our souls by this seasonable and blessed change of weather; We praise
and glorify thy holy Name for this thy mercy, and will always declare
thy loving-kindness from generation to generation; through Jesus Christ
our Lord. \emph{Amen}.

\begin{quote}

\begin{quote}

\emph{For
Plenty.}

\end{quote}

\end{quote}

O
MOST merciful Father, who of thy gracious goodness hast heard the devout
prayers of. thy Church, and turned our dearth and scarcity into plenty;
We give thee humble thanks for this thy special bounty; beseeching thee
to continue thy loving-kindness unto us, that our land may yield us her
fruits of increase, to thy glory and our comfort; through Jesus Christ
our Lord. \emph{Amen}.

\emph{For
Peace, and Deliverance from our Enemies.}

O
ALMIGHTY God, who art a strong tower of defence unto thy servants against
the face of their enemies; We yield thee praise and thanksgiving for our
deliverance from those great and apparent dangers wherewith we were compassed.
We acknowledge it thy goodness that we were not delivered over as a prey
unto them; beseeching thee still to continue such thy mercies towards
us, that all the world may know that thou art our Saviour and mighty Deliverer;
through Jesus Christ our Lord. \emph{Amen}.

\begin{quote}

\begin{quote}

\emph{For
Restoring Public Peace at Home.}

\end{quote}

\end{quote}

O
ETERNAL God, our heavenly Father, who alone makest men to be of one mind
in a house, and stillest the outrage of a violent and unruly people; We
bless thy holy Name, that it hath pleased thee to appease the seditious
tumults which have been lately raised up amongst us; most humbly beseeching
thee to grant to all of us grace, that we may henceforth obediently walk
in thy holy commandments; and, leading a quiet and peaceable life in all
godliness and honesty, may continually offer unto thee our sacrifice of
praise and thanksgiving for these thy mercies towards us; through Jesus
Christ our Lord. \emph{Amen}.

\begin{quote}

\begin{quote}

\emph{For
Recovery from Sickness.}

\end{quote}

\end{quote}

O
GOD, who art the giver of life, of health, and of safety; We bless thy
Name, that thou hast been pleased to deliver from \emph{his} bodily
sickness \emph{this} thy \emph{servant}, who now \emph{desireth}
to return thanks unto thee, in the presence of all thy people. Gracious
art thou, O Lord, and full of compassion to the children of men. May \emph{his
heart} be duly impressed with a sense of thy merciful goodness, and
may \emph{he} devote the residue of \emph{his} days to an humble,
holy, and obedient walking before thee; through Jesus Christ our Lord.
\emph{Amen}.

\begin{quote}

\begin{quote}

\emph{For
a Child’s Recovery from Sickness.}

\end{quote}

\end{quote}

ALMIGHTY
God and heavenly Father, we give thee humble thanks for that thou hast
been graciously pleased to deliver from \emph{his} bodily sickness the
\emph{child} in whose behalf we bless and praise thy Name, in the presence
of all thy people. Grant, we beseech thee, O gracious Father, that. \emph{he},
through thy help, may both faithfully live in this world according to
thy will, and also may be \emph{partaker} of everlasting glory in the
life to come; through Jesus Christ our Lord. \emph{Amen}.

\begin{quote}

\begin{quote}

\emph{For
a Safe Return from a Journey.}

\end{quote}

\end{quote}

MOST
gracious Lord, whose mercy is over all thy works; We praise thy holy Name
that thou hast been pleased to conduct in safety, through the perils of
\emph{the great deep}, [\emph{his way},] \emph{this} thy \emph{servant},
who now \emph{desireth} to return \emph{his} thanks unto thee in thy
holy Church. May \emph{he} be duly sensible of thy merciful providence
towards \emph{him}, and ever express \emph{his} thankfulness by a
holy trust in thee, and obedience to thy laws; through Jesus Christ our
Lord. \emph{Amen}.

\begin{quote}

\emph{Return to the U.
S. Book of Common Prayer (1928)}

\end{quote}

\xchapter{The Collects, Epistles, and Gospels}

\textbf{The Book of Common Prayer (1928)}

A
Saint’s Day.

\emph{The
Collect.}

ALMIGHTY
and everlasting God, who dost enkindle the name of thy love in the hearts
of the Saints; Grant to us, thy humble servants, the same faith and power
of love; that, as we rejoice in their triumphs, we may profit by their
examples; through Jesus Christ our Lord. \emph{Amen}.

¶
\emph{Or this}.

O
ALMIGHTY God, who hast called us to faith in thee, and hast compassed
us about with so great a cloud of witnesses; Grant that we, encouraged
by the good exam-ples of thy Saints, and especially of thy servant [Saint——],
may persevere in running the race that is set before us, until at length,
through thy mercy, we, with them, attain to thine eternal joy; through
him who is the author and finisher of our faith, thy Son Jesus Christ
our Lord. \emph{Amen}.

\emph{The Epistle.} Hebrews xii. 1.

SEEING
we also are compassed about with so great a cloud of witnesses, let us
lay aside every weight, and the sin which doth so easily beset us, and
let us run with patience the race that is set before us, looking unto
Jesus the author and finisher of our faith; who for the joy that was set
before him endured the cross, despising the shame, and is set down at
the right hand of the throne of God.

\emph{The Gospel}. St. Matthew xxv. 31.

WHEN
the Son of man shall come in his glory, and all the holy angels with him,
then shall he sit upon the throne of his glory: and before him shall be
gathered all nations: and he shall separate them one from another, as
a shepherd divideth his sheep from the goats: and he shall set the sheep
on his right hand, but the goats on the left. Then shall the King say
unto them on his right hand, Come, ye blessed of my Father, inherit the
kingdom pre-pared for you from the foundation of the world: for I was
an hungred, and ye gave me meat: I was thirsty, and ye gave me drink:
I was a stranger, and ye took me in: naked, and ye clothed me: I was sick,
and ye visited me: I was in prison, and ye came unto me. Then shall the
righteous answer him, saying, Lord, when saw we thee an hungred, and fed
thee? or thirsty, and gave thee drink? When saw we thee a stranger, and
took thee in? or naked, and clothed thee? Or when saw we thee sick, or
in prison, and came unto thee? And the King shall answer and say unto
them, Verily I say unto you, Inasmuch as ye have done it unto one of the
least of these my brethren, ye have done it unto me.

Feast
of the Dedication of a Church.

\emph{The
Collect.}

O
GOD, whom year by year we praise for the dedication of this church; Hear,
we beseech thee, the prayers of thy people, and grant that whosoever shall
worship before thee in this place, may obtain thy merciful aid and protection;
through Jesus Christ our Lord. \emph{Amen}.

\emph{The Epistle}. 1 St. Peter ii. 1.

LAYING
aside all malice, and all guile, and hypocrisies, and envies, and all
evil speakings, as newborn babes, desire the sincere milk of the word,
that ye may grow thereby: if so be ye have tasted that the Lord is gracious.
To whom coming, as unto a living stone, disallowed indeed of men, but
chosen of God, and precious, ye also, as lively stones, are built up a
spiritual house, an holy priesthood, to offer up spiritual sacrifices,
acceptable to God by Jesus Christ.

\emph{The Gospel}. St. Matthew xxi. 12.

JESUS
went into the temple of God, and cast out all them that sold and bought
in the temple, and overthrew the tables of the money-changers, and the
seats of them that sold doves, and said unto them, It is written, My house
shall be called the house of prayer; but ye have made it a den of thieves.
And the blind and the lame came to him in the temple; and he healed them.
And when the chief priests and scribes saw the wonderful things that he
did, and the children crying in the temple, and saying, Hosanna to the
son of David; they were sore displeased, and said unto him, Hearest thou
what these say? And Jesus saith unto them, Yea; have ye never read, Out
of the mouth of babes and sucklings thou hast perfected praise?

This
page gives propers for occasions not provided for in earlier books. These
were simply added on to the end of the propers for Holy Days, with no
separate heading or designation

The Ember Days

At the Four Seasons.

\emph{The Collect.}

O
ALMIGHTY God, who hast committed to the hands of men the ministry of reconciliation;
We humbly beseech thee, by the inspiration of thy Holy Spirit, to put
it into the hearts of many to offer themselves for this ministry; that
thereby mankind may be drawn to thy blessed kingdom; through Jesus Christ
our Lord. \emph{Amen}.

\emph{For the Epistle}. Acts xiii. 44.

THE
next sabbath day came almost the whole city together to hear the word
of God. But when the Jews saw the multitudes, they were filled with envy,
and spake against those things which were spoken by Paul, contradicting
and blaspheming. Then Paul and Barnabas waxed bold, and said, It was necessary
that the word of God should first have been spoken to you: but seeing
ye put it from you, and judge yourselves unworthy of everlasting life,
lo, we turn to the Gentiles. For so hath the Lord commanded us, saying,
I have set thee to be a light of the Gentiles, that thou shouldest be
for salvation unto the ends of the earth. And when the Gentiles heard
this, they were glad, and glorified the word of the Lord: and as many
as were ordained to eternal life believed. And the word of the Lord was
published throughout all the region.

\emph{The Gospel}. St. Luke iv. 16.

JESUS
came to Nazareth, where he had been brought up: and, as his custom was,
he went into the synagogue on the sabbath day, and stood up for to read.
And there was delivered unto him the book of the prophet Esaias. And when
he had opened the book, he found the place where it was written, The Spirit
of the Lord is upon me, because he hath anointed me to preach the gospel
to the poor; he hath sent me to heal the brokenhearted, to preach deliverance
to the captives, and recovering of sight to the blind, to set at liberty
them that are bruised, to preach the acceptable year of the Lord. And
he closed the book, and he gave it again to the minister, and sat down.
And the eyes of all them that were in the synagogue were fastened on him.
And he began to say unto them, This day is this scripture fulfilled in
your ears.

The Rogation
Days

Being
the Three Days before Ascension Day.

\emph{The
Collect.}

ALMIGHTY
God, Lord of heaven and earth; We beseech thee to pour forth thy blessing
upon this land, and to give us a fruitful season; that we, constantly
receiving thy bounty, may evermore give thanks unto thee in thy holy Church;
through Jesus Christ our Lord. Amen.

\emph{For the Epistle}. Ezekiel xxxiv. 25.

I WILL make with them a
covenant of peace, and will cause the evil beasts to cease out of the
land: and they shall dwell safely in the wilderness, and sleep in the
woods. And I will make them and the places round about my hill a blessing;
and I will cause the shower to come down in his season; there shall be
showers of blessing. And the tree of the field shall yield her fruit,
and the earth shall yield her increase, and they shall be safe in their
land, and shall know that I am the Lord, when I have broken the bands
of their yoke, and delivered them out of the hand of those that served
themselves of them. And they shall no more be a prey to the heathen, neither
shall the beast of the land devour them; but they shall dwell safely,
and none shall make them afraid. And I will raise up for them a plant
of renown, and they shall be no more consumed with hunger in the land,
neither bear the shame of the heathen any more. Thus shall they know that
I the Lord their God am with them, and that they, even the house of Israel,
are my people, saith the Lord God. And ye my flock, the flock of my pasture,
are men, and I am your God, saith the Lord God.

\emph{The Gospel.} St. Luke xi. 5.

JESUS said unto them, Which
of you shall have a friend, and shall go unto him at midnight, and say
unto him, Friend, lend me three loaves; for a friend of mine in his journey
is come to me, and I have nothing to set before him? and he from within
shall answer and say, Trouble me not: the door is now shut, and my children
are with me in bed; I cannot rise and give thee. I say unto you, Though
he will not rise and give him, because he is his friend, yet because of
his importunity he will rise and give him as many as he needeth. And I
say unto you, Ask, and it shall be given you; seek, and ye shall find;
knock, and it shall be opened unto you. For every one that asketh receiveth;
and he that seeketh findeth; and to him that knocketh it shall be opened.
If a son shall ask bread of any of you that is a father, will he give
him a stone? or if he ask a fish, will he for a fish give him a serpent?
or if he shall ask an egg, will he offer him a scorpion? If ye then, being
evil, know how to give good gifts unto your children: how much more shall
your heavenly Father give the Holy Spirit to them that ask him?

Independence Day.

[July 4.]

\emph{The Collect.}

O
ETERNAL God, through whose mighty power our fathers won their liberties
of old; Grant, we beseech thee, that we and all the people of this land
may have grace to maintain these liberties in righteousness and peace;
through Jesus Christ our Lord. \emph{Amen.}

\emph{For the Epistle}. Deuteronomy x. 17.

THE Lord your God is God
of gods, and Lord of lords, a great God, a mighty, and a terrible, which
regardeth not persons, nor taketh reward: he doth execute the judg-ment
of the fatherless and widow, and loveth the stranger, in giving him food
and raiment. Love ye therefore the stranger: for ye were strangers in
the land of Egypt. Thou shalt fear the Lord thy God; him shalt thou serve,
and to him shalt thou cleave, and swear by his name. He is thy praise,
and he is thy God, that hath done for thee these great and terrible things,
which thine eyes have seen.

\emph{The Gospel}. St. Matthew v. 43.

JESUS said, Ye have heard
that it hath been said, Thou shalt love thy neighbour, and hate thine
enemy. But I say unto you, Love your enemies, bless them that curse you,
do good to them that hate you, and pray for them which despitefully use
you, and persecute you; that ye may be the children of your Father which
is in heaven: for he maketh his sun to rise on the evil and on the good,
and sendeth rain on the just and on the unjust. For if ye love them which
love you, what reward have ye? do not even the publicans the same? And
if ye salute your brethren only, what do ye more than others? do not even
the publicans so? Be ye therefore perfect, even as your Father which is
in heaven is perfect.

Thanksgiving Day

¶ \emph{Instead of the Venite, the
following shall be said or sung.}

O PRAISE the Lord, for
it is a good thing to sing praises unto our God; * yea, a joyful and pleasant
thing it is to be thankful.

The Lord doth build up Jerusalem, * and gather together
the outcasts of Israel.

He healeth those that are broken in heart, * and giveth
medicine to heal their sickness.

O sing unto the Lord with thanksgiving; * sing praises
upon the harp unto our God:

Who covereth the heaven with clouds, and prepareth
rain for the earth; * and maketh the grass to grow upon the mountains,
and herb for the use of men;

Who giveth fodder unto the cattle, * and feedeth the
young ravens that call upon him.

Praise the Lord, O Jerusalem; * praise thy God, O Sion.

For he hath made fast the bars of thy gates, * and
hath blessed thy children within thee.

He maketh peace in thy borders, * and filleth thee
with the flour of wheat.

Glory be to the Father, and to the Son, * and to the
Holy Ghost;

As it was in the beginning, is now, and ever shall
be, * world without end. Amen.

\emph{The Collect.}

O MOST merciful Father,
who hast blessed the labours of the husbandman in the returns of the fruits
of the earth; We give thee humble and hearty thanks for this thy bounty;
beseeching thee to continue thy loving-kindness to us, that our land may
still yield her increase, to thy glory and our comfort; through Jesus
Christ our Lord. \emph{Amen}.

\emph{The Epistle.} St. James i. 16.

DO not err, my beloved
brethren. Every good gift and every perfect gift is from above, and cometh
down from the Father of lights, with whom is no variableness, neither
shadow of turning. Of his own will begat he us with the word of truth,
that we should be a kind of firstfruits of his creatures. Wherefore, my
beloved brethren, let every man be swift to hear, slow to speak, slow
to wrath: for the wrath of man worketh not the righteousness of God. Wherefore
lay apart all filthiness and superfluity of naughtiness, and receive with
meekness the engrafted word, which is able to save your souls. But be
ye doers of the word, and not hearers only, deceiving your own selves.
For if any be a hearer of the word, and not a doer, he is like unto a
man beholding his natural face in a glass: for he beholdeth himself, and
goeth his way, and straightway forgetteth what manner of man he was. But
whoso looketh into the perfect law of liberty, and continueth therein,
he being not a forgetful hearer, but a doer of the work, this man shall
be blessed in his deed. If any man among you seem to be religious, and
bridleth not his tongue, but deceiveth his own heart, this man's religion
is vain. Pure religion and undefiled before God and the Father is this,
To visit the fatherless and widows in their affliction, and to keep himself
unspotted from the world.

\emph{The Gospel}. St. Matthew vi. 25.

JESUS
said, Be not anxious for your life, what ye shall eat, or what ye shall
drink; nor yet for your body, what ye shall put on. Is not the life more
than food, and the body than raiment? Behold the fowls of the air: for
they sow not, neither do they reap, nor gather into barns; yet your heavenly
Father feedeth them. Are ye not much better than they? Which of you by
being anxious can add one cubit unto the measure of his life? And why
are ye anxious for raiment? Consider the lilies of the field, how they
grow; they toil not, neither do they spin: and yet I say unto you, That
even Solomon in all his glory was not arrayed like one of these. Wherefore,
if God so clothe the grass of the field, which to-day is, and to-morrow
is cast into the oven, shall he not much more clothe you, O ye of little
faith? There-fore be not anxious, saying, What shall we eat? or, What
shall we drink? or, Wherewithal shall we be clothed? (For after all these
things do the Gentiles seek:) for your heavenly Father knoweth that ye
have need of all these things. But seek ye first the kingdom of God, and
his righteousness; and all these things shall be added unto you. Be not
therefore anxious for the morrow: for the morrow shall take thought for
the things of itself. Sufficient unto the day is the evil thereof.

At a Marriage.

\emph{The Collect.}

O
ETERNAL God, we humbly beseech thee, favourably to behold these thy servants
now (\emph{or} about to be) joined in wedlock according to thy holy
ordinance; and grant that they, seeking first thy kingdom and thy right-eousness,
may obtain the manifold blessings of thy grace; through Jesus Christ our
Lord. \emph{Amen.}

\emph{The Epistle.} Ephesians v. 20.

GIVE thanks always for
all things unto God and the Father in the name of our Lord Jesus Christ;
submit-ting yourselves one to another in the fear of God. Wives, submit
yourselves unto your own husbands, as unto the Lord. For the husband is
the head of the wife, even as Christ is the head of the church: and he
is the saviour of the body. Therefore as the church is subject unto Christ,
so let the wives be to their own husbands in every thing. Husbands, love
your wives, even as Christ also loved the church, and gave himself for
it; that he might sanctify and cleanse it with the washing of water by
the word, that he might present it to himself a glorious church, not having
spot, or wrinkle, or any such thing; but that it should be holy and without
blemish. So ought men to love their wives as their own bodies. He that
loveth his wife loveth himself. For no man ever yet hated his own flesh;
but nourisheth and cherisheth it, even as the Lord the church: for we
are members of his body, of his flesh, and of his bones. For this cause
shall a man leave his father and mother, and shall be joined unto his
wife, and they two shall be one flesh. This is a great mystery: but I
speak concerning Christ and the church. Nevertheless let every one of
you in particular so love his wife even as himself; and the wife see that
she reverence her husband.

\emph{The Gospel}. St. Matthew xix. 4.

JESUS answered and said
unto them, Have ye not read, that he which made them at the beginning
made them male and female, and said, For this cause shall a man leave
father and mother, and shall cleave to his wife: and they twain shall
be one flesh? Wherefore they are no more twain, but one flesh. What therefore
God hath joined together, let not man put asunder.

At the Burial of the Dead.

\emph{The Collect.}

O ETERNAL Lord God, who
holdest all souls in life; Vouchsafe, we beseech thee, to thy whole Church
in paradise and on earth, thy light and thy peace; and grant that we,
following the good examples of those who have served thee here and are
now at rest, may at the last enter with them into thine unending joy;
through Jesus Christ our Lord. Amen.

\emph{¶ Or this.}

O GOD, whose mercies cannot
be numbered; Accept our prayers on behalf of the soul of thy servant departed,
and grant him an entrance into the land of light and joy, in the fellowship
of thy saints; through Jesus Christ our Lord. \emph{Amen}.

\emph{The Epistle}. 1 Thessalonians iv. 13.

I WOULD not have you to
be ignorant, brethren, concern-ing them which are asleep, that ye sorrow
not, even as others which have no hope. For if we believe that Jesus died
and rose again, even so them also which sleep in Jesus will God bring
with him. For this we say unto you by the word of the Lord, that we which
are alive and remain unto the coming of the Lord shall not prevent them
which are asleep. For the Lord himself shall descend from heaven with
a shout, with the voice of the archangel, and with the trump of God: and
the dead in Christ shall rise first: then we which are alive and remain
shall be caught up together with them in the clouds, to meet the Lord
in the air: and so shall we ever be with the Lord. Wherefore comfort one
another with these words.

\emph{The Gospel.}St. John vi. 37.

JESUS said unto them, All
that the Father giveth me shall come to me; and him that cometh to me
I will in no wise cast out. For I came down from heaven, not to do mine
own will, but the will of him that sent me. And this is the Father’s
will which hath sent me, that of all which he hath given me I should lose
nothing, but should raise it up again at the last day. And this is the
will of him that sent me, that every one which seeth the Son, and believeth
on him, may have everlasting life: and I will raise him up at the last
day.

\begin{quote}

\emph{Return to the U.
S. Book of Common Prayer (1928)}

\end{quote}

\xchapter{The Order for the Holy Communion}

\textbf{The Book of Common Prayer (1928)}

\begin{quote}

\begin{quote}

The Order for

The
Administration of the Lord's Supper

or

Holy
Communion

\end{quote}

\end{quote}

\emph{¶
At the Communion time the Holy Table shall have upon it a fair white linen
cloth, and the Priest, standing reverently before the Holy Table, shall
say the}Lord’s Prayer\emph{and the}Collect\emph{following,
the people kneeling; but the}Lord’s Prayer\emph{may be omitted
at the discretion of the Priest.}

OUR
Father, who art in heaven, Hallowed be thy Name. Thy kingdom come. Thy
will be done on earth, As it is in heaven. Give us this day our daily
bread. And forgive us our trespasses, As we forgive those who trespass
against us. And lead us not into temptation; But deliver us from evil.
Amen.

\begin{quote}

\begin{quote}

\emph{The
Collect.}

\end{quote}

\end{quote}

ALMIGHTY
God, unto whom all hearts are open, all desires known, and from whom no
secrets are hid; Cleanse the thoughts of our hearts by the inspiration
of thy Holy Spirit, that we may perfectly love thee, and worthily magnify
thy holy Name; through Christ our Lord. \emph{Amen}.

¶
\emph{Then shall the Priest, turning to the People, rehearse distinctly
The Ten Commandments; and the People, still kneeling, shall, after every
Commandment, ask God mercy for their transgressions for the time past,
and grace to keep the law for the time to come.}

\emph{¶
And}NOTE,\emph{that in rehearsing the Ten Commandments, the Priest may omit that part
of the Commandment which is inset.}

\emph{¶
The Decalogue may be omitted, provided it be said at least one Sunday
in each month. But}NOTE,\emph{That when ever it is omitted, the Minister shall say the Summary of the
Law, beginning,}Hear what our Lord Jesus Christ saith.

GOD
spake these words, and said:

\textbf{}I am the LORD
thy God; Thou shalt have none other gods but me.

\emph{Lord, have mercy upon us, and incline our hearts
to keep this law.}

Thou shalt not make to thyself any graven image, nor the likeness of any
thing that is in heaven above, or in the earth beneath, or in the water
under the earth; thou shalt not bow down to them, nor worship them:

for I the Lord thy God am a jealous
God, and visit the sins of the fathers upon the children, unto the third
and fourth generation of them that hate me; and show mercy unto thousands
in them that love me and keep my commandments.

\emph{Lord,
have mercy upon us, and incline our hearts to keep this law.}

Thou
shalt not take the Name of the Lord thy God in vain;

for the Lord will not hold him guiltless,
that taketh his Name in vain.

\emph{Lord, have mercy upon vs, and incline our hearts
to keep this law.}

Remember that thou keep holy the Sabbath-day.

Six days shalt thou labour, and do
all that thou hast to do; but the seventh day is the Sabbath of the
Lord thy God. In it thou shalt do no manner of work; thou, and thy son,
and thy daughter, thy man-servant, and thy maid-servant, thy cattle,
and the stranger that is within thy gates. For in six days the Lord
made heaven and earth, the sea, and all that in them is, and rested
the seventh day: wherefore the Lord blessed the seventh day, and hallowed
it.

\emph{Lord, have mercy upon us, and incline our hearts
to keep this law.}

Honour thy father and thy mother;

that thy days may be long in the land
which the Lord thy God giveth thee.

\emph{Lord, have mercy upon us, and incline our’
hearts to keep this law.}

Thou shalt do no murder.

\emph{Lord, have mercy upon us, and incline our hearts
to keep this law.}

Thou shalt not commit adultery.

\emph{Lord, have mercy upon us, and incline our hearts
to keep this law.}

Thou shalt not steal.

\emph{Lord, have mercy upon us, and incline our hearts
to keep this law.}

Thou shalt not bear false witness against thy neighbor.

\emph{Lord, have mercy upon us, and incline our hearts
to keep this law.}

Thou shalt not covet

thy neighbor’s house, thou shalt
not covet thy neighbour’s wife, nor his servant, nor his maid,
nor his ox, nor his ass, nor any thing that is his.

\emph{Lord, have mercy upon us, and write all these
thy laws in our hearts, we beseech thee.}

\emph{¶
Then may the Priest say,}
Hear
what our Lord Jesus Christ saith.
THOU
shalt love the Lord thy God with all thy heart, and with all thy soul,
and with all thy mind. This is the first and great commandment. And the
second is like unto it; Thou shalt love thy neighbor as thyself. On these
two commandments hang all the Law and the Prophets.

\emph{¶
Here, if the Decalogue hath been omitted, shall be said,}

Lord,
have mercy upon us.

\emph{Christ, have mercy upon us.}

Lord, have mercy upon us.

\emph{¶ Then the Priest may say,}

O
ALMIGHTY Lord, and everlasting God, vouchsafe, we beseech thee, to direct,
sanctify, and govern, both our hearts and bodies, in the ways of thy laws,
and in the works of thy commandments; that, through thy most mighty protection,
both here and ever, we may be preserved in body and soul; through our
Lord and Saviour Jesus Christ. \emph{Amen}.

\emph{¶
Here shall be said,}

\emph{Answer. Minister.} The Lord be with
you. \\
And
with thy spirit. 

Let us pray.

\emph{¶
Then shall the Priest say the Collect of the Day. And after the Collect
the Minister appointed shall read the Epistle, first saying,}The
Epistle is written in the — Chapter of —, beginning at the
— Verse. \emph{The Epistle ended, he shall say,}Here endeth the
Epistle.

\emph{¶
Here may be sung a Hymn or an Anthem.}

\emph{¶
Then, all the People standing, the Minister appointed shall read the Gospel,
first saying,}\emph{}The Holy Gospel is written in the — Chapter of —, beginning
at the — Verse.

\begin{quote}

\begin{quote}

\emph{¶
Here shall be said,}

Glory
be to thee, O Lord.

\emph{¶ And after the Gospel may be said,}

Praise
be to thee, O Christ.

\end{quote}

\end{quote}

\emph{¶ Then shall be said the Creed commonly called the Nicene, or
else the Apostles’ Creed; but the Creed may be omitted, if it hath
been said immediately before in Morning Prayer; Provided, That the Nicene
creed shall be said on Christmas-day, Easter-day, Ascension-day, Whitsunday,
and Trinity-Sunday.}

I
BELIEVE
in one God the Father Almighty, Maker of heaven and earth, And of all
things visible and invisible:

And in one Lord Jesus Christ, the only-begotten Son
of God; Begotten of his Father before all worlds, God of God, Light of
Light, Very God of very God; Begotten, not made; Being of one substance
with the Father; By whom all things were made: Who for us men and for
our salvation came down from heaven, And was incarnate by the Holy Ghost
of the Virgin Mary, And was made man: And was crucified also for us under
Pontius Pilate; He suffered and was buried: And the third day he rose
again according to the Scriptures: And ascended into heaven, And sitteth
on the right hand of the Father: And he shall come again, with glory,
to judge both the quick and the dead; Whose kingdom shall have no end.

And I believe in the Holy Ghost, The Lord, and Giver
of Life, Who proceedeth from the Father and the Son; Who with the Father
and the Son together is worshipped and glorified; Who spake by the Prophets:
And I believe one Catholic and Apostolic Church: I acknowledge one Baptism
for the remission of sins: And I look for the Resurrection of the dead:
And the Life of the world to come. Amen.

\emph{¶ Then shall be declared unto the People what Holy-days, or Fasting-days,
are in the week following to be observed; and (if occasion be) shall Notice
be given of the Communion, and of the Banns of Matrimony, and other matters
to be published.}

\emph{¶
Here, or immediately after the Creed, may be said the Bidding Prayer,
or other authorized Prayers and intercessions.}

\emph{¶
Then followeth the Sermon. After which, the Priest, when there is a Communion,
shall return to the Holy Table, and begin the Offertory, saying one or
more of these Sentences following, as he thinketh most convenient.}

REMEMBER
the words of the Lord Jesus, how he said, It is more blessed to give than
to receive. \emph{Acts} xx. 35.

Let your light so shine before men, that they may see
your good works, and glorify your Father which is in heaven. \emph{St.
Matt}. v. 16.

Lay not up for yourselves treasures upon earth, where
moth and rust doth corrupt, and where thieves break through and steal:
but lay up for yourselves treasures in heaven, where neither moth nor
rust doth corrupt, and where thieves do not break through nor steal. \emph{St.
Matt}. vi. 19, 20.

Not every one that saith unto me, Lord, Lord, shall
enter into the kingdom of heaven; but he that doeth the will of my Father
which is in heaven. \emph{St. Matt.} vii. 21.

He that soweth little shall reap little; and he that
soweth plenteously shall reap plenteously. Let every man do according
as he is disposed in his heart, not grudgingly, or of necessity; for God
loveth a cheerful giver. 2 \emph{Cor}. ix. 6, 7.

While we have time, let us do good unto all men; and
especially unto them that are of the household of faith. \emph{Gal}.
vi. 10.

God is not unrighteous, that he will forget your works,
and labour that proceedeth of love; which love ye have showed for his
Name’s sake, who have ministered unto the saints, and yet do minister.
\emph{Heb}. vi. 10.

To do good, and to distribute, forget not; for with
such sacrifices God is well pleased. \emph{Heb}. xiii. 16.

Whoso hath this world’s goods, and seeth his
brother have need, and shutteth up his compassion from him, how dwelleth
the love of God in him’? 1 \emph{St. John} iii. 17.

Be merciful after thy power. If thou hast much, give
plenteously; if thou hast little, do thy diligence gladly to give of that
little: for so gatherest thou thyself a good reward in the day of necessity.
\emph{Tobit}iv. 8,9.

And the King shall answer and say unto them, Verily
I say unto you, Inasmuch as ye have done it unto one of the least of these
my brethren, ye have done it unto me. \emph{St. Matt.} xxv. 40.

How then shall they call on him in whom they have not
believed’? and how shall they believe in him of whom they have not
heard’? and how shall they hear without a preacher’? and how
shall they preach, except they be sent’? \emph{Rom}. x.
14, 15.

Jesus said unto them, The harvest truly is plenteous,
but the labourers are few; pray ye therefore the Lord of the harvest,
that he would send forth labourers into his harvest. \emph{St. Luke}
x. 2.

Ye shall not appear before the LORD
empty; every man shall give as he is able, according to the blessing of
the LORD thy God which he hath given thee. \emph{Deut.}
xvi. 16, 17.

Thine, O LORD, is the greatness,
and the power, and the glory, and the victory, and the majesty: for all
that is in the heaven and in the earth is thine; thine is the kingdom,
O LORD, and thou art exalted as head above all.
1 \emph{Chron}. xxix. 11.

All things come of thee, O LORD,
and of thine own have we given thee. 1 \emph{Chron}. xxix. 14.

\emph{¶ And}NOTE,\emph{That these Sentences
may be used on any other occasion of Public Worship when the offerings
of the People are to be received.}

\emph{¶
The Deacons, Church-wardens, or other fit persons appointed for that purpose,
shall receive the Alms for the Poor, and other Offerings of the People,
in a decent Basin to be provided by the Parish; and reverently bring it
to the Priest, who shall humbly present and place it upon the Holy Table.}

\emph{¶
And the Priest shall then offer, and shall place upon the Holy Table,
the Bread and the Wine.}

\emph{¶
And when the Alms and Oblations are being received and presented, there
may be sung a Hymn, or an Offertory Anthem in the words of Holy Scripture
or of the Book of Common Prayer, under the direction of the Priest.}

\emph{¶
Here the Priest may ask the secret intercessions of the Congregation for
any who have desired the prayers of the Church.}

\emph{¶
Then shall the Priest say,}

Let us
pray for the whole state of Christ’s Church.

ALMIGHTY
and everliving God, who by thy holy Apostle hast taught us to make prayers,
and supplications, and to give thanks for all men; We humbly beseech thee
most mercifully to accept our [alms and] oblations, and to receive these
our prayers, which we offer unto thy Divine Majesty; beseeching thee to
inspire continually the Universal Church with the spirit of truth, unity,
and concord: And grant that all those who do confess thy holy Name may
agree in the truth of thy holy Word, and live in unity and godly love.

We beseech thee also, so to direct and dispose the
hearts of all Christian Rulers, that they may truly and impartially administer
justice, to the punishment of wickedness and vice, and to the maintenance
of thy true religion, and virtue.

Give grace, O heavenly Father, to all Bishops and other
Ministers, that they may, both by their life and doctrine, set forth thy
true and lively Word, and rightly and duly administer thy holy Sacraments.

And to all thy People give thy heavenly grace; and
especially to this congregation here present; that, with meek heart and
due reverence, they may hear, and receive thy holy Word; truly serving
thee in holiness and righteousness all the days of their life.

And we most humbly beseech thee, of thy goodness, O
Lord, to comfort and succour all those who, in this transitory life, are
in trouble, sorrow, need, sickness, or any other adversity.

And we also bless thy holy Name for all thy servants
departed this life in thy faith and fear; beseeching thee to grant them
continual growth in thy love and service, and to give us grace so to follow
their good examples, that with them we may be partakers of thy heavenly
kingdom. Grant this, O Father, for Jesus Christ’s sake, our only
Mediator and Advocate. \emph{Amen}.

\emph{¶
Then shall the Priest say to those who come to receive the Holy Communion,}

YE
who do truly and earnestly repent you of your sins, and are in love and
charity with your neighbours, and intend to lead a new life, following
the commandments of God, and walking from henceforth in his holy ways;
Draw near with faith, and take this holy Sacrament to your comfort; and
make your humble confession to Almighty God, devoutly kneeling.

\emph{¶
Then shall this}General Confession\emph{be made, by the Priest and
all those who are minded to receive the Holy Communion, humbly kneeling.}

ALMIGHTY
God, Father of our Lord Jesus Christ, Maker of all things, Judge of all
men; We acknowledge and bewail our manifold sins and wickedness, Which
we, from time to time, most grievously have committed, By thought, word,
and deed, Against thy Divine Majesty, Provoking most justly thy wrath
and indignation against us. We do earnestly repent, And are heartily sorry
for these our misdoings; The remembrance of them is grievous unto us;
The burden of them is intolerable. Have mercy upon us, Have mercy upon
us, most merciful Father; For thy Son our Lord Jesus Christ’s sake,
Forgive us all that is past; And grant that we may ever hereafter Serve
and please thee In newness of life, To the honour and glory of thy Name;
Through Jesus Christ our Lord. Amen.

\emph{¶
Then shall the Priest (the Bishop if he be present) stand up, and turning
to the People, say,}

ALMIGHTY
God, our heavenly Father, who of his great mercy hath promised forgiveness
of sins to all those who with hearty repentance and true faith turn unto
him; Have mercy upon you; pardon and deliver you from all your sins; confirm
and strengthen you in all goodness; and bring you to everlasting life;
through Jesus Christ our Lord. \emph{Amen}.

\emph{¶
Then shall the Priest say,}

Hear what comfortable words our Saviour Christ saith unto all who truly
turn to him.

COME
unto me, all ye that travail and are heavy laden, and I will refresh you.
\emph{St. Matt}. xi. 28.

So God loved the world, that he gave his only-begotten Son, to the end
that all that believe in him should not perish, but have everlasting life.
\emph{St. John} iii. 16.

Hear also what Saint Paul saith.

This is a true saying, and worthy of all men to be received, That Christ
Jesus came into the world to save sinners. 1 \emph{Tim}. i. 15.

Hear also what Saint John saith.

If any man sin, we have an Advocate with the Father, Jesus Christ the
righteous; and he is the Propitiation for our sins. 1 \emph{St. John}
ii. 1, 2.

\emph{¶
After which the Priest shall proceed, saying,}

\emph{Answer.}\emph{Priest. Answer.}

Lift
up your hearts.

We lift them up unto the Lord. 

Let us give thanks unto our Lord God. 

It is meet and right so to do.

\emph{¶
Then shall the Priest turn to the Holy Table, and say,}

IT
is very meet, right, and our bounden duty, that we should at all times,
and in all places, give thanks unto thee, O Lord, Holy Father, Almighty,
Everlasting God.

\emph{¶ Here shall follow the Proper Preface, according to the time,
if there be any specialty appointed; or else immediately shall be said
or sung by the Priest,}

THEREFORE
with Angels and Archangels, and with all the company of heaven, we laud
and magnify thy glorious Name; evermore praising thee, and saying,

HOLY, HOLY, HOLY, Lord God of, hosts, Heaven and earth are full of thy
glory: Glory be to thee, O Lord Most High. Amen.

\begin{quote}

\begin{quote}

PROPER
PREFACES.

CHRISTMAS.

\end{quote}

\end{quote}

\emph{Upon Christmas Day, and seven days after.}

BECAUSE
thou didst give Jesus Christ, thine only Son, to be born as at this time
for us; who, by the operation of the Holy Ghost, was made very man, of
the substance of the Virgin Mary his mother; and that without spot of
sin, to make us clean from all sin. Therefore with Angels, etc.

\begin{quote}

\begin{quote}

EPIPHANY.

\end{quote}

\end{quote}

\emph{Upon
The Epiphany, and seven days after.}

THROUGH
Jesus Christ our Lord, who, in substance of our mortal flesh, manifested
forth his glory; that he might bring us out of darkness into his own glorious
light. 

Therefore with Angels, \emph{etc}.

PURIFICATION, ANNUNCIATION,
AND TRANSFIGURATION.

\emph{Upon
the}Purification, Annunciation,\emph{and}Transfiguration.

BECAUSE
in the Mystery of the Word made flesh, thou hast caused a new light to
shine in our hearts, to give the knowledge of thy glory in the face of
thy Son, Jesus Christ our Lord.

Therefore with Angels, \emph{etc}.

EASTER.

\emph{Upon}Easter-day,\emph{and seven days after.}

BUT chiefly are we bound to praise thee
for the glorious Resurrection of thy Son Jesus Christ our Lord: for he
is the very Paschal Lamb, which was offered for us, and hath taken away
the sin of the world; who by his death hath destroyed death, and by his
rising to life again hath restored to us everlasting life. 

Therefore with Angels, \emph{etc}.

ASCENSION.

\emph{Upon}Ascension-day, \emph{and seven days after.}

THROUGH
thy most dearly beloved Son Jesus Christ our Lord; who, after his most
glorious Resurrection, manifestly appeared to all his Apostles, and in
their sight ascended up into heaven, to prepare a place for us; that where
he is, thither we might also ascend, and reign with him in glory. 

Therefore with Angels,\emph{etc}.

WHITSUNTIDE.

\emph{Upon
the Feast of}Whitsunday,\emph{and six days after.}

THROUGH
Jesus Christ our Lord; according to whose most true promise, the Holy
Ghost came down as at this time from heaven, lighting upon the Apostles,
to teach them, and to lead them into all truth; giving them boldness with
fervent zeal constantly to preach the Gospel unto all nations; whereby
we have been brought out of darkness and error into the clear light and
true knowledge of thee, and of thy Son Jesus Christ. 

Therefore with Angels, \emph{etc}.

TRINITY SUNDAY.

\emph{Upon
the Feast of}Trinity\emph{only.}

WHO,
with thine only-begotten Son, and the Holy Ghost, art one God, one Lord,
in Trinity of Persons and in Unity of Substance. For that which we believe
of thy glory, O Father, the same we believe of the Son, and of the Holy
Ghost, without any difference of inequality.

Therefore with Angels, etc.

\emph{¶
Or this.}

FOR
the precious death and merits of thy Son Jesus Christ our Lord, and for
the sending to us of the Holy Ghost, the Comforter; who are one with thee
in thy Eternal Godhead. 

Therefore with Angels, \emph{etc}.

ALL SAINTS.

\emph{Upon}All Saints’ Day,\emph{and seven days after.}

WHO,
in the multitude of thy saints, hast compassed us about with so great
a cloud of witnesses that we, rejoicing in their fellowship, may run with
patience the race that is set before us, and together with them may receive
the crown of glory that fadeth not away.

Therefore with Angels and Archangels, and with all
the company of heaven, we laud and magnify thy glorious Name; evermore
praising thee, and saying,

HOLY, HOLY, HOLY, Lord God of, hosts, Heaven and earth are full of thy
glory: Glory be to thee, O Lord Most High. Amen.

\emph{¶
When the Priest, standing before the Table, hath so ordered the Bread
and Wine, that he may with the more readiness and decency break the Bread
before the People, and take the Cup into his hands, he shall say the}Prayer
of Consecration,\emph{as followeth.}

ALL
glory be to thee, Almighty God, our heavenly Father, for that thou, of
thy tender mercy, didst give thine only Son Jesus Christ to suffer death
upon the Cross for our redemption; who made there (by his one oblation
of himself once offered) a full, perfect, and sufficient sacrifice, oblation,
and satisfaction, for the sins of the whole world; and did institute,
and in his holy Gospel command us to continue, a
perpetual memory of that his precious death and sacrifice, until
his coming again: For in the night in which he was betrayed, (\emph{a}) he took Bread; and when he had given thanks,
(\emph{b}) he brake it, and gave it to his disciples, saying, Take,
eat, (\emph{c}) this is my Body, which is given for you; Do this in
remembrance of me. Likewise, after supper, (\emph{d}) he took the Cup;
and when he had given thanks, he gave it to them, saying, Drink ye all
of this; for (\emph{e}) this is my Blood of the New Testament, which
is shed for you, and for many, for the remission of sins; Do this, as
oft as ye shall drink it, in remembrance of me.

WHEREFORE
O Lord and heavenly Father, according to the institution of thy dearly
beloved Son our Saviour Jesus Christ, we, thy humble servants, do celebrate
and make here before thy Divine Majesty, with these thy holy gifts, which
we now offer unto thee, the memorial thy Son hath commanded us to make;
having in remembrance his blessed passion and precious death, his mighty
resurrection and glorious ascension; rendering unto thee most hearty thanks
for the innumerable benefits procured unto us by the same.

AND
we most humbly beseech thee, O merciful Father, to hear us; and, of thy
almighty goodness, vouchsafe to bless and sanctify, with thy Word and
Holy Spirit, these thy gifts and creatures of bread and wine; that we,
receiving them according to thy Son our Saviour Jesus Christ’s holy
institution, in remembrance of his death and passion, may be partakers
of his most blessed Body and Blood.

AND
we earnestly desire thy fatherly goodness, mercifully to accept this our
sacrifice of praise and thanksgiving; most humbly beseeching thee to grant
that, by the merits and death of thy Son Jesus Christ, and through faith
in his blood, we, and all thy whole Church, may obtain remission of our
sins, and all other benefits of his passion. And here we offer and present
unto thee, O Lord, our selves, our souls and bodies, to be a reasonable,
holy, and living sacrifice unto thee; humbly beseeching thee, that we,
and all others who shall be partakers of this Holy Communion, may worthily
receive the most precious Body and Blood of thy Son Jesus Christ, be filled
with thy grace and heavenly benediction, and made one body with him, that
he may dwell in us, and we in him. And although we are unworthy, through
our manifold sins, to offer unto thee any sacrifice; yet we beseech thee
to accept this our bounden duty and service; not weighing our merits,
but pardoning our offences, through Jesus Christ our Lord; by whom, and
with whom, in the unity of the Holy Ghost, all honour and glory be unto
thee, O Father Almighty, world without end. \emph{Amen}.

And now, as our Saviour Christ hath taught us, we are bold to say,

OUR
Father, who art in heaven, Hallowed be thy Name. Thy kingdom come. Thy
will be done, on earth, As it is in heaven. Give us this day our daily
bread. And forgive us our trespasses, As we forgive those who trespass
against us. And lead us not into temptation, But deliver us from evil.
For thine is the kingdom, and the power, and the glory, for ever and ever.
Amen.

\emph{¶
Then shall the Priest, kneeling down at the Lord’s Table, say, in
the name of all those who shall receive the Communion, this Prayer following.}

WE
do not presume to come to this thy Table, O merciful Lord, trusting in
our own righteousness, but in thy manifold and great mercies. We are not
worthy so much as to gather up the crumbs under thy Table. But thou art
the same Lord, whose property is always to have mercy: Grant us therefore,
gracious Lord, so to eat the flesh of thy dear Son Jesus Christ, and to
drink his blood, that our sinful bodies may be made clean by his Body,
and our souls washed through his most precious Blood, and that we may
evermore dwell in him, and he in us. \emph{Amen}.
\emph{¶ Here
may be sung a Hymn.}

\emph{¶
Then shall the Priest first receive the Holy Communion in both kinds
himself, and proceed to deliver the same to the Bishops, Priests,
and Deacons, in like manner, (if any be present,) and, after
that, to the People also in order, into their hands, all devoutly
kneeling. And sufficient opportunity shall be given to those
present to communicate. And when he delivereth the Bread, he
shall say,}

THE
Body of our Lord Jesus Christ, which was given for thee, preserve thy
body and soul unto everlasting life. Take and eat this in remembrance
that Christ died for thee, and feed on him in thy heart by faith, with
thanksgiving.

\emph{¶
And the Minister who delivereth the Cup shall say,}

THE
Blood of our Lord Jesus Christ, which was shed for thee, preserve thy
body and soul unto everlasting life. Drink this in remembrance that Christ’s
Blood was shed for thee, and be thankful.

\emph{¶
If the consecrated Bread or Wine be spent before all have communicated,
the Priest is to consecrate more, according to the Form before prescribed;
beginning at,}All glory be to thee, Almighty God, \emph{and ending
with these words,} partakers of his most blessed Body and Blood.

\emph{¶
When all have communicated, the Priest shall return to the Lord’s
Table, and reverently place upon it what remaineth of the consecrated
Elements, covering the same with a fair linen cloth.}

\emph{¶
Then shall the Priest say,}

\begin{quote}

\begin{quote}

Let
us pray.

\end{quote}

\end{quote}

ALMIGHTY
and everliving God, we most heartily thank thee, for that thou dost vouchsafe
to feed us who have duly received these holy mysteries with the spiritual
food of the most precious Body and Blood of thy Son our Saviour Jesus
Christ; and dost assure us thereby of thy favour and goodness towards
us; and that we are very members incorporate in the mystical body of thy
Son, which is the blessed company of all faithful people; and are also
heirs through hope of thy everlasting kingdom, by the merits of his most
precious death and passion. And we humbly beseech thee, O heavenly
Father, so to assist us with thy grace, that we may continue in that holy
fellowship, and do all such good works as thou hast prepared for us to
walk in; through Jesus Christ our Lord, to whom, with thee and the Holy
Ghost, be all honour and glory, world without end. \emph{Amen}.

\emph{¶
Then shall be said Gloria in excelsis, all standing, or some proper Hymn.}

GLORY
be to God on high, and on earth peace, good will towards men. We praise
thee, we bless thee, we worship thee, we glorify thee, we give thanks
to thee for thy great glory, O Lord God, heavenly King, God the Father
Almighty.

O Lord, the only-begotten Son, Jesus Christ; O Lord
God, Lamb of God, Son of the Father, that takest away the sins of the
world, have mercy upon us. Thou that takest away the sins of the world,
receive our prayer. Thou that sittest at the right hand of God the Father,
have mercy upon us.

For thou only art holy; thou only art the Lord; thou
only, O Christ, with the Holy Ghost, art most high in the glory of God
the Father. Amen.

\emph{¶
Then, the people kneeling, the Priest (the Bishop if he be present) shall
let them depart with this Blessing.}

THE
Peace of God, which passeth all understanding, keep your hearts and minds
in the knowledge and love of God, and of his Son Jesus Christ our Lord:
And the Blessing of God Almighty, the Father, the Son, and the Holy Ghost,
be amongst you, and remain with you always. \emph{Amen}.

\begin{quote}

\begin{quote}

GENERAL RUBRICS

\end{quote}

\end{quote}

\emph{¶ In the absence of a Priest, a Deacon may say all that is before
appointed unto the end of the Gospel.}

\emph{¶
Upon the Sundays and other Holy-days, (though there be no Sermon or Communion,)
may be said all that is appointed at the Communion, unto the end of the
Gospel, concluding with the Blessing.}

\emph{¶
And if any of the consecrated Bread and Wine remain after the Communion,
it shall not be carried out of the Church; but the Minister and other
Communicants shall, immediately after the Blessing, reverently eat and
drink the same.}

\emph{¶
If among those who come to be partakers of the Holy Communion, the Minister
shall know any to be an open and notorious evil liver, or to have done
any wrong to his neighbours by word or deed, so that the Congregation
be thereby offended; he shall advertise him, that he presume not to come
to the Lord’s Table, until he have openly declared himself to have
truly repented and amended his former evil life, that the Congregation
may thereby be satisfied; and that he hath recompensed the parties to
whom he hath done wrong; or at least declare himself to be in full purpose
so to do, as soon as he conveniently may.}

\emph{¶
The same order shall the Minister use with those, betwixt whom he perceiveth
malice and hatred to reign; not suffering them to be partakers of the
Lord’s Table, until he know them to be reconciled. And if one of
the parties, so at variance, be content to forgive from the bottom of
his heart all that the other hath trespassed against him, and to make
amends for that wherein he himself hath offended; and the other party
will not be persuaded to a godly unity, but remain still in his forwardness
and malice; the Minister in that case ought to admit the penitent person
to the Holy Communion, and not him that is obstinate.}Provided,\emph{That every Minister so repelling any, as is herein specified, shall be
obliged to give an account of the same to the Ordinary, within fourteen
days after, at the farthest.}

\begin{quote}

\begin{quote}

EXHORTATIONS

\end{quote}

\end{quote}

\emph{¶
At the time of the celebration of the Communion, the Priest may say this
Exhortation. And Note, that the Exhortation shall be said on the First
Sunday in Advent, the First Sunday in Lent, and Trinity Sunday.}

DEARLY
beloved in the Lord, ye who mind to come to the holy Communion of the
Body and Blood of our Saviour Christ, must consider how Saint Paul exhorteth
all persons diligently to try and examine themselves, before they presume
to eat of that Bread, and drink of that Cup. For as the benefit is great,
if with a true penitent heart and lively faith we receive that holy Sacrament;
so is the danger great, if we receive the same unworthily. Judge therefore
yourselves, brethren, that ye be not judged of the Lord; repent you truly
for your sins past; have a lively and steadfast faith in Christ our Saviour;
amend your lives, and be in perfect charity with all men; so shall ye
be meet partakers of those holy mysteries. And above all things ye must
give most humble and hearty thanks to God, the Father, the Son, and the
Holy Ghost, for the redemption of the world by the death and passion of
our Saviour Christ, both God and man; who did humble himself, even to
the death upon the Cross, for us, miserable sinners, who lay in darkness
and the shadow of death; that he might make us the children of God, and
exalt us to everlasting life. And to the end that we should always remember
the exceeding great love of our Master, and only Saviour, Jesus Christ,
thus dying for us, and the innumerable benefits which by his precious
blood-shedding he hath obtained for us; he hath instituted and ordained
holy mysteries, as pledges of his love, and for a continual remembrance
of his death, to our great and endless comfort. To him therefore, with
the Father and the Holy Ghost, let us give (as we are most bounden) continual
thanks; submitting ourselves wholly to his holy will and pleasure, and
studying to serve him in true holiness and righteousness all the days
of our life. Amen.

\emph{¶
When the Minister giveth warning for the Celebration of the Holy Communion,
(which he shall always do upon the Sunday, or some Holy-day, immediately
preceding,) he shall read this Exhortation following, or so much thereof
as, in his discretion, he may think convenient.}

DEARLY
beloved, on——day next I purpose, through God's assistance,
to administer to all such as shall be religiously and devoutly disposed
the most comfortable Sacrament of the Body and Blood of Christ; to be
by them received in remembrance of his meritorious Cross and Pas-sion;
whereby alone we obtain remission of our sins, and are made partakers
of the Kingdom of heaven. Wherefore it is our duty to render most humble
and hearty thanks to Almighty God, our heavenly Father, for that he hath
given his Son our Saviour Jesus Christ, not only to die for us, but also
to be our spiritual food and sustenance in that holy Sacrament. Which
being so divine and comfortable a thing to them who receive it worthily,
and so dangerous to those who will presume to receive it unworthily; my
duty is to exhort you, in the mean season to consider the dignity of that
holy mystery, and the great peril of the unworthy receiving thereof; and
so to search and examine your own consciences, and that not lightly, and
after the manner of dissemblers with God; but so that ye may come holy
and clean to such a heavenly Feast, in the marriage-garment required by
God in holy Scripture, and be received as worthy partakers of that holy
Table.

The way and means thereto is: First, to examine your
lives and conversations by the rule of God’s command-ments; and
whereinsoever ye shall perceive yourselves to have offended, either by
will, word, or deed, there to bewail your own sinfulness, and to confess
yourselves to Almighty God, with full purpose of amendment of life. And
if ye shall perceive your offences to be such as are not only against
God, but also against your neighbours; then ye shall reconcile yourselves
unto them; being ready to make restitution and satisfaction, according
to the uttermost of your powers, for all injuries and wrongs done by you
to any other; and being likewise ready to forgive others who have offended
you, as ye would have forgiveness of your offences at God's hand: for
otherwise the receiving of the holy Communion doth nothing else but increase
your condemnation. There-fore, if any of you be a blasphemer of God, an
hinderer or slanderer of his Word, an adulterer, or be in malice, or envy,
or in any other grievous crime; repent you of your sins, or else come
not to that holy Table.

And because it is requisite that no man should come
tothe holy Communion, but with a full trust in God’s mercy, and
with a quiet conscience; therefore, if there be any of you, who by this
means cannot quiet his own conscience herein, but requireth further comfort
or counsel, let him come to me, or to some other Minister of God’s
Word, and open his grief; that he may receive such godly counsel and advice,
as may tend to the quieting of his conscience, and the removing of all
scruple and doubtfulness. 

\emph{¶
Or, in case he shall see the People negligent to come to the Holy Communion,
instead of the former, he may use this Exhortation.}

DEARLY
beloved brethren, on——I intend, by God’s grace, to celebrate
the Lord's Supper: unto which, in God's behalf, I bid you all who are
here present; and beseech you, for the Lord Jesus Christ's sake, that
ye will not refuse to come thereto, being so lovingly called and bidden
by God himself. Ye know how grievous and unkind a thing it is, when a
man hath prepared a rich feast, decked his table with all kind of provision,
so that there lacketh nothing but the guests to sit down; and yet they
who are called, without any cause, most unthankfully refuse to come. Which
of you in such a case would not be moved? Who would not think a great
injury and wrong done unto him? Wherefore, most dearly beloved in Christ,
take ye good heed, lest ye, withdrawing yourselves from this holy Supper,
provoke God's indignation against you. It is an easy matter for a man
to say, I will not communicate, because I am otherwise hindered with worldly
business. But such excuses are not so easily accepted and allowed before
God. If any man say, I am a grievous sinner, and therefore am afraid to
come: wherefore then do ye not repent and amend? When God calleth you,
are ye not ashamed to say ye will not come? When ye should return to God,
will ye excuse yourselves, and say ye are not ready? Considerearnestly
with yourselves how little such feigned excuses will avail before God.
Those who refused the feast in the Gospel, because they had bought a farm,
or would try their yokes of oxen, or because they were married, were not
so excused, but counted unworthy of the heavenly feast. Wherefore, according
to mine office, I bid you in the Name of God, I call you in Christ's behalf,
I exhort you, as ye love your own salvation, that ye will be partakers
of this holy Communion. And as the Son of God did vouchsafe to yield up
his soul by death upon the Cross for your salvation; so it is your duty
to receive the Communion in remembrance of the sacrifice of his death,
as he himself hath commanded: which if ye shall neglect to do, consider
with yourselves how great is your ingratitude to God, and how sore punishment
hangeth over your heads for the same; when ye wilfully abstain from the
Lord's Table, and separate from your brethren, who come to feed on the
banquet of that most heavenly food. These things if ye earnestly consider,
ye will by God’s grace return to a better mind: for the obtaining
whereof we shall not cease to make our humble petitions unto Almighty
God, our heavenly Father.

\begin{quote}

\emph{Return to the U.
S. Book of Common Prayer (1928)}

\end{quote}

\xchapter{The Ministration of Holy Baptism}

\textbf{The Book of Common Prayer (1928)}

The Ministration
of

Holy
Baptism

\emph{¶
The Minister of every Parish shall often admonish the People, that they
defer not the Baptism of their Children, and that it is most convenient
that Baptism should be administered upon Sundays and other Holy Days.
Nevertheless, if necessity so require, Baptism may be administered upon
any other day. And also he shall warn them that, except for urgent cause,
they seek not to have their Children baptized in their houses.}

\emph{¶
There shall be for every Male-child to be baptized. when they can be had,
two Godfathers and one Godmother; and for every Female, one Godfather
and two Godmothers; and Parents shall be admitted as Sponsors, if it be
desired.}

\emph{¶
When there are Children to be baptized, the Parents or Sponsors shall
give knowledge thereof to the Minister. And then the Godfathers and
Godmothers, and the People with the Children, must be ready at the Font,
either immediately after the Second Lesson at Morning or Evening Prayer,
or at such other time as the Minister shall appoint.}

\emph{¶
When any such Persons as are of riper years are to be baptized, timely
notice shall be given to the Minister, that so due care may be taken
for their examination, whether they be sufficiently instructed in the
Principles of the Christian Religion; and that they may be exhorted
to prepare themselves, with Prayers and Fasting, for the receiving of
this holy Sacrament. And}NOTE\emph{, That
at the time of the Baptism of an Adult, there shall be present with
him at the Font at least two Witnesses.}

\emph{}
\emph{¶
The Minister, having come to the Font, which is then to be filled with pure
Water, shall say as followeth, the People all standing,} HATH
this Child (Person) been already baptized, or no? 

\emph{¶ If they answer,}No:\emph{then shall
the Minister proceed as followeth.}

DEARLY
beloved, forasmuch as our Saviour Christ saith, None can enter into the
Kingdom of God, except he be regenerate and born anew of Water and of
the Holy Ghost; I beseech you to call upon God the Father, through our
Lord Jesus Christ, that of his bounteous mercy he will grant to \emph{this
Child} (or \emph{Person}) that which by nature \emph{he} cannot
have; that \emph{he} may be baptized with Water and the Holy Ghost,
and received into Christ’s holy Church, and be made \emph{a} living
\emph{member} of the same.

\begin{quote}

\begin{quote}

\emph{¶ Then shall the Minister say,}

Let
us pray.

\end{quote}

\end{quote}

ALMIGHTY
and immortal God, the aid of all who need, the helper of all who flee
to thee for succour, the life of those who believe, and the resurrection
of the dead; We call upon thee for \emph{this Child} (or \emph{this}
thy \emph{Servant}), that \emph{he}, coming to thy holy Baptism, may
receive remission of sin, by spiritual regeneration. Receive \emph{him},
O Lord, as thou hast promised by thy well-beloved Son, saying, Ask, arid
ye shall have; seek, and ye shall find; knock, and it shall be opened
unto you. So give now unto us who ask; let us ‘who seek, find; open
the gate unto us who knock; that \emph{this Child} (or \emph{this}
thy \emph{Servant}) may enjoy the everlasting benediction of thy heavenly
washing, and may come to the eternal kingdom which thou hast promised
by Christ our Lord. \emph{Amen}.

\begin{quote}

\begin{quote}

\emph{¶
Then the Minister shall say as followeth:}

\end{quote}

\end{quote}

Hear the words of the Gospel, written by St. Mark, in the tenth Chapter,
at the thirteenth Verse.

THEY
brought young children to Christ, that he should touch them: and his disciples
rebuked those that brought them. But when Jesus saw it, he was much displeased,
and said unto them, Suffer the little children to come unto me, and forbid
them not: for of such is the kingdom of God. Verily I say unto you, Whosoever
shall not receive the kingdom of God as a little child, he shall not enter
therein. And he took them up in his arms, putS his hands upon them, and
blessed them.

\begin{quote}

\begin{quote}

\emph{¶
Or this.}

\end{quote}

\end{quote}

Hear the words of the Gospel, written by St. John, in the third Chapter,
at the first Verse.

THERE
was a man of the Pharisees, named Nicodemus, a ruler of the Jews: the
same came to Jesus by night, and said unto him, Rabbi, we know that thou
art a teacher come from God: for no man can do these miracles that thou
doest, except God be with him. Jesus answered and said unto him, Verily,
verily, I say unto thee, Except a man be born again, he cannot see the
kingdom of God. Nicodemus saith unto him, How can a man be born when he
is old? can he enter the second time into his mother’s womb, and
be born? Jesus answered, Verily, verily, I say unto thee, Except a man
be born of water and of the Spirit, he cannot enter into the kingdom of
God. That which is born of the flesh is flesh; and that which is born
of the Spirit is spirit. Marvel not that I said unto thee, Ye must be
born again. The wind bloweth where it listeth, and thou hearest the sound
thereof, but. canst not tell whence it cometh, and whither it goeth: so
is every one that is born of the Spirit.

\begin{quote}

\begin{quote}

\emph{¶ Or this.}

\end{quote}

\end{quote}

Hear the words of the Gospel, written by St. Matthew, in the twenty-eighth
Chapter, at the eighteenth Verse.

JESUS
came and spake unto them, saying, All power is given unto me in heaven
and in earth. Go ye therefore, and teach all nations, baptizing them in
the name of the Father, and of the Son, and of the Holy Ghost: teaching
them to observe all things whatsoever I have commanded you: and, lo, I
am with you alway, even unto the end of the world.

\begin{quote}

\begin{quote}

\emph{¶
Then shall the Minister say:}

\end{quote}

\end{quote}

AND
now, we being persuaded of the good will of our heavenly Father toward
this \emph{Child} (\emph{this Person}) declared by his Son Jesus Christ;
let us faithfully and devoutly give thanks unto him and say,

\emph{¶
Minister and People}

ALMIGHTY
and everlasting God, heavenly Father, We give thee humble thanks, That
thou hast vouchsafed to call us to the knowledge of thy grace, and faith
in thee: Increase this knowledge, And confirm this faith, in us evermore.
Give thy Holy Spirit to \emph{this Child} (or \emph{this} thy \emph{Servant}),
That \emph{he} may be born again, And be made \emph{an heir} of everlasting
salvation; Through our Lord Jesus Christ, Who liveth and reigrieth with
thee and the Holy Spirit, Now and forever. \emph{Amen}.

\emph{¶
When the Office is used for Children, the Minister shall speak unto
the Godfathers and Godmothers on this wise.}

DEARLY
beloved, ye have brought this Child here to be baptized; ye have prayed
that our Lord Jesus Christ would vouchsafe to receive him, to release
him from sin, to sanctify him with the Holy Ghost, to give him the kingdom
of heaven, and everlasting life.

Dost thou, therefore, in the name of this Child, renounce the devil and
all his works, the vain pomp and glory of the world, with all covetous
desires of the same, and the sinful desires of the flesh, so that thou
wilt not follow, nor be led by them?

\emph{Answer}. I renounce them all; and, by God’s help, will endeavor
not to follow, nor be led by them.

\emph{Minister}. Dost thou believe all the Articles of the Christian Faith,
as contained in the Apostles’ Creed?

\emph{Answer}. I do.

\emph{Minister}. Wilt thou be baptized in this Faith?

\emph{Answer}. That is my desire.

\emph{Minister}. Wilt thou then obediently keep God’s
holy will and commandments, and walk in the same all the days of thy life?

\emph{Answer}. I will, by God’s help.

\emph{Minister}. Having now, in the name of \emph{this Child}, made these
promises, will ye also on your part take heed that \emph{he} learn the
Creed, the Lord’s Prayer, and the Ten Commandments, and all other
things which a Christian ought to know and believe, to his soul’s
health?

\emph{Answer}. I will, by God’s help.

\emph{Minister}. Will ye take heed that this Child,
so soon as sufficiently
instructed, be brought to the Bishop to be confirmed by him?

\emph{Answer.} I will, God being my helper.

\emph{¶
When the Office is used for Adults, the Minister shall address them
on this wise, the Persons to be baptized answering the questions for
themselves.}

WELL-BELOVED,
you have come hither desiring to receive holy Baptism. We have prayed
that our Lord Jesus Christ would vouchsafe to receive you, to release
you from sin, to sanctify you with the Holy Ghost, to give you the kingdom
of heaven, and everlasting life.

DOST
thou renounce the devil and all his works, the vain pomp and glory of
the world, with all covetous desires of the same, and the sinful desires
of the flesh, so that thou wilt not follow, nor be led by them? 

\emph{Answer}. I renounce them all; and, by God's
help, will endeavour not to follow, nor be led by them.

\emph{Minister}.
Dost thou believe in Jesus the Christ, the Son of the Living God?

\emph{Answer}. I do.

\emph{Minister}. Dost thou accept him, and desire
to follow him as thy Saviour and Lord?

\emph{Answer}. I do.

\emph{Minister}. Dost thou believe all the Articles
of the Christian Faith, as contained in the Apostles' Creed?

\emph{Answer}. I do.

\emph{Minister}. Wilt thou be baptized in this Faith?

\emph{Answer}. That is my desire.

\emph{Minister}. Wilt thou then obediently keep God's
holy will and commandments, and walk in the same all the days of thy life?

\emph{Answer}. I will, by God's help. 

\begin{quote}

\begin{quote}

\emph{¶
Then shall the Minister say,}

\end{quote}

\end{quote}

O
MERCIFUL God, grant that like as Christ died and rose again, so \emph{this
Child} (\emph{this}thy \emph{Servant}) may die to sin and rise
to newness of life. \emph{Amen}.

Grant that all sinful affections may die in \emph{him}, and that all
things belonging to the Spirit may live and grow in \emph{him. Amen.}

Grant that \emph{he} may have power and strength
to have victory, and to triumph, against the devil, the world, and the
flesh. \emph{Amen}.

Grant that whosoever is here dedicated to thee by our
office and ministry, may also be endued with heavenly virtues, and everlastingly
rewarded, through thy mercy, O blessed Lord God, who dost live, and govern
all things, world without end. \emph{Amen}.

\emph{Minister.} The
Lord be with you.

\emph{Answer.} And with thy spirit.

\emph{Minister.} Lift up your hearts.

\emph{Answer.} We
lift them up unto the Lord.

\emph{Minister.} Let
us give thanks unto our Lord God.

\emph{Answer.} It
is meet and right so to do.

\emph{¶
Then the Minister shall say,}

IT
is very meet, right, and our bounden duty, that we should give thanks
unto thee, O Lord, Holy Father, Al-mighty, Everlasting God, for that
thy dearly beloved Son Jesus Christ, for the forgiveness of our sins,
did shed out of his most precious side both water and blood; and gave
commandment to his disciples, that they should go teach all nations,
and baptize them In the Name of the Father, and of the Son, and of the
Holy Ghost. Regard, we beseech thee, the supplications of thy congregation;
sanctify this Water to the mystical washing away of sin; and grant that
\emph{this Child} (\emph{this} thy \emph{Servant}), now to be
baptized therein, may receive the fulness of thy grace, and ever remain
in the number of thy faithful children; through the same Jesus Christ
our Lord, to whom, with thee, in the unity of the Holy Spirit, be all
honour and glory, now and evermore. \emph{Amen}.

\emph{¶
Then the Minister shall take the Child into his hands, and shall say to
the Godfathers and Godmothers,}

Name this
Child.

\emph{¶
And then, naming the child after them, he shall dip}him\emph{in the
Water discreetly, or shall pour Water upon}him\emph{, saying,}

\emph{N}.
I baptize thee In the Name of the Father, and of the Son, and of the Holy
Ghost. Amen.

\emph{¶
But}NOTE,\emph{That if the Person to be baptized
be an Adult, the Minister shall take him by the hand, and shall ask
the Witnesses the Name; and then shall dip him in the Water, or pour
Water upon him, using the same form of words.}

\begin{quote}

\begin{quote}

\emph{¶
Then shall the Minister say,}

\end{quote}

\end{quote}

WE
receive this Child
(\emph{or} person) into the congregation of Christ’s flock; and
do *sign \emph{him} with the sign of the Cross, in token that hereafter
\emph{he} shall not be ashamed to confess the faith of Christ crucified,
and manfully to fight under his banner, against sin, the world, and the
devil; and to continue Christ’s faithful soldier and servant unto
\emph{his} life’s end. Amen.

\begin{quote}

\begin{quote}

\emph{¶ Then shall the Minister say,}

\end{quote}

\end{quote}

SEEING
now, dearly beloved brethren, that \emph{this Child} (or \emph{this Person})
is regenerate, and grafted into the body of Christ’s Church, let
us give thanks unto Almighty God for these benefits; and with one accord
make our prayers unto him, that \emph{this Child} (or \emph{this Person})
may lead the rest of \emph{his} life according to this beginning.

\begin{quote}

\begin{quote}

\emph{¶
Then shall be said,}

\end{quote}

\end{quote}

OUR
Father, who art in heaven, hallowed be thy Name. Thy kingdom come. Thy
will be done on earth, As it is in heaven. Give us this day our daily
bread. And forgive us our trespasses, As we forgive those who trespass
against us. And lead us not into temptation; But deliver us from evil.
For thine is the kingdom, and the power, and the glory, for ever and ever.
Amen.

\begin{quote}

\begin{quote}

\emph{¶
Then shall the Minister say,}

\end{quote}

\end{quote}

WE
yield thee hearty thanks, most merciful Father, that it hath pleased thee
to regenerate \emph{this Child} (or \emph{this} thy \emph{Servant})
with thy Holy Spirit, to receive \emph{him} for thine own \emph{Child},
and to incorporate \emph{him} into thy holy Church. And humbly we beseech
thee to grant, that \emph{he}, being dead unto sin, may live unto righteousness,
and being buried with Christ in his death, may also be \emph{partaker}
of his resurrection; so that finally, with the residue of thy holy Church,
\emph{he} may be an inheritor of thine everlasting kingdom; through
Christ our Lord. \emph{Amen}.

\emph{¶
Then the Minister shall add,}

THE
Almighty God, the Father of our Lord Jesus Christ, of whom the whole family
in heaven and earth is named; Grant you to be strengthened with might
by his Spirit in the inner man; that Christ dwelling in your hearts by
faith, ye may be filled with all the fulness of God. \emph{Amen}.

\emph{¶
It is expedient that every Adult, thus baptized, should be confirmed
by the Bishop, so soon after his Baptism as conveniently may be; that
so he may be admitted to the Holy Communion.}

PRIVATE
BAPTISM

\emph{¶
When, in consideration of extreme sickness, necessity may require, then
the following form shall suffice:}

\emph{¶The Child (or Person) being named by some one who is present,
the Minister shall pour Water upon him, saying these words:}

\emph{N}.
I baptize thee In the Name of the Father, and of the Son, and of the Holy
Ghost. Amen.

\emph{¶ After which shall be said the}Lord’s Prayer,\emph{and the}Thanksgiving\emph{from the Office, beginning,}We yield
thee hearty thanks,\emph{etc.}

\emph{¶
But}NOTE,\emph{That in the case of an Adult,
the Minister shall first ask the questions provided in this Office for
the Baptism of Adults.}

\emph{¶ In cases of extreme sickness, or any imminent peril, if a Minister
cannot be procured, then any baptized person present may administer
holy Baptism, using the foregoing form. Such Baptism shall be promptly
reported to the Parish authorities.}

\emph{}

THE
RECEIVING OF ONE PRIVATELY
BAPTIZED.

\emph{¶
It is expedient that a Child or Person so baptized be afterward brought
to the Church, at which time these parts of the foregoing service shall
be used: The Gospel, the Questions (omitting the question}Wilt
thou be baptized in this Faith? \emph{and the answer thereto), the Declaration,
We receive this Child}(\emph{or}Person),\emph{etc., and the remainder
of the Office.}

CONDITIONAL
BAPTISM.

\emph{¶
If there be reasonable doubt whether any Person were baptized with Water,}In the Name of the Father, and of the Son, and of the Holy Ghost\emph{(which are essential parts of Baptism), such Person may be baptized in
the manner herein appointed; saving that, at the immersion or the pouring
of water, the Minister shall use this Form of words.}

IF
thou art not already baptized, \emph{N}., I baptize thee In the Name
of the Father, and of the Son, and of the Holy Ghost. Amen.

\begin{quote}

\emph{Return to the U.
S. Book of Common Prayer (1928)}

\end{quote}

\xchapter{The Offices of Instruction, with the Catechism}

\textbf{The Book of Common Prayer (1928)}

A
Catechism

that
is to say, an Instruction,

to
be Learned by Every Person before he

be
brought to be Confirmed

by
the Bishop.

\emph{Q}\emph{UESTION}.
What is your Name?

\emph{Answer. N. or N. N.}

\emph{Question}. Who gave you this Name?

\emph{Answer}. My Sponsors in Baptism; wherein I
was made a member of Christ, the child of God, and an inheritor of the
kingdom of heaven.

\emph{Question}. What did your Sponsors then for
you?

\emph{Answer}. They did promise and vow three things
in my name: First, that I should renounce the devil and all his works,
the pomps and vanity of this wicked world, and all the sinful lusts of
the flesh; Secondly, that I should believe all the Articles of the Christian
Faith; And Thirdly, that I should keep God's holy will and commandments,
and walk in the same all the days of my life.

\emph{Question}. Dost thou not think that thou art
bound to believe, and to do, as they have promised for thee?

\emph{Answer}. Yes, verily; and by God's help so
I will. And I heartily thank our heavenly Father, that he hath called
me to this state of salvation, through Jesus Christ our Saviour. And I
pray unto God to give me his grace, that I may continue in the same unto
my life's end.

\emph{Catechist}. Rehearse the Articles of thy Belief.

\emph{Answer}. I believe in God the Father Almighty,
Maker of heaven and earth:

And in Jesus Christ his only Son our Lord: Who was
conceived by the Holy Ghost, Born of the Virgin Mary: Suffered under Pontius
Pilate, Was crucified, dead, and buried: He descended into hell; The third
day he rose again from the dead: He ascended into heaven, And sitteth
on the right hand of God the Father Almighty: From thence he shall come
to judge the quick and the dead.

I believe in the Holy Ghost: The holy Catholic Church;
The Communion of Saints: The Forgiveness of sins: The Resurrection of
the body: And the Life everlasting. Amen.

\emph{Question}. What dost thou chiefly learn in
these Articles of thy Belief?

\emph{Answer}. First, I learn to believe in God the
Father, who hath made me, and all the world.

Secondly, in God the Son, who hath redeemed me, and
all mankind.

Thirdly, in God the Holy Ghost, who sanctifieth me,
and all the people of God.

\emph{Question}. You said that your Sponsors did
promise for you, that you should keep God's Commandments. Tell me how
many there are?

\emph{Answer}. Ten.

\emph{Question}. Which are they?

\emph{Answer}. The same which God spake in the
twentieth Chapter of Exodus, saying, I am the Lord thy God, who brought
thee out of the land of Egypt, out of the house of bondage.

I. Thou shalt have none other gods but me.

II. Thou shalt not make to thyself any graven image,
nor the likeness of any thing that is in heaven above, or in the earth
beneath, or in the water under the earth; thou shalt not bow down to them,
nor worship them; for I the LORD thy God am a jealous
God, and visit the sins of the fathers upon the children, unto the third
and fourth generation of them that hate me; and show mercy unto thousands
in them that love me and keep my commandments.

III. Thou shalt not take the Name of the Lord thy God
in vain; for the LORD will not hold him guiltless,
that taketh his Name in vain; for the LORD will
not hold him guiltless, that taketh his name in vain.

IV. Remember that thou keep holy the Sabbath-day. Six
days shalt thou labour, and do all that thou hast to do; but the seventh
day is the Sabbath of the Lord thy God. In it thou shalt do no manner
of work; thou, and thy son, and thy daughter, thy man-servant, and thy
maid-servant, thy cattle, and the stranger that is within thy gates. For
in six days the Lord made heaven and earth, the sea, and all that in them
is, and rested the seventh day: wherefore the Lord blessed the seventh
day, and hallowed it.

V. Honour thy father and thy mother; that thy days may
be long in the land which the LORD thy God giveth
thee.

VI. Thou shalt do no murder.

VII. Thou shalt not commit adultery.

VIII. Thou shalt not steal.

IX. Thou shalt not bear false witness against thy neighbour.

X. Thou shalt not covet thy neighbour's house, thou
shalt not covet thy neighbour's wife, nor his servant, nor his maid, nor
his ox, nor his ass, nor any thing that is his.

\emph{Question}. What dost thou chiefly learn
by these Commandments?

\emph{Answer}. I learn two things; my duty towards
God, and my duty towards my Neighbour.

\emph{Question}. What is thy duty towards God?

\emph{Answer.} My duty towards God is To believe
in him, to fear him, And to love him with all my heart, with all my mind,
with all my soul, and with all my strength: To worship him, to give him
thanks: To put my whole trust in him, to call upon him: To honour his
holy Name and his Word: And to serve him truly all the days of my life.

\emph{Question}. What is thy duty towards thy
Neighbour?

\emph{Answer}. My duty towards my Neighbour
is To love him as myself, and to do to all men as I would they should
do unto me: To love, honour, and succour my father and mother: To honour
and obey the civil authority: To submit myself to all my governors, teachers,
spiritual pastors and masters: To order myself lowly and reverently to
all my betters: To hurt nobody by word or deed: To be true and just in
all my dealings: To bear no malice nor hatred in my heart: To keep my
hands from picking and stealing, and my tongue from evil speaking, lying,
and slandering: To keep my body in temperance, soberness, and chastity:
Not to covet nor desire other men's goods; But to learn and labour truly
to get mine own living, And to do my duty in that state of life unto which
it shall please God to call me.

\emph{Catechist}. My good Child, know this; that
thou art not able to do these things of thyself, nor to walk in the Com-
mandments of God, and to serve him, without his special grace; which thou
must learn at all times to call for by diligent prayer. Let me hear, therefore,
if thou canst say the Lord’s Prayer.

\emph{Answer}. Our Father, who art in heaven,
Hallowed be thy Name. Thy kingdom come. Thy will be done, On earth as
it is in heaven. Give us this day our daily bread. And forgive us our
trespasses, As we forgive those who trespass against us. And lead us not
into temptation, But deliver us from evil. Amen.

\emph{Question}. What desirest thou of God in
this Prayer?

\emph{Answer}. I desire my Lord God, our heavenly
Father, who is the giver of all goodness, to send his grace unto me, and
to all people; that we may worship him, serve him, and obey him, as we
ought to do. And I pray unto God, that he will send us all things that
are needful both for our souls and bodies; and that he will be merciful
unto us, and forgive us our sins; and that it will please him to save
and defend us in all dangers both of soul and body; and that he will keep
us from all sin and wickedness, and from our spiritual enemy, and from
everlasting death. And this I trust he will do of his mercy and goodness,
through our Lord Jesus Christ. And therefore I say, Amen, So be it.

\emph{Question}. How many Sacraments hath Christ
ordained in his Church?

\emph{Answer}. Two only, as generally necessary
to salvation; that is to say, Baptism, and the Supper of the Lord.

\emph{Question}. What meanest thou by this word
Sacrament?

\emph{Answer}. I mean an outward and visible
sign of an inward and spiritual grace given unto us; ordained by Christ
himself, as a means whereby we receive the same, and a pledge to assure
us thereof.

\emph{Question}. How many parts are there in
a Sacrament?

\emph{Answer}. Two; the outward visible sign,
and the inward spiritual grace.

\emph{Question}. What is the outward visible
sign or form in Baptism?

\emph{Answer}. Water; wherein the person is
baptized, In the Name of the Father, and of the Son, and of the Holy Ghost.

\emph{Question}. What is the inward and spiritual
grace?

\emph{Answer}. A death unto sin, and a new birth
unto righteousness: for being by nature born in sin, and the children
of wrath, we are hereby made the children of grace.

\emph{Question}. What is required of persons
to be baptized?

\emph{Answer}. Repentance, whereby they forsake
sin; and Faith, whereby they stedfastly believe the promises of God made
to them in that Sacrament.

\emph{Question}. Why then are Infants baptized,
when by reason of their tender age they cannot perform them?

\emph{Answer}. Because they promise them both
by their Sureties; which promise, when they come to age, themselves are
bound to perform.

\emph{Question}. Why was the Sacrament of the
Lord’s Supper ordained?

\emph{Answer}. For the continual remembrance
of the sacrifice of the death of Christ, and of the benefits which we
receive thereby.

\emph{Question}. What is the outward part or
sign of the Lord’s Supper?

\emph{Answer}. Bread and Wine, which the Lord
hath commanded to be received.

\emph{Question}. What is the inward part, or
thing signified?

\emph{Answer}. The Body and Blood of Christ,
which are spiritually taken and received by the faithful in the Lord's
Supper.

\emph{Question}. What are the benefits whereof
we are partakers thereby?

\emph{Answer.} The strengthening and refreshing
of our souls by the Body and Blood of Christ, as our bodies are by the
Bread and Wine.

\emph{Question}. What is required of those who
come to the Lord's Supper?

\emph{Answer}. To examine themselves, whether
they repent them truly of their former sins, stedfastly purposing to lead
a new life; have a lively faith in God's mercy through Christ, with a
thankful remembrance of his death; and be in charity with all men. 

¶
\emph{The Minister of every Parish shall diligently, upon Sundays and
Holy Days, or on some other convenient occasions, openly in the Church,
instruct or examine so many Children of his Parish, sent unto him, as
he shall think convenient, in some part of this Catechism.}

¶ \emph{And all Fathers, Mothers, Masters, and Mistresses, shall
cause their Children, Servants, and Apprentices, who have not learned
their Catechism, to come to the Church at the time appointed, and obediently
to hear and to be ordered by the Minister, until such time as they have
learned all that is here appointed for them to learn.}

¶ \emph{So soon as Children are come to a competent age, and can
say the Creed, the Lord's Prayer, and the Ten Commandments, and can
answer to the other questions of this short Catechism, they shall be
brought to the Bishop.}

¶\emph{And whensoever the Bishop shall give knowledge for Children
to be brought unto him for their Confirmation, the Minister of every
Parish shall either bring, or send in writing, with his hand subscribed
thereunto, the Names of all such Persons within his Parish, as he shall
think it to be presented to the Bishop to be confirmed.}

\begin{quote}

\emph{Return to the U.
S. Book of Common Prayer (1928)}

\end{quote}

\xchapter{The Order of Confirmation}

\textbf{The Book of Common Prayer (1928)}

Offices
of Instruction

FIRST OFFICE.

¶ \emph{After the singing of a Hymn, shall be said by the Minister
and People together, all kneeling, the following Prayer, the Minister
first pronouncing,}

\emph{Answer.} The Lord be with
you. 

And with thy spirit.

Let us
pray.

LORD
of all power and might, Who art the author and giver of all good things;
Graft in our hearts the love of thy Name, Increase in us true religion,
Nourish us with all goodness, And of thy great mercy keep us in the same;
Through Jesus Christ our Lord. Amen.

¶
\emph{Then, the People being seated, the Minister shall ask them the Questions
which follow, the People reading or repeating the Answers as appointed.}

\emph{Question}. What is your Christian Name? 

\emph{Answer}. My Christian Name is ——.

\emph{Question}. Who gave you this Name? 

\emph{Answer}. My Sponsors gave me this Name in Baptism;
wherein I was made a member of Christ, the child of God, and an inheritor
of the kingdom of heaven.

\emph{Question}. What did your Sponsors then promise
for you?

\emph{Answer}. My Sponsors did promise and vow three
things in my name: First, that I should renounce the devil and all his
works, the pomps and vanity of this wicked world, and all the sinful lusts
of the flesh; Secondly, that I should believe all the Articles of the
Christian Faith; And Thirdly, that I should keep God’s holy will
and commandments, and walk in the same all the days of my life.

\emph{Question.} Do you not think that you are bound
so to do? 

\emph{Answer}. Yes, verily; and by God's help so
I will. And I heartily thank our heavenly Father, that he hath called
me to this state of salvation, through Jesus Christ our Saviour. And I
pray unto God to give me his grace, that I may continue in the same unto
my life’s end.

¶
\emph{Then the Minister shall say,}

YOU
said that your Sponsors promised and vowed that you should believe all
the Articles of the Christian Faith. Recite the Articles of the Christian
Faith as con-tained in the Apostles’ Creed.

¶
\emph{Then, all standing, shall be said the Apostles' Creed by the Minister
and People.}

I
BELIEVE in God the Father Almighty, Maker of heaven and earth: 

And in Jesus Christ his only Son our Lord: Who was
conceived by the Holy Ghost, Born of the Virgin Mary: Suffered under Pontius
Pilate, Was crucified, dead, and buried: He descended into hell; The third
day he rose again from the dead: He ascended into heaven, And sitteth
on the right hand of God the Father Almighty: From thence he shall come
to judge the quick and the dead. 

I believe in the Holy Ghost: The holy Catholic Church;
The Communion of Saints: The Forgiveness of sins: The Resurrection of
the body: And the Life everlasting. Amen. 

¶
\emph{Then, the Minister, turning to the People, shall ask the Question
following, the People responding.}

\emph{Question}. What do you chiefly learn in these Articles of your
Belief?

\emph{Answer}. First, I learn to believe in God the
Father, who hath made me, and all the world. 

Secondly, in God the Son, who hath redeemed me, and
all mankind. 

Thirdly, in God the Holy Ghost, who sanctifieth me,
and all the people of God. 

And this Holy Trinity, One God, I praise and magnify,
saying,

¶
\emph{Minister and People.}

GLORY
be to the Father, and to the Son, and to the Holy Ghost; 

As it was in the beginning, is now,
and ever shall be, world without end. Amen.

¶
\emph{Here may be sung a Hymn, after which the Minister, turning to the
People, shall say,}

YOU
said that your Sponsors promised and vowed that you should keep God’s
holy will and commandments. Tell me how many Commandments there are. 

\emph{Answer}. There are Ten Commandments, given
in old time by God to the people of Israel.

¶
\emph{Then shall the Minister say,}

Let us
ask God's help to know and to keep them.

\emph{Answer.} The Lord be with
you. 

And with thy spirit.

Let us
pray.

¶
\emph{Then shall be said this Prayer by the Minister and People together,
all kneeling.}

O
ALMIGHTY God, Who alone canst order the unruly wills and affections of
sinful men; Grant unto thy people, That they may love the thing which
thou com-mandest, And desire that which thou dost promise; That so, among
the sundry and manifold changes of the world, Our hearts may surely there
be fixed, Where true joys are to be found; Through Jesus Christ our Lord.
Amen. 

¶
\emph{Then shall the Minister repeat the Ten Commandments, and after every
Commandment the People shall say the response. But} NOTE,
\emph{That where so instructed, the People may repeat the Commandments,
the Minister saying theresponse. And}NOTE\emph{further, That the part of the Commandment which is inset may be omitted.}

I.
Thou shalt have none other gods but me.

\emph{Lord, have mercy upon us, and incline our hearts to keep this law.}

II. Thou shalt not make to thyself any graven image, nor the likeness
of any thing that is in heaven above, or in the earth beneath, or in the
water under the earth; thou shalt not bow down to them, nor worship them;

\begin{quote}

for
I the LORD thy God am a jealous God, and visit
the sins of the fathers upon the children, unto the third and fourth
generation of them that hate me; and show mercy unto thou-sands in them
that love me and keep my commandments.

\end{quote}

\emph{Lord, have mercy upon us, and incline our hearts to keep this law.}

III. Thou shalt not take the Name of the LORD thy
God in vain

\begin{quote}

for
the LORD will not hold him guiltless, that taketh
his Name in vain.

\end{quote}

\emph{Lord, have mercy upon us, and incline our hearts to keep this law.}

IV. Remember that thou keep holy the Sabbath-day.

\begin{quote}

Six
days shalt thou labour, and do all that thou hast to do; but the seventh
day is the Sabbath of the LORD thy God. In it
thou shalt do no manner of work; thou, and thy son, and thy daughter,
thy man-servant, and thy maid-servant, thy cattle, and the stranger
that is within thy gates. For in six days the LORD
made heaven and earth, the sea, and all that in them is, and rested
the seventh day: wherefore the LORD blessed the
seventh day, and hallowed it.

\end{quote}

\emph{Lord, have mercy upon us, and incline our hearts to keep this law.}

V.
Honour thy father and thy mother;

\begin{quote}

that
thy days may be long in the land which the Lord thy God giveth thee.

\end{quote}

\emph{Lord, have mercy upon us, and incline our hearts to keep this law.}

VI. Thou shalt do no murder.

\emph{Lord, have mercy upon us, and incline our hearts to keep this law.}

VII. Thou shalt not commit adultery.

\emph{Lord, have mercy upon us, and incline our hearts to keep this law.}

VIII. Thou shalt not steal.

\emph{Lord, have mercy upon us, and incline our hearts to keep this law.}

IX. Thou shalt not bear false witness against thy neighbour.

\emph{Lord, have mercy upon us, and incline our hearts to keep this law.}

X. Thou shalt not covet

\begin{quote}

thy
neighbour's house, thou shalt not covet thy neighbour's wife, nor his
servant, nor his maid, nor his ox, nor his ass, nor any thing that is
his.

\end{quote}

\emph{Lord, have mercy upon us, and write all these thy laws in our hearts,
we beseech thee.}

¶
\emph{Then shall the Minister say,}

GRANT
to us, Lord, we beseech thee, the spirit to think and do always such things
as are right; that we, who cannot do any thing that is good without thee,
may by thee be enabled to live according to thy will; through Jesus Christ
our Lord. \emph{Amen}.

¶
\emph{After this, the People being seated, the Minister, turning to them,
shall ask the Questions which follow, the People reading or repeating
the Answers.}

\emph{Question.} What does our Lord Jesus Christ teach us about these
Commandments?

\emph{Answer.} Our Lord Jesus Christ teaches us that
they are summed up in these two Commandments: Thou shalt love the Lord
thy God with all thy heart, with all thy mind, with all thy soul, and
with all thy strength; this is the first and great Commandment. And the
second is: Thou shalt love thy neighbour as thyself.

\emph{Question}. What then do you chiefly learn from
the Ten Commandments?

\emph{Answer}. I learn two things from these Commandments;
my duty towards God, and my duty towards my Neighbour.

\emph{Question}. What is your duty towards God?

\emph{Answer}. My duty towards God is To believe
in him, to fear him, And to love him with all my heart, with all my mind,
with all my soul, and with all my strength:

I., II. To worship him, to give him thanks: To put my whole trust in him,
to call upon him:

III. To honour his holy Name and his Word:

IV. And to serve him truly all the days of my life.

\emph{Question}. What is your duty towards your Neighbour?

\emph{Answer.} My duty towards my Neighbour is To
love him as myself, and to do to all men as I would they should do unto
me:

V. To love, honour, and help my father and mother: To honour and obey
the civil authority: To submit myself to all my governors, teachers, spiritual
pastors and masters: And to order myself in that lowliness and reverence
which becometh a servant of God:

VI. To hurt nobody by word or deed: To bear no malice
nor hatred in my heart:

VII. To keep my body in temperance, soberness, and
chastity:

VIII. To keep my hands from picking and stealing: To
be true and just in all my dealings:

IX. To keep my tongue from evil speaking, lying, and
slandering:

X. Not to covet nor desire other men's goods; But to
learn and labour truly to earn mine own living, And to do my duty in that
state of life unto which it shall please God to call me.

¶
\emph{Then shall be sung a Hymn, after which the Minister shall say as followeth.}

KNOW
this; that you are not able to do these things of yourself, nor to walk
in the Commandments of God, and to serve him, without his special grace;
which you must learn at all times to call for by diligent prayer. What
is the prayer that our Lord taught us to pray?

\emph{Answer}. The Lord's Prayer.

\emph{Minister}. Let us pray, as our Saviour Christ
hath taught us, and say,

¶
\emph{Then shall be said by the Minister and People together, all kneeling,}

OUR
Father, who art in heaven, Hallowed be thy Name. Thy kingdom come. Thy
will be done, On earth as it is in heaven. Give us this day our daily
bread. And forgive us our trespasses, As we forgive those who trespass
against us. And lead us not into temptation, But deliver us from evil.
For thine is the kingdom, and the power, and the glory, for ever and ever.
Amen.

THE
grace of our Lord Jesus Christ, and the love of God, and the fellowship
of the Holy Ghost, be with us all evermore. \emph{Amen}.

SECOND
OFFICE.

¶\emph{After the singing of a Hymn, there shall be said the following Sentence
by the Minister and People together.}

COME
ye, and let us walk in the light of the Lord. And he will teach us of
his ways, and we will walk in his paths.

\emph{Minister.}Show thy servants thy work; 

\emph{People}. And their children thy glory. 

\emph{Minister}. Let thy merciful kindness, O Lord,
be upon us; 

\emph{People}. As we do put our trust in thee. 

\emph{Minister}. Not unto us, O Lord, not unto us,

\emph{People}. But unto thy Name be the praise. 

\emph{Minister}. Lord, hear our prayer. 

\emph{People}. And let our cry come unto thee. 

\emph{Minister}. The Lord be with you. 

\emph{People}. And with thy spirit. 

\emph{Minister}. Let us pray.

O
ALMIGHTY God, who hast built thy Church upon the foundation of the Apostles
and Prophets, Jesus Christ himself being the head corner-stone; Grant
us so to be joined together in unity of spirit by their doctrine, that
we may be made an holy temple acceptable unto thee; through the same Jesus
Christ our Lord. Amen.

¶
\emph{Here may be sung a Hymn, after which, the People being seated, the
Minister shall ask them the Questions concerning the Church which follow,
the People responding.}

WHEN
were you made a member of the Church?

\emph{Answer}. I was made a member of the Church when I was baptized.

\emph{Question}. What is the Church?

\emph{Answer}. The Church is the Body of which Jesus
Christ is the Head, and all baptized people are the members.

\emph{Question}. How is the Church described in the
Apostles' and Nicene Creeds?

\emph{Answer}. The Church is described in the Creeds
as One, Holy, Catholic, and Apostolic.

\emph{Question}. What do we mean by these words?

\emph{Answer}. We mean that the Church is 

One; because it is one Body under one Head; 

Holy; because the Holy Spirit dwells in it, and sanctifies
its members; 

Catholic; because it is universal, holding earnestly
the Faith for all time, in all countries, and for all people; and is sent
to preach the Gospel to the whole world; 

Apostolic; because it continues stedfastly in the Apostles'
teaching and fellowship.

\emph{Question}. What is your bounden duty as a member
of the Church?

\emph{Answer}. My bounden duty is to follow Christ,
to worship God every Sunday in his Church; and to work and pray and give
for the spread of his kingdom.

\emph{Question.} What special means does the Church
provide to help you to do all these things?

\emph{Answer}. The Church provides the Laying on
of Hands, or Confirmation, wherein, after renewing the promises and vows
of my Baptism, and declaring my loyalty and devotion to Christ as my Master,
I receive the strengthening gifts of the Holy Spirit.

\emph{Question}. After you have been confirmed, what
great privilege doth our Lord provide for you?

\emph{Answer}. Our Lord provides the Sacrament of
the Lord’s Supper, or Holy Communion, for the continual streng-thening
and refreshing of my soul. 

¶
\emph{After another Hymn, the Minister shall proceed with the Questions
on the Sacraments, as followeth.}

HOW
many Sacraments hath Christ ordained in his Church?

\emph{Answer}. Christ hath ordained two Sacraments
only, as generally necessary to salvation; that is to say, Baptism, and
the Supper of the Lord.

\emph{Question}. What do you mean by this word Sacrament?

\emph{Answer}. I mean by this word Sacrament an outward
and visible sign of an inward and spiritual grace given unto us; ordained
by Christ himself, as a means whereby we receive this grace, and a pledge
to assure us thereof.

\emph{Question}. How many parts are there in a Sacrament?

\emph{Answer.} There are two parts in a Sacrament;
the out-ward and visible sign, and the inward and spiritual grace.

\emph{Question}. What is the outward and visible
sign or form in Baptism?

\emph{Answer}. The outward and visible sign or form
in Baptism is Water; wherein the person is baptized, In the Name of the
Father, and of the Son, and of the Holy Ghost.

\emph{Question.} What is the inward and spiritual
grace in Baptism?

\emph{Answer}. The inward and spiritual grace in
Baptism is a death unto sin, and a new birth unto righteousness; whereby
we are made the children of grace.

\emph{Question}. What is required of persons to be
baptized?

\emph{Answer}. Repentance, whereby they forsake sin;
and Faith, whereby they stedfastly believe the promises of God to them
in that Sacrament.

\emph{Question.} Why then are infants baptized, when
by reason of their tender age they cannot perform them?

\emph{Answer}. Because, by the faith of their Sponsors,
infants are received into Christ's Church, become the recipients of his
grace, and are trained in the household of faith.

\emph{Question}. Why was the Sacrament of the Lord's
Supper ordained?

\emph{Answer}. The Sacrament of the Lord's Supper
was ordained for the continual remembrance of the sacrifice of the death
of Christ, and of the benefits which we receive thereby.

\emph{Question}. What is the outward part or sign
of the Lord's Supper?

\emph{Answer.} The outward part or sign of the Lord's
Supper is, Bread and Wine, which the Lord hath commanded to be received.

\emph{Question}. What is the inward part, or thing
signified?

\emph{Answer}. The inward part, or thing signified,
is the Body and Blood of Christ, which are spiritually taken and received
by the faithful in the Lord's Supper.

\emph{Question}. What are the benefits whereof we
are partakers in the Lord's Supper?

\emph{Answer}. The benefits whereof we are partakers
in the Lord’s Supper are the strengthening and refreshing of our
souls by the Body and Blood of Christ, as our bodies are strengthened
and refreshed by the Bread and Wine.

\emph{Question}. What is required of those who come
to the Lord's Supper?

\emph{Answer.} It is required of those who come to
the Lord's Supper to examine themselves, whether they repent them truly
of their former sins, with stedfast purpose to lead a new life; to have
a lively faith in God's mercy through Christ, with a thankful remembrance
of his death; and to be in charity with all men.

¶
\emph{Here may be sung a Hymn, after which the Minister shall ask the
People the following Questions concerning the Ministry, the People responding.}

WHAT
orders of Ministers are there in the Church?

\emph{Answer}. Bishops, Priests, and Deacons; which
orders have been in the Church from the earliest times.

\emph{Question}. What is the office of a Bishop?

\emph{Answer}. The office of a Bishop is, to be
a chief pastor in the Church; to confer Holy Orders; and to administer
Confirmation.

\emph{Question}. What is the office of a Priest?

\emph{Answer}. The office of a Priest is, to minister
to the people committed to his care; to preach the Word of God; to baptize;
to celebrate the Holy Communion; and to pronounce Absolution and Blessing
in God's Name.

\emph{Question}. What is the office of a Deacon?

\emph{Answer}. The office of a Deacon is, to assist
the Priest in Divine Service, and in his other ministrations, under
the direction of the Bishop.

¶ \emph{Then shall the Minister add,}

\emph{Answer.} The Lord be with
you. 

And with thy spirit.

Let us pray.

GRANT, O Lord, that they
who shall renew the promises and vows of their Baptism, and be confirmed
by the Bishop, may receive such a measure of thy Holy Spirit, that they
may grow in grace unto their life's end; through Jesus Christ our Lord.
\emph{Amen}.

GRANT, O Father, that when
we receive the blessed Sacrament of the Body and Blood of Christ, coming
to those holy mysteries in faith, and love, and true repentance, we may
receive remission of our sins, and befilled with thy grace and heavenly
benediction; through Jesus Christ our Lord. \emph{Amen}. 

THE grace of our Lord Jesus Christ, and
the love of God, and the fellowship of the Holy Ghost, be with us all
evermore. \emph{Amen.}

¶ \emph{The Minister
of every Parish shall diligently, upon Sundays and Holy Days, or on some
other convenient occasions, openly in the Church, instruct or examine
the Youth of his Parish.}

¶ \emph{And all Fathers,
Mothers, Guardians, and Sponsors shall bring those, for whose religious
nurture they are responsible, to the Church at the time appointed, to
receive instruction by the Minister.}

¶
\emph{So soon as Children are come to a competent age, and can say the Creed,
the Lord’s Prayer, and the Ten Commandments and are sufficiently
instructed in the matter contained in these Offices they shall be brought
to the Bishop to be confirmed by him.}

The
Order of Confirmation

Or Laying on of Hands upon Those 

that are Baptized, and come to 

Years of Discretion.

¶\emph{Upon the day appointed, all that are to be confirmed shall stand in
order before the Bishop, sitting in his chair near to the Holy Table,
the People all standing until the Lord's Prayer; and the Minister shall
say,}

REVEREND
Father in God, I present unto you these persons to receive the Laying
on of Hands.

\emph{¶ Then the Bishop, or some Minister
appointed by him, may say,}

Hear the words of the Evangelist Saint
Luke, in the eighth Chapter of the Acts of the Apostles.

WHEN the apostles which
were at Jerusalem heard that Samaria had received the word of God, they
sent unto them Peter and John: who, when they were come down, prayed for
them, that they might receive the Holy Ghost: for as yet he was fallen
upon none of them: only they were baptized in the name of the Lord Jesus.
Then laid they their hands on them, and they received the Holy Ghost.

¶ \emph{Then shall the Bishop say,}

DO ye here, in the presence
of God, and of this con- gregation, renew the solemn promise and vow that
ye made, or that was made in your name, at your Baptism; ratifying and
confirming the same; and acknowledging your-selves bound to believe and
to do all those things which ye then undertook, or your Sponsors then
undertook for you?

¶
\emph{And every one shall audibly answer,}

I do.

¶ \emph{Then shall the Bishop say,}

DO ye promise to follow
Jesus Christ as your Lord and Saviour?

¶
\emph{And every one shall answer,}

I do.

\emph{Bishop.} Our help is in the
Name of the Lord;

\emph{Answer.} Who hath made heaven and earth.

\emph{Bishop}. Blessed be the Name of the Lord;

\emph{Answer}. Henceforth, world without end.

\emph{Bishop}. Lord, hear our prayer.

\emph{Answer.} And let our cry come unto thee.

\emph{Bishop}. Let us pray.

ALMIGHTY and everliving
God, who hast vouchsafed to regenerate these thy servants by Water and
the Holy Ghost, and hast given unto them forgiveness of all their sins;
Strengthen them, we beseech thee, O Lord, with the Holy Ghost, the Comforter,
and daily increase in them thy manifold gifts of grace: the spirit of
wisdom and under-standing, the spirit of counsel and ghostly strength,
the spirit of knowledge and true godliness; and fill them, O Lord, with
the spirit of thy holy fear, now and for ever. \emph{Amen}. 

¶
\emph{Then all of them in order kneeling before the Bishop, he shall lay
his hand upon the head of every one severally, saying,}

DEFEND,
O Lord, this thy Child with thy heavenly grace; that he may continue
thine for ever; and daily increase in thy Holy Spirit more and more,
until he come unto thy everlasting kingdom. \emph{Amen}.

¶ \emph{Then shall the Bishop say,}

\emph{Answer.} The Lord be with
you. 

And with thy spirit.

\emph{Bishop}.
Let us pray.

¶ \emph{Then shall the Bishop say the Lord's Prayer, the People kneeling
and repeating it with him.}

OUR
Father, who art in heaven, Hallowed be thy Name. Thy kingdom come. Thy
will be done, On earth as it is in heaven. Give us this day our daily
bread. And forgive us our trespasses, As we forgive those who trespass
against us. And lead us not into temptation, But deliver us from evil.
For thine is the kingdom, and the power, and the glory, for ever and
ever. Amen.

¶ \emph{Then shall the Bishop say,}

ALMIGHTY and everliving
God, who makest us both to will and to do those things which are good,
and accept-able unto thy Divine Majesty; We make our humble suppli-cations
unto thee for these thy servants, upon whom, after the example of thy
holy Apostles, we have now laid our hands, to certify them, by this sign,
of thy favour and gra-cious goodness towards them. Let thy fatherly hand,
we beseech thee, ever be over them; let thy Holy Spirit ever be with them;
and so lead them in the knowledge and obedience of thy Word, that in the
end they may obtain everlasting life; through our Lord Jesus Christ, who
with thee and the same Holy Spirit liveth and reigneth ever, one God,
world without end. \emph{Amen}.

O ALMIGHTY Lord, and everlasting
God, vouchsafe, we beseech thee, to direct, sanctify, and govern, both
our hearts and bodies, in the ways of thy laws, and in the works of thy
commandments; that, through thy most mighty pro-tection, both here and
ever, we may be pre served in body and soul; through our Lord and Saviour
Jesus Christ. Amen.

¶ \emph{Then the Bishop shall bless
them, saying thus,}

THE Blessing of God Almighty,
the Father, the Son, and the Holy Ghost, be upon you, and remain with
you for ever. \emph{Amen}.

¶ \emph{The Minister
shall not omit earnestly to move the Persons confirmed to come, without
delay, to the Lord’s Supper.}

¶ \emph{And there shall
none be admitted to the Holy Communion, until such time as he be confirmed,
or be ready and desirous to be confirmed.}

\begin{quote}

\emph{Return to the U.
S. Book of Common Prayer (1928)}

\end{quote}

\xchapter{The Solemnization of Matrimony}

\textbf{The Book of Common Prayer (1928)}

The Form of

Solemnization of Matrimony

¶
\emph{At the day and time appointed for Solemnization of Matrimony, the
Persons to be married shall come into the body of the Church, or shall
be ready in some proper house, with their friends and neighbours; and
there standing together, the Man on the right hand, and the Woman on the
left, the Minister shall say,}

DEARLY
beloved, we are gathered together here in the sight of God, and in the
face of this company, to join together this Man and this Woman in holy
Matrimony; which is an honourable estate, instituted of God, signifying
unto us the mystical union that is betwixt Christ and his Church: which
holy estate Christ adorned and beautified with his presence and first
miracle that he wrought in Cana of Galilee, and is commended of Saint
Paul to be honour-able among all men: and therefore is not by any to be
entered into unadvisedly or lightly; but reverently, dis-creetly, advisedly,
soberly, and in the fear of God. Into this holy estate these two persons
present come now to be joined. If any man can show just cause, why they
may not lawfully be joined together, let him now speak, or else hereafter
for ever hold his peace.

¶
\emph{And also speaking unto the Persons who are to be married, he shall
say,}

I
REQUIRE and charge you both, as ye will answer at the dreadful day of
judgment when the secrets of all hearts shall be disclosed, that if either
of you know any impedi-ment, why ye may not be lawfully joined together
in Matrimony, ye do now confess it. For be ye well assured, that if any
persons are joined together otherwise than as God’s Word doth allow,
their marriage is not lawful.

¶
\emph{The Minister, if he shall have reason to doubt of the lawfulness of
the proposed Marriage, may demand sufficient surety for his indemnification:
but if no impediment shall be alleged, or suspected, the Minister shall
say to the Man,}

\emph{N.}
WILT thou have this Woman to thy wedded wife, to live together after God’s
ordinance in the holy estate of Matrimony? Wilt thou love her, comfort
her, honour, and keep her in sickness and in health; and, forsaking all
others, keep thee only unto her, so long as ye both shall live?

¶
\emph{The Man shall answer,}

I will.

¶
\emph{Then shall the Minister say unto the Woman,}

\emph{N.}
WILT thou have this Man to thy wedded husband, to live together after
God’s ordinance in the holy estate of Matrimony? Wilt thou love
him, comfort him, honour, and keep him in sickness and in health; and,
forsaking all others, keep thee only unto him, so long as ye both shall
live?

¶
\emph{The Woman shall answer,}

I will.

¶
\emph{Then shall the Minister say,}

Who giveth this
Woman to be married to this Man?

¶
\emph{Then shall they give their troth to each other in this manner. The
Minister, receiving the Woman at her father's or friend's hands, shall
cause the Man with his right hand to take the Woman by her right hand,
and to say after him as followeth.}

I
\emph{N.} take thee \emph{N.} to my wedded Wife, to have and to hold
from this day forward, for better for worse, for richer for poorer, in
sickness and in health, to love and to cherish, till death us do part,
according to God's holy ordinance; and thereto I plight thee my troth.


¶ \emph{Then shall they loose their hands; and the Woman with her right
hand taking the Man by his right hand, shall likewise say after the Minister,}

I
\emph{N}. take thee \emph{N}. to my wedded Husband, to have and to
hold from this day forward, for better for worse, for richer for poorer,
in sickness and in health, to love and to cherish, till death us do part,
according to God's holy ordinance; and thereto I give thee my troth.

¶
\emph{Then shall they again loose their hands; and the Man shall give unto
the Woman a Ring on this wise: the Minister taking the Ring shall deliver
it unto the Man, to put it upon the fourth finger of the Woman’s
left hand. And the Man holding the Ring there, and taught by the Minister,
shall say,}

WITH
this Ring I thee wed: In the Name of the Father, and of the Son, and of
the Holy Ghost. Amen. 

¶\emph{And, before delivering the Ring to the Man, the Minister may say as
followeth.}

BLESS,
O Lord, this Ring, that he who gives it and she who wears it may abide
in thy peace, and continue in thy favour, unto their life's end; through
Jesus Christ our Lord. Amen.

¶\emph{Then, the Man leaving the Ring upon the fourth finger of the Woman's left
hand, the Minister shall say,}

Let us
pray.

¶
\emph{Then shall the Minister and the People, still standing, say the Lord's
Prayer.}

OUR
Father, who art in heaven, Hallowed be thy Name. Thy kingdom come. Thy
will be done, On earth as it is in heaven. Give us this day our daily
bread. And forgive us our trespasses, As we forgive, those who trespass
against us. And lead us not into temptation, But deliver us from evil.
For thine is the kingdom, and the power, and the glory, for ever and ever.
Amen.

¶
\emph{Then shall the Minister add,}

O
ETERNAL God, Creator and Preserver of all mankind, Giver of all spiritual
grace, the Author of everlasting life; Send thy blessing upon these thy
servants, this man and this woman, whom we bless in thy Name; that they,
living faithfully together, may surely perform and keep the vow and covenant
betwixt them made, (whereof this Ring given and received is a token and
pledge,) and may ever remain in perfect love and peace together, and live
accord-ing to thy laws; through Jesus Christ our Lord. \emph{Amen}.

¶
\emph{The Minister may add one or both of the following prayers.}

O
ALMIGHTY God, Creator of mankind, who only art the well-spring of life;
Bestow upon these thy servants, if it be thy will, the gift and heritage
of children; and grant that they may see their children brought up in
thy faith and fear, to the honour and glory of thy Name; through Jesus
Christ our Lord. \emph{Amen}.

O
GOD, who hast so consecrated the state of Matrimony that in it is represented
the spiritual marriage and unity betwixt Christ and his Church; Look mercifully
upon these thy servants, that they may love, honour, and cherish each
other, and so live together in faithfulness and patience, in wisdom and
true godliness, that their home may be a haven of blessing and of peace;
through the same Jesus Christ our Lord, who liveth and reigneth with thee
and the Holy Spirit ever, one God, world without end. \emph{Amen}.

¶
\emph{Then shall the Minister join their right hands together, and say,}

Those whom God hath joined together let no man put asunder.

¶
\emph{Then shall the Minister speak unto the company.}

FORASMUCH
as \emph{N}. and \emph{N}. have consented together in holy wedlock,
and have witnessed the same before God and this company, and thereto have
given and pledged their troth, each to the other, and have declared the
same by giving and receiving a Ring, and by joining hands; I pronounce
that they are Man and Wife, In the Name of the Father, and of the Son,
and of the Holy Ghost. Amen.

¶
\emph{The Man and Wife kneeling, the Minister shall add this Blessing.}

GOD
the Father, God the Son, God the Holy Ghost, bless, preserve, and keep
you; the Lord mercifully with his favour look upon you, and fill you with
all spiritual benediction and grace; that ye may so live together in this
life, that in the world to come ye may have life everlasting.\emph{Amen.}

¶
\emph{The laws respecting Matrimony, whether by publishing the Banns in
Churches, or by Licence, being different in the several States, every
Minister is left to the direction of those laws, in every thing that regards
the civil contract between the parties.}

¶ \emph{And when the Banns are published, it shall be in the following
form:}I publish the Banns of Marriage between\emph{N.} of —,
and\emph{N.}of —. If any of you know cause, or just impediment,
why these two persons should not be joined together in holy Matrimony,
ye are to declare it. This is the first [second\emph{or}third] time
of asking.

The
Thanksgiving of Women 

after Child-birth

Commonly
called the Churching of Women.

¶
\emph{This Service, or the concluding prayer alone, as it stands among
the Occasional Prayers and Thanksgivings, may be used at the discretion
of the Minister.}

¶ \emph{The Woman, at the usual time after her delivery, shall come
into the Church decently apparelled, and there shall kneel down in some
convenient place as hath been accustomed, or as the Ordinary shall direct.}

¶ \emph{The Minister shall then say unto her,}

FORASMUCH
as it hath pleased Almighty God, of his goodness, to give you safe deliverance,
and to preserve you in the great danger of Child-birth; you shall therefore
give hearty thanks unto God, and say,

¶
\emph{Then shall be said by both of them the following Hymn, the Woman still
kneeling.}

\emph{Dilexi,
quoniam}. Psalm cxvi.

MY
delight is in the Lord; because he hath heard the voice of my prayer;

Because he hath inclined his ear unto me; therefore
will I call upon him as long as I live.

I found trouble and heaviness; then called I upon the
Name of the Lord; O Lord, I beseech thee, deliver my soul.

Gracious is the Lord, and righteous; yea, our God is
merciful.

What reward shall I give unto the Lord for all the
benefits that he hath done unto me?

I will receive the cup of salvation, and call upon
the Name of the Lord.

I will pay my vows now in the presence of all his people,
inthe courts of the Lord’s house; even in the midst of thee, O
Jerusalem. Praise the Lord.

Glory be to the Father, and to the Son, and to the Holy Ghost;

As it was in the beginning, is now, and ever shall
be, world without end. Amen.

¶
\emph{Then shall the Minister say the Lord's Prayer, with what followeth:
but the Lord's Prayer may be omitted, if this be used with the Morning
or Evening Prayer.}

OUR
Father, who art in heaven, Hallowed be thy Name. Thy kingdom come. Thy
will be done, On earth as it is in heaven. Give us this day our daily
bread. And forgive us our trespasses, As we forgive those who trespass
against us. And lead us not into temptation, But deliver us from evil.
Amen. 

\emph{Minister}. O Lord, save this woman thy servant;

\emph{Answer}. Who putteth her trust in thee. 

\emph{Minister}. Be thou to her a strong tower; 

\emph{Answer}. From the face of her enemy. 

\emph{Minister.} Lord, hear our prayer. 

\emph{Answer}. And let our cry come unto thee.

\emph{Minister.}
Let us pray.

O
ALMIGHTY God, we give thee humble thanks for that thou hast been graciously
pleased to preserve, through the great pain and peril of child-birth,
this woman, thy servant, who desireth now to offer her praises and thanksgivings
unto thee. Grant, we beseech thee, most merciful Father, that she, through
thy help, may faithfully live according to thy will in this life, and
also may be partaker of everlasting glory in the life to come; through
Jesus Christ our Lord. Amen.

¶
\emph{Then may be said,}

GRANT,
we beseech thee, O heavenly Father, that the child of this thy servant
may daily increase in wisdom and stature, and grow in thy love and service,
until he come to thy eternal joy; through Jesus Christ our Lord. \emph{Amen}.

¶
\emph{The Woman, that cometh to give her Thanks, must offer accustomed offerings,
which shall be applied by the Minister and the Church-wardens to the relief
of distressed women in child-bed; and if there be a Communion, it is convenient
that she receive the Holy Communion.}

\begin{quote}

\emph{Return to the U.
S. Book of Common Prayer (1928)}

\end{quote}

\xchapter{The Visitation of the Sick, and the Communion of the Sick}

\textbf{The Book of Common Prayer (1928)}

The
Order for 

the Visitation of the Sick

¶
\emph{The following Service, or any part thereof, may be used at the discretion
of the Minister.}

¶ \emph{When any person is sick, notice shall be given thereof to
the Minister of the Parish; who, coming into the sick person's presence,
shall say,}

PEACE be to this house,
and to all that dwell in it.

¶\emph{After which he shall say the Antiphon following, and then, according
to his discretion, one of the Penitential Psalms.}

\emph{Antiphon}. Remember not,
Lord, our iniquities; Nor the iniquities of our forefathers.

¶\emph{Then the Minister shall say,}

Let us pray.

\begin{quote}

\begin{quote}

\begin{quote}

Lord,
have mercy upon us. 

\emph{Christ, have mercy upon us.}

Lord, have mercy upon us.

\end{quote}

\end{quote}

\end{quote}

OUR
Father, who art in heaven, Hallowed be thy Name. Thy kingdom come. Thy
will be done, On earth as it is in heaven. Give us this day our daily
bread. And forgive us our trespasses, As we forgive those who trespass
against us. And lead us not into temptation, But deliver us from evil.
Amen.

\emph{Minister}. O Lord, save thy
servant; 

\emph{Answer}. Who putteth \emph{his} trust in
thee. 

\emph{Minister}. Send \emph{him}help from thy
holy place; 

\emph{Answer}. And evermore mightily defend \emph{him.}

\emph{Minister}. Let the enemy have no advantage
of \emph{him};

\emph{Answer.} Nor the wicked approach to hurt \emph{him}.

\emph{Minister.} Be unto \emph{him}, O Lord, a
strong tower; 

\emph{Answer}. From the face of \emph{his} enemy.

\emph{Minister}. O Lord, hear our prayer.

\emph{Answer}. And let our cry come unto thee.

\emph{Minister.}

O LORD, look down from
heaven, behold, visit, and relieve this thy servant. Look upon \emph{him}
with the eyes of thy mercy, give \emph{him} comfort and sure confidence
in thee, defend \emph{him} in all danger, and keep \emph{him} in perpetual
peace and safety; through Jesus Christ our Lord.\emph{Amen}.

¶ \emph{Here may be said
any one or more of the Psalms following, with Antiphon and Collect.}

\emph{Antiphon.} I did call upon
the Lord with my voice; And he heard me out of his holy hill.

\emph{Domine, quid multiplicati?} Psalm iii.

LORD, how are they increased
that trouble me! many are they that rise against me.

Many one there be that say of my soul, There is no
help for him in his God.

But thou, O LORD, art my defender;
thou art my worship, and the lifter up of my head.

I did call upon the LORD with
my voice, and he heard me out of his holy hill.

I laid me down and slept, and rose up again; for the
LORD sustained me.

Salvation belongeth unto the LORD;
and thy blessing is upon thy people.

\emph{The Collect.}

HEAR us, Almighty and most
merciful God and Saviour; extend thy accustomed goodness to this thy servant
who is grieved with sickness. Visit \emph{him}, O Lord, with thy loving
mercy, and so restore \emph{him} to \emph{his}former health, that
\emph{he} may give thanks unto thee in thy holy Church; through Jesus
Christ our Lord. \emph{Amen}.

\emph{Antiphon}. I will go unto the altar of God; Unto the God of my
joy and gladness.

\emph{Judica
me, Deus}. Psalm xliii.

GIVE
sentence with me, O God, and defend my cause against the ungodly people;
O deliver me from the deceitful and wicked man.

For thou art the God of my strength; why hast thou
put me from thee? and why go I so heavily, while the enemy oppresseth
me?

O send out thy light and thy truth, that they may
lead me, and bring me unto thy holy hill, and to thy dwelling;

And that I may go unto the altar of God, even unto
the God of my joy and gladness; and upon the harp will I give thanks
unto thee, O God, my God.

Why art thou so heavy, O my soul? and why art thou
so disquieted within me?

O put thy trust in God; for I will yet give him thanks,
which is the help of my countenance, and my God.

\emph{The
Collect.}

SANCTIFY,
we beseech thee, O Lord, the sickness of this thy servant; that the
sense of \emph{his} weakness may add strength to \emph{his} faith,
and seriousness to \emph{his} repentance; and grant that \emph{he}
may dwell with thee in life everlasting; through Jesus Christ our Lord.
\emph{Amen}.

\emph{Antiphon.} I have considered the days of old; And the years
that are past.

\emph{Voce
mea ad Dominum}. Psalm lxxvii.

I
WILL cry unto God with my voice; even unto God will I cry with my voice,
and he shall hearken unto me.

In the time of my trouble I sought the Lord: I stretched
forth my hands unto him, and ceased not in the night season; my soul
refused comfort.

When I am in heaviness, I will think upon God; when
my heart is vexed, I will complain.

Thou holdest mine eyes waking: I am so feeble that
I cannot speak.

I have considered the days of old, and the years
that are past.

I call to remembrance my song, and in the night I
commune with mine own heart, and search out my spirit.

Will the Lord absent himself for ever? and will he
be no more intreated?

Is his mercy clean gone for ever? and is his promise
come utterly to an end for evermore?

Hath God forgotten to be gracious? and will he shut
up his loving-kindness in displeasure?

And I said, It is mine own infirmity; but I will
remember the years of the right hand of the Most Highest.

\emph{The
Collect.}

HEAR,
O Lord, we beseech thee, these our prayers, as we call upon thee on
behalf of this thy servant; and bestow upon \emph{him} the help of
thy merciful consolation; through Jesus Christ our Lord. \emph{Amen}.

\emph{Antiphon.} Though I walk in the midst of trouble; Yet shalt
thou refresh me.

\emph{Confitebor
tibi}. Psalm cxxxviii.

I
WILL give thanks unto thee, O Lord, with my whole heart; even before
the gods will I sing praise unto thee.

I will worship toward thy holy temple, and praise
thy Name, because of thy loving-kindness and truth; for thou hast magnified
thy Name, and thy word, above all things.

When I called upon thee, thou heardest me; and enduedst
my soul with much strength.

Though I walk in the midst of trouble, yet shalt
thou refresh me; thou shalt stretch forth thy hand upon the furiousness
of mine enemies, and thy right hand shall save me.

The LORD shall make good his
loving-kindness toward me; yea, thy mercy, O Lord, endureth for ever;
despise not then the works of thine own hands.

\emph{The
Collect.}

O
GOD, the strength of the weak and the comfort of sufferers; Mercifully
accept our prayers, and grant to thy servant the help of thy power,
that \emph{his} sickness may be turned into health, and our sorrow
into joy; through Jesus Christ our Lord. \emph{Amen}.

\emph{Antiphon}.
The LORD saveth thy life from destruction; And
crowneth thee with mercy and loving-kindness.

\emph{Benedic,
anima mea}. Psalm ciii.

PRAISE
the LORD, O my soul; and all that is within me,
praise his holy Name.

Praise the LORD, O my soul,
and forget not all his benefits:

Who forgiveth all thy sin, and healeth all thine
infirmities;

Who saveth thy life from destruction, and crowneth
thee with mercy and loving-kindness.

O praise the LORD, ye angels
of his, ye that excel in strength; ye that fulfil his commandment, and
hearken unto the voice of his word.

O praise the LORD, all ye
his hosts; ye servants of his that do his pleasure.

O speak good of the LORD,
all ye works of his, in all places of his dominion: praise thou the
Lord, O my soul.

\emph{The
Collect.}

ACCEPT,
we beseech thee, merciful Lord, the devout praise of thy humble servant;
and grant \emph{him} an abiding sense of thy loving-kindness; through
Jesus Christ our Lord. \emph{Amen}.

¶\emph{Any
of the following Psalms, 20, 27, 42, 91, 121, 146, may, at the discretion
of the Minister, be substituted for any of those given above.}

¶\emph{Here shall be said,}

O
SAVIOUR of the world, who by thy Cross and precious Blood hast redeemed
us; Save us, and help us, we humbly beseech thee, O Lord.

¶
\emph{As occasion demands, the Minister shall address the sick person
on the meaning and use of the time of sickness, and the opportunity
it affords for spiritual profit.}

¶ \emph{Here may the Minister inquire of the sick person as to his
acceptance of the Christian Faith, and as to whether he repent him truly
of his sins, and be in charity with all the world; exhorting him to
forgive, from the bottom of his heart, all persons that have offended
him; and if he hath offended any other, to ask them forgiveness; and
where he hath done injury or wrong to any man, that he make amends to
the uttermost of his power.}

¶ \emph{Then shall the sick person be moved to make a special confession
of his sins, if he feel his conscience troubled with any matter; after
which confession, on evidence of his repentance, the Minister shall
assure him of God's mercy and forgiveness.}

¶ Then the Minister shall say,

Let
us pray.

O
MOST merciful God, who, according to the multitude of thy mercies, dost
so put away the sins of those who truly repent, that thou rememberest
them no more; Open thine eye of mercy upon this thy servant, who most
earnestly desireth pardon and forgiveness. Renew in \emph{him}, most
loving Father, whatsoever hath been decayed by the fraud and malice
of the devil, or by \emph{his} own carnal will and frailness; preserve
and continue this sick member in the unity of the Church; consider \emph{his}
contrition, accept \emph{his} tears, assuage \emph{his} pain, as
shall seem to thee most expedient for \emph{him}. And forasmuch as
\emph{he} putteth \emph{his} full trust only in thy mercy, impute
not unto \emph{him} his former sins, but strengthen \emph{him} with
thy blessed Spirit; and, when thou art pleased to take \emph{him}
hence, take \emph{him} unto thy favour; through the merits of thy
most dearly beloved Son, Jesus Christ our Lord. Amen.

¶
\emph{Then shall the Minister say,}

THE
Almighty Lord, who is a most strong tower to all those who put their
trust in him, to whom all things in heaven, in earth, and under the
earth, do bow and obey; Be now and evermore thy defence; and make thee
know and feel, that there is none other Name under heaven given to man,
in whom, and through whom, thou mayest receive health and salvation,
but only the Name of our Lord Jesus Christ. \emph{Amen}.

¶
\emph{Here the Minister may use any part of the service of the Book, which,
in his discretion, he shall think convenient to the occasion; and after
that shall say,}

UNTO
God's gracious mercy and protection we commit thee. The Lord bless thee,
and keep thee. The Lord make his face to shine upon thee, and be gracious
unto thee. The Lord lift up his countenance upon thee, and give thee
peace, both now and evermore. \emph{Amen}.

PRAYERS.

¶ \emph{These Prayers may be said with the foregoing Service, or any
part thereof, at the discretion of the Minister.}

\emph{A
Prayer for Recovery.}

O
GOD of heavenly powers, who, by the might of thy command, drivest away
from men's bodies all sickness and all infirmity; Be present in thy goodness
with this thy servant, that \emph{his} weakness may be banished and
his strength recalled; that \emph{his} health being thereupon restored,
\emph{he} may bless thy holy Name; through Jesus Christ our Lord. \emph{Amen}.

\emph{A
Prayer for Healing.}

O
ALMIGHTY God, who art the giver of all health, and the aid of them that
turn to thee for succour; We entreat thy strength and goodness in behalf
of this thy servant, that \emph{he} may be healed of his infirmities,
to thine honour and glory; through Jesus Christ our Lord. \emph{Amen}.

\emph{A
Thanksgiving for the Beginning of a Recovery.}

GREAT
and mighty God, who bringest down to the grave, and bringest up again;
We bless thy wonderful goodness, for having turned our heaviness into
joy and our mourning into gladness, by restoring this our \emph{brother}
to some degree of his former health. Blessed be thy Name that thou didst
not forsake \emph{him} in \emph{his} sickness; but didst visit \emph{him}
with comforts from above; didst support him in patience and submission
to thy will; and at last didst send \emph{him} seasonable relief. Perfect,
we beseech thee, this thy mercy towards \emph{him}; and prosper the
means which shall be made use of for \emph{his} cure: that, being restored
to health of body, vigour of mind, and cheerfulness of spirit, \emph{he}
may be able to go to thine house, to offer thee an oblation with great
gladness, and to bless thy holy Name for all thy goodness towards \emph{him};
through Jesus Christ our Saviour, to whom, with thee and the Holy Spirit,
be all honour and glory, world without end. \emph{Amen}.

\emph{A
Prayer for a Sick Person, when there appeareth but small hope of Recovery.}

O
FATHER of mercies, and God of all comfort, our only help in time of need;
We fly unto thee for succour in behalf of this thy servant, here lying
in great weakness of body. Look graciously upon \emph{him}, O Lord;
and the more the outward man decayeth, strengthen \emph{him}, we beseech
thee, so much the more continually with thy grace and Holy Spirit in the
inner man. Give \emph{him} unfeigned repentance for all the errors of
\emph{his} life past, and stedfast faith in thy Son Jesus; that \emph{his}
sins may be done away by thy mercy, and \emph{his} pardon sealed in
heaven; through the same thy Son, our Lord and Saviour. \emph{Amen}.

\emph{A
Prayer for the Despondent.}

COMFORT,
we beseech thee, most gracious God, this thy servant, cast down and faint
of heart amidst the sorrows and difficulties of the world; and grant that,
by the power of thy Holy Spirit, \emph{he} may be enabled to go upon
\emph{his} way rejoicing, and give thee continual thanks for thy sustaining
providence; through Jesus Christ our Lord. \emph{Amen}.

\emph{A Prayer which may be said by the Minister in behalf of all present
at the Visitation.}

O
GOD, whose days are without end, and whose mercies cannot be numbered;
Make us, we beseech thee, deeply sensible of the shortness and uncertainty
of human life; and let thy Holy Spirit lead us in holiness and right-eousness,
all our days: that, when we shall have served thee in our generation,
we may be gathered unto our fathers, having the testimony of a good conscience;
in the communion of the Catholic Church; in the confidence of a certain
faith; in the comfort of a reasonable, religious, and holy hope; in favour
with thee our God, and in perfect charity with the world. All which we
ask through Jesus Christ our Lord. \emph{Amen}.

\emph{A
Commendatory Prayer for a Sick Person at the point of Departure.}

O
ALMIGHTY God, with whom do live the spirits of just men made perfect,
after they are delivered from their earthly prisons; We humbly commend
the soul of this thy servant, our dear \emph{brother}, into thy hands,
as into the hands of a faithful Creator, and most merciful Saviour; beseeching
thee, that it may be precious in thy sight. Wash it, we pray thee, in
the blood of that immaculate Lamb, that was slain to take away the sins
of the world; that whatsoever defilements it may have contracted, through
the lusts of the flesh or the wiles of Satan, being purged and done away,
it may be presented pure and without spot before thee; through the merits
of Jesus Christ thine only Son our Lord. \emph{Amen}.

LITANY FOR THE DYING.

O GOD the Father; 

\emph{Have mercy upon the soul of thy servant.}

O God the Son; 

\emph{Have mercy upon the soul of thy servant.}

O God the Holy Ghost; 

\emph{Have mercy upon the soul of thy servant.}

O holy Trinity, one God; 

\emph{Have mercy upon the soul of thy servant.}

From all evil, from all sin, from all tribulation;

\emph{Good Lord, deliver}him. 

By thy holy Incarnation, by thy Cross and Passion,
by thy precious Death and Burial; 

\emph{Good Lord, deliver} him. 

By thy glorious Resurrection and Ascension, and by
the coming of the Holy Ghost; 

\emph{Good Lord, deliver} him.

We sinners do beseech thee to hear us, O Lord God; That it may please
thee to deliver the soul of thy servant from the power of the evil one,
and from eternal death; 

\emph{We beseech thee to hear us, good Lord.}

That it may please thee mercifully to pardon all \emph{his}sins.

\emph{We beseech thee to hear us, good Lord.}

That it may please thee to grant \emph{him} a place
of refreshment and everlasting blessedness; 

\emph{We beseech thee to hear us, good Lord.}

That it may please thee to give \emph{him} joy and gladness in thy
kingdom, with thy saints in light; 

\emph{We beseech thee to hear us, good Lord.}

O Lamb of God, who takest away the sins of the world; 

\emph{Have mercy upon} him. 

O Lamb of God, who takest away the sins of the world;

\emph{Have mercy upon him.}

O Lamb of God, who takest away the sins of the world;

\emph{Grant him thy peace.}

\begin{quote}

\begin{quote}

\begin{quote}

Lord, have mercy. 

\emph{Christ, have mercy.}

Lord, have mercy.

\end{quote}

\end{quote}

\end{quote}

OUR
Father, who art in heaven, Hallowed be thy Name. Thy kingdom come. Thy
will be done, On earth as it is in heaven. Give us this day our daily
bread. And forgive us our trespasses, As we forgive those who trespass
against us. And lead us not into temptation, But deliver us from evil.
Amen.

Let
us pray.

O
SOVEREIGN Lord, who desirest not the death of a sinner; We beseech thee
to loose the spirit of this thy servant from every bond, and set \emph{him}
free from all evil; that \emph{he} may rest with all thy saints in
the eternal habitations; through Jesus Christ our Lord, who liveth and
reigneth with thee and the Holy Ghost, one God, world without end. \emph{Amen}.

\emph{An
Absolution to be said by the Priest.}

THE
Almighty and merciful Lord grant thee pardon and remission of all thy
sins, and the grace and comfort of the Holy Spirit. \emph{Amen.}

\emph{A
Commendation.}

DEPART,
O Christian soul, out of this world,

In the Name of God the Father Almighty who created thee. 

In the Name of Jesus Christ who redeemed thee. 

In the Name of the Holy Ghost who sanctifieth thee.

May thy rest be this day in peace, and thy dwellingplace
in the Paradise of God.

\emph{A
Commendatory Prayer when the Soul is Departed.}

INTO
thy hands, O merciful Saviour, we commend the soul of thy servant, now
departed from the body. Acknowledge, we humbly beseech thee, a sheep
of thine own fold, a lamb of thine own flock, a sinner of thine own
redeeming. Receive \emph{him} into the arms of thy mercy, into the
blessed rest of everlasting peace, and into the glorious company of
the saints in light. \emph{Amen}.

UNCTION OF THE SICK.

\emph{¶
When any sick person shall in humble faith desire the ministry of healing
through Anointing or Laying on of Hands, the Minister may use such portion
of the foregoing Office as he shall think fit, and the following:}

O
BLESSED Redeemer, relieve, we beseech thee, by thy indwelling power,
the distress of this thy servant; release \emph{him} from sin, and
drive away all pain of soul and body, that being restored to soundness
of health, \emph{he} may offer thee praise and thanksgiving; who livest
and reignest with the Father and the Holy Ghost, one God, world without
end. \emph{Amen}.

I
ANOINT thee with oil (\emph{or} I lay my hand upon thee), In the Name
of the Father, and of the Son, and of the Holy Ghost; beseeching the
mercy of our Lord Jesus Christ, that all thy pain and sickness of body
being put to flight, the blessing of health may be restored unto thee.
Amen.

\emph{¶
The Minister is ordered, from time to time, to advise the People, whilst
they are in health, to make Wills arranging for the disposal of their
temporal goods, and, when of ability, to leave Bequests for religious
and charitable uses.}

\emph{}

The
Communion of the Sick

¶ \emph{Forasmuch as all mortal men are subject to many sudden perils,
diseases and sicknesses, and ever uncertain what time they shall depart
out of this life; therefore, to the intent they may be always in readiness
to die, whensoever it shall please Almighty God to call them, the Ministers
shall diligently from time to time, but especially in the time of pestilence,
or other infectious sickness, exhort their parishioners to the often
receiving of the Holy Communion of the Body and Blood of our Saviour
Christ, when it shall be publicly administered in the Church; that so
doing, they may, in case of sudden visitation, have the less cause to
be disquieted for lack of the same. But if the sick person be not able
to come to the Church, and yet is desirous to receive the Communion
in his house; then he must give timely notice to the Minister, signifying
also how many there are to communicate with him; and all things necessary
being prepared, the Minister shall there celebrate the Holy Communion,
beginning with the Collect, Epistle, and Gospel, here following.}

\emph{The
Collect.}

ALMIGHTY,
everliving God, Maker of mankind, who dost correct those whom thou dost
love, and chastise every one whom thou dost receive; We beseech thee
to have mercy upon this thy servant visited with thine hand, and to
grant that \emph{he} may take \emph{his} sickness patiently, and
recover \emph{his} bodily health, if it be thy gracious will; and
that, whensoever \emph{his} soul shall depart from the body, it may
be without spot presented unto thee; through Jesus Christ our Lord.
\emph{Amen}.

\emph{The
Epistle}. Hebrews xii. 5.

MY
son, despise not thou the chastening of the Lord, nor faint when thou
art rebuked of him: for whom the Lord loveth he chasteneth, and scourgeth
every son whom he receiveth.

\emph{The
Gospel}. St. John v. 24.

VERILY,
verily, I say unto you, He that heareth my word, and believeth on him
that sent me, hath everlasting life, and shall not come into condemnation;
but is passed from death unto life.

¶
\emph{Or the following Collect, Epistle, and Gospel.}

\emph{The
Collect.}

O
LORD, holy Father, by whose loving-kindness our souls and bodies are
renewed; Mercifully look upon this thy servant, that, every cause of
sickness being removed, \emph{he} may be restored to soundness of
health; through Jesus Christ our Lord. \emph{Amen}.

\emph{The
Epistle}. 1 St. John v. 13.

THESE
things have I written unto you that believe on the name of the Son of
God; that ye may know that ye have eternal life, and that ye may believe
on the name of the Son of God. And this is the confidence that we have
in him, that, if we ask any thing according to his will, he heareth
us: and if we know that he hear us, whatsoever we ask, we know that
we have the petitions that we desired of him.

\emph{The
Gospel}. St. John vi. 47.

JESUS
said, Verily, verily, I say unto you, He that believeth on me hath everlasting
life. I am that bread of life. Your fathers did eat manna in the wilderness,
and are dead. This is the bread which cometh down from heaven, that
a man may eat thereof, and not die. I am the living bread which came
down from heaven: if any man eat of this bread, he shall live for ever:
and the bread that I will give is my flesh, which I will give for the
life of the world.

¶
\emph{After which the Minister shall proceed according to the form before
prescribed for the Holy Communion, beginning at these words,}Ye
who do truly,\emph{etc.}

¶ \emph{At the time of the distribution of the holy Sacrament, the
Minister shall first receive the Communion himself, and after minister
unto those who are ap-pointed to communicate with the sick, and last
of all to the sick person.}

¶\emph{When circumstances render it expedient to shorten the Service, the following
form shall suffice: The Confession and the Absolution;} Lift up your
hearts, \emph{etc., through the Sanctus; The Prayer of Consecration, ending
with these words,} partakers of his most blessed Body and Blood;
\emph{The Prayer of Humble Access; The Communion; The Lord's Prayer; The
Blessing. And}NOTE, \emph{That for the Confession
and Absolution the following may be used.}

\emph{The
Confession.}

O
ALMIGHTY Father, Lord of heaven and earth, we confess that we have sinned
against thee in thought, word, and deed. Have mercy upon us, O God,
after thy great goodness; according to the multitude of thy mercies,
do away our offences and cleanse us from our sins; for Jesus Christ's
sake. Amen.

\emph{The
Absolution.}

THE
Almighty and merciful Lord grant you Absolution and Remission of all
your sins, true repentance, amendment of life, and the grace and consolation
of his Holy Spirit. \emph{Amen}.

¶\emph{But if a man, either by reason of extremity of sickness, or for want
of warning in due time to the Minister, or by any other just impediment,
do not receive the Sacrament of Christ's Body and Blood, the Minister
shall instruct him, that if he do truly repent him of his sins, and
stedfastly believe that Jesus Christ hath suffered death upon the Cross
for him, and shed his Blood for his redemption, earnestly remembering
the benefits he hath thereby, and giving him hearty thanks therefor,
he doth eat and drink the Body and Blood of our Saviour Christ profitably
to his soul's health, although he do not receive the Sacrament with
his mouth.}

¶ \emph{This Office may be used with aged and bed-ridden persons,
or such as are not able to attend the public Ministration in Church,
substituting the Collect, Epistle, and Gospel for the Day, for those
appointed above.}

\begin{quote}

\emph{Return to the U.
S. Book of Common Prayer (1928)}

\end{quote}

\xchapter{The Burial of the Dead}

\textbf{The Book of Common Prayer (1928)}

\begin{quote}

\begin{quote}

The Order
for 

The
Burial of the Dead

\end{quote}

\end{quote}

\emph{¶
The Minister, meeting the Body, and going before it, either into the Church
or towards the Grave, shall say or sing,}

I
AM the resurrection and the life, saith the Lord: he that believeth in
me, though he were dead, yet shall he live: and whosoever liveth and believeth
in me, shall never die.

I
know that my redeemer liveth, and that he shall stand at the latter day
upon the earth: and though this body be destroyed, yet shall I see God:
whom I shall see for myself, and mine eyes shall behold, and not as a
stranger.

We
brought nothing into this world, and it is certain we can carry nothing
out. The LORD gave, and the LORD
hath taken away; blessed be the name of the LORD.

\emph{¶ After they are come into the Church, shall be said one or more
or the following Selections, taken from the Psalms. The}Gloria Patri\emph{may be omitted except at the end of the whole Portion or Selection from
the Psalter.}

\emph{Dixi,
custodiam}. Psalm xxxix.

LORD,
let me know mine end, and the number of my days; * that I may be certified
how long I have to live.

Behold, thou hast made my days as it were a span long,
and mine age is even as nothing in respect of thee; * and verily every
man living is altogether vanity.

For man walketh in a vain shadow, and disquieteth himself
in vain; * he heapeth up riches, and cannot tell who shall gather them.

And now, Lord, what is my hope? * truly my hope is
even in thee.

Deliver me from all mine offences; * and make me not
a rebuke unto the foolish.When thou with rebukes dost chasten man for
sin, thou makest his beauty to consume away, like as it were a moth fretting
a garment: * every man therefore is but vanity.

Hear my prayer, O Lord, and with thine ears consider
my calling; * hold not thy peace at my tears;

For I am a stranger with thee, and a sojourner, * as
all my fathers were.

O spare me a little, that I may recover my strength,
* before I go hence, and be no more seen.

\emph{Domine,
refugium}. Psalm xc.

LORD,
thou hast been our refuge, * from one generation to another.

Before the mountains were brought forth, or ever the
earth and the world were made, * thou art God from everlasting, and world
without end.

Thou turnest man to destruction; * again thou sayest,
Come again, ye children of men.

For a thousand years in thy sight are but as yesterday,
when it is past, * and as a watch in the night.

As soon as thou scatterest them they are even as a
sleep; * and fade away suddenly like the grass.

In the morning it is green, and groweth up; * but in
the evening it is cut down, dried up, and withered.

For we consume away in thy displeasure, * and are afraid
at thy wrathful indignation.

Thou hast set our misdeeds before thee; * and our secret
sins in the light of thy countenance.

For when thou art angry all our days are gone: * we
bring our years to an end, as it were a tale that is told.

The days of our age are threescore years and ten; and
though men be so strong that they come to fourscore years, * yet is their
strength then but labour and sorrow; so soon passeth it away, and we are
gone.

So teach us to number our days, * that we may apply
our hearts unto wisdom.

\emph{Dominus
illuininatio.}
Psalm xxvii.

THE
Lord is my light and my salvation; whom then shall I fear? * the Lord
is the strength of my life; of whom then shall I be afraid?

One thing have I desired of the Lord, which I will
require; * even that I may dwell in the house of the Lord all the days
of my life, to behold the fair beauty of the Lord, and to visit his
temple.

For in the time of trouble he shall hide me in his
tabernacle; * yea, in the secret place of his dwelling shall he hide
me, and set me up upon a rock of stone.

And now shall he lift up mine head * above mine enemies
round about me.

Therefore will I offer in his dwelling an oblation,
with great gladness: * I will sing and speak praises unto the Lord.

Hearken unto my voice, O Lord, when I cry unto thee;
* have mercy upon me, and hear me.

My heart hath talked of thee, Seek ye my face: *
Thy face, Lord, will I seek.

O hide not thou thy face from me, * nor cast thy
servant away in displeasure.

Thou hast been my succour; * leave me not, neither
forsake me, O God of my salvation.

I should utterly have fainted, * but that I believe
verily to see the goodness of the Lord in the land of the living.

O tarry thou the Lord’s leisure; * be strong,
and he shall comfort thine heart; and put thou thy trust in the Lord.

\emph{Deus noster refugium}. Psalm xlvi.

GOD is our hope and strength,
* a very present help in trouble.

Therefore will we not fear, though the earth be moved,
* and though the hills be carried into the midst of the sea;

Though the waters thereof rage and swell, * and though
the mountains shake at the tempest of the same.

There is a river, the streams whereof make glad the
city of God; * the holy place of the tabernacle of the Most Highest.

God is in the midst of her, therefore shall she not
be removed; * God shall help her, and that right early.

Be still then, and know that I am God: * I will be
exalted among the nations, and I will be exalted in the earth.

The Lord of hosts is with us; * the God of Jacob is
our refuge.

\emph{Levavi oculos}. Psalm cxxi.

I WILL lift up mine eyes
unto the hills; * from whence cometh my help?

My help cometh even from the Lord, * who hath made
heaven and earth.

He will not suffer thy foot to be moved; * and he that
keepeth thee will not sleep.

Behold, he that keepeth Israel * shall neither slumber
nor sleep.

The Lord himself is thy keeper; * the Lord is thy defence
upon thy right hand;

So that the sun shall not burn thee by day, * neither
the moon by night.

The Lord shall preserve thee from all evil; * yea,
it is even he that shall keep thy soul.

The Lord shall preserve thy going out, and thy coming
in, * from this time forth for evermore.

\emph{De profundis}. Psalm cxxx.

OUT of the deep have I
called unto thee, O Lord; * Lord, hear my voice.

O let thine ears consider well * the voice of my complaint.

If thou, Lord, wilt be extreme to mark what is done
amiss, * O Lord, who may abide it?

For there is mercy with thee; * therefore shalt thou
be feared.

I look for the Lord; my soul doth wait for him; * in
his word is my trust.

My soul fleeth unto the Lord before the morning watch;
* I say, before the morning watch.

O Israel, trust in the Lord, for with the Lord there
is mercy, * and with him is plenteous redemption.

And he shall redeem Israel * from all his sins. 

\emph{¶
Then shall follow the Lesson, taken out of the fifteenth Chapter of the
first Epistle of St. Paul to the Corinthians.}

1 Corinthians
xv. 20.

NOW
is Christ risen from the dead, and become the firstfruits of them that
slept. For since by man came death, by man came also the resurrection
of the dead. For as in Adam all die, even so in Christ shall all be made
alive. But every man in his own order: Christ the firstfruits; afterward
they that are Christ's at his coming. Then cometh the end, when he shall
have delivered up the kingdom to God, even the Father; when he shall have
put down all rule and all authority and power. For he must reign, till
he hath put all enemies under his feet. The last enemy that shall be destroyed
is death. For he hath put all things under his feet. But when he saith
all things are put under him, it is manifest that he is excepted, which
did put all things under him. And when all things shall be subdued unto
him, then shall the Son also himself be subject unto him that put all
things under him, that God may be all in all.

But some man will say, How are the dead raised up?
and with what body do they come? Thou foolish one, that which thou sowest
is not quickened, except it die: and that which thou sowest, thou sowest
not that body that shall be, but bare grain, it may chance of wheat, or
of some other grain: but God giveth it a body as it hath pleased him,
and to every seed its own body. All flesh is not the same flesh: but there
is one kind of flesh of men, another flesh of beasts, another of fishes,
and another of birds. There are also celestial bodies, and bodies terrestrial:
but the glory of the celestial is one, and the glory of the terrestrial
is another. There is one glory of the sun, and another glory of the moon,
and another glory of the stars: for one star differeth from another star
in glory. So also is the resurrection of the dead. It is sown in corruption;
it is raised in incorruption: it is sown in dishonour; it is raised in
glory: it is sown in weakness; it is raised in power: it is sown a natural
body; it is raised a spiritual body. There is a natural body, and there
is a spiritual body. And so it is written, The first man Adam was made
a living soul; the last Adam was made a quickening spirit. Howbeit that
was not first which is spiritual, but that which is natural; and afterward
that which is spiritual. The first man is of the earth, earthy: the second
man is the Lord from heaven. As is the earthy, such are they also that
are earthy: and as is the heavenly, such are they also that are heavenly.
And as we have borne the image of the earthy, we shall also bear the image
of the heavenly.

Now this I say, brethren, that flesh and blood cannot
inherit the kingdom of God; neither doth corruption inherit incorruption.
Behold, I shew you a mystery; We shall not all sleep, but we shall all
be changed, in a moment, in the twinkling of an eye, at the last trump:
for the trumpet shall sound, and the dead shall be raised incorruptible,
and we shall be changed. For this corruptible must put on incor-ruption,
and this mortal must put on immortality. So when this corruptible shall
have put on incorruption, and this mortal shall have put on immortality,
then shall be brought to pass the saying that is written, Death is swallowed
up in victory. O death, where is thy sting? O grave, where is thy victory?
The sting of death is sin; and the strength of sin is the law. But thanks
be to God, which giveth us the victory through our Lord Jesus Christ.
Therefore, my beloved brethren, be ye stedfast, unmoveable, always abounding
in the work of the Lord, forasmuch as ye know that your labour is not
in vain in the Lord. 

\emph{¶
Or this.}

Rom. viii.
14.

AS
many as are led by the Spirit of God, they are the sons of God. For ye
have not received the spirit of bondage again to fear; but ye have received
the Spirit of adoption, whereby we cry, Abba, Father. The Spirit himself
beareth witness with our spirit, that we are the children of God: and
if children, then heirs; heirs of God, and joint heirs with Christ; if
so be that we suffer with him, that we may be also glorified together.
For I reckon that the sufferings of this present time are not worthy to
be compared with the glory which shall be revealed in us. For the earnest
expectation of the creature waiteth for the manifestation of the sons
of God. We know that all things work together for good to them that love
God, to them who are the called according to his purpose. What shall we
then say to these things? If God be for us, who can be against us? He
that spared not his own Son, but delivered him up for us all, how shall
he not with him also freely give us all things? Who is he that condemneth?
It is Christ that died, yea rather, that is risen again, who is even at
the right hand of God, who also maketh intercession for us. Who shall
separate us from the love of Christ? shall tribulation, or distress, or
persecution, or famine, or nakedness, or peril, or sword? Nay, in all
these things we are more than conquerors through him that loved us. For
I am persuaded, that neither death, nor life, nor angels, nor principalities,
nor powers, nor things present, nor things to come, nor height, nor depth,
nor any other creature, shall be able to separate us from the love of
God, which is in Christ Jesus our Lord.

\emph{¶
Or this.}

St. John
xiv. 1.

JESUS
said, Let not your heart be troubled: ye believe in God, believe also
in me. In my Father's house are many mansions: if it were not so, I would
have told you. I go to prepare a place for you. And if I go and prepare
a place for you, I will come again, and receive you unto myself; that
where I am, there ye may be also. And whither I go ye know, and the way
ye know. Thomas saith unto him, Lord, we know not whither thou goest;
and how can we know the way? Jesus saith unto him, I am the way, the truth,
and the life: no man cometh unto the Father, but by me.

\emph{¶
Here may be sung a Hymn or Anthem; and at the discretion of the Minister,
the Creed, the Lord’s Prayer, the prayer which followeth,
and such other fitting Prayers as are elsewhere provided in this Book,
ending with, the Blessing; the Minister, before the Prayers, first pronouncing,}

\begin{quote}

\begin{quote}

The
Lord be with you.

\emph{Answer:} And with
thy spirit.

Let
us pray.

\end{quote}

\end{quote}

REMEMBER
thy servant, O Lord, according to the favour which thou bearest unto thy
people, and grant that, increasing in knowledge and love of thee, \emph{he}
may go from strength to strength, in the life of perfect service, in thy
heavenly kingdom; through Jesus Christ our Lord, who liveth and reigneth
with thee and the Holy Ghost, ever, one God, world without end. \emph{Amen}.

UNTO
God’s gracious mercy and protection we commit you. The Lord bless
you and keep you. The Lord make his face to shine upon you, and be gracious
unto you. The Lord lift up his countenance upon you, and give you peace,
both now and evermore. \emph{Amen}.

\emph{}

AT THE GRAVE

¶
\emph{When they come to the Grave, while the Body is made ready to be
laid into the earth, shall be sung or said,}

MAN,
that is born of a woman, hath but a short time to live, and is full of
misery. He cometh up, and is cut down, like a flower; he fleeth as it
were a shadow, and never continueth in one stay.

In the midst of life we are in death; of whom may we
seek for succour, but of thee, O Lord, who for our sins art justly displeased?

Yet, O Lord God most holy, O Lord most mighty, O holy
and most merciful Saviour, deliver us not into the bitter pains of eternal
death.

Thou knowest, Lord, the secrets of our hearts; shut
not thy merciful ears to our prayer; but spare us, Lord most holy, O God
most mighty, O holy and merciful Saviour, thou most worthy Judge eternal,
suffer us not, at our last hour, for any pains of death, to fall from
thee.

\begin{quote}

\begin{quote}

\emph{¶
Or this.}

\end{quote}

\end{quote}

ALL
that the Father giveth me shall come to me: and him that cometh to me
I will in no wise cast out.

He that raised up Jesus from the dead: will also quicken
your mortal bodies by the spirit which dwelleth in you.

Wherefore my heart is glad, and my glory rejoiceth:
my flesh also shall rest in hope.

Thou shalt show me the path of life; in thy presence
is the fulness of joy: and at thy right hand there is pleasure for evermore.

\emph{¶
Then, while the earth shall be cast upon the Body by some standing by,
the Minister shall say,}

UNTO
Almighty God we commend the soul of our \emph{brother} departed, and
we commit \emph{his} body to the ground; earth to earth, ashes to ashes,
dust to dust; in sure and certain hope of the Resurrection unto eternal
life, through our Lord Jesus Christ, at whose coming in glorious majesty
to judge the world, the earth and the sea shall give up their dead; and
the corruptible bodies of those who sleep in him shall be changed, and
made like unto his own glorious body; according to the mighty working
whereby he is able to subdue all things unto himself.

\emph{¶
Then shall be said or sung,}

I
HEARD a voice from heaven, saying unto me, Write, From henceforth blessed
are the dead who die in the Lord: even so saith the Spirit; for they rest
from their labours.

\begin{quote}

\begin{quote}

\emph{¶
Then the Minister shall say,}

\end{quote}

\end{quote}

The
Lord be with you.

\emph{Answer:} And with
thy spirit.

\begin{quote}

\begin{quote}

Let
us pray.

\end{quote}

Lord,
have mercy upon us.

\emph{Christ, have mercy upon us.}

Lord, have mercy upon us.

\end{quote}

OUR
Father, who art in heaven, Hallowed be thy Name. Thy kingdom come. Thy
will be done on earth, As it is in heaven. Give us this day our daily
bread. And forgive us our trespasses, As we forgive those who trespass
against us. And lead us not into temptation; But deliver us from evil.
Amen.

\emph{¶
Then the Minister shall say one or more of the following Prayers, at his
discretion.}

O
GOD, whose mercies cannot be numbered; Accept our prayers on behalf of
the soul of thy servant departed, and grant him an entrance into the land
of light and joy, in the fellowship of thy saints; through Jesus Christ
our Lord. \emph{Amen}.

ALMIGHTY
God, with whom do live the spirits of those who depart hence in the Lord,
and with whom the souls of the faithful, after they are delivered from
the burden of the flesh, are in joy and felicity; We give thee hearty
thanks for the good examples of all those thy servants, who, having finished
their course in faith, do now rest from their labours. And we beseech
thee, that we, with all those who are departed in the true faith of thy
holy Name, may have our perfect consummation and bliss, both in body and
soul, in thy eternal and everlasting glory; through Jesus Christ our Lord.
\emph{Amen}.

O
MERCIFUL God, the Father of our Lord Jesus Christ, who is the Resurrection
and the Life; in whom whosoever believeth, shall live, though he die;
and whoso-ever liveth, and believeth in him, shall not die eternally;
who also hath taught us, by his holy Apostle Saint Paul, not to be sorry,
as men without hope, for those who sleep in him; We humbly beseech thee,
O Father, to raise us from the death of sin unto the life of righteousness;
that, when we shall depart this life, we may rest in him; and that, at
the general Resurrection in the last day, we may be found acceptable in
thy sight; and receive that blessing, which thy well-beloved Son shall
then pronounce to all who love and fear thee, saying, Come, ye blessed
children of my Father, receive the kingdom prepared for you from the beginning
of the world. Grant this, we beseech thee, O merciful Father, through
Jesus Christ, our Mediator and Redeemer. \emph{Amen}.

THE
God of peace, who brought again from the dead our Lord Jesus Christ, the
great Shepherd of the sheep, through the blood of the everlasting covenant;
Make you perfect in every good work to do his will, working in you that
which is well pleasing in his sight; through Jesus Christ, to whom be
glory for ever and ever. \emph{Amen}.

¶
\emph{The Minister, at his discretion, may also use any of the following
Prayers before the final Blessing.}

O
ALMIGHTY God, the God of the spirits of all flesh, who by a voice from
heaven didst proclaim, Blessed are the dead who die in the Lord; Multiply,
we beseech thee, to those who rest in Jesus, the manifold blessings of
thy love, that the good work which thou didst begin in them may be perfected
unto the day of Jesus Christ. And of thy mercy, O heavenly Father, vouchsafe
that we, who now serve thee here on earth, may at last, together with
them, be found meet to be partakers of the inheritance of the saints in
light; for the sake of the same thy Son Jesus Christ our Lord. Amen.

MOST
merciful Father, who hast been pleased to take unto thyself the soul of
this thy servant (or this thy child); Grant to us who are still in our
pilgrimage, and who walk as yet by faith, that having served thee with
constancy on earth, we may be joined hereafter with thy blessed saints
in glory everlasting; through Jesus Christ our Lord. \emph{Amen}.

O
LORD Jesus Christ, who by thy death didst take away the sting of death;
Grant unto us thy servants so to follow in faith where thou hast led the
way, that we may at length fall asleep peacefully in thee, and awake up
after thy likeness; through thy mercy, who livest with the Father and
the Holy Ghost, one God, world without end. \emph{Amen}.

ALMIGHTY
and everliving God, we yield unto thee most high praise and hearty thanks,
for the wonderful grace and virtue declared in all thy saints, who have
been the choice vessels of thy grace, and the lights of the world in their
several generations; most humbly beseeching thee to give us grace so to
follow the example of their stedfast-ness in thy faith, and obedience
to thy holy command-ments, that at the day of the general Resurrection,
we, with all those who are of the mystical body of thy Son, may be set
on his right hand, and hear that his most joyful voice: Come, ye blessed
of my Father, inherit the kingdom prepared for you from the foundation
of the world. Grant this, O Father, for the sake of the same, thy Son
Jesus Christ, our only Mediator and Advocate. \emph{Amen.}

\emph{¶
Inasmuch as it may sometimes be expedient to say under shelter of the
Church the whole or a part of the Service appointed to be said at the
Grave, the same is hereby allowed for weighty cause.}

\emph{¶
It is to be noted that this Office is appropriate to be used only for
the faithful departed in Christ, provided that in any other case the
Minister may, at his discretion, use such part of this Office, or such
devotions taken from other parts of this Book, as may be fitting.}

\begin{quote}

\begin{quote}

\emph{At
the Burial of the Dead at Sea.}

\end{quote}

\end{quote}

\emph{¶
The same office may be used; but instead of the Sentence of Committal,
the Minister shall say,}

UNTO
Almighty God we commend the soul of our \emph{brother} departed, and
we commit \emph{his} body to the deep; in sure and certain hope of the
Resurrection unto eternal life, through our Lord Jesus Christ; at whose
coming in glorious majesty to judge the world, the sea shall give up her
dead; and the corruptible bodies of those who sleep in him shall be changed,
and made like unto his glorious body; according to the mighty working
whereby he is able to subdue all things unto himself.

\begin{quote}

\begin{quote}

\end{quote}

\end{quote}

\emph{}
\textbf{AT
THE BURIAL OF A CHILD.}

\emph{¶
The Minister, meeting the Body, and going before it, either into the Church
or towards the Grace, shall say:}

I
AM the resurrection and the life, saith the Lord: he that believeth in
me, though he were dead, yet shall he live: and whosoever liveth and believeth
in me, shall never die.

JESUS
called them unto him and said, Suffer the little children to come unto
me, and forbid them not: for of such is the kingdom of God.

HE
shall feed his flock like a shepherd: he shall gather the lambs with his
arms, and carry them in his bosom.

¶
\emph{When they are come into the Church, shall be said the following Psalms;
and at the end of each Psalm shall be said the Gloria Patri;}

\emph{Dominus
regit me}. Psalm xxiii.

THE
Lord is my shepherd; * therefore can I lack nothing.

He shall feed me in a green pasture, * and lead me
forth beside the waters of comfort.

He shall convert my soul, * and bring me forth in the
paths of righteousness for his Name's sake.

Yea, though I walk through the valley of the shadow
of death, I will fear no evil; * for thou art with me; thy rod and thy
staff comfort me.

Thou shalt prepare a table before me in the presence
of them that trouble me; * thou hast anointed my head with oil, and my
cup shall be full.

Surely thy loving-kindness and mercy shall follow me
all the days of my life; * and I will dwell in the house of the Lord for
ever.

\emph{Levavi
oculos}. Psalm cxxi.

I
WILL lift up mine eyes unto the hills; * from whence cometh my help?

My help cometh even from the Lord, * who hath made
heaven and earth.

He will not suffer thy foot to be moved; * and he that
keepeth thee will not sleep.

Behold, he that keepeth Israel * shall neither slumber
nor sleep.

The LORD himself is thy keeper;
* the Lord is thy defence upon thy right hand;

So that the sun shall not burn thee by day, * neither
the moon by night.

The LORD shall preserve thee
from all evil; * yea, it is even he that shall keep thy soul.

The LORD shall preserve thy
going out, and thy coming in, * from this time forth for evermore.

\emph{¶
Then shall follow the Lesson taken out of the eighteenth Chapter of
the Gospel according to St. Matthew.}

AT
the same time came the disciples unto Jesus, saying, Who is the greatest
in the kingdom of heaven? And Jesus called a little child unto him,
and set him in the midst of them, and said, Verily I say unto you, Except
ye be converted, and become as little children, ye shall not enter into
the kingdom of heaven. Whosoever therefore shall humble himself as this
little child, the same is greatest in the kingdom of heaven. And whoso
shall receive one such little child in my name receiveth me. Take heed
that ye despise not one of these little ones; for I say unto you, That
in heaven their angels do always behold the face of my Father which
is in heaven. 

\emph{¶
Here may be sung a Hymn or an Anthem; then shall the Minister say,}

The
Lord be with you.

\emph{Answer:} And with
thy spirit.

Let
us pray.

\begin{quote}

Lord,
have mercy upon us.

\emph{Christ, have mercy upon us.}

Lord, have mercy upon us.

\end{quote}

\emph{¶
Then shall be said by the Minister and People,}

OUR
Father, who art in heaven, Hallowed be thy Name. Thy kingdom come. Thy
will be done on earth, As it is in heaven. Give us this day our daily
bread. And forgive us our trespasses, As we forgive those who trespass
against us. And lead us not into temptation; But deliver us from evil.
Amen.

\emph{Minister. Answer. Minister. Answer. Minister. Answer.} Blessed
are the pure in heart;

For they shall see God.

Blessed be the name of the Lord;

Henceforth, world without end.

Lord, hear our prayer;

And let our cry come unto thee.

\emph{¶
Here shall be said the following Prayers, or other fitting Prayers from
this Book.}

O
MERCIFUL Father, whose face the angels of thy little ones do always behold
in heaven; Grant us stedfastly to believe that this thy child hath been
taken into the safe keeping of thine eternal love; through Jesus Christ
our Lord. \emph{Amen}.

ALMIGHTY
and merciful Father, who dost grant to children an abundant entrance into
thy kingdom; Grant us grace so to conform our lives to their innocency
andperfect faith, that at length, united with them, we may stand in thy
presence in fulness of joy; through Jesus Christ our Lord. \emph{Amen.}

THE
grace of our Lord Jesus Christ, and the love of God, and the fellowship
of the Holy Ghost, be with us all evermore. \emph{Amen}.

¶\emph{When they are come to the Grave shall be said or sung,}

JESUS
saith to his disciples, Ye now therefore have sorrow: but I will see you
again, and your heart shall rejoice, and your joy no man taketh from you.

¶
\emph{While the earth is being cast upon the Body, the Minister shall say,}

IN
sure and certain hope of the Resurrection to eternal life through our
Lord Jesus Christ, we commit the body of this child to the ground. The
LORD bless him and keep \emph{him}, the LORD
make his face to shine upon \emph{him} and be gracious unto \emph{him},
the LORD lift up his countenance upon \emph{him},
and give \emph{him} peace, both now and evermore.

¶
\emph{Then shall be said or sung,}

THEREFORE
are they before the throne of God, and serve him day and night in his
temple: and he that sitteth on the throne shall dwell among them.

They shall hunger no more, neither thirst any more;
neither shall the sun light on them, nor any heat.

For the Lamb which is in the midst of the throne shall
feed them, and shall lead them unto living fountains of waters: and God
shall wipe away all tears from their eyes.

¶
\emph{Then shall the Minister say,}

The
Lord be with you.

\emph{Answer:} And with thy
spirit.

Let us
pray.

O
GOD, whose most dear Son did take little children into his arms and bless
them; Give us grace, we beseech thee, to entrust the soul of this child
to thy neverfailing care and love, and bring us all to thy heavenly kingdom;
through the same thy Son, Jesus Christ our Lord. \emph{Amen}.

ALMIGHTY
God, Father of mercies and giver of all comfort; Deal graciously, we pray
thee, with all those who mourn, that, casting every care on thee, they
may know the consolation of thy love; through Jesus Christ our Lord. \emph{Amen}.

MAY
Almighty God, the Father, the Son, and the Holy Ghost, bless you and keep
you, now and for evermore. \emph{Amen.}

\begin{quote}

\emph{Return to the U.
S. Book of Common Prayer (1928)}

\end{quote}

\xchapter{The Psalter, Books I and II}

\textbf{The Psalter from the 1789, 1892 and 1928 U. S. Books of Common
Prayer}

The
Psalter 

or Psalms of David

BOOK
I.

The
First Day.

Morning
Prayer.

Psalm
1. \emph{Beatus vir qui non abiit.}

BLESSED
is the man that hath not walked in the counsel of the ungodly, nor stood
in the way of sinners, * and hath not sat in the seat of the scornful.

2 But his delight is in the law of the LORD;
* and in his law will he exercise himself day and night.

3 And he shall be like a tree planted by the water-side,
* that will bring forth his fruit in due season.

4 His leaf also shall not wither; * and look, whatsoever
he doeth, it shall prosper.

5 As for the ungodly, it is not so with them; * but
they are like the chaff, which the wind scattereth away from the face
of the earth.

6 Therefore the ungodly shall not be able to stand
in the judgment, * neither the sinners in the congregation of the righteous.

7 But the LORD knoweth the way
of the righteous; * and the way of the ungodly shall perish.

Psalm
2. \emph{Quare fremuerunt gentes?}

WHY
do the heathen so furiously rage together? * and why do the people imagine
a vain thing?

2 The kings of the earth stand up, and the rulers take
counsel together * against the LORD, and against
his Anointed:

3 Let us break their bonds asunder, * and cast away
their cords from us.

4 He that dwelleth in heaven shall laugh them to scorn:
* the Lord shall have them in derision.

5 Then shall he speak unto them in his wrath, * and
vex them in his sore displeasure:

6 Yet have I set my King * upon my holy hill of Sion.

Go
to Psalms 51-100; 

101-150

Return to the U. S. Book of Common Prayer: 1789;
1892; 1928

7 I will rehearse the decree; * the LORD hath said
unto me, Thou art my Son, this day have I begotten thee.

8 Desire of me, and I shall give thee the nations1
for thine inheritance, * and the utmost parts of the earth for thy possession.

9 Thou shalt bruise them with a rod of iron, * and break
them in pieces like a potter's vessel.

7 I will preach the law,
whereof the LORD ... until
1928.

1 the heathen until
1928.

10 Be wise now therefore, O ye kings; * be instructed2, ye
that are judges of the earth.

11 Serve the LORD in fear, *
and rejoice unto him with reverence.

12 Kiss the Son, lest he be angry, and so ye perish
from the right way, if his wrath be kindled, yea but a little. * Blessed3
are all they that put their trust in him.

2 be
learned until
1928.

3blessed (i.
e., not capitalized) until 1845.

Psalm
3. \emph{Domine, quid multiplicati?}

LORD,
how are they increased that trouble me! * many are they that rise against
me.

2 Many one there be that say of my soul, * There is
no help for him in his God.

3 But thou, O LORD, art my defender;
* thou art my worship, and the lifter up of my head.

4 I did call upon the LORD with
my voice, * and he heard me out of his holy hill.

5 I laid me down and slept, and rose up again; * for
the LORD sustained me.

6 I will not be afraid for ten thousands of the people,
* that have set themselves against me round about.

7 Up, LORD, and help me, O my
God, * for thou smitest all mine enemies upon the cheekbone; thou hast
broken the teeth of the ungodly.

8 Salvation belongeth unto the LORD;
* and thy blessing is upon thy people. 

Psalm
4. \emph{Cum invocarem.}

HEAR
me when I call, O God of my righteousness: * thou hast set me at liberty
when I was in trouble; have mercy upon me, and hearken unto my prayer.

2 O ye sons of men, how long will ye blaspheme mine
honour, * and have such pleasure in vanity, and seek after falsehood1?

No
exclamation point before 1928.

1 leasing
in Engl. Book

3 Know this also, that the LORD
hath chosen to himself the man that is godly; * when I call upon the LORD
he will hear me.

4 Stand in awe, and sin not; * commune with your own
heart, and in your chamber, and be still.

5 Offer the sacrifice of righteousness, * and put your
trust in the LORD.

6 There be many that say, * Who will show us any good?

7 LORD, lift thou up* the light
of thy countenance upon us.

8 Thou hast put gladness in my heart; * yea, more than
when their corn and wine and oil increase.

9 I will lay me down in peace, and take my rest; *
for it is thou, LORD, only, that makest me dwell
in safety.

8 ... heart : since the time that
their corn ... until
1928.

Psalm
5. \emph{Verba mea auribus.}

PONDER
my words, O LORD, * consider my meditation.

2 O hearken thou unto the voice of my calling, my King
and my God: * for unto thee will I make my prayer.

3 My voice shalt thou hear betimes, O LORD;
* early in the morning will I direct my prayer unto thee, and will look
up.

4 For thou art the God that hast no pleasure in wickedness;
* neither shall any evil dwell with thee.

5 Such as be foolish shall not stand in thy sight;
* for thou hatest all them that work iniquity.

6 Thou shalt destroy them that speak lies1:
* the LORD will abhor both the blood-thirsty and
deceitful man. 

1 leasing
in Engl. Book

7 But as for me, in the multitude of thy mercy I will come into thine
house; * and in thy fear will I worship toward thy holy temple.

8 Lead me, O LORD, in thy righteousness,
because of mine enemies; * make thy way plain before my face.

9 For there is no faithfulness in their mouth2;
* their inward parts are very wickedness.

10 Their throat is an open sepulchre; * they flatter
with their tongue.

11 Destroy thou them, O God; let them perish through
their own imaginations; * cast them out in the multitude of their ungodliness;
for they have rebelled against thee.

12 And let all them that put their trust in thee rejoice:
* they shall ever be giving of thanks, because thou defendest them; they
that love thy Name shall be joyful in thee;

13 For thou, LORD, wilt give
thy blessing unto the righteous, * and with thy favourable kindness wilt
thou defend him as with a shield.

Evening
Prayer.

Psalm
6. \emph{Domine, ne in furore.}

O
LORD, rebuke me not in thine indignation, * neither chasten me in thy
displeasure.

2 Have mercy upon me, O LORD,
for I am weak; * O LORD, heal me, for my bones
are vexed.

3 My soul also is sore troubled: * but, LORD,
how long wilt thou punish me?

4 Turn thee, O LORD, and deliver
my soul; * O save me, for thy mercy's1 sake.

5 For in death no man remembereth thee; * and who will
give thee thanks in the pit?

6 I am weary of my groaning; * every night wash I my
bed, and water my couch with my tears.

7 My beauty is gone for very trouble, * and worn away
because of all mine enemies.

7 But as for me, I will come into thine
house, even upon the multitude of thy mercy until
1928.

2 his mouth until
1928.

1 mercies until
1822; mercies' 1822-45

8 Away from me, all ye that work iniquity2; * for the LORD
hath heard the voice of my weeping.

9 The LORD hath heard my petition;
* the LORD will receive my prayer.

10 All mine enemies shall be confounded, and sore vexed;
* they shall be turned back, and put to shame suddenly.

Psalm 7. \emph{Domine, Deus meus}.

O LORD my God, in thee have I put my trust:
* save me from all them that persecute me, and deliver me;

2 Lest he devour my soul like a lion, and tear it in
pieces, * while there is none to help.

3 O LORD my God, if I have done
any such thing; * or if there be any wickedness in my hands;
2 vanity until
1928.

4 If I have rewarded evil unto him that dealt friendly with
me; * (yea, I have delivered him that without any cause is mine enemy;)

5 Then let mine enemy persecute my soul, and take me; * yea,
let him tread my life down upon the earth, and lay mine honour in the dust.

6 Stand up, O LORD, in thy wrath, and
lift up thyself, because of the indignation of mine enemies; * arise up
for me in the judgment that thou hast commanded.

no
parentheses until 1928

7 And so shall the congregation of the peoples1 come about
thee: * for their sakes therefore lift up thyself again.

8 The LORD shall judge the peoples1:
give sentence with me, O LORD, * according to my
righteousness, and according to the innocency that is in me.

9 O let the wickedness of the ungodly come to an end;
* but guide thou the just.

10 For the righteous God * trieth the very hearts and
reins.

11 My help cometh of God, * who preserveth them that
are true of heart.

12 God is a righteous Judge, strong, and patient; * and
God is provoked every day.

13 If a man will not turn, he will whet his sword; *
he hath bent his bow, and made it ready.

14 He hath prepared for him the instruments of death;
* he ordaineth his arrows against the persecutors.

1 people
until 1928

15 Behold, the ungodly travaileth with iniquity; * he hath conceived
mischief, and brought forth falsehood.

16 He hath graven and digged up a pit, * and is fallen
himself into the destruction that he made for other.

17 For his travail shall come upon his own head, * and
his wickedness shall fall on his own pate.

18 I will give thanks unto the LORD,
according to his righteousness; * and I will praise the Name of the LORD
Most High.

Psalm
8. \emph{Domine, Dominus noster.}

O
LORD our Governor, how excellent is thy Name in all the world; * thou
that hast set thy glory above the heavens!

2 Out of the mouth of very babes and sucklings hast
thou ordained strength, because of thine enemies, * that thou mightest
still the enemy and the avenger.

15 Behold, he travaileth with
mischief : he hath conceived sorrow, and brought forth ungodliness.
until 1928

3 When I consider thy heavens, even the work of thy fingers; * the
moon and the stars which thou hast ordained;

4 What is man, that thou art mindful of him? * and the
son of man, that thou visitest him?

5 Thou madest him lower than the angels, * to crown him
with glory and worship.

6 Thou makest him to have dominion of the works of thy
hands; * and thou hast put all things in subjection under his feet:

7 All sheep and oxen; * yea, and the beasts of the field;

8 The fowls of the air, and the fishes of the sea; *
and whatsoever walketh through the paths of the seas.

9 O LORD our Governor, * how excellent
is thy Name in all the world!

The Second Day.

Morning
Prayer.

Psalm
9. \emph{Confitebor tibi.}

I
WILL give thanks unto thee, O LORD, with my whole
heart; * I will speak of all thy marvellous works.

2 I will be glad and rejoice in thee; * yea, my songs
will I make of thy Name, O thou Most Highest.

3 While mine enemies are driven back, * they shall
fall and perish at thy presence.

4 For thou hast maintained my right and my cause; *
thou art set in the throne that judgest right.

5 Thou hast rebuked the heathen, and destroyed the
ungodly; * thou hast put out their name for ever and ever. 3 For I will consider thy
heavens
until 1928

6 O thou enemy, thy1 destructions are come to a perpetual
end; * even as the cities which thou hast destroyed, whose memorial is perished
with them.

7 But the LORD shall endure for
ever; * he hath also prepared his seat for judgment.

8 For he shall judge the world in righteousness, * and
minister true judgment unto the people.

9 The LORD also will be a defence
for the oppressed, * even a refuge in due time of trouble.

10 And they that know thy Name will put their trust in
thee; * for thou, LORD, hast never failed them that
seek thee.

11 O praise the LORD which dwelleth
in Sion; * show the people of his doings.

12 For when he maketh inquisition for blood, he remembereth
them, * and forgetteth not the complaint of the poor.

13 Have mercy upon me, O LORD;
consider the trouble which I suffer of them that hate me, * thou that liftest
me up from the gates of death;

14 That I may show all thy praises within the gates2
of the daughter of Sion: * I will rejoice in thy salvation.

15 The heathen are sunk down in the pit that they made;
* in the same net which they hid privily is their foot taken.

16 The LORD is known to execute
judgment; * the ungodly is trapped in the work of his own hands.

1 thy
introduced in 1928.

2 ports until
1928

17 The wicked shall be turned to destruction3,
* and all the people that forget God.

18 For the poor shall not alway be forgotten; * the patient
abiding of the meek shall not perish for ever.

19 Up, LORD, and let not man have
the upper hand; * let the heathen be judged in thy sight.

20 Put them in fear, O LORD, *
that the heathen may know themselves to be but men.


Psalm 10. \emph{Ut quid,
Domine ?}

WHY
standest thou so far off, O LORD, * and hidest
thy face in the needful time of trouble?

2 The ungodly, for his own lust, doth persecute the
poor: * let them be taken in the crafty wiliness that they have imagined.

3 For the ungodly hath made boast of his own heart's
desire, * and speaketh good of the covetous, whom the LORD
abhorreth.

4 The ungodly is so proud, that he careth not for God,
* neither is God in all his thoughts.

3 hell until
1928

5 His ways are alway1 grievous; * thy judgments are far above
out of his sight, and therefore defieth he all his enemies.

6 For he hath said in his heart, Tush, I shall never
be cast down, * there shall no harm happen unto me.

7 His mouth is full of cursing, deceit, and fraud; *
under his tongue is ungodliness and vanity.

8 He sitteth lurking in the thievish corners of the streets,
* and privily in his lurking dens doth he murder the innocent; his eyes
are set against the poor.

9 For he lieth waiting secretly; even as a lion lurketh
he in his den, * that he may ravish the poor.

10 He doth ravish the poor, * when he getteth him into
his net.

11 He falleth down, and humbleth himself, * that the
congregation of the poor may fall into the hands of his captains.

12 He hath said in his heart, Tush, God hath forgotten;
* he hideth away his face, and he will never see it.

13 Arise, O LORD God, and lift
up thine hand; * forget not the poor.

14 Wherefore should the wicked blaspheme God, * while
he doth say in his heart, Tush, thou God carest not for it? 1 always
in the English Book

15 Surely thou hast seen it; * for thou beholdest ungodliness
and wrong, that thou mayest take the matter into thy hand.

16 The poor committeth himself unto thee; * for thou
art the helper of the friendless.

17 Break thou the power of the ungodly and malicious;
* search out his ungodliness, until thou find none.

18 The LORD is King for ever and
ever, * and the heathen are perished out of the land.

19 LORD, thou hast heard the desire
of the poor; * thou preparest their heart, and thine ear hearkeneth2;

20 To help the fatherless and poor unto their right,
* that the man of the earth be no more exalted against them.

Psalm
11. \emph{In Domino confido.}

IN
the LORD put I my trust; * how say ye then to my
soul, that she should flee as a bird unto the hill?

2 For lo, the ungodly bend their bow, and make ready
their arrows within the quiver, * that they may privily shoot at them
which are true of heart.

15 Surely thou hast seen it : for thou
beholdest ungodliness and wrong, 

16 That thou mayest take the matter into thy hand :
the poor committeth himself unto thee; for thou art the helper of the
friendless

17 Break thou the power of the ungodly and malicious
: take away his ungodliness, and thou shalt find none. until
1928

2 hearkeneth thereto; until
1928

3 If the foundations be destroyed, * what can the righteous
do?

4 The LORD is in his holy temple;
* the LORD’S seat is in heaven.

5 His eyes consider the poor, * and his eyelids try the
children of men.

6 The LORD approveth1
the righteous: * but the ungodly, and him that delighteth in wickedness,
doth his soul abhor.

7 Upon the ungodly he shall rain snares, fire and brimstone,
storm and tempest: * this shall be their portion to drink.

8 For the righteous LORD loveth
righteousness; * his countenance will behold the thing that is just.

Evening
Prayer.

Psalm
12. \emph{Salvum me fac.}

HELP
me,LORD, for there is not one godly man left; *
for the faithful are minished from among the children of men.

2 They talk of vanity every one with his neighbour; * they do but flatter
with their lips, and dissemble in their double heart.

3 The LORD shall root out all deceitful lips, *
and the tongue that speaketh proud things;

4 Which have said, With our tongue will we prevail; * we are they that
ought to speak; who is lord over us?

5 Now, for the comfortless troubles' sake of the needy, * and because
of the deep sighing of the poor,

6 I will up, saith the LORD; * and will help every
one from him that swelleth against him, and will set him at rest.

7 The words of the LORD are pure words; * even
as the silver which from the earth is tried, and purified seven times
in the fire.

8 Thou shalt keep them, O LORD; * thou shalt preserve
them from this generation for ever.

9 The ungodly walk on every side: * when they are exalted, the children
of men are put to rebuke.

3 For the foundations will be cast
down : and what hath the righteous done? until
1928

1 alloweth until
1928

Psalm
13. \emph{Usquequo, Domine?}

HOW
long wilt thou forget me, O LORD; for ever? * how
long wilt thou hide thy face from me?

2 How long shall I seek counsel in my soul, and be
so vexed in my heart? * how long shall mine enemy1 triumph
over me?

3 Consider, and hear me, O LORD
my God; * lighten mine eyes, that I sleep not in death;

4 Lest mine enemy say, I have prevailed against him:
* for if I be cast down, they that trouble me will rejoice at it.

5 But my trust is in thy mercy, * and my heart is joyful
in thy salvation.

6 I will sing of the LORD, because
he hath dealt so lovingly with me; * yea, I will praise the Name of the
Lord Most Highest.

1 enemies until
1928

Psalm
14. \emph{Dixit insipiens.}

THE
fool hath said in his heart, * There is no God.

2 They are corrupt, and become abominable in their doings; * there is
none that doeth good, no not one.

3 The LORD looked down from
heaven upon the children of men, * to see if there were any that would
understand, and seek after God.

4 But they are all gone out of the way, they are altogether
become abominable; * there is none that doeth good, no not one.

5 Have they no knowledge, that they are all such workers
of mischief, * eating up my people as it were bread, and call not upon
the LORD?

6 There were they brought in great fear, even where
no fear was; * for God is in the generation of the righteous.

7 As for you, ye have made a mock at the counsel of
the poor; * because he putteth his trust in the LORD.

8 Who shall give salvation unto Israel out of Sion?
* When the LORD turneth the captivity of his people,
then shall Jacob rejoice, and Israel shall be glad.

The
following additional verses were present prior to 1928: 

5. Their throat is an open sepulchre, with their tongues have they deceived
: the poison of asps is under their lips.\\
6. Their
mouth is full of cursing and bitterness : their feet are swift to shed blood.

7. Destruction and unhappiness is in their ways, and
the way of peace have they not known ; there is no fear of God before their
eyes. \\
giving
this Psalm 11 verses.

The
Third Day.

Morning Prayer.

Psalm
15. \emph{Domine, quis habitabit?}

LORD,
who shall dwell in thy tabernacle? * or who shall rest upon thy holy hill?

2 Even he that leadeth an uncorrupt life, * and doeth
the thing which is right, and speaketh the truth from his heart.

3 He that hath used no deceit in his tongue, nor done
evil to his neighbour, * and hath not slandered his neighbour.

4 He that setteth not by himself, but is lowly in his
own eyes, * and maketh much of them that fear the LORD.

5 He that sweareth unto his neighbour, and disappointeth
him not, * though it were to his own hindrance.

6 He that hath not given his money upon usury, * nor
taken reward against the innocent.

7 Whoso doeth these things * shall never fall.

Psalm
16. \emph{Conserva me, Domine.}

PRESERVE
me, O God; * for in thee have I put my trust.

2 O my soul, thou hast said unto the LORD,
* Thou art my God; I have no good like unto thee.

3 All my delight is upon the saints that are in the
earth, * and upon such as excel in virtue.

4 But they that run after another god * shall have
great trouble.

5 Their drink-offerings of blood will I not offer,
* neither make mention of their names within my lips.

6 The LORD himself is the portion
of mine inheritance, and of my cup; * thou shalt maintain my lot.

7 The lot is fallen unto me in a fair ground; * yea,
I have a goodly heritage.

8 I will thank the LORD for
giving me warning; * my reins also chasten me in the night season.

2 ... : Thou art my God; my goods are nothing unto
thee. until
1928

9 I have set the LORD alway before me; * for
he is on my right hand, therefore I shall not fall.

10 Wherefore my heart is glad, and my glory rejoiceth1:
* my flesh also shall rest in hope.

11 For why? thou shalt not leave my soul in hell; * neither
shalt thou suffer thy Holy One to see corruption.

12 Thou shalt show me the path of life: in thy presence
is the fulness of joy, * and at thy right hand there is pleasure for evermore.

Psalm
17. \emph{Exaudi, Domine}.

HEAR
the right, O LORD, consider my complaint, * and
hearken unto my prayer, that goeth not out of feigned lips.

2 Let my sentence come forth from thy presence; * and
let thine eyes look upon the thing that is equal.

3 Thou hast proved and visited mine heart in the night
season; thou hast tried me, and shalt find no wickedness in me; * for
I am utterly purposed that my mouth shall not offend.

9 I have set GOD
always until
1892; I have set GOD alway 1892-1928

1 was rejoiced until
1928

4 As for the works of men, * by the word of thy lips I have kept
me from the ways of the destroyer.

5 O hold thou up my goings in thy paths, * that my footsteps
slip not.

6 I have called upon thee, O God, for thou shalt hear
me: * incline thine ear to me, and hearken unto my words.

7 Show thy marvellous loving-kindness, thou that art
the Saviour of them which put their trust in thee, * from such as resist
thy right hand.

8 Keep me as the apple of an eye; * hide me under the
shadow of thy wings,

9 From the ungodly, that trouble me; * mine enemies compass
me round about, to take away my soul.

10 They are inclosed in their own fat, * and their mouth
speaketh proud things.

11 They lie waiting in our way on every side, * watching
to cast us down to the ground;

12 Like as a lion that is greedy of his prey, * and as
it were a lion’s whelp lurking in secret places.

4 Because of men's works that
are done against the words of thy lips : until
1928

11 ... : turning their eyes down
to the ground;
until 1928

13 Up, LORD, disappoint him, and cast him down; *
deliver my soul from the ungodly, by thine own sword; 

14 Yea, by thy hand, O LORD; from
the men of the evil world; * which have their portion in this life, whose
bellies thou fillest with thy hid treasure.

15 They have children at their desire, * and leave the
rest of their substance for their babes.

16 But as for me, I shall behold thy presence in righteousness;
* and when I awake up after thy likeness, I shall be satisfied. 

13 ... from the ungodly, which
is a sword of thine;\\
14 From the men of the hand, O LORD,
from the men, I say, and from the evil world : until
1928

with it added
to end of verse 16, until 1928

Evening Prayer.

Psalm 18. \emph{Diligam
te, Domine.}

I
WILL love thee, O LORD, my strength. * The LORD
is my stony rock, and my defence;

2 My Saviour, my God, and my might, in whom I will
trust; * my buckler, the horn also of my salvation, and my refuge.

3 I will call upon the LORD,
which is worthy to be praised; * so shall I be safe from mine enemies.

4 The sorrows of death compassed me, * and the overflowings
of ungodliness made me afraid.

5 The pains of hell came about me; * the snares of
death overtook me.

6 In my trouble I called upon the LORD,
* and complained unto my God:

7 So he heard my voice out of his holy temple, * and
my complaint came before him; it entered even into his ears.

8
The earth trembled and quaked, * the very foundations also of the hills
shook, and were removed, because he was wroth.

9 There went a smoke out in his presence, * and a consuming
fire out of his mouth, so that coals were kindled at it.

10 He bowed the heavens also, and came down, * and
it was dark under his feet.

Vs.
1 \& 2 were combined as one verse prior to 1928 (and the rest through
v. 44 numbered one lower)

11 He rode upon the Cherubim1, and did fly; * he came flying
upon the wings of the wind.

12 He made darkness his secret place, * his pavilion
round about him with dark water, and thick clouds to cover him.

13 At the brightness of his presence his clouds removed;
* hailstones and coals of fire.

14 The LORD also thundered out
of heaven, and the Highest gave his thunder; * hailstones and coals of fire.

15 He sent out his arrows, and scattered them; * he cast
forth lightnings, and destroyed them.

16 The springs of waters were seen, and the foundations
of the round world were discovered, * at thy chiding, O LORD,
at the blasting of the breath of thy displeasure.

17 He sent down from on high to fetch me, * and took
me out of many waters.

18 He delivered me from my strongest enemy, and from
them which hate me; * for they were too mighty for me. 1 Cherubins in
the English Book; Cherubims 1790-1793

19 They came upon me in the day of my trouble; * but
the LORD was my upholder.

20 He brought me forth also into a place of liberty;
* he brought me forth, even because he had a favour unto me.

21
The LORD rewarded me after my righteous dealing,
* according to the cleanness of my hands did he recompense me.

22 Because I have kept the ways of the LORD,
* and have not forsaken my God, as the wicked doth.

23 For I have an eye unto all his laws, * and will
not cast out his commandments from me.

24 I was also uncorrupt before him, * and eschewed
mine own wickedness.

25 Therefore the LORD rewarded
me after my righteous dealing, * and according unto the cleanness of my
hands in his eyesight.

26 With the holy thou shalt be holy, * and with a perfect
man thou shalt be perfect.

[18] They prevented me in
the day until
1928

27 With the clean thou shalt be clean, * and with the
froward thou shalt be froward.

28 For thou shalt save the people that are in adversity,
and shalt bring down the high looks of the proud.

29 Thou also shalt light my candle; * the LORD
my God shall make my darkness to be light.

30 For in thee I shall discomfit an host of men, *
and with the help of my God I shall leap over the wall.

31 The way of God is an undefiled way: * the word of
the LORD also is tried in the fire; he is the defender
of all them that put their trust in him.

32 For who is God, but the LORD?
* or who hath any strength, except our God?

33 It is God that girdeth me with strength of war,
* and maketh my way perfect.34 He maketh my feet like harts' feet, *
and setteth me up on high.

35 He teacheth mine hands to fight, * and mine arms
shall bend even a bow of steel.

36
Thou hast given me the defence of thy salvation; * thy right hand also
shall hold me up, and thy loving correction shall make me great.

37 Thou shalt make room enough under me for to go,
* that my footsteps shall not slide.

38 I will follow upon mine enemies, and overtake them;
* neither will I turn again till I have destroyed them.

39 I will smite them, that they shall not be able to
stand, * but fall under my feet.

40 Thou hast girded me with strength unto the battle;
* thou shalt throw down mine enemies under me.

41 Thou hast made mine enemies also to turn their backs
upon me, * and I shall destroy them that hate me.

42 They shall cry, but there shall be none to help
them; * yea, even unto the LORD shall they cry,
but he shall not hear them.

43 I will beat them as small as the dust before the
wind: * I will cast them out as the clay in the streets.

[26] ... with the froward
thou shalt learn frowardness until
1928

44 Thou shalt deliver me from the strivings of the people,
* and thou shalt make me the head of the nations; a people whom I have not
known shall serve me.

45 As soon as they hear of me, they shall obey me; *
the strangers shall feign obedience unto me.

46 The strangers shall fail, * and come trembling out
of their strongholds.

47 The LORD liveth; and blessed
be my strong helper, * and praised be the God of my salvation;

48 Even the God that seeth that I be avenged, * and subdueth
the people unto me.

49 It is he that delivereth me from my cruel enemies,
and setteth me up above mine adversaries: * thou shalt rid me from the wicked
man.

50 For this cause will I give thanks unto thee, O LORD,
among the Gentiles, * and sing praises unto thy Name.

51 Great prosperity giveth he unto his King, * and showeth
loving-kindness unto David his anointed, and unto his seed for evermore.

43. Thou shalt deliver me
from the strivings of the people : and thou shalt make me the head of the
heathen.

44. A people whom I have not known : shall serve me.

45. As soon as they hear of me, they shall obey me : 

but the strange children shall dissemble with me.

46. The strange children shall fail : and be afraid out of their prisons.
until
1928

The
Fourth Day.

Morning Prayer.

Psalm
19. \emph{Caeli enarrant.}

THE
heavens declare the glory of God; * and the firmament showeth his handy-work.

2 One day telleth another; * and one night certifieth
another.

3 There is neither speech nor language; * but their
voices are heard among them.

4 Their sound is gone out into all lands; * and their
words into the ends of the world.

5 In them hath he set a tabernacle for the sun; * which
cometh forth as a bridegroom out of his chamber, and rejoiceth as a giant
to run his course.

6 It goeth forth from the uttermost part of the heaven,
and runneth about unto the end of it again; * and there is nothing hid
from the heat thereof.

7 The law of the LORD is an
undefiled law, converting the soul; * the testimony of the LORD
is sure, and giveth wisdom unto the simple.

8 The statutes of the LORD are
right, and rejoice the heart; * the commandment of the LORD
is pure, and giveth light unto the eyes.

9 The fear of the LORD is clean,
and endureth for ever; * the judgments of the LORD
are true, and righteous altogether.

10 More to be desired are they than gold, yea, than
much fine gold; * sweeter also than honey, and the honeycomb.

11 Moreover, by them is thy servant taught; * and in
keeping of them there is great reward.

12 Who can tell how oft he offendeth? * O cleanse thou
me from my secret faults.

13 Keep thy servant also from presumptuous sins, lest
they get the dominion over me; * so shall I be undefiled, and innocent
from the great offence.

14 Let the words of my mouth, and the meditation of
my heart, be alway acceptable in thy sight, * O LORD,
my strength and my redeemer.

Last
verse was two until 1928:\\
14 Let the words
of my mouth, and the meditation of my heart, be alway acceptable in thy
sight,\\
15 O LORD, my strength
and my redeemer.

Psalm
20. \emph{Exaudiat te Dominus.}

THE
LORD hear thee in the day of trouble; * the Name
of the God of Jacob defend thee:

2 Send thee help from the sanctuary, * and strengthen
thee out of Sion:

3 Remember all thy offerings, * and accept thy burnt-sacrifice:

4 Grant thee thy heart’s desire, * and fulfil
all thy mind.

5 We will rejoice in thy salvation, and triumph in
the Name of the LORD our God: * the LORD
perform all thy petitions.

6 Now know I that the LORD helpeth
his anointed, and will hear him from his holy heaven, * even with the
wholesome strength of his right hand.

7 Some put their trust in chariots, and some in horses;
* but we will remember the Name of the LORD our
God.

8 They are brought down and fallen; * but we are risen
and stand upright.

9 Save,LORD; and hear us, O
King of heaven, * when we call upon thee.

Psalm
21. \emph{Domine, in virtute tua.}

THE
King shall rejoice in thy strength, O LORD; * exceeding
glad shall he be of thy salvation.

2 Thou hast given him his heart’s desire, * and
hast not denied him the request of his lips.

3 For thou shalt meet1 him with the blessings
of goodness, * and shalt set a crown of pure gold upon his head.

4 He asked life of thee; and thou gavest him a long life,
* even for ever and ever.

5 His honour is great in thy salvation; * glory and great
worship shalt thou lay upon him.

6 For thou shalt give him everlasting felicity, * and
make him glad with the joy of thy countenance.

7 And why? because the King putteth his trust in the
LORD; * and in the mercy of the Most Highest he shall
not miscarry. 1 prevent
until 1928

8 All thine enemies shall feel thine2 hand;
* thy right hand shall find out them that hate thee.

9 Thou shalt make them like a fiery oven in time of
thy wrath: * the LORD shall destroy them in his
displeasure, and the fire shall consume them.

10 Their fruit shalt thou root out of the earth, *
and their seed from among the children of men.

11 For they intended mischief against thee, * and imagined
such a device as they are not able to perform.

12 Therefore shalt thou put them to flight, * and the
strings of thy bow shalt thou make ready against the face of them.

13 Be thou exalted, LORD, in
thine own strength; * so will we sing, and praise thy power.

Evening
Prayer.

Psalm
22. \emph{Deus, Deus meus.}

MY
God, my God, look upon me; why hast thou forsaken me? * and art so far
from my health, and from the words of my complaint?

2 O my God, I cry in the day-time, but thou hearest
not; * and in the night season also I take no rest.

3 And thou continuest holy, * O thou Worship of Israel.

4 Our fathers hoped in thee; * they trusted in thee,
and thou didst deliver them.

5 They called upon thee, and were holpen; * they put
their trust in thee, and were not confounded.

6 But as for me, I am a worm, and no man; * a very
scorn of men, and the outcast of the people.

7 All they that see me laugh me to scorn; * they shoot
out their lips, and shake their heads, saying,

8 He trusted in the LORD, that
he would deliver him; * let him deliver him, if he will have him.

9 But thou art he that took me out of my mother's womb;
* thou wast my hope, when I hanged yet upon my mother’s breasts.

10 I have been left unto thee ever since I was born;
* thou art my God even from my mother's womb.

11 O go not from me; for trouble is hard at hand, *
and there is none to help me.

2 thy
until 1892

12 Many oxen are come about me; * fat bulls of Bashan1 close
me in on every side.

13 They gape upon me with their mouths, * as it were
a ramping and a roaring lion.

14 I am poured out like water, and all my bones are out
of joint; * my heart also in the midst of my body is even like melting wax.

15 My strength is dried up like a potsherd, and my tongue
cleaveth to my gums, * and thou bringest me into the dust of death. \\
1 Basan
until 1928

16 For many dogs are come about me, * and the council2
of the wicked layeth siege against me.

17 They pierced my hands and my feet: I may tell all
my bones: * they stand staring and looking upon me.

18 They part my garments among them, * and cast lots
upon my vesture.

19 But be not thou far from me, O LORD;
* thou art my succour, haste thee to help me.

20
Deliver my soul from the sword, * my darling from the power of the dog.

21 Save me from the lion's mouth; * thou hast heard
me also from among the horns of the unicorns.

22 I will declare thy Name unto my brethren; * in the
midst of the congregation will I praise thee.

23 O praise the LORD, ye that
fear him: * magnify him, all ye of the seed of Jacob; and fear him, all
ye seed of Israel.

24 For he hath not despised nor abhorred the low estate
of the poor; * he hath not hid his face from him; but when he called unto
him he heard him.

25 My praise is of thee in the great congregation;
* my vows will I perform in the sight of them that fear him.

26 The poor shall eat, and be satisfied; they that
seek after the LORD shall praise him: * your heart
shall live for ever.

27 All the ends of the world shall remember themselves,
and be turned unto the LORD; * and all the kindreds
of the nations shall worship before him.

2 counsel 1845-1892

28 For the kingdom is the LORD’S,
* and he is the Governor among the nations3.

29 All such as be fat upon earth * have eaten, and worshipped.

30 All they that go down into the dust shall kneel before
him; * and no man hath quickened his own soul.

31 My seed shall serve him: * they shall be counted unto
the Lord for a generation.

32 They shall come, and shall declare his righteousness
* unto a people that shall be born, whom the Lord hath made.

3 people
until 1928

32 ... and the heavens shall declare
until 1928

Psalm
23. \emph{Dominus regit me.}

THE
LORD is my shepherd; * therefore can I lack nothing.

2 He shall feed me in a green pasture, * and lead me
forth beside the waters of comfort.

3 He shall convert my soul, * and bring me forth in
the paths of righteousness for his Name’s sake.

4 Yea, though I walk through the valley of the shadow
of death, I will fear no evil; * for thou art with me; thy rod and thy
staff comfort me.

5 Thou shalt prepare a table before me in the presence
of them that trouble me; * thou hast anointed my head with oil, and my
cup shall be full.

6 Surely1 thy loving-kindness and mercy
shall follow me all the days of my life; * and I will dwell in the house
of the LORD for ever. 

5 ... before me against them that
trouble me :
until 1928

1 But
until 1928

The Fifth Day.

Morning
Prayer.

Psalm
24. \emph{Domini est terra.}

THE
earth is the LORD'S, and all that therein is; *
the compass of the world, and they that dwell therein.

2 For he hath founded it upon the seas, * and stablished
it upon the floods.

3 Who shall ascend into the hill of the LORD?
* or who shall rise up in his holy place?

4 Even he that hath clean hands, and a pure heart;
* and that hath not lift up his mind unto vanity, nor sworn to deceive
his neighbour.

5 He shall receive the blessing from the LORD,
* and righteousness from the God of his salvation.

6 This is the generation of them that seek him; * even
of them that seek thy face, O God of Jacob.

6 ... O Jacob
until 1928

7 Lift up your heads, O ye gates; and be ye lift up,
ye everlasting doors; * and the King of glory shall come in.

8 Who is this1 King of glory? * It is the
LORD strong and mighty, even the LORD
mighty in battle.

9 Lift up your heads, O ye gates; and be ye lift up,
ye everlasting doors; * and the King of glory shall come in.

10 Who is this1 King of glory? * Even the
LORD of hosts, he is the King of glory.

1 the
until 1892

Psalm 25. \emph{Ad te, Domine, levavi.}

UNTO
thee, O LORD, will I lift up my soul; my God, I
have put my trust in thee: * O let me not be confounded, neither let
mine enemies triumph over me.

2 For all they that hope in thee shall not be ashamed;
* but such as transgress without a cause shall be put to confusion.

3 Show me thy ways, O LORD,
* and teach me thy paths.

4 Lead me forth in thy truth, and learn me: * for thou
art the God of my salvation; in thee hath been my hope all the day long.

5 Call to remembrance, O LORD,
thy tender mercies, * and thy loving-kindnesses, which have been ever
of old.

6 O remember not the sins and offences of my youth;
* but according to thy mercy think thou upon me, O LORD,
for thy goodness.

7 Gracious and righteous is the LORD;
* therefore, will he teach sinners in the way.

8 Them that are meek shall he guide in judgment; *
and such as are gentle, them shall he learn his way.

9 All the paths of the LORD
are mercy and truth, * unto such as keep his covenant and his testimonies.

10 For thy Name's sake, O LORD,
* be merciful unto my sin; for it is great.

11 What man is he that feareth the LORD?
* him shall he teach in the way that he shall choose.

12 His soul shall dwell at ease, * and his seed shall
inherit the land.

13 The secret of the LORD is
among them that fear him; * and he will show them his covenant.

14 Mine eyes are ever looking unto the LORD;
* for he shall pluck my feet out of the net.

15 Turn thee unto me, and have mercy upon me; * for
I am desolate, and in misery.

16 The sorrows of my heart are enlarged: * O bring
thou me out of my troubles.

17 Look upon my adversity and misery, * and forgive
me all my sin.

18 Consider mine enemies, how many they are; * and
they bear a tyrannous hate against me.

19 O keep my soul, and deliver me: * let me not be
confounded, for I have put my trust in thee.

20 Let perfectness and righteous dealing wait upon
me; * for my hope hath been in thee.

21 Deliver Israel, O God, * out of all his troubles.

Psalm
26. \emph{Judica me, Domine.}

BE
thou my Judge, O LORD, for I have walked innocently:
* my trust hath been also in the LORD, therefore
shall I not fall.

2 Examine me, O LORD, and prove
me; * try out my reins and my heart.

3 For thy loving-kindness is ever before mine eyes;
* and I will walk in thy truth.

4 I have not dwelt with vain persons; * neither will
I have fellowship with the deceitful.

5 I have hated the congregation of the wicked; * and
will not sit among the ungodly.

6 I will wash my hands in innocency, O LORD;
* and so will I go to thine altar;

7 That I may show the voice of thanksgiving, * and
tell of all thy wondrous works.

8 LORD, I have loved the habitation
of thy house, * and the place where thine honour dwelleth.

9 O shut not up my soul with the sinners, * nor my
life with the blood-thirsty;

10 In whose hands is wickedness, * and their right
hand is full of gifts.

11 But as for me, I will walk innocently: * O deliver
me, and be merciful unto me.

12 My foot standeth right: * I will praise the LORD
in the congregations.

Evening
Prayer.

Psalm
27. \emph{Dominus illuminatio.}

THE
LORD is my light and my salvation; whom then shall
I fear? * the LORD is the strength of my life;
of whom then shall I be afraid?

2 When the wicked, even mine enemies and my foes, came
upon me to eat up my flesh, * they stumbled and fell.

3 Though an host of men were laid against me, yet shall
not my heart be afraid; * and though there rose up war against me, yet
will I put my trust in him.

4 One thing have I desired of the LORD,
which I will require; * even that I may dwell in the house of the LORD
all the days of my life, to behold the fair beauty of the LORD,
and to visit his temple.

5 For in the time of trouble he shall hide me in his
tabernacle; * yea, in the secret place of his dwelling shall he hide me,
and set me up upon a rock of stone.

6 And now shall he lift up mine head * above mine enemies
round about me.

7 Therefore will I offer in his dwelling an oblation,
with great gladness: * I will sing and speak praises unto the LORD.

8 Hearken unto my voice, O LORD,
when I cry unto thee; * have mercy upon me, and hear me.

9 My heart hath talked of thee, Seek ye my face: *
Thy face, LORD, will I seek.

10 O hide not thou thy face from me, * nor cast thy
servant away in displeasure.

11 Thou hast been my succour; * leave me not, neither
forsake me, O God of my salvation.

12 When my father and my mother forsake me, * the LORD
taketh me up.

13 Teach me thy way, O LORD,
* and lead me in the right way, because of mine enemies.

14 Deliver me not over into the will of mine adversaries:
* for there are false witnesses risen up against me, and such as speak
wrong.

15 I should utterly have fainted, * but that I believe
verily to see the goodness of the LORD in the land
of the living.

16 O tarry thou the LORD’S
leisure; * be strong, and he shall comfort thine heart; and put thou thy
trust in the LORD.

Psalm
28. \emph{Ad te, Domine.}

UNTO
thee will I cry, O LORD, my strength: * think no
scorn of me; lest, if thou make as though thou hearest not, I become like
them that go down into the pit.

2 Hear the voice of my humble petitions, when I cry
unto thee; * when I hold up my hands towards the mercyseat of thy holy
temple.

3 O pluck me not away, neither destroy me with the
ungodly and wicked doers, * which speak friendly to their neighbours,
but imagine mischief in their hearts.

4 Reward them according to their deeds, * and according
to the wickedness of their own inventions.

5 Recompense them after the work of their hands; *
pay them that they have deserved.

6 For they regard not in their mind the works of the
LORD, nor the operation of his hands; * therefore
shall he break them down, and not build them up.

7 Praised be the LORD; * for
he hath heard the voice of my humble petitions.

8 The LORD is my strength, and
my shield; my heart hath trusted in him, and I am helped; * therefore
my heart danceth for joy, and in my song will I praise him.

9 The LORD is my strength, *
and he is the wholesome defence of his anointed.

10 O save thy people, and give thy blessing unto thine
inheritance: * feed them, and set them up for ever.

Psalm
29. \emph{Afferte Domino.}

ASCRIBE
unto the LORD, O ye mighty, * ascribe unto the
LORD worship and strength.

2 Ascribe unto the LORD the
honour due unto his Name; * worship the LORD with
holy worship.

3 The voice of the LORD is upon
the waters; * it is the glorious God that maketh the thunder.

4 It is the LORD that ruleth
the sea; the voice of the LORD is mighty in operation;
* the voice of the LORDd is a glorious voice.

5 The voice of the LORD breaketh
the cedar-trees; * yea, the LORD breaketh the cedars
of Lebanon1.

6 He maketh them also to skip like a calf; * Lebanon1
also, and Sirion, like a young unicorn.

BRING unto the LORD,
O ye mighty, bring young rams unto the LORD : ascribe
unto the LORD worship and strength.

2. Give the LORD the honour
due unto his Name : worship the LORD with holy
worship.

3. It is the LORD that commandeth
the waters : it is the glorious God that maketh the thunder. until
1928.

1 Libanus until
1928.

7 The voice of the LORD divideth
the flames of fire; the voice of the LORD shaketh
the wilderness; * yea, the LORD shaketh the wilderness
of Kadesh2.

8 The voice of the LORD maketh
the hinds to bring forth young, and strippeth bare the forests: * in his
temple doth every thing speak of his honour.

9 The LORD sitteth above the
water-flood, * and the LORD remaineth a King for
ever.

10 The LORD shall give strength
unto his people; * the LORD shall give his people
the blessing of peace.

The
Sixth Day.

Morning
Prayer.

Psalm
30. \emph{Exaltabo te, Domine.}

I
WILL magnify thee, O LORD; for thou hast set me
up, * and not made my foes to triumph over me.

2 O LORD my God, I cried unto
thee; * and thou hast healed me.

2 Cades until
1928.

8 ... young, and discovereth the thick bushes : in
his temple until
1928.

3 Thou, LORD, hast brought my
soul out of hell: * thou hast kept my life, that I should not go down
into the pit.

4 Sing praises unto the LORD,
O ye saints of his; * and give thanks unto him, for a remembrance of his
holiness.

5 For his wrath endureth but the twinkling of an eye,
and in his pleasure is life; * heaviness may endure for a night, but joy
cometh in the morning.

6 And in my prosperity I said, I shall never be removed:
* thou, LORD, of thy goodness, hast made my hill
so strong.

7 Thou didst turn thy face from me, * and I was troubled.

8 Then cried I unto thee, O LORD;
* and gat me to my LORD right humbly.

9 What profit is there in my blood, * when I go down
into the pit?

10 Shall the dust give thanks unto thee? * or shall
it declare thy truth?

11 Hear, O LORD, and have mercy
upon me; * LORD, be thou my helper.

12 Thou hast turned my heaviness into joy; * thou hast
put off my sackcloth, and girded me with gladness:

13 Therefore shall every good man sing of thy praise
without ceasing. * O my God, I will give thanks unto thee for ever.

Psalm
31. \emph{In te, Domine, speravi.}

IN
thee, O LORD, have I put my trust; let me never
be put to confusion; * deliver me in thy righteousness.

2 Bow down thine ear to me; * make haste to deliver
me.

3 And be thou my strong rock, and house of defence,
* that thou mayest save me.

4 For thou art my strong rock, and my castle: * be
thou also my guide, and lead me for thy Name's sake.

5 Draw me out of the net that they have laid privily
for me; * for thou art my strength.

6 Into thy hands I commend my spirit; * for thou hast
redeemed me, O LORD, thou God of truth.

3 ... my life from them that
go down into the pit. until
1928.

7 I have hated them that hold of lying1 vanities, * and
my trust hath been in the LORD.

8 I will be glad and rejoice in thy mercy; * for thou
hast considered my trouble, and hast known my soul in adversities.

9 Thou hast not shut me up into the hand of the enemy;
* but hast set my feet in a large room.

10 Have mercy upon me, O LORD,
for I am in trouble, * and mine eye is consumed for very heaviness; yea,
my soul and my body.

11 For my life is waxen old with heaviness, * and my
years with mourning.

12 My strength faileth me, because of mine iniquity,
* and my bones are consumed.

1 superstitious until
1928

13 I became a reproach2 among all mine enemies,
but especially among my neighbours; * and they of mine acquaintance were
afraid of me; and they that did see me without, conveyed themselves from
me.

14 I am clean forgotten as a dead man out of mind; *
I am become like a broken vessel.

15 For I have heard the blasphemy of the multitude, and
fear is on every side; * while they conspire together against me, and take
their counsel to take away my life.

16 But my hope hath been in thee, O LORD;
* I have said, Thou art my God.

2 reproof until
1928

17 My times are3 in thy hand; deliver me from
the hand of mine enemies, * and from them that persecute me.

18 Show thy servant the light of thy countenance, * and
save me for thy mercy's sake.

19 Let me not be confounded, O LORD,
for I have called upon thee; * let the ungodly be put to confusion, and
be put to silence in the grave.

20 Let the lying lips be put to silence, * which cruelly,
disdainfully, and despitefully speak against the righteous.

21 O how plentiful is thy goodness, which thou hast laid
up for them that fear thee, * and that thou hast prepared for them that
put their trust in thee, even before the sons of men!

22 Thou shalt hide them in the covert of thine own presence
from the plottings of men: * thou shalt keep them secretly in thy tabernacle
from the strife of tongues.

23 Thanks be to the LORD; * for
he hath showed me marvellous great kindness in a strong city.

3 My time is until
1928

22 Thou shalt hide them privily
by thine own presence from the provoking of all men : until
1928

24 But in my haste I said, * I am cast out of the sight
of thine eyes.

25 Nevertheless, thou heardest the voice of my prayer,
* when I cried unto thee.

26 O love the LORD, all ye his
saints; * for the LORD preserveth them that are
faithful, and plenteously rewardeth the proud doer.

27 Be strong, and he shall establish your heart, *
all ye that put your trust in the LORD.

Evening
Prayer.

Psalm
32. \emph{Beati quorum.}

BLESSED
is he whose unrighteousness is forgiven, * and whose sin is covered.

2 Blessed is the man unto whom the LORD
imputeth no sin, * and in whose spirit there is no guile.

3 For whilst I held my tongue, * my bones consumed
away through my daily complaining.

4 For thy hand was heavy upon me day and night, * and
my moisture was like the drought in summer.

5 I acknowledged my sin unto thee; * and mine unrighteousness
have I not hid.

6 I said, I will confess my sins unto the LORD;
* and so thou forgavest the wickedness of my sin.

24 And when I made haste,
I said : until
1928

7 For this shall every one that is godly make his prayer
unto thee, in a time when thou mayest be found; * surely1 the
great water-floods shall not come nigh him.

8 Thou art a place to hide me in; thou shalt preserve
me from trouble; * thou shalt compass me about with songs of deliverance.

9 I will inform thee, and teach thee in the way wherein
thou shalt go; * and I will guide thee with mine eye.

1 but until
1928

10 Be ye not like to horse and mule, which have no
understanding; * whose mouths must be held with bit and bridle, else they
will not obey thee.

11 Great plagues remain for the ungodly; * but whoso
putteth his trust in the LORD, mercy embraceth
him on every side.

12 Be glad, O ye righteous, and rejoice in the LORD;
* and be joyful, all ye that are true of heart.

Psalm
33. \emph{Exultate, justi.}

REJOICE
in the LORD, O ye righteous; * for it becometh
well the just to be thankful.

2 Praise the LORD with harp;
* sing praises unto him with the lute, and instrument of ten strings.

3 Sing unto the LORD a new song;
* sing praises lustily unto him with a good courage.

4 For the word of the LORD is
true; * and all his works are faithful.

5 He loveth righteousness and judgment; * the earth
is full of the goodness of the LORD.

10 ... bit and bridle, lest they fall upon thee. until
1928

6 By the word of the LORD were
the heavens made; * and all the host1 of them by the breath of
his mouth.

7 He gathereth the waters of the sea together, as it
were upon an heap; * and layeth up the deep, as in a treasure-house.

8 Let all the earth fear the LORD:
* stand in awe of him, all ye that dwell in the world.

9 For he spake, and it was done; * he commanded, and
it stood fast.

10 The LORD bringeth the counsel
of the heathen to nought, * and maketh the devices of the people to be of
none effect, and casteth out the counsels of princes.

11 The counsel of the LORD shall
endure for ever, * and the thoughts of his heart from generation to generation.

12 Blessed are the people whose God is the Lord JEHOVAH;
* and blessed are the folk that he hath chosen to him, to be his inheritance.

13 The LORD looketh down from
heaven, and beholdeth all the children of men; * from the habitation of
his dwelling, he considereth all them that dwell on the earth.

14 He fashioneth all the hearts of them, * and understandeth
all their works.

15 There is no king that can be saved by the multitude
of an host; * neither is any mighty man delivered by much strength.

16 A horse is counted but a vain thing to save a man;
* neither shall he deliver any man by his great strength.

17 Behold, the eye of the LORD
is upon them that fear him, * and upon them that put their trust in his
mercy;

18 To deliver their soul from death, * and to feed them
in the time of dearth.

19 Our soul hath patiently tarried for the LORD;
* for he is our help and our shield.

20 For our heart shall rejoice in him; * because we have
hoped in his holy Name.

21 Let thy merciful kindness, O LORD,
be upon us, * like as we do put our trust in thee.

1 hosts until
1928

Psalm
34. \emph{Benedicam Dominum.}

I
WILL alway give thanks unto the LORD; * his praise
shall ever be in my mouth.

2 My soul shall make her boast in the LORD;
* the humble shall hear thereof, and be glad.

3 O praise the LORD with me,
* and let us magnify his Name together.

4 I sought the LORD, and he
heard me; * yea, he delivered me out of all my fear.

5 They had an eye unto him, and were lightened; * and
their faces were not ashamed.

6 Lo, the poor crieth, and the LORD
heareth him; * yea, and saveth him out of all his troubles.

7 The angel of the LORD tarrieth
round about them that fear him, * and delivereth them.

8 O taste, and see, how gracious the LORD
is: * blessed is the man that trusteth in him.

9 O fear the LORD, ye that are
his saints; * for they that fear him lack nothing.

10 The lions do lack, and suffer hunger; * but they
who seek the LORD shall want no manner of thing
that is good.

11 Come, ye children, and hearken unto me; * I will
teach you the fear of the LORD.

12 What man is he that lusteth to live, * and would
fain see good days?

13 Keep thy tongue from evil, * and thy lips, that
they speak no guile.

14 Eschew evil, and do good; * seek peace, and ensue
it.

15 The eyes of the LORD are
over the righteous, * and his ears are open unto their prayers.

16 The countenance of the LORD
is against them that do evil, * to root out the remembrance of them from
the earth.

17 The righteous cry, and the LORD
heareth them, * and delivereth them out of all their troubles.

18 The LORD is nigh unto them
that are of a contrite heart, * and will save such as be of an humble
spirit.

19 Great are the troubles of the righteous; * but the
LORD delivereth him out of all.

20 He keepeth all his bones, * so that not one of them
is broken.

21 But misfortune shall slay the ungodly; * and they
that hate the righteous shall be desolate.

22 The LORD delivereth the souls
of his servants; * and all they that put their trust in him shall not
be destitute.
The
Seventh Day.

Morning
Prayer.

\emph{Benedicam Domino}until
1892

Psalm
35. \emph{Judica, Domine.}

PLEAD
thou my cause, O LORD, with them that strive with
me, * and fight thou against them that fight against me.

2 Lay hand upon the shield and buckler, * and stand
up to help me. \emph{Judica me, Domine} 1822-1892

3 Bring forth the spear, and stop the way against them
that pursue1 me: * say unto my soul, I am thy salvation.

4 Let them be confounded, and put to shame, that seek
after my soul; * let them be turned back, and brought to confusion, that
imagine mischief for me.

5 Let them be as the dust before the wind, * and the
angel of the LORD scattering them.

6 Let their way be dark and slippery, * and let the
angel of the LORD pursue them.

7 For they have privily laid their net to destroy me
without a cause; * yea, even without a cause have they made a pit for
my soul.

8 Let a sudden destruction come upon him unawares,
and his net that he hath laid privily catch himself; * that he may fall
into his own mischief.

1 pursecute until
1928

9 And my soul shall be joyful in the LORD;
* it shall rejoice in his salvation.

10 All my bones shall say, LORD,
who is like unto thee, who deliverest the poor from him that is too strong
for him; * yea, the poor, and him that is in misery, from him that spoileth
him?

11 False witnesses did rise up: * they laid to my charge
things that I knew not.

12 They rewarded me evil for good, * to the great discomfort
of my soul.

13 Nevertheless, when they were sick, I put on sack-cloth,
and humbled my soul with fasting; * and my prayer shall turn into mine
own bosom.

14 I behaved myself as though it had been my friend
or my brother; * I went heavily, as one that mourneth for his mother.

15 But in mine adversity they rejoiced, and gathered
themselves together; * yea, the very abjects came together against me
unawares, making mouths at me, and ceased not.

16 With the flatterers were busy mockers, * who gnashed
upon me with their teeth.

17 Lord, how long wilt thou look upon this? * O deliver
my soul from the calamities which they bring on me, and my darling from
the lions.

18 So will I give thee thanks in the great congregation;
* I will praise thee among much people.

19 O let not them that are mine enemies triumph over
me ungodly; * neither let them wink with their eyes, that hate me without
a cause.

20 And why? their communing is not for peace; * but
they imagine deceitful words against them that are quiet in the land.

21 They gaped upon me with their mouths, and said,
* Fie on thee! fie on thee! we saw it with our eyes.

22 This thou hast seen, O LORD;
* hold not thy tongue then; go not far from me, O Lord.

23 Awake, and stand up to judge my quarrel; * avenge
thou my cause, my God and my Lord.

24 Judge me, O LORD my God,
according to thy righteousness; * and let them not triumph over me.

25 Let them not say in their hearts, There! there!
so would we have it; * neither let them say, We have devoured him.

26 Let them be put to confusion and shame together,
that rejoice at my trouble; * let them be clothed with rebuke and dishonour,
that boast themselves against me.

27 Let them be glad and rejoice, that favour my righteous
dealing; * yea, let them say alway, Blessed be the LORD,
who hath pleasure in the prosperity of his servant.

28 And as for my tongue, it shall be talking of thy
righteousness, * and of thy praise, all the day long.

Psalm
36. \emph{Dixit injustus.}

MY
heart showeth me the wickedness of the ungodly, * that there is no fear
of God before his eyes.

2 For he flattereth himself in his own sight, * until
his abominable sin be found out.

3 The words of his mouth are unrighteous and full of
deceit: * he hath left off to behave himself wisely, and to do good.

4 He imagineth mischief upon his bed, and hath set
himself in no good way; * neither doth he abhor any thing that is evil.

5 Thy mercy, O LORD, reacheth
unto the heavens, * and thy faithfulness unto the clouds.

6 Thy righteousness standeth like the strong mountains:
* thy judgments are like the great deep.

7 Thou, LORD, shalt save both
man and beast: how excellent is thy mercy, O God! * and the children of
men shall put their trust under the shadow of thy wings.

8 They shall be satisfied with the plenteousness of
thy house; * and thou shalt give them drink of thy pleasures, as out of
the river.

9 For with thee is the well of life; * and in thy light
shall we see light.

10 O continue forth thy loving-kindness unto them that
know thee, * and thy righteousness unto them that are true of heart.

11 O let not the foot of pride come against me; * and
let not the hand of the ungodly cast me down.

12 There are they fallen, all that work wickedness;
* they are cast down, and shall not be able to stand.

Evening
Prayer.

Psalm
37. \emph{Noli aemulari.}

FRET
not thyself because of the ungodly; * neither be thou envious against
the evil doers.

2 For they shall soon be cut down like the grass, *
and be withered even as the green herb.

3 Put thou thy trust in the LORD,
and be doing good; * dwell in the land, and verily thou shalt be fed.

4 Delight thou in the LORD,
* and he shall give thee thy heart's desire.

5 Commit thy way unto the LORD,
and put thy trust in him, * and he shall bring it to pass.

6 He shall make thy righteousness as clear as the light,
* and thy just dealing as the noon-day.

7 Hold thee still in the LORD,
and abide patiently upon him: * but grieve not thyself at him whose way
doth prosper, against the man that doeth after evil counsels.

8 Leave off from wrath, and let go displeasure: * fret
not thyself, else shalt thou be moved to do evil.

9 Wicked doers shall be rooted out; * and they that
patiently abide the LORD, those shall inherit the
land.

10 Yet a little while, and the ungodly shall be clean
gone: * thou shalt look after his place, and he shall be away.

11 But the meek-spirited shall possess the earth, *
and shall be refreshed in the multitude of peace.

12 The ungodly seeketh counsel against the just, *
and gnasheth upon him with his teeth.

13 The Lord shall laugh him to scorn; * for he hath
seen that his day is coming. 9 And, my soul, be joyful
... until
1928

14 The ungodly have drawn out the sword, and have bent
their bow, * to cast down the poor and needy, and to slay such as be upright
in their ways.

15 Their sword shall go through their own heart, * and
their bow shall be broken.

16 A small thing that the righteous hath, * is better
than great riches of the ungodly.

17 For the arms of the ungodly shall be broken, * and
the LORD upholdeth the righteous.

18 The LORD knoweth the days of
the godly; * and their inheritance shall endure for ever.

19 They shall not be confounded in the perilous time;
* and in the days of dearth they shall have enough.

20 As for the ungodly, they shall perish, and the enemies
of the LORD shall consume as the fat of lambs: *
yea, even as the smoke shall they consume away.

21 The ungodly borroweth, and payeth not again; * but
the righteous is merciful and liberal.

22 Such as are blessed of God, shall possess the land;
* and they that are cursed of him, shall be rooted out.

23 The LORD ordereth a good man's
going, * and maketh his way acceptable to himself.

24 Though he fall, he shall not be cast away; * for the
LORD upholdeth him with his hand.

25 I have been young, and now am old; * and yet saw I
never the righteous forsaken, nor his seed begging their bread.

26 The righteous is ever merciful, and lendeth; * and
his seed is blessed.

27 Flee from evil, and do the thing that is good; * and
dwell for evermore.

28 For the LORD loveth the thing
that is right; * he forsaketh not his that be godly, but they are preserved
for ever.

29 The unrighteous shall be punished; * as for the seed
of the ungodly, it shall be rooted out.

30 The righteous shall inherit the land, * and dwell
therein for ever.

31 The mouth of the righteous is exercised in wisdom,
* and his tongue will be talking of judgment.

32 The law of his God is in his heart, * and his goings
shall not slide.

14 ... slay such as are of a right conversation. until
1928

33 The ungodly watcheth1 the righteous,
* and seeketh occasion to slay him.

34 The LORD will not leave him
in his hand, * nor condemn him when he is judged.

35 Hope thou in the LORD, and
keep his way, and he shall promote thee, that thou shalt possess the land:
* when the ungodly shall perish, thou shalt see it.

36 I myself have seen the ungodly in great power, *
and flourishing like a green bay-tree.

37 I went by, and lo, he was gone: * I sought him,
but his place could no where be found.

38 Keep innocency, and take heed unto the thing that
is right; * for that shall bring a man peace at the last.

39 As for the transgressors, they shall perish together;
* and the end of the ungodly is, they shall be rooted out at the last.

40 But the salvation of the righteous cometh of the
LORD; * who is also their strength in the time
of trouble.

41 And the LORD shall stand
by them, and save them: * he shall deliver them from the ungodly, and
shall save them, because they put their trust in him.

The Eighth
Day

Morning
Prayer.

Psalm
38. \emph{Domine, ne in furore.}

PUT
me not to rebuke, O LORD, in thine anger; * neither
chasten me in thy heavy displeasure:

2 For thine arrows stick fast in me, * and thy hand
presseth me sore.

3 There is no health in my flesh, because of thy displeasure;
* neither is there any rest in my bones, by reason of my sin.

4 For my wickednesses are gone over my head, * and
are like a sore burden, too heavy for me to bear.

5 My wounds stink, and are corrupt, * through my foolishness.

6 I am brought into so great trouble and misery, *
that I go mourning all the day long.

7 For my loins are filled with a sore disease, * and
there is no whole part in my body.

8 I am feeble and sore smitten; * I have roared for
the very disquietness of my heart.

9 Lord, thou knowest all my desire; * and my groaning
is not hid from thee. 1 seeth
until
1928

10 My heart panteth, my strength hath failed me, *
and the light1 of mine eyes is gone from me.

11 My lovers and my neighbours did stand looking upon
my trouble, * and my kinsmen stood afar off.

12 They also that sought after my life laid snares
for me; * and they that went about to do me evil talked of wickedness,
and imagined deceit all the day long.

13 As for me, I was like a deaf man, and heard not;
* and as one that is dumb, who doth not open his mouth.

14 I became even as a man that heareth not, * and in
whose mouth are no reproofs.

15 For in thee, O LORD, have
I put my trust; * thou shalt answer for me, O Lord my God.

16 I have required that they, even mine enemies, should
not triumph over me; * for when my foot slipt, they rejoiced greatly against
me.

17 And I truly am set in the plague, * and my heaviness
is ever in my sight.

18 For I will confess my wickedness, * and be sorry
for my sin.

19 But mine enemies live, and are mighty; * and they
that hate me wrongfully are many in number.

20 They also that reward evil for good are against
me; * because I follow the thing that good is.

21 Forsake me not, O LORD my
God; * be not thou far from me.

22 Haste thee to help me, * O Lord God of my salvation.

Psalm
39. \emph{Dixi, Custodiam.}

I
SAID, I will take heed to my ways, * that I offend not in my tongue.

2 I will keep my mouth as it were with a bridle, *
while the ungodly is in my sight.

3 I held my tongue, and spake nothing: * I kept silence,
yea, even from good words; but it was pain and grief to me.

4 My heart was hot within me: and while I was thus
musing the fire kindled, * and at the last I spake with my tongue: \\
1 sight until
1892.

5 LORD, let me know mine1
end, and the number of my days; * that I may be certified how long I have
to live.

6 Behold, thou hast made my days as it were a span
long, and mine age is even as nothing in respect of thee; * and verily
every man living is altogether vanity.

7 For man walketh in a vain shadow, and disquieteth
himself in vain; * he heapeth up riches, and cannot tell who shall gather
them.

8 And now, Lord, what is my hope? * truly my hope is
even in thee.

9 Deliver me from all mine offences; * and make me
not a rebuke unto the foolish.

10 I became dumb, and opened not my mouth; * for it
was thy doing.

11 Take thy plague away from me: * I am even consumed
by the means of thy heavy hand.

12 When thou with rebukes dost chasten man for sin,
thou makest his beauty to consume away, like as it were a moth fretting
a garment: * every man therefore is but vanity.

13 Hear my prayer, O LORD, and
with thine ears consider my calling; * hold not thy peace at my tears;

14 For I am a stranger with thee, and a sojourner,
* as all my fathers were.

15 O spare me a little, that I may recover my strength,
* before I go hence, and be no more seen.

Psalm
40. \emph{Expectans expectavi.}

I
WAITED patiently for the LORD, * and he inclined
unto me, and heard my calling.

2 He brought me also out of the horrible pit, out of
the mire and clay, * and set my feet upon the rock, and ordered my goings.

3 And he hath put a new song in my mouth, * even a
thanksgiving unto our God.

4 Many shall see it, and fear, * and shall put their
trust in the LORD.

5 Blessed is the man that hath set his hope in the
LORD, * and turned not unto the proud, and to such
as go about with lies.

6 O LORD my God, great are the
wondrous works which thou hast done, like as be also thy thoughts, which
are to usward; * and yet there is no man that ordereth them unto thee.

7 If I should declare them, and speak of them, * they
should be more than I am able to express. 1 my until
1892.

8 Sacrifice and offering1 thou wouldest
not, * but mine ears hast thou opened.

9 Burnt-offering and sacrifice for sin hast thou not
required: * then said I, Lo, I come;

10 In the volume of the book it is written of me, that
I should fulfil thy will, O my God: * I am content to do it; yea, thy
law is within my heart.

11 I have declared thy righteousness in the great congregation:
* lo, I will not refrain my lips, O LORD, and that
thou knowest.

12 I have not hid thy righteousness within my heart;
* my talk hath been of thy truth, and of thy salvation.

13 I have not kept back thy loving mercy and truth
* from the great congregation.

14 Withdraw not thou thy mercy from me, O LORD;
* let thy loving-kindness and thy truth alway preserve me.

15 For innumerable troubles are come about me; my sins
have taken such hold upon me, that I am not able to look up; * yea, they
are more in number than the hairs of my head, and my heart hath failed
me.

16 O LORD, let it be thy pleasure
to deliver me; * make haste, O LORD, to help me.

17
Let them be ashamed, and confounded together, that seek after my soul
to destroy it; * let them be driven backward, and put to rebuke, that
wish me evil.

18 Let them be desolate, and rewarded with shame, *
that say unto me, Fie upon thee! fie upon thee!

19 Let all those that seek thee, be joyful and glad
in thee; * and let such as love thy salvation, say alway, The LORD
be praised.

20 As for me, I am poor and needy; * but the Lord careth
for me.

21 Thou art my helper and redeemer; * make no long
tarrying, O my God.

Evening
Prayer.

Psalm
41. \emph{Beatus qui intelligit.}

BLESSED
is he that considereth the poor and needy; * the LORD
shall deliver him in the time of trouble.

2 The LORD preserve him, and
keep him alive, that he may be blessed upon earth; * and deliver not thou
him into the will of his enemies.

3 The LORD comfort him when
he lieth sick upon his bed; * make thou all his bed in his sickness.

4 I said, LORD, be merciful
unto me; * heal my soul, for I have sinned against thee.

5 Mine enemies speak evil of me, * When shall he die,
and his name perish?

6 And if he come to see me, he speaketh vanity, * and
his heart conceiveth falsehood within himself; and when he cometh forth,
he telleth it.

7 All mine enemies whisper together against me; * even
against me do they imagine this evil.

1 meat-offering until
1928

8 An evil disease, say they, cleaveth fast unto him;
* and now that he lieth, he shall rise up no more.

9 Yea, even mine own familiar friend whom I trusted,
* who did also eat of my bread, hath laid great wait for me.

10 But be thou merciful unto me, O LORD;
* raise thou me up again, and I shall reward them.

11 By this I know thou favourest me, * that mine enemy
doth not triumph against me. 8 Let the sentence of guiltiness
proceed against him : until
1928

12 And in my innocency thou upholdest me, * and shalt
set me before thy face for ever.

13
Blessed be the LORD God of Israel, * world without
end. Amen.

BOOK
II.

Psalm
42. \emph{Quemadmodum.}

LIKE
as the hart desireth the water-brooks, * so longeth my soul after thee,
O God.

2 My soul is athirst for God, yea, even for the living
God: * when shall I come to appear before the presence of God?

3 My tears have been my meat day and night, * while
they daily say unto me, Where is now thy God?

4 Now when I think thereupon, I pour out my heart by
myself; * for I went with the multitude, and brought them forth into the
house of God;

5 In the voice of praise and thanksgiving, * among
such as keep holy-day.

6 Why art thou so full of heaviness, O my soul? * and
why art thou so disquieted within me?

12 And when I am in my health,
thou upholdest me : until
1928

7
O put thy trust in God; * for I will yet thank him, which is the help
of my countenance, and my God.

8
My soul is vexed within me; * therefore will I remember thee from the
land of Jordan, from Hermon and the little hill.

9
One deep calleth another, because of the noise of thy water-floods1;
* all thy waves and storms are gone over me.

10
The LORD will grant his loving-kindness in the
daytime; * and in the night season will I sing of him, and make my prayer
unto the God of my life.

11
I will say unto the God of my strength, Why hast thou forgotten me? *
why go I thus heavily, while the enemy oppresseth me?

12
My bones are smitten asunder as with a sword, * while mine enemies that
trouble me cast me in the teeth;

13
Namely, while they say daily unto me, * Where is now thy God?

14
Why art thou so vexed, O my soul? * and why art thou so disquieted within
me?

15
O put thy trust in God; * for I will yet thank him, which is the help
of my countenance, and my God.

Psalm
43. \emph{Judica me, Deus.}

GIVE
sentence with me, O God, and defend my cause against the ungodly people;
* O deliver me from the deceitful and wicked man.

2
For thou art the God of my strength; why hast thou put me from thee? *
and why go I so heavily, while the enemy oppresseth me?

3
O send out thy light and thy truth, that they may lead me, * and bring
me unto thy holy hill, and to thy dwelling;

4
And that I may go unto the altar of God, even unto the God of my joy and
gladness; * and upon the harp will I give thanks unto thee, O God, my
God.

5
Why art thou so heavy, O my soul? * and why art thou so disquieted within
me?

6
O put thy trust in God; * for I will yet give him thanks, which is the
help of my countenance, and my God.

The
Ninth Day.

Morning
Prayer.

Psalm
44. \emph{Deus, auribus.}

WE
have heard with our ears, O God, our fathers have told us * what thou
hast done in their time of old:

7. Put thy trust in God :
for I will yet give him thanks for the help of his countenance.

8. My God, my soul is vexed within me : therefore will I remember thee concerning
the land of Jordan, and the little hill of Hermon.
until 1928\\
1 the water-pipes until
1892; thy water pipes in
the 1892 book

10. The LORD hath granted
his loving-kindness in the day-time : and in the night-season did I sing
of him, and made my prayer unto the God of my life. until
1928

2 How thou hast driven out the heathen with thy hand,
and planted our fathers in; * how thou hast destroyed the nations, and made
thy people to flourish.

3 For they gat not the land in possession through their
own sword, * neither was it their own arm that helped them;

4 But thy right hand, and thine arm, and the light of
thy countenance; * because thou hadst a favour unto them.

5 Thou art my King, O God; * send help unto Jacob.

6 Through thee will we overthrow our enemies, * and in
thy Name will we tread them under that rise up against us.

7 For I will not trust in my bow, * it is not my sword
that shall help me;

8 But it is thou that savest us from our enemies, * and
puttest them to confusion that hate us.

9 We make our boast of God all day long, * and will praise
thy Name for ever.

10 But now thou art far off, and puttest us to confusion,
* and goest not forth with our armies.

11 Thou makest us to turn our backs upon our enemies,
* so that they which hate us spoil our goods.

12 Thou lettest us be eaten up like sheep, * and hast
scattered us among the heathen.

13 Thou sellest thy people for nought, * and takest no
money for them.

14 Thou makest us to be rebuked of our neighbours, *
to be laughed to scorn, and had in derision of them that are round about
us.

2. ... and planted them in
: how thou hast destroyed the nations and cast them out. until
1928

15 Thou makest us to be a byword among the nations1,
* and that the peoples2 shake their heads at us.

16 My confusion is daily before me, * and the shame
of my face hath covered me;

17 For the voice of the slanderer and blasphemer, *
for the enemy and avenger.

18 And though all this be come upon us, yet do we not
forget thee, * nor behave ourselves frowardly in thy covenant.

19 Our heart is not turned back, * neither our steps
gone out of thy way;

20 No, not when thou hast smitten us into the place
of dragons, * and covered us with the shadow of death.

21 If we have forgotten the Name of our God, and holden
up our hands to any strange god, * shall not God search it out? for he
knoweth the very secrets of the heart.

22 For thy sake also are we killed all the day long,
* and are counted as sheep appointed to be slain.

23 Up, Lord, why sleepest thou? * awake, and be not
absent from us for ever.

24 Wherefore hidest thou thy face, * and forgettest
our misery and trouble?

25 For our soul is brought low, even unto the dust;
* our belly cleaveth unto the ground.

26 Arise, and help us, * and deliver us, for thy mercy's
sake.

1 heathen until
1928\\
2 people until
1928

Psalm
45. \emph{Eructavit cor meum.}

MY
heart overfloweth with a good matter; I speak the things which I have
made concerning the King. * My tongue is the pen of a ready writer.

2 Thou art fairer than the children of men; * full
of grace are thy lips, because God hath blessed thee for ever.

3 Gird thee with thy sword upon thy thigh, O thou Most
Mighty, * according to thy worship and renown.

4 Good luck have thou with thine honour: * ride on,
because of the word of truth, of meekness, and righteousness; and thy
right hand shall teach thee terrible things.

MY HEART is inditing of a good matter :
I speak of the things which I have made unto the King.

2. My tongue is the pen : of a ready writer. until
1928; and remaining verses numbered one higher.

5 Thy arrows are very sharp in the heart of the King's
enemies, * and the people shall be subdued unto thee.

6 Thy seat, O God, endureth for ever; * the sceptre of
thy kingdom is a right sceptre.

7 Thou hast loved righteousness, and hated iniquity;
* wherefore God, even thy God, hath anointed thee with the oil of gladness
above thy fellows.

8 All thy garments smell of myrrh, aloes, and cassia;
* out of the ivory palaces, whereby they have made thee glad.

9 Kings’ daughters are1 among thy honourable
women; * upon thy right hand doth2 stand the queen in a vesture
of gold, wrought about with divers colours.

10 Hearken, O daughter, and consider; incline thine ear;
* forget also thine own people, and thy father's house. [6.] Thy arrows are very sharp,
and the people shall be subdued unto thee : even in the midst among the
King's enemies. until
1928

1 were until
1928

2 did until
1928

11 So shall the King have pleasure in thy beauty; * for
he is thy Lord3, and worship thou him.

12 And the daughter of Tyre shall be there with a gift;
* like as the rich also among the people shall make their supplication before
thee.

13 The King’s daughter is all glorious within;
* her clothing is of wrought gold.

14 She shall be brought unto the King in raiment of needlework:
* the virgins that be her fellows shall bear her company, and shall be brought
unto thee.

15 With joy and gladness shall they be brought, * and
shall enter into the King's palace.

16 Instead of thy fathers, thou shalt have children,
* whom thou mayest make princes in all lands.

3 Lord God, until
1928

17 I will make thy Name to be remembered from one generation to another;
* therefore shall the people give thanks unto thee, world without end.

Psalm
46. \emph{Deus noster refugium.}

GOD
is our hope and strength, * a very present help in trouble.

2 Therefore will we not fear, though the earth be moved,
* and though the hills be carried into the midst of the sea;

3 Though the waters thereof rage and swell, * and though
the mountains shake at the tempest of the same. [18] I will remember thy name
from one generation ...
until 1928

4 There is a river, the streams whereof make glad the
city of God; * the holy place of the tabernacle of the Most Highest.

5 God is in the midst of her, therefore shall she not
be removed; * God shall help her, and that right early.

4 The rivers of the
flood thereof shall make glad ...
until 1928

6 The nations1 make much ado, and the kingdoms
are moved; * but God hath showed his voice, and the earth shall melt away.

7 The Lord of hosts is with us; * the God of Jacob is
our refuge.

8 O come hither, and behold the works of the Lord, *
what destruction he hath brought upon the earth.

9 He maketh wars to cease in all the world; * he breaketh
the bow, and knappeth the spear in sunder, and burneth the chariots in the
fire. 1 heathen
until 1928

10 Be still then, and know that I am God: * I will
be exalted among the nations1, and I will be exalted in the
earth.

11 The Lord of hosts is with us; * the God of Jacob
is our refuge.

Evening
Prayer.

Psalm
47. \emph{Omnes gentes, plaudite.}

CLAP
your hands together, all ye peoples: * O sing unto God with the voice
of melody.

2 For the Lord is high, and to be feared; * he is the
great King upon all the earth.

1 heathen
until 1928

3 He shall subdue the peoples1 under us, * and the nations under
our feet.

4 He shall choose out an heritage for us, * even the
excellency2 of Jacob, whom he loved.

5 God is gone up with a merry noise, * and the Lord with
the sound of the trump.

6 O sing praises, sing praises unto our God; * O sing
praises, sing praises unto our King.

7 For God is the King of all the earth: * sing ye praises
with understanding.

1 people
until 1928

2 worship
until 1928

8 God reigneth over the nations3; * God
sitteth upon his holy seat.

9 The princes of the peoples4 are joined
unto the people of the God of Abraham; * for God, which is very high exalted,
doth defend the earth, as it were with a shield.

Psalm
48. \emph{Magnus Dominus.}

GREAT
is the Lord, and highly to be praised * in the city of our God, even upon
his holy hill.

2 The hill of Sion is a fair place, and the joy of
the whole earth; * upon the north side lieth the city of the great King:
God is well known in her palaces as a sure refuge.

3 For lo, the kings of the earth * were gathered, and
gone by together.

4 They marvelled to see such things; * they were astonished,
and suddenly cast down.

5 Fear came there upon them; and sorrow, * as upon
a woman in her travail.

6 Thou dost break the ships of the sea * through the
east-wind.

7 Like as we have heard, so have we seen in the city
of the Lord of hosts, in the city of our God; * God upholdeth the same
for ever.

8 We wait for thy loving-kindness, O God, * in the
midst of thy temple.

9 O God, according to thy Name, so is thy praise unto
the world's end; * thy right hand is full of righteousness.

10 Let the mount Sion rejoice, and the daughters of
Judah be glad, * because of thy judgments.

11 Walk about Sion, and go round about her; * and tell
the towers thereof.

3 heathen
until 1928

4 people
until 1928

12 Mark well her bulwarks, consider her palaces, *
that ye may tell them that come after.

13 For this God is our God for ever and ever: * he
shall be our guide unto death.

Psalm
49. \emph{Audite haec, omnes.}

O
HEAR ye this, all ye people; * ponder it with your ears, all ye that dwell
in the world;

2 High and low, rich and poor, * one with another.

3 My mouth shall speak of wisdom, * and my heart shall
muse of understanding.

4 I will incline mine ear to the parable, * and show
my dark speech upon the harp. 12 ... bulwarks, set up her
houses :
until 1928

5 Wherefore should I fear in the days of evil, * when
wickedness at my heels compasseth me round about?

6 There be some that put their trust in their goods,
* and boast themselves in the multitude of their riches.

7 But no man may deliver his brother, * nor give a ransom
unto God for him, 5 ... days of wickedness :
and when the wickedness
until 1928

8 (For it cost more to redeem their souls, * so that
he must let that alone for ever;)

9 That he shall live alway, * and not see the grave.

10 For he seeth that wise men also die and perish together,
* as well as the ignorant and foolish, and leave their riches for other.

11 And yet they think that their houses shall continue
for ever, and that their dwelling-places shall endure from one generation
to another; * and call the lands after their own names.

Parentheses
added in 1928.

9 Yea, though he live long :
until 1928.

12 Nevertheless, man being in honour abideth not, * seeing he may be compared
unto the beasts that perish;

13 This their way is very foolishness; * yet their posterity
praise their saying.

14 They lie in the grave like sheep; death is their shepherd;
and the righteous shall have dominion over them in the morning: * their
beauty shall consume in the sepulchre, and have no abiding.

15 But God hath delivered my soul from the power of the
grave; * for he shall receive me.

16 Be not thou afraid, though one be made rich, * or
if the glory of his house be increased;

17 For he shall carry nothing away with him when he dieth,
* neither shall his pomp follow him.

18 For while he lived, he counted himself an happy man;
* and so long as thou doest well unto thyself, men will speak good of thee.

19 He shall follow the generation of his fathers, * and
shall never see light.\\
12. Nevertheless, man will
not abide in honour : seeing he may be compared unto the beasts that perish;
this is the way of them.\\
13. This is their foolishness :
and their posterity praise their saying.\\
14. They
lie in the hell like sheep, death gnaweth upon them, and the righteous shall
have domination [dominion in
the 1789 Book] over them in the morning : their beauty shall consume
in the sepulchre out of their dwelling.\\
15. But God
hath delivered my soul from the place of hell : for he shall receive me.
until
1928.

20 Man that is in honour but hath no understanding
* is compared unto the beasts that perish.

The
Tenth Day.

Morning
Prayer.

Psalm
50. \emph{Deus deorum.}

THE
Lord, even the Most Mighty God, hath spoken, * and called the world, from
the rising up of the sun unto the going down thereof.

2 Out of Sion hath God appeared * in perfect beauty.

3 Our God shall come, and shall not keep silence; *
there shall go before him a consuming fire, and a mighty tempest shall
be stirred up round about him.

4 He shall call the heaven from above, * and the earth,
that he may judge his people.

5 Gather my saints together unto me; * those that have
made a covenant with me with sacrifice.

6 And the heavens shall declare his righteousness;
* for God is Judge himself.

7 Hear, O my people, and I will speak; * I myself will
testify against thee, O Israel; for I am God, even thy God.

8 I will not reprove thee because of thy sacrifices;
* as for thy burnt-offerings, they are alway before me.

20. Man being in honour hath
no understanding : but is compared ...
until 1928.

9 I will take no bullock out of thine house, * nor he-goats1
out of thy folds.

10 For all the beasts of the forest are mine, * and so
are the cattle upon a thousand hills.

11 I know all the fowls upon the mountains, * and the
wild beasts of the field are in my sight.

12 If I be hungry, I will not tell thee; * for the whole
world is mine, and all that is therein.

13 Thinkest thou that I will eat bulls' flesh, * and
drink the blood of goats?

14 Offer unto God thanksgiving, * and pay thy vows unto
the Most Highest.

15 And call upon me in the time of trouble; * so will
I hear thee, and thou shalt praise me. 

1 he-goat
until 1892

16 But unto the ungodly saith2 God, * Why
dost thou preach my laws, and takest my covenant in thy mouth;

17 Whereas thou hatest to be reformed, * and hast cast
my words behind thee?

18 When thou sawest a thief, thou consentedst unto him;
* and hast been partaker with the adulterers.

19 Thou hast let thy mouth speak wickedness, * and with
thy tongue thou hast set forth deceit.

20 Thou sattest and spakest against thy brother; * yea,
and hast slandered thine own mother's son.

21 These things hast thou done, and I held my tongue,
and thou thoughtest wickedly, that I am even such a one as thyself; * but
I will reprove thee, and set before thee the things that thou hast done.

22 O consider this, ye that forget God, * lest I pluck
you away, and there be none to deliver you.

2 said
until 1928.

23 Whoso offereth me thanks and praise, he honoureth
me; * and to him that ordereth his way aright, will I show the salvation
of God. 23 ... ordereth his conversation
right, ...
until 1928.

\begin{quote}

\emph{Psalms 51-100;
101-150 Return to the U. S. Book of Common Prayer: 1789;
1892; 1928}

\end{quote}

\xchapter{The Psalter, Books III and IV}

\textbf{The Psalter from the 1789, 1892 and 1928 U. S. Books of Common
Prayer}

Psalm 51. \emph{Miserere mei, Deus.}

HAVE mercy upon me, O God, after thy great
goodness; * according to the multitude of thy mercies do away mine offences.

2 Wash me throughly from my wickedness, * and cleanse
me from my sin.

3 For I acknowledge my faults, * and my sin is ever
before me.

4 Against thee only have I sinned, and done this evil
in thy sight; * that thou mightest be justified in thy saying, and clear
when thou shalt judge.

5 Behold, I was shapen in wickedness, * and in sin
hath my mother conceived me.

6 But lo, thou requirest truth in the inward parts,
* and shalt make me to understand wisdom secretly.

7 Thou shalt purge me with hyssop, and I shall
be clean; * thou shalt wash me, and I shall be whiter than snow.

8 Thou shalt make me hear of joy and gladness, * that
the bones which thou hast broken may rejoice.

9 Turn thy face from my sins, * and put out all my
misdeeds.

10 Make me a clean heart, O God, * and renew a right
spirit within me.

11 Cast me not away from thy presence, * and take not
thy holy Spirit from me.

12 O give me the comfort of thy help again, * and stablish
me with thy free Spirit.

13 Then shall I teach thy ways unto the wicked, * and
sinners shall be converted unto thee.

14 Deliver me from blood-guiltiness, O God, thou that
art the God of my health; * and my tongue shall sing of thy righteousness.

15 Thou shalt open my lips, O Lord, * and my mouth
shall show thy praise.

16 For thou desirest no sacrifice, else would I give
it thee; * but thou delightest not in burnt-offerings.

17 The sacrifice of God is a troubled spirit: * a broken
and contrite heart, O God, shalt thou not despise.

18 O be favourable and gracious unto Sion; * build
thou the walls of Jerusalem.

19 Then shalt thou be pleased with the sacrifice of
righteousness, with the burnt-offerings and oblations; * then shall they
offer young bullocks upon thine altar.

Psalm 52. \emph{Quid gloriaris?}

WHY boastest thou thyself, thou tyrant,
* that thou canst do mischief;

2 Whereas the goodness of God * endureth yet daily?

3 Thy tongue imagineth wickedness, * and with lies
thou cuttest like a sharp razor.

Go
to Psalms 1-50; 

101-150

Return to the U. S. Book of Common Prayer: 1789;
1892; 1928

4 Thou hast loved unrighteousness more than goodness, * and falsehood more
than righteousness.

5 Thou hast loved to speak all words that may do hurt,
* O thou false tongue.

6 Therefore shall God destroy thee for ever; * he shall
take thee, and pluck thee out of thy dwelling, and root thee out of the
land of the living.

7 The righteous also shall see this, and fear, * and
shall laugh him to scorn:

8 Lo, this is the man that took not God for his strength;
* but trusted unto the multitude of his riches, and strengthened himself
in his wickedness.

9 As for me, I am like a green olive-tree in the house
of God; * my trust is in the tender mercy of God for ever and ever.

4 ... and talk
of lies more than righteousness.
until 1928

10 I will alway1 give thanks unto thee for that thou hast done;
* and I will hope in thy Name, for thy saints like it well.

Evening Prayer.

Psalm
53. \emph{Dixit insipiens.}

THE
foolish body hath said in his heart, * There is no God.

2 Corrupt are they, and become abominable in their
wickedness; * there is none that doeth good.

3 God looked down from heaven upon the children of
men, * to see if there were any that would understand, and seek after
God.

4 But they are all gone out of the way, they are altogether
become abominable; * there is also none that doeth good, no not one.

5 Are not they without understanding that work wickedness,
* eating up my people as if they would eat bread? they have not called
upon God.

6 They were afraid where no fear was; * for God hath
broken the bones of him that besieged thee; thou hast put them to confusion,
because God hath despised them.

7 O that the salvation were given unto Israel out of
Sion! * O that the Lord would deliver his people out of captivity!

8 Then should Jacob rejoice, * and Israel should be
right glad.

Psalm
54. \emph{Deus, in Nomine.}

SAVE
me, O God, for thy Name's sake, * and avenge me in thy strength.

2 Hear my prayer, O God, * and hearken unto the words
of my mouth.

3 For strangers are risen up against me; * and tyrants,
which have not God before their eyes, seek after my soul.

4 Behold, God is my helper; * the Lord is with them
that uphold my soul.

5 He shall reward evil unto mine enemies: * destroy
thou them in thy truth.

6 An offering of a free heart will I give thee, and
praise thy Name, O LORD; * because it is so comfortable.

7 For he hath delivered me out of all my trouble; *
and mine eye hath seen his desire upon mine enemies.

Psalm
55. \emph{Exaudi, Deus.}

HEAR
my prayer, O God, * and hide not thyself from my petition.

2 Take heed unto me, and hear me, * how I mourn in
my prayer, and am vexed;

3 The enemy crieth so, and the ungodly cometh on so
fast; * for they are minded to do me some mischief, so maliciously are
they set against me.

4 My heart is disquieted within me, * and the fear
of death is fallen upon me.

5 Fearfulness and trembling are come upon me, * and
an horrible dread hath overwhelmed me.

6 And I said, O that I had wings like a dove! * for
then would I flee away, and be at rest.

7 Lo, then would I get me away far off, * and remain
in the wilderness.8 I would make haste to escape, * because of the stormy
wind and tempest.

9 Destroy their tongues, O Lord, and divide them; *
for I have spied unrighteousness and strife in the city.

10 Day and night they go about within the walls thereof:
* mischief also and sorrow are in the midst of it.

1 always until
1892

11 Wickedness is therein; * deceit and guile go not out of her1
streets.

12 For it is not an open enemy that hath done me this
dishonour; * for then I could have borne it;

13 Neither was it mine adversary that did magnify himself
against me; * for then peradventure I would have hid myself from him;

14 But it was even thou, my companion, * my guide, and
mine own familiar friend.

15 We took sweet counsel together, * and walked in the
house of God as friends.

1 their until
1892

16 Let death come hastily upon them, and let them go down alive into the
pit; * for wickedness is in their dwellings, and among them.

17 As for me, I will call upon God, * and the LORD shall
save me.

18 In the evening, and morning, and at noon-day will
I pray, and that instantly; * and he shall hear my voice. \\
16 ... go down
quick into hell
until 1928

19 It is he that hath delivered my soul in peace from the battle
that was against me; * for there were many that strove with me.

20 Yea, even God, that endureth for ever, shall hear
me, and bring them down; * for they will not turn, nor fear God.

21 He laid his hands upon such as be at peace with
him, * and he brake his covenant.

22 The words of his mouth were softer than butter,
having war in his heart; * his words were smoother than oil, and yet be
they very swords.

23 O cast thy burden upon the LORD, and he shall nourish
thee, * and shall not suffer the righteous to fall for ever.

24 And as for them, * thou, O God, shalt bring them
into the pit of destruction.

25 The blood-thirsty and deceitful men shall not live
out half their days: * nevertheless, my trust shall be in thee, O Lord.

The
Eleventh Day.

Morning
Prayer.

Psalm
56. \emph{Miserere mei, Deus.}

BE
merciful unto me, O God, for man goeth about to devour me; * he is daily
fighting, and troubling me.

19 ... were many with me. until
1928

2 Mine enemies are daily at hand to swallow me up; * for they be many that
fight against me, O thou Most Highest.

3 Nevertheless, though I am sometime afraid, * yet put
I my trust in thee.

4 I will praise God, because of his word: * I have put
my trust in God, and will not fear what flesh can do unto me.

5 They daily mistake my words; * all that they imagine
is to do me evil.

6 They hold all together, and keep themselves close,
* and mark my steps, when they lay wait for my soul.

7 Shall they escape for their wickedness? * thou, O God,
in thy displeasure shalt cast them down. 2. Mine enemies are daily
in hand to swallow me up : for they be many that fight against me, O thou
most Highest.
until 1928

8 Thou tellest my wanderings1; put my tears into thy
bottle: * are not these things noted in thy book?

9 Whensoever I call upon thee, then shall mine enemies
be put to flight: * this I know; for God is on my side.

10 In God's word will I rejoice; * in the LORD’S
word will I comfort me.

11 Yea, in God have I put my trust; * I will not be
afraid what man can do unto me.

12 Unto thee, O God, will I pay my vows; * unto thee
will I give thanks.

13 For thou hast delivered my soul from death, and
my feet from falling, * that I may walk before God in the light of the
living.

Psalm
57. \emph{Miserere mei, Deus.}

BE
merciful unto me, O God, be merciful unto me; for my soul trusteth in
thee; * and under the shadow of thy wings shall be my refuge, until this
tyranny be overpast.

2 I will call unto the Most High God, * even unto the
God that shall perform the cause which I have in hand.

3 He shall send from heaven, * and save me from the
reproof of him that would eat me up.

4 God shall send forth his mercy and truth: * my soul
is among lions;

5 And I lie even among the children of men, that are
set on fire, * whose teeth are spears and arrows, and their tongue a sharp
sword.

6 Set up thyself, O God, above the heavens; * and thy
glory above all the earth.

7 They have laid a net for my feet, and pressed down
my soul; * they have digged a pit before me, and are fallen into the midst
of it themselves.

8 My heart is fixed, O God, my heart is fixed; * I
will sing and give praise.

9 Awake up, my glory; awake, lute and harp: * I myself
will awake right early. 1 flittings
in Engl. Book

10 I will give thanks unto thee, O Lord, among the peoples1;
* and I will sing unto thee among the nations.

11 For the greatness of thy mercy reacheth unto the
heavens, * and thy truth unto the clouds.

12 Set up thyself, O God, above the heavens; * and
thy glory above all the earth.

Psalm
58. \emph{Si vere utique.}

ARE
your minds set upon righteousness, O ye congregation? * and do ye judge
the thing that is right, O ye sons of men?

2 Yea, ye imagine mischief in your heart upon the earth,
* and your hands deal with wickedness.

3 The ungodly are froward, even from their mother's
womb; * as soon as they are born, they go astray, and speak lies.

4 They are as venomous as the poison of a serpent,
* even like the deaf adder, that stoppeth her ears;

5 Which refuseth to hear the voice of the charmer,
* charm he never so wisely.

1 people
until 1928

6 Break their teeth, O God, in their mouths; * smite the jaw-bones
of the lions, O LORD.

7 Let them fall away like water that runneth apace;
* when they shoot their arrows, let them be rooted out.

8 Let them consume away like a snail, and be like the
untimely fruit of a woman; * and let them not see the sun.

9 Or ever your pots be made hot with thorns, * he shall
take them away with a whirlwind, the green and the burning alike.

10 The righteous shall rejoice when he seeth the vengeance;
* he shall wash his footsteps in the blood of the ungodly.

11 So that a man shall say, Verily there is a reward
for the righteous; * doubtless there is a God that judgeth the earth.

Evening
Prayer.

Psalm
59. \emph{Eripe me de inimicis.}

DELIVER
me from mine enemies, O God; * defend me from them that rise up against
me.

2 O deliver me from the wicked doers, * and save me
from the blood-thirsty men.

3 For lo, they lie waiting for my soul; * the mighty
men are gathered against me, without any offence or fault of me, O LORD.

4 They run and prepare themselves without my fault;
* arise thou therefore to help me, and behold.

5 Stand up, O LORD God of hosts, thou God of Israel,
to visit all the heathen, * and be not merciful unto them that offend
of malicious wickedness.

6 They go to and fro in the evening, * they grin like
a dog, and run about through the city.

7 Behold, they speak with their mouth, and swords are
in their lips; * for who doth hear?

8 But thou, O LORD, shalt have them in derision, *
and thou shalt laugh all the heathen to scorn.

9 My strength will I ascribe unto thee; * for thou
art the God of my refuge.

10 God showeth me his goodness plenteously; * and God
shall let me see my desire upon mine enemies.

11 Slay them not, lest my people forget it; * but scatter
them abroad among the people, and put them down, O Lord our defence.

6. Break their teeth, O
God, in their mouths; smite the jaw-bones of the lions, O LORD
: let them fall away like water that runneth apace; and when they

shoot their arrows let them be rooted out.
and remaining verses numbered one higher, until 1928.

[8.] Or ever your pots be made hot with thorns
: so let 

indignation vex him, even as a thing that is raw.
until 1928

12 For the sin of their mouth, and for the words of their lips, they shall
be taken in their pride: * and why? their talk1 is of cursing
and lies.

13 Consume them in thy wrath, consume them, that they
may perish; * and know that it is God that ruleth in Jacob, and unto the
ends of the world.

14 And in the evening they will return, * grin like
a dog, and will go about the city.

15 They will run here and there for meat, * and grudge
if they be not satisfied.

16 As for me, I will sing of thy power, and will praise
thy mercy betimes in the morning; * for thou hast been my defence and
refuge in the day of my trouble.

17 Unto thee, O my strength, will I sing; * for thou,
O God, art my refuge, and my merciful God.

Psalm
60. \emph{Deus, repulisti nos.}

O
GOD, thou hast cast us out, and scattered us abroad; * thou hast also
been displeased: O turn thee unto us again.

2 Thou hast moved the land, and divided it: * heal
the sores thereof, for it shaketh.

3 Thou hast showed thy people heavy things; * thou
hast given us a drink of deadly wine.

4 Thou hast given a token for such as fear thee, *
that they may triumph because of the truth.

5 Therefore were thy beloved delivered; * help me with
thy right hand, and hear me.

1 preaching
until 1928

6 God hath spoken in his holiness, I will rejoice, and divide Shechem1,
* and mete out the valley of Succoth.

7 Gilead is mine, and Manasseh2 is mine;
* Ephraim also is the strength of my head; Judah is my law-giver;

8 Moab is my wash-pot; over Edom will I cast out my
shoe; * Philistia, be thou glad of me.

9 Who will lead me into the strong city? * who will
bring me into Edom?

10 Hast not thou cast us out, O God? * wilt not thou,
O God, go out with our hosts?

11 O be thou our help in trouble; * for vain is the
help of man.

12 Through God will we do great acts; * for it is he
that shall tread down our enemies.

Psalm
61. \emph{Exaudi, Deus.}

HEAR
my crying, O God, * give ear unto my prayer.

2 From the ends of the earth will I call upon thee,
* when my heart is in heaviness.

3 O set me up upon the rock that is higher than I;
* for thou hast been my hope, and a strong tower for me against the enemy.

4 I will dwell in thy tabernacle for ever, * and my
trust shall be under the covering of thy wings.

5 For thou, O Lord, hast heard my desires, * and hast
given an heritage unto those that fear thy Name.

6 Thou shalt grant the King a long life, * that his
years may endure throughout all generations.

7 He shall dwell before God for ever: * O prepare thy
loving mercy and faithfulness, that they may preserve him.

1 Sichem
until 1928

2 Manasses
until 1928

8 So will I alway1 sing praise unto thy Name, * that I may
daily perform my vows. 

The
Twelfth Day.

Morning
Prayer.

Psalm
62. \emph{Nonne Deo?}

MY
soul truly waiteth still upon God; * for of him cometh my salvation.

2 He verily is my strength and my salvation; * he is
my defence, so that I shall not greatly fall.

3 How long will ye imagine mischief against every man?
* Ye shall be slain all the sort of you; yea, as a tottering wall shall
ye be, and like a broken hedge.

4 Their device is only how to put him out whom God
will exalt; * their delight is in lies; they give good words with their
mouth, but curse with their heart.

5 Nevertheless, my soul, wait thou still upon God;
* for my hope is in him.

6 He truly is my strength and my salvation; * he is
my defence, so that I shall not fall.

7 In God is my health and my glory; * the rock of my
might; and in God is my trust.

8 O put your trust in him alway, ye people; * pour
out your hearts before him, for God is our hope.
1 always
until 1892

9 As for the children of men, they are but vanity; the children of men
are deceitful; * upon the weights they are altogether lighter than vanity
itself.

10 O trust not in wrong and robbery; give not yourselves
unto vanity: * if riches increase, set not your heart upon them.

11 God spake once, and twice I have also heard the
same, * that power belongeth unto God;

12 And that thou, Lord, art merciful; * for thou rewardest
every man according to his work.

Psalm
63. \emph{Deus, Deus meus.}

O
GOD, thou art my God; * early will I seek thee.

2 My soul thirsteth for thee, my flesh also longeth
after thee, * in a barren and dry land where no water is. 9. As for the children of
men, they are but vanity : the children of men are deceitful upon the weights,
they are altogether lighter than vanity itself.
until 1928

3 Thus have I looked for thee in the sanctuary, * that I might behold thy
power and glory.

4 For thy loving-kindness is better than the life itself:
* my lips shall praise thee. 3 ... for thee in holiness
: until
1928

5 As long as I live will I magnify thee in1 this manner, *
and lift up my hands in thy Name.

6 My soul shall be satisfied, even as it were with
marrow and fatness, * when my mouth praiseth thee with joyful lips.

7 Have I not remembered thee in my bed, * and thought
upon thee when I was waking?

8 Because thou hast been my helper; * therefore under
the shadow of thy wings will I rejoice.

9 My soul hangeth upon thee; * thy right hand hath
upholden me.

10 These also that seek the hurt of my soul, * they
shall go under the earth.

11 Let them fall upon the edge of the sword, * that
they may be a portion for foxes.

12 But the King shall rejoice in God; all they also
that swear by him shall be commended; * for the mouth of them that speak
lies shall be stopped.

Psalm
64. \emph{Exaudi, Deus.}

HEAR
my voice, O God, in my prayer; * preserve my life from fear of the enemy.

2 Hide me from the gathering together of the froward,
* and from the insurrection of wicked doers;

3 Who have whet their tongue like a sword, * and shoot
out their arrows, even bitter words;

4 That they may privily shoot at him that is perfect:
* suddenly do they hit him, and fear not.

5 They encourage themselves in mischief, * and commune
among themselves, how they may lay snares; and say, that no man shall
see them.

6 They imagine wickedness, and practise it; * that
they keep secret among themselves, every man in the deep of his heart.

7 But God shall suddenly shoot at them with a swift
arrow, * that they shall be wounded.

8 Yea, their own tongues shall make them fall; * insomuch
that whoso seeth them shall laugh them to scorn.

9 And all men that see it shall say, This hath God
done; * for they shall perceive that it is his work.

10 The righteous shall rejoice in the LORD, and put
his trust in him; * and all they that are true of heart shall be glad.

Evening
Prayer.

Psalm
65. \emph{Te decet hymnus.}

THOU,
O God, art praised in Sion; * and unto thee shall the vow be performed
in Jerusalem.

2 Thou that hearest the prayer, * unto thee shall all
flesh come.

3 My misdeeds prevail against me: * O be thou merciful
unto our sins.

4 Blessed is the man whom thou choosest, and receivest
unto thee: * he shall dwell in thy court, and shall be satisfied with
the pleasures of thy house, even of thy holy temple.

5 Thou shalt show us wonderful things in thy righteousness,
O God of our salvation; * thou that art the hope of all the ends of the
earth, and of them that remain in the broad sea.

6 Who in his strength setteth fast the mountains, *
and is girded about with power.

7 Who stilleth the raging of the sea, * and the noise
of his waves, and the madness of the peoples.

8 They also that dwell in the uttermost parts of the
earth shall be afraid at thy tokens, * thou that makest the outgoings
of the morning and evening to praise thee.

9 Thou visitest the earth, and blessest it; * thou
makest it very plenteous.

10 The river of God is full of water: * thou preparest
their corn, for so thou providest for the earth.

11 Thou waterest her furrows; thou sendest rain into
the little valleys thereof; * thou makest it soft with the drops of rain,
and blessest the increase of it.

12 Thou crownest the year with thy goodness; * and
thy clouds drop fatness.

13 They shall drop upon the dwellings of the wilderness;
* and the little hills shall rejoice on every side.

14 The folds shall be full of sheep; * the valleys
also shall stand so thick with corn, that they shall laugh and sing.

Psalm
66. \emph{Jubilate Deo.}

O
BE joyful in God, all ye lands; * sing praises unto the honour of his
Name; make his praise to be glorious. 1
on in
Engl. Book

2 Say unto God, O how wonderful art thou in thy works! * through the greatness
of thy power shall thine enemies bow down unto thee.

3 For all the world shall worship thee, * sing of thee,
and praise thy Name.

4 O come hither, and behold the works of God; * how wonderful
he is in his doing toward the children of men.

5 He turned the sea into dry land, * so that they went
through the water on foot; there did we rejoice thereof. \\
2 ... enemies be found liars
unto thee.
until 1928

6 He ruleth with his power for ever; his eyes behold the nations1:
* and such as will not believe shall not be able to exalt themselves.

7 O praise our God, ye peoples, * and make the voice
of his praise to be heard;

8 Who holdeth our soul in life; * and suffereth not
our feet to slip.

9 For thou, O God, hast proved us; * thou also hast
tried us, like as silver is tried.

10 Thou broughtest us into the snare; * and laidest
trouble upon our loins.

11 Thou sufferedst men to ride over our heads; * we
went through fire and water, and thou broughtest us out into a wealthy
place.

12 I will go into thine house with burnt-offerings,
and will pay thee my vows, * which I promised with my lips, and spake
with my mouth, when I was in trouble.

13 I will offer unto thee fat burnt-sacrifices, with
the incense of rams; * I will offer bullocks and goats.

14 O come hither, and hearken, all ye that fear God;
* and I will tell you what he hath done for my soul.

15 I called unto him with my mouth, * and gave him
praises with my tongue.

16 If I incline unto wickedness with mine heart, *
the Lord will not hear me.

17 But God hath heard me; * and considered the voice
of my prayer.

18 Praised be God, who hath not cast out my prayer,
* nor turned his mercy from me.

Psalm
67. \emph{Deus misereatur.}

GOD
be merciful unto us, and bless us, * and show us the light of his countenance,
and be merciful unto us;

2 That thy way may be known upon earth, * thy saving
health among all nations.

3 Let the peoples praise thee, O God; * yea, let all
the peoples praise thee.

4 O let the nations rejoice and be glad; * for thou
shalt judge the folk righteously, and govern the nations upon earth.

1
people
until 1928

5 Let the peoples praise thee, O God; * yea1, let all the peoples
praise thee.

6 Then shall the earth bring forth her increase; *
and God, even our own God, shall give us his blessing.

7 God shall bless us; * and all the ends of the world
shall fear him. 

The
Thirteenth Day.

Morning
Prayer.

Psalm
68. \emph{Exsurgat Deus.}

LET
God arise, and let his enemies be scattered; * let them also that hate
him flee before him.

2 Like as the smoke vanisheth, so shalt thou drive
them away; * and like as wax melteth at the fire, so let the ungodly perish
at the presence of God.

3 But let the righteous be glad, and rejoice before
God; * let them also be merry and joyful.

1 yea
added, 1822

4 O sing unto God, and sing praises unto his Name; magnify him that rideth
upon the heavens; * praise him in his Name JAH, and rejoice before him.

5 He is a Father of the fatherless, and defendeth the
cause of the widows; * even God in his holy habitation.

6 He is the God that maketh men to be of one mind in
an house, and bringeth the prisoners out of captivity; * but letteth the
runagates continue in scarceness.

7 O God, when thou wentest forth before the people; *
when thou wentest through the wilderness,

8 The earth shook, and the heavens dropped at the presence
of God; * even as Sinai also was moved at the presence of God, who is the
God of Israel.

9 Thou, O God, sentest a gracious rain upon thine inheritance,
* and refreshedst it when it was weary.

10 Thy congregation shall dwell therein; * for thou,
O God, hast of thy goodness prepared for the poor. 4. O sing unto God, and sing
praises unto his Name : magnify him that rideth upon the heavens, as it
were upon an horse; praise him in his Name JAH, and rejoice before him.
until
1928

11 The Lord gave the word; * great was the company of women that bare the
tidings.

12 Kings with their armies did flee, and were discomfited,
* and they of the household divided the spoil. 11 ... great was the company
of preachers.
until 1928

13 Though ye have lain1 among the sheep-folds2, yet
shall ye be as the wings of a dove * that is covered with silver wings,
and her feathers like gold.

14 When the Almighty scattered kings for their sake,
* then were they as white as snow in Salmon. 1 lien in
Engl. Book\\
2 pots
until 1928

15 As the hill of Bashan3, so is God's hill; * even an high hill,
as the hill of Bashan3.

16 Why mock4 ye so, ye high hills? this is
God's hill, in the which it pleaseth him to dwell; * yea, the LORD will
abide in it for ever.

17 The chariots of God are twenty thousand, even thousands
of angels; * and the Lord is among them as in the holy place of Sinai.

3 Basan
until 1928

4 hop
until 1928

18 Thou art gone up on high, thou hast led captivity captive, and received
gifts from5 men; * yea, even from5 thine enemies,
that the LORD God might dwell among them.

19 Praised be the Lord daily, * even the God who helpeth
us, and poureth his benefits upon us.

20 He is our God, even the God of whom cometh salvation:
* GOD is the Lord, by whom we escape death.

21 God shall wound the head of his enemies, * and the
hairy scalp of such a one as goeth on still in his wickedness.

5 for
until 1928

22 The Lord hath said, I will bring my people again, as I did from Bashan6;
* mine own will I bring again, as I did sometime from the deep of the sea.

23 That thy foot may be dipped in the blood of thine
enemies, * and that the tongue of thy dogs may be red through the same.

24 It is well seen, O God, how thou goest; * how thou,
my God and King, goest in the sanctuary. \\
6 Basan
until 1928

25 The singers go before, the minstrels follow after, * in the midst of7
the damsels playing with the timbrels.

26 Give thanks unto God the Lord in the congregation,
* ye that are of the fountain of Israel.

27 There is little Benjamin their ruler, and the princes
of Judah their council8; * the princes of Zebulon9,
and the princes of Naphthali.

28 Thy God hath sent forth strength for thee; * stablish
the thing, O God, that thou hast wrought in us,

29 For thy temple's sake at Jerusalem; * so shall kings
bring presents unto thee.

7 are
until 1928

26. Give thanks, O Israel, unto God the Lord in the

congregations : from the ground of the heart.
until 1928

8 counsel in
Engl. Book, 1845-1892 

9
Zabulon until
1793, and in 1892 Book

30 Rebuke thou the dragon and the bull, with the leaders of the heathen,
so that they humbly bring pieces of silver; * scatter thou the peoples
that delight in war;

31 Then shall the princes come out of Egypt; * the
Morians’ land shall soon stretch out her hands unto God.

32 Sing unto God, O ye kingdoms of the earth; * O sing
praises unto the Lord;

33 Who sitteth in the heavens over all, from the beginning:
* lo, he doth send out his voice; yea, and that a mighty voice.

34 Ascribe ye the power to God over Israel; * his worship
and strength is in the clouds.

35 O God, wonderful art thou in thy holy places: *
even the God of Israel, he will give strength and power unto his people.
Blessed be God. 

Evening
Prayer.

Psalm
69. \emph{Salvum me fac.}

SAVE
me, O God; * for the waters are come in, even unto my soul.

2 I stick fast in the deep mire, where no ground is;
* I am come into deep waters, so that the floods run over me.

3 I am weary of crying; my throat is dry; * my sight
faileth me for waiting so long upon my God.

4 They that hate me without a cause are more than the
hairs of my head; * they that are mine enemies, and would destroy me guiltless,
are mighty.

5 I paid them the things that I never took: * God,
thou knowest my simpleness, and my faults are not hid from thee.

6 Let not them that trust in thee, O Lord GOD of hosts,
be ashamed for my cause; * let not those that seek thee be confounded
through me, O Lord God of Israel.

7 And why? for thy sake have I suffered reproof; *
shame hath covered my face.

8 I am become a stranger unto my brethren, * even an
alien unto my mother's children.

9 For the zeal of thine house hath even eaten me; *
and the rebukes of them that rebuked thee are fallen upon me.

10 I wept, and chastened myself with fasting, * and
that was turned to my reproof.

11 I put on sackcloth also, * and they jested upon
me.

12 They that sit in the gate speak against me, * and
the drunkards make songs upon me.

13 But, LORD, I make my prayer unto thee * in an acceptable
time.

14 Hear me, O God, in the multitude of thy mercy, *
even in the truth of thy salvation.

15 Take me out of the mire, that I sink not; * O let
me be delivered from them that hate me, and out of the deep waters.

16 Let not the water-flood drown me, neither let the
deep swallow me up; * and let not the pit shut her mouth upon me.

17 Hear me, O LORD, for thy loving-kindness is comfortable;
* turn thee unto me according to the multitude of thy mercies.

18 And hide not thy face from thy servant; for I am
in trouble: * O haste thee, and hear me.

19 Draw nigh unto my soul, and save it; * O deliver
me, because of mine enemies. 30. When the company of the
spearmen, and multitude of the mighty are scattered abroad among the beasts
of the people, so that they humbly bring pieces of silver : and when he
hath scattered the people that delight in war;
until 1928

20 Thou hast known my reproach1, my shame, and my dishonour:
* mine adversaries are all in thy sight.

21 Reproach hath broken my heart; I am full of heaviness:
* I looked for some to have pity on me, but there was no man, neither
found I any to comfort me.

22 They gave me gall to eat; * and when I was thirsty
they gave me vinegar to drink.

23
Let their table be made a snare to take themselves withal; * and let the
things that should have been for their wealth be unto them an occasion
of falling.

24 Let their eyes be blinded, that they see not; *
and ever bow thou down their backs.

25 Pour out thine indignation upon them, * and let
thy wrathful displeasure take hold of them.

26 Let their habitation be void, * and no man to dwell
in their tents.

27 For they persecute him whom thou hast smitten; *
and they talk how they may vex them whom thou hast wounded.

28 Let them fall from one wickedness to another, *
and not come into thy righteousness.

29 Let them be wiped out of the book of the living,
* and not be written among the righteous.

30 As for me, when I am poor and in heaviness, * thy
help, O God, shall lift me up.

31 I will praise the Name of God with a song, * and
magnify it with thanksgiving.

32 This also shall please the LORD
* better than a bullock that hath horns and hoofs.

33 The humble shall consider this, and be glad: * seek
ye after God, and your soul shall live.

34 For the LORD heareth the
poor, * and despiseth not his prisoners.

35 Let heaven and earth praise him: * the sea, and
all that moveth therein.

36 For God will save Sion, and build the cities of
Judah, * that men may dwell there, and have it in possession.

37 The posterity also of his servants shall inherit
it; * and they that love his Name shall dwell therein.

Psalm
70. \emph{Deus, in adjutorium.}

HASTE
thee, O God, to deliver me; * make haste to help me, O LORD.

2 Let them be ashamed and confounded that seek after
my soul; * let them be turned backward and put to confusion that wish
me evil.

3 Let them for their reward be soon brought to shame,
* that cry over me, There! there!

4 But let all those that seek thee be joyful and glad
in thee: * and let all such as delight in thy salvation say alway, The
Lord be praised.

5 As for me, I am poor and in misery: * haste thee
unto me, O God.

6 Thou art my helper, and my redeemer: * O LORD, make
no long tarrying.

The
Fourteenth Day.

Morning
Prayer.

Psalm
71. \emph{In te, Domine, speravi.}

IN
thee, O LORD, have I put my trust; let me never be put to confusion, *
but rid me and deliver me in thy righteousness; incline thine ear unto
me, and save me.

2 Be thou my stronghold, whereunto I may alway resort:
* thou hast promised to help me, for thou art my house of defence, and
my castle.

3 Deliver me, O my God, out of the hand of the ungodly,
* out of the hand of the unrighteous and cruel man.

4 For thou, O Lord GOD, art the thing that I long for:
* thou art my hope, even from my youth.

1 reproof
until 1928

21 Thy rebuke hath broken ...
until 1928

5 Through thee have I been holden up ever since I was born: * thou art he
that took me out of my mother's womb: my praise shall be alway1
of thee.

6 I am become as it were a monster unto many, * but my
sure trust is in thee.

7 O let my mouth be filled with thy praise, * that I
may sing of thy glory and honour all the day long.

8 Cast me not away in the time of age; * forsake me not
when my strength faileth me.

1 always in
1789 Book

9 For mine enemies speak against me; * and they that lay wait for my soul
take their counsel together, saying,

10 God hath forsaken him; * persecute him, and take him,
for there is none to deliver him.

11 Go not far from me, O God; * my God, haste thee to
help me.

12 Let them be confounded and perish that are against
my soul; * let them be covered with shame and dishonour that seek to do
me evil.

13 As for me, I will patiently abide alway, * and will
praise thee more and more.

14 My mouth shall daily speak of thy righteousness and
salvation; * for I know no end thereof.

15 I will go forth in the strength of the Lord GOD, * and will make mention of thy righteousness only.

16 Thou, O God, hast taught me from my youth up until
now; * therefore will I tell of thy wondrous works.

17 Forsake me not, O God, in mine old age, when I am
gray-headed, * until I have showed thy strength unto this generation, and
thy power to all them that are yet for to come.

18 Thy righteousness, O God, is very high, * and great
things are they that thou hast done: O God, who is like unto thee!

19 O what great troubles and adversities hast thou showed
me ! and yet didst thou turn and refresh me; * yea, and broughtest me from
the deep of the earth again.

20 Thou hast brought me to great honour, * and comforted
me on every side:

21 Therefore will I praise thee, and thy faithfulness,
O God, playing upon an instrument of music: * unto thee will I sing upon
the harp, O thou Holy One of Israel.

9. For mine enemies speak
against me, and they that lay wait for my soul take their counsel together,
saying : God hath forsaken him; persecute him, and take him, for there is
none to deliver him.
until 1928, and remaining verses numbered one lower

22 My lips will be glad2 when I sing unto thee; * and so will
my soul whom thou hast delivered.

23 My tongue also shall talk of thy righteousness all
the day long; * for they are confounded and brought unto shame that seek
to do me evil.

Psalm
72. \emph{Deus, judicium.}

GIVE
the King thy judgments, O God, * and thy righteousness unto the King's
son.

2 Then shall he judge thy people according unto right,
* and defend the poor.

3 The mountains also shall bring peace, * and the little
hills righteousness unto the people.

4 He shall keep the simple folk by their right, * defend
the children of the poor, and punish the wrong doer.

5 They shall fear thee, as long as the sun and moon
endureth, * from one generation to another.

2 fain
until 1928

6 He shall come down like the rain upon the mown grass, * even as the drops
that water the earth.

7 In his time shall the righteous flourish; * yea, and
abundance of peace, so long as the moon endureth. 6 ... like rain unto a fleece
of wool :
until 1928

8 His dominion shall be also from the one sea to the other, * and from the
River3 unto the world's end.

9 They that dwell in the wilderness shall kneel before
him; * his enemies shall lick the dust. \\
3 flood
until 1928

10 The kings of Tarshish4 and of the isles shall give presents;
* the kings of Arabia and Saba shall bring gifts.

11 All kings shall fall down before him; * all nations
shall do him service.

12 For he shall deliver the poor when he crieth; * the
needy also, and him that hath no helper.

13 He shall be favourable to the simple and needy, *
and shall preserve the souls of the poor.

14 He shall deliver their souls from falsehood and wrong;
* and dear shall their blood be in his sight.

15 He shall live, and unto him shall be given of the
gold of Arabia; * prayer shall be made ever unto him, and daily shall he
be praised.

4 Tharsis
until 1928

16 There shall be an heap of corn in the earth, high upon the hills ;
the fruit thereof shall shake like Lebanon: * and they of the city shall
flourish like grass upon the earth.

17 His Name shall endure for ever; his Name shall remain
under the sun among4 the posterities, which shall be blessed
in5 him; * and all the nations6 shall praise him.

18 Blessed be the Lord God, even the God of Israel,
* which only doeth wondrous things;

19 And blessed be the Name of his majesty for ever:
* and all the earth shall be filled with his majesty. Amen, Amen.

BOOK
III. 

Evening Prayer.

Psalm
73. \emph{Quam bonus Israel!}

TRULY
God is loving unto Israel: * even unto such as are of a clean heart.

2 Nevertheless, my feet were almost gone, * my treadings
had well-nigh slipt.

3 And why? I was grieved at the wicked: * I do also
see the ungodly in such prosperity.

4 For they are in no peril of death; * but are lusty
and strong.

5 They come in no misfortune like other folk; * neither
are they plagued like other men.

16. ... shall shake like Libanus, and
shall be green in the city like grass upon the earth.
until 1928

4 amongst 1793
- 1892

5 through
until 1928 

6
heathen
until 1928

6 And this is the cause that they are so holden with pride, * and
cruelty covereth them as a garment.

7 Their eyes swell with fatness, * and they do even what
they lust. 6 ... with pride : and overwhelmed
with cruelty.
until 1928

8 They corrupt other, and speak of wicked blasphemy; * their talking is
against the Most1 High.

9 For they stretch forth their mouth unto the heaven,
* and their tongue goeth through the world.

10 Therefore fall the people unto them, * and thereout
suck they no small advantage.

11 Tush, say they, how should God perceive it? * is there
knowledge in the Most High? \\
1 most (i.
e., not capitalized) in the 1892 Book

12 Lo, these are the ungodly, * these prosper in the world, and these have
riches in possession:

13 And I said, Then have I cleansed my heart in vain,
* and washed my2 hands in innocency.

14 All the day long have I been punished, * and chastened
every morning.

15 Yea, and I had almost said even as they; * but lo,
then I should have condemned the generation of thy children.

16 Then thought I to understand this; * but it was too
hard for me,

17 Until I went into the sanctuary of God: * then understood
I the end of these men;

18 Namely, how thou dost set them in slippery places,
* and castest them down, and destroyest them.

19 O how suddenly do they consume, * perish, and come
to a fearful end!

20 Yea, even like as a dream when one awaketh; * so shalt
thou make their image to vanish out of the city.

21 Thus my heart was grieved, * and it went even through
my reins.

22 So foolish was I, and ignorant, * even as it were
a beast before thee.

23 Nevertheless, I am alway by thee; * for thou hast
holden me by my right hand.

24 Thou shalt guide me with thy counsel, * and after
that receive me with glory.

25 Whom have I in heaven but thee? * and there is none
upon earth that I desire in comparison of thee.

26 My flesh and my heart faileth; * but God is the strength
of my heart, and my portion for ever.

12. Lo, these are the ungodly,
these prosper in the world, and these have riches in possession : and I
said, Then have I cleansed my heart in vain, and washed my2 hands
in innocency.
until 1928, and remaining verses numbered one lower.

2 mine until
1892

27 For
lo, they that forsake thee shall perish; * thou hast destroyed all them
that are unfaithful unto thee.

28 But it is good for me to hold me fast by God, to
put my trust in the Lord GOD, * and to speak of all thy works in the gates
of the daughter of Sion.

Psalm
74. \emph{Ut quid, Deus?}

O
GOD, wherefore art thou absent from us so long? * why is thy wrath so
hot against the sheep of thy pasture?

2 O think upon thy congregation, * whom thou hast purchased,
and redeemed of old.

3 Think upon the tribe of thine inheritance, * and
Mount Sion, wherein thou hast dwelt.

4 Lift up thy feet, that thou mayest utterly destroy
every enemy, * which hath done evil in thy sanctuary.

5 Thine adversaries roar in the midst of thy congregations,
* and set up their banners for tokens.

6 He that hewed timber afore out of the thick trees,
* was known to bring it to an excellent work.

7 But now they break down all the carved work thereof
* with axes and hammers.

8 They have set fire upon thy holy places, * and have
defiled the dwelling-place of thy Name, even unto the ground.

9 Yea, they said in their hearts, Let us make havoc
of them altogether: * thus have they burnt up all the houses of God in
the land.

10 We see not our tokens; there is not one prophet
more; * no, not one is there among us, that understandeth any more.

[26] ... them that commit
fornication against thee.
until 1928

11 O
God, how long shall the adversary do this dishonour? * shall the enemy
blaspheme thy Name for ever?

12 Why withdrawest thou thy hand? * why pluckest thou
not thy right hand out of thy bosom to consume the enemy?

13 For God is my King of old; * the help that is done
upon earth, he doeth it himself.

14 Thou didst divide the sea through thy power; * thou
brakest the heads of the dragons in the waters.

15 Thou smotest the heads of leviathan in pieces, *
and gavest him to be meat for the people of the wilderness.

16 Thou broughtest out fountains and waters out of
the hard rocks; * thou driedst up mighty waters.

17 The day is thine, and the night is thine; * thou
hast prepared the light and the sun.

18 Thou hast set all the borders of the earth; * thou
hast made summer and winter.

19 Remember this, O LORD, how the enemy hath rebuked;
* and how the foolish people hath blasphemed thy Name.

20 O deliver not the soul of thy turtle-dove unto the
multitude of the enemies; * and forget not the congregation of the poor
for ever.

21 Look upon the covenant; * for all the earth is full
of darkness and cruel habitations.

22 O let not the simple go away ashamed; * but let
the poor and needy give praise unto thy Name.

23 Arise, O God, maintain thine own cause; * remember
how the foolish man blasphemeth thee daily.

24 Forget not the voice of thine enemies: * the presumption
of them that hate thee increaseth ever more and more. 

The
Fifteenth Day.

Morning
Prayer.

Psalm
75. \emph{Confitebimur tibi.}

UNTO
thee, O God, do we give thanks; * yea, unto thee do we give thanks.

2 Thy Name also is so nigh; * and that do thy wondrous
works declare. 11 ... this dishonour : how
long shall the enemy blaspheme thy Name? for ever?
until 1928

3 In the appointed time, saith God, * I shall judge according unto right.

4 The earth is weak, and all the inhabiters thereof:
* I bear up the pillars of it.

5 I said unto the fools, Deal not so madly; * and to
the ungodly, Set not up your horn.

6 Set not up your horn on high, * and speak not with
a stiff neck.

7 For promotion cometh neither from the east, nor from
the west, * nor yet from the south.

8 And why? God is the Judge; * he putteth down one, and
setteth up another.

9 For in the hand of the LORD there is a cup, and the
wine is red; * it is full mixt, and he poureth out of the same.

10 As for the dregs thereof, * all the ungodly of the
earth shall drink them, and suck them out.

11 But I will talk of the God of Jacob, * and praise
him for ever.

12 All the horns of the ungodly also will I break, *
and the horns of the righteous shall be exalted.

3 When I receive the congregation
: I shall judge ...
until 1928

Psalm
76. \emph{Notus in Judæa.}

IN
Judah is God known; * his Name is great in Israel.

2 At Salem is his tabernacle, * and his dwelling in
Sion.

3 There brake he the arrows of the bow, * the shield,
the sword, and the battle.

4 Thou art glorious in might, * when thou comest from the hills of the
robbers.

5 The proud are robbed, they have slept their sleep;
* and all the men whose hands were mighty have found nothing.

6 At thy rebuke, O God of Jacob, * both the chariot
and horse are fallen.

7 Thou, even thou art to be feared; * and who may stand
in thy sight when thou art angry?

8 Thou didst cause thy judgment to be heard from heaven;
* the earth trembled, and was still,

9 When God arose to judgment, * and to help all the
meek upon earth.

10 The fierceness of man shall turn to thy praise;
* and the fierceness of them shalt thou refrain.

11 Promise unto the LORD your God, and keep it, all
ye that are round about him; * bring presents unto him that ought to be
feared.

12 He shall refrain the spirit of princes, * and is
wonderful among the kings of the earth.

Psalm
77. \emph{Voce mea ad Dominum.}

I
WILL cry unto God with my voice; * even unto God will I cry with my voice,
and he shall hearken unto me. 4. Thou art of more honour
and might : than the hills of the robbers.
until 1928

2 In the time of my trouble I sought the Lord: * I stretched forth my hands
unto him, and ceased not in the night season; my soul refused comfort.

3 When I am in heaviness, I will think upon God; * when
my heart is vexed, I will complain.

4 Thou holdest mine eyes waking: * I am so feeble that
I cannot speak.

5 I have considered the days of old, * and the years
that are past. 2. ... the Lord : my sore
ran and ceased not in the night-season ; ...
until 1928

6 I call to remembrance my song, * and in the night I commune with mine
own heart, and search out my spirit1.

7 Will the Lord absent himself for ever? * and will he
be no more intreated?

8 Is his mercy clean gone for ever? * and is his promise
come utterly to an end for evermore?

9 Hath God forgotten to be gracious? * and will he shut
up his loving-kindness in displeasure?

10 And I said, It is mine own infirmity; * but I will
remember the years of the right hand of the Most Highest.

11 I will remember the works of the LORD, * and call
to mind thy wonders of old time.

12 I will think also of all thy works, * and my talking
shall be of thy doings.

13 Thy way, O God, is holy: * who is so great a God as
our God?

1 spirits
until 1845

14 Thou art the God that doest2 wonders, * and hast declared
thy power among the peoples3.

15 Thou hast mightily delivered thy people, * even
the sons of Jacob and Joseph.

16 The waters saw thee, O God, the waters saw thee,
and were afraid; * the depths also were troubled.

17 The clouds poured out water, the air thundered,
* and thine arrows went abroad.

18 The voice of thy thunder was heard round about:
* the lightnings shone upon the ground; the earth was moved, and shook
withal.

19 Thy way is in the sea, and thy paths in the great
waters, * and thy footsteps are not known.

20 Thou leddest thy people like sheep, * by the hand
of Moses and Aaron. 

Evening
Prayer.

Psalm
78. \emph{Attendite, popule.}

HEAR
my law, O my people; * incline your ears unto the words of my mouth.

2 I will open my mouth in a parable; * I will declare
hard sentences of old;

3 Which we have heard and known, * and such as our
fathers have told us;

4 That we should not hide them from the children of
the generations to come; * but to show the honour of the LORD,
his mighty and wonderful works that he hath done.

5 He made a covenant with Jacob, and gave Israel a
law, * which he commanded our forefathers to teach their children;

6 That their posterity might know it, * and the children
which were yet unborn;

7 To the intent that when they came up, * they might
show their children the same;

8 That they might put their trust in God; * and not
to forget the works of God, but to keep his commandments;

9 And not to be as their forefathers, a faithless and
stubborn generation; * a generation that set not their heart aright, and
whose spirit clave not stedfastly unto God;

10 Like as the children of Ephraim; * who being harnessed,
and carrying bows, turned themselves back in the day of battle.

11 They kept not the covenant of God, * and would not
walk in his law;

12 But forgat what he had done, * and the wonderful
works that he had showed for them.

13 Marvellous things did
he in the sight of our forefathers, in the land of Egypt, * even in the
field of Zoan.

14 He divided the sea, and let them go through; * he
made the waters to stand on an heap.

15 In the day-time also he led them with a cloud, *
and all the night through with a light of fire.

16 He clave the hard rocks in the wilderness, * and
gave them drink thereof, as it had been out of the great depth.

17 He brought waters out of the stony rock, * so that
it gushed out like the rivers.

18 Yet for all this they sinned more against him, *
and provoked the Most Highest in the wilderness.

19 They tempted God in their hearts, * and required
meat for their lust.

20 They spake against God also, saying, * Shall God
prepare a table in the wilderness?

21 He smote the stony rock indeed, that the water gushed
out, and the streams flowed withal; * but can he give bread also, or provide
flesh for his people?

22 When the LORD heard this,
he was wroth; * so the fire was kindled in Jacob, and there came up heavy
displeasure against Israel;

23 Because they believed not in God, * and put not
their trust in his help.

24 So he commanded the clouds above, * and opened the
doors of heaven.

25 He rained down manna also upon them for to eat,
* and gave them food from heaven.

26 So man did eat angels' food; * for he sent them
meat enough.

27 He caused the east-wind to blow under heaven; *
and through his power he brought in the southwest-wind.

28 He rained flesh upon them as thick as dust, * and
feathered fowls like as the sand of the sea.

29 He let it fall among their tents, * even round about
their habitation.

30 So they did eat, and were well filled; for he gave
them their own desire: * they were not disappointed of their lust.

31 But while the meat was yet in their mouths, the
heavy wrath of God came upon them, and slew the wealthiest of them; *
yea, and smote down the chosen men that were in Israel.

32 But for all this they sinned yet more, * and believed
not his wondrous works.

33 Therefore their days did he consume in vanity, *
and their years in trouble.

34 When he slew them, they sought him, * and turned
them early, and inquired after God.

35 And they remembered that God was their strength,
* and that the High God was their redeemer.

36 Nevertheless, they did but flatter him with their
mouth, * and dissembled with him in their tongue.

37 For their heart was not whole with him, * neither
continued they stedfast in his covenant.

38 But he was so merciful, that he forgave their misdeeds,
* and destroyed them not.

39 Yea, many a time turned he his wrath away, * and
would not suffer his whole displeasure to arise.

40 For he considered that they were but flesh, * and
that they were even a wind that passeth away, and cometh not again.

41 Many a time did they
provoke him in the wilderness, * and grieved him in the desert.

2 doeth
until 1845\\
3 people
until 1928

42 They turned back, and tempted God, * and provoked1 the Holy
One in Israel.

43 They thought not of his hand, * and of the day when
he delivered them from the hand of the enemy;

44 How he had wrought his miracles in Egypt, * and his
wonders in the field of Zoan.

45 He turned their waters into blood, * so that they
might not drink of the rivers. \\
1 moved
until 1928

46 He sent flies2 among them, and devoured them up; * and frogs
to destroy them.

47 He gave their fruit unto the caterpillar, * and their
labour unto the grasshopper.

48 He destroyed their vines with hailstones, * and their
mulberry-trees with the frost.

49 He smote their cattle also with hailstones, * and
their flocks with hot thunderbolts.

50 He cast upon them the furiousness of his wrath, anger,
displeasure, and trouble: * and sent evil angels among them.

51 He made a way to his indignation, and spared not their
soul from death; * but gave their life over to the pestilence;

52 And smote all the firstborn in Egypt, * the most principal
and mightiest in the dwellings of Ham.

53 But as for his own people, he led them forth like
sheep, * and carried them in the wilderness like a flock.

54 He brought them out safely, that they should not fear,
* and overwhelmed their enemies with the sea.

55 And brought them within the borders of his sanctuary,
* even to this mountain, which he purchased with his right hand.

56 He cast out the heathen also before them, * caused
their land to be divided among them for an heritage, and made the tribes
of Israel to dwell in their tents.

2 lice
until 1928

57 Yet3 they tempted and displeased the Most High God, * and
kept not his testimonies.

58 They4 turned their backs, and fell away
like their forefathers; * starting aside like a broken bow.

59 For they grieved him with their hill-altars, * and
provoked him to displeasure with their images.

60 When God heard this, he was wroth, * and took sore
displeasure at Israel;

3 So
until 1928

4 But
until 1928

61 So that he forsook the tabernacle in Shiloh5, * even the tent
that he had pitched among men.

62 He delivered their power into captivity, * and their
beauty into the enemy's hand.

63 He gave his people over also unto the sword, * and
was wroth with his inheritance. 5 Silo
until 1928

64 The fire consumed their young men, * and their maidens were not given
in6 marriage.

65 Their priests were slain with the sword, * and there
were no widows to make lamentation.

66 So the Lord awaked as one out of sleep, * and like
a giant refreshed with wine. \\
6 to
until 1928

67 He drave his enemies backward, * and put them to a perpetual shame.

68 He refused the tabernacle of Joseph, * and chose not
the tribe of Ephraim;

69 But chose the tribe of Judah, * even the hill of Sion
which he loved.

70 And there he built his temple on high, * and laid
the foundation of it like the ground which he hath made continually.

71 He chose David also his servant, * and took him away
from the sheep-folds: 67 He smote his enemies in
the hinder parts : and put them ...
until 1928

72 As he was following the ewes with their young he took him, * that he
might feed Jacob his people, and Israel his inheritance.

73 So he fed them with a faithful and true heart, *
and ruled them prudently with all his power. 

The
Sixteenth Day.

Morning
Prayer.

Psalm
79. \emph{Deus, venerunt.}

O
GOD, the heathen are come into thine inheritance; * thy holy temple have
they defiled, and made Jerusalem an heap of stones.

2 The dead bodies of thy servants have they given to
be meat unto the fowls of the air, * and the flesh of thy saints unto
the beasts of the land.

3 Their blood have they shed like water on every side
of Jerusalem, * and there was no man to bury them.

4 We are become an open shame to our enemies, * a very
scorn and derision unto them that are round about us.

5 LORD, how long wilt thou be angry? * shall thy jealousy
burn like fire for ever?

6 Pour out thine indignation upon the heathen that
have not known thee; * and upon the kingdoms that have not called upon
thy Name.

7 For they have devoured Jacob, * and laid waste his
dwelling-place.

8 O remember not our old sins, but have mercy upon
us, and that soon; * for we are come to great misery.

9 Help us, O God of our salvation, for the glory of
thy Name: * O deliver us, and be merciful unto our sins, for thy Name's
sake.

10 Wherefore do the heathen say, * Where is now their
God?

11 O let the vengeance of thy servants' blood that
is shed, * be openly showed upon the heathen, in our sight.

12 O let the sorrowful sighing of the prisoners come
before thee; * according to the greatness of thy power, preserve thou
those that are appointed to die.

13 And for the blasphemy wherewith our neighbours have
blasphemed thee, * reward thou them, O Lord, sevenfold into their bosom.

14 So we, that are thy people, and sheep of thy pasture,
shall give thee thanks for ever, * and will alway be showing forth thy
praise from generation to generation.

72 ... the ewes great with
their young ones, he took ...
until 1928

Psalm
80. \emph{Qui regis Israel.}

HEAR,
O thou Shepherd of Israel, thou that leadest Joseph like a flock1;
* show thyself also, thou that sittest upon the Cherubim2.

2 Before Ephraim, Benjamin, and Manasseh3,
* stir up thy strength, and come and help us.

3 Turn us again, O God; * show the light of thy countenance,
and we shall be whole.

4 O LORD God of hosts, * how long wilt thou be angry
with thy people that prayeth?

5 Thou feedest them with the bread of tears, * and
givest them plenteousness of tears to drink.

6 Thou hast made us a very strife unto our neighbours,
* and our enemies laugh us to scorn.

7 Turn us again, thou God of hosts; * show the light
of thy countenance, and we shall be whole.

8 Thou hast brought a vine out of Egypt; * thou hast
cast out the heathen, and planted it.

9 Thou madest room for it; * and when it had taken
root, it filled the land.

10 The hills were covered with the shadow of it, *
and the boughs thereof were like the goodly cedar-trees.

1 sheep
until 1928

2 Cherubims
until 1793

3
Mannases
until 1928

11 She stretched out her branches unto the sea, * and her boughs unto
the River4.

12 Why hast thou then broken down her hedge, * that
all they that go by pluck off her grapes?

13 The wild boar out of the wood doth root it up, *
and the wild beasts of the field devour it.

14 Turn thee again, thou God of hosts, look down from
heaven, * behold, and visit this vine;

15 And the place of the vineyard that thy right hand
hath planted, * and the branch that thou madest so strong for thyself.

16 It is burnt with fire, and cut down; * and they
shall perish at the rebuke of thy countenance.

17 Let thy hand be upon the man of thy right hand,
* and upon the son of man, whom thou madest so strong for thine own self.

18 And so will not we go back from thee: * O let us
live, and we shall call upon thy Name.

19 Turn us again, O LORD God of hosts; * show the light
of thy countenance, and we shall be whole.

Psalm
81. \emph{Exultate Deo.}

SING
we merrily unto God our strength; * make a cheerful noise unto the God
of Jacob.

2 Take the psalm, bring hither the tabret, * the merry
harp with the lute.

3 Blow up the trumpet in the new moon, * even in the
time appointed, and upon our solemn feast-day.

4 For this was made a statute for Israel, * and a law
of the God of Jacob.

5 This he ordained in Joseph for a testimony, * when
he came out of the land of Egypt, and had heard a strange language.

6 I eased his shoulder from the burden, * and his hands
were delivered from making the pots.

7 Thou calledst upon me in troubles, and I delivered
thee; * and heard thee what time as the storm fell upon thee.

8 I proved thee also * at the waters of strife.

9 Hear, O my people; and I will assure thee, O Israel,
* if thou wilt hearken unto me,

10 There shall no strange god be in thee, * neither
shalt thou worship any other god.

11 I am the LORD thy God, who brought thee out of the
land of Egypt: * open thy mouth wide, and I shall fill it.

12 But my people would not hear my voice; * and Israel
would not obey me;

13 So I gave them up unto their own hearts' lusts,
* and let them follow their own imaginations.

14 O that my people would have hearkened unto me! *
for if Israel had walked in my ways,

15 I should soon have put down their enemies, * and
turned my hand against their adversaries. \\
4 river (i.
e., not capitallized) until 1928

16 The haters of the LORD should have submitted themselves unto him; *
but their time should have endured for ever.

17 I would have fed them also with the finest wheat-flour;
* and with honey out of the stony rock would I have satisfied thee. 

Evening
Prayer.

Psalm
82. \emph{Deus stetit.}

GOD
standeth in the congregation of princes; * he is a Judge among gods.

2 How long will ye give wrong judgment, * and accept
the persons of the ungodly?

3 Defend the poor and fatherless; * see that such as
are in need and necessity have right.

4 Deliver the outcast and poor; * save them from the
hand of the ungodly.

16 ... should have been found liars
: but their time ...
until 1928

17. He should have fed them ... stony rock should
I have satisfied thee.
until 1928

5 They know not, neither do they understand, but walk on still in darkness:
* all the foundations of the earth are out of course.

6 I have said, Ye are gods, * and ye are all the children
of the Most Highest.

7 But ye shall die like men, * and fall like one of the
princes. 5 They will not be learned,
nor understand, but ...
until 1928

8 Arise, O God, and judge thou the earth; * for thou shalt take all nations1
to thine inheritance.

Psalm
83. \emph{Deus, quis similis?}

HOLD
not thy tongue, O God, keep not still silence: * refrain not thyself,
O God.

2 For lo, thine enemies make a murmuring; * and they
that hate thee have lift up their head.

3 They have imagined craftily against thy people, *
and taken counsel against thy secret ones.

4 They have said, Come, and let us root them out, that
they be no more a people, * and that the name of Israel may be no more
in remembrance.

5 For they have cast their heads together with one
consent, * and are confederate against thee: \\
1 heathen
until 1928

6 The tabernacles of the Edomites, and the Ishmaelites1; * the
Moabites, and Hagarenes2;

7 Gebal, and Ammon, and Amalek; * the Philistines, with
them that dwell at Tyre. \\
1 Ismaelites in
the Engl. Book

2 Hagarens
until 1793

8 Assyria3 also is joined with them; * they have holpen the children
of Lot.

9 But do thou to them as unto the Midianites4;
* unto Sisera, and unto Jabin at the brook of Kishon5;

10 Who perished at Endor, * and became as the dung of
the earth.

3 Assur
until 1928

4 Madianites
until 1822, and in 1892 Book

5 Kison
until 1928

11 Make them and their princes like Oreb and Zeëb; * yea, make all
their princes like as Zebah and Zalmunna;

12 Who say, Let us take to ourselves * the houses of
God in possession. 11 ... Oreb and Zeb : ...
Zeba and Salmana
until 1928

13 O my God, make them like unto the whirling dust, * and as the stubble
before the wind;

14 Like as the fire that burneth up the forest6,
* and as the flame that consumeth the mountains;

15 Pursue7 them even so with thy tempest,
* and make them afraid with thy storm.

16 Make their faces ashamed, O LORD, * that they may
seek thy Name.

17 Let them be confounded and vexed ever more and more;
* let them be put to shame, and perish.

18 And they shall know that thou, whose Name is JEHOVAH,
* art only the Most Highest over all the earth.

Psalm
84. \emph{Quam dilecta!}

O
HOW amiable are thy dwellings, * thou LORD of hosts!

2 My soul hath a desire and longing to enter into the
courts of the LORD; * my heart and my flesh rejoice in the living God.

3 Yea, the sparrow hath found her an house, and the
swallow a nest, where she may lay her young; * even thy altars, O LORD
of hosts, my King and my God.

4 Blessed are they that dwell in thy house; * they
will be alway praising thee.

5 Blessed is the man whose strength is in thee; * in
whose heart are thy ways.

6 Who going through the vale of misery use it for a
well; * and the pools are filled with water.

7 They will go from strength to strength, * and unto
the God of gods appeareth every one of them in Sion.

8 O LORD God of hosts, hear my prayer; * hearken, O
God of Jacob.

9 Behold, O God our defender, * and look upon the face
of thine anointed.

10 For one day in thy courts * is better than a thousand.

11 I had rather be a door-keeper in the house of my
God, * than to dwell in the tents of ungodliness.

12 For the LORD God is a light and defence; * the LORD
will give grace and worship; and no good thing shall he withhold from
them that live a godly life.

13 O LORD God of hosts, * blessed is the man that putteth
his trust in thee.

Psalm
85. \emph{Benedixisti, Domine.}

LORD,
thou art become gracious unto thy land; * thou hast turned away the captivity
of Jacob.

2 Thou hast forgiven the offence of thy people, * and
covered all their sins.

3 Thou hast taken away all thy displeasure, * and turned
thyself from thy wrathful indignation.

4 Turn us then, O God our Saviour, * and let thine
anger cease from us.

5 Wilt thou be displeased at us for ever? * and wilt
thou stretch out thy wrath from one generation to another?

6 Wilt thou not turn again, and quicken us, * that
thy people may rejoice in thee?

7 Show us thy mercy, O LORD, * and grant us thy salvation.

13 ... like unto a wheel : and as ...
until
1928 

6 wood
until 1928

7 Perscute
until
1928

8 I will hearken what the LORD God will say; * for he shall speak peace
unto his people, and to his saints, that they turn not again unto foolishness.

9 For his salvation is nigh them that fear him; * that
glory may dwell in our land.

10 Mercy and truth are met together: * righteousness
and peace have kissed each other.

11 Truth shall flourish out of the earth, * and righteousness
hath looked down from heaven.

12 Yea, the LORD shall show loving-kindness; * and
our land shall give her increase.

13 Righteousness shall go before him, * and shall direct
his going in the way. 

The
Seventeenth Day.

Morning
Prayer.

Psalm
86. \emph{Inclina, Domine.}

BOW
down thine ear, O LORD, and hear me; * for I am poor, and in misery.

2 Preserve thou my soul, for I am holy: * my God, save
thy servant that putteth his trust in thee.

3 Be merciful unto me, O Lord; * for I will call daily
upon thee.

4 Comfort the soul of thy servant; * for unto thee,
O Lord, do I lift up my soul.

5 For thou, Lord, art good and gracious, * and of great
mercy unto all them that call upon thee.

6 Give ear, LORD, unto my prayer, * and ponder the
voice of my humble desires.

7 In the time of my trouble I will call upon thee;
* for thou hearest me.

8 Among the gods there is none like unto thee, O Lord;
* there is not one that can do as thou doest.

9 All nations whom thou hast made shall come and worship
thee, O Lord; * and shall glorify thy Name.

10 For thou art great, and doest wondrous things: *
thou art God alone.

11 Teach me thy way, O LORD, and I will walk in thy
truth: * O knit my heart unto thee, that I may fear thy Name.

12 I will thank thee, O Lord my God, with all my heart;
* and will praise thy Name for evermore.

13 For great is thy mercy toward me; * and thou hast
delivered my soul from the nethermost hell. 8. ... will say concerning
me : for he shall speak peace unto his people, and to his saints, that they
turn not again.
until 1928

14 O God, the proud are risen against me; * and the congregations of violent1
men have sought after my soul, and have not set thee before their eyes.

15 But thou, O Lord God, art full of compassion and
mercy, * long-suffering, plenteous in goodness and truth.

16 O turn thee then unto me, and have mercy upon me;
* give thy strength unto thy servant, and help the son of thine handmaid.

17 Show some token upon me for good; that they who
hate me may see it, and be ashamed, * because thou, LORD, hast holpen
me, and comforted me.

1 naughty
until 1928

Psalm
87. \emph{Fundamenta ejus.}

HER
foundations are upon the holy hills: * the LORD loveth the gates of Sion
more than all the dwellings of Jacob.

2 Very excellent things are spoken of thee, * thou
city of God.

3 I will make mention of Egypt and Babylon, * among them that know me.

4 Behold1, Philistia also; and Tyre, with
Ethiopia; * lo, in Sion were they born.

5 Yea, of Sion it shall be reported, this one and that
one were born in her; * and the Most High shall stablish her.

6 The LORD shall record it, when he writeth up the
peoples; * lo, in Sion were they born.

7 The singers also and trumpeters shall make answer:
* All my fresh springs are in thee.

Psalm
88. \emph{Domine, Deus.}

O
LORD God of my salvation, I have cried day and night before thee: * O
let my prayer enter into thy presence, incline thine ear unto my calling;

3. I will think upon Rahab and Babylon : with them
that know me.

4. Behold, yea1
the Philistines also : and they of Tyre, with the Morians; lo, there was
he born.

5. And of Sion it shall be reported that he was born
in her : and the most High shall stablish her.

6. The LORD shall rehearse it
when he writeth up the people : that he was born there.

7. The singers also and trumpeters shall he rehearse
: All my fresh springs shall be in thee.
until 1928 

1
Behold ye
in the 1789 Book

2 For my soul is full of trouble, * and my life draweth nigh unto the
grave.

3 I am counted as one of them that go down into the
pit, * and I am even as a man that hath no strength;

4 Cast off among the dead, like unto them that are
slain, and lie in the grave, * who are out of remembrance, and are cut
away from thy hand.

5 Thou hast laid me in the lowest pit, * in a place
of darkness, and in the deep.

6 Thine indignation lieth hard upon me, * and thou
hast vexed me with all thy storms.

7 Thou hast put away mine acquaintance far from me,
* and made me to be abhorred of them.

8 I am so fast in prison * that I cannot get forth.

9 My sight faileth for very trouble; * LORD, I have
called daily upon thee, I have stretched forth my hands unto thee.

10 Dost thou show wonders among the dead? * or shall
the dead rise up again, and praise thee?

11 Shall thy loving-kindness be showed in the grave?
* or thy faithfulness in destruction?

12 Shall thy wondrous works be known in the dark? *
and thy righteousness in the land where all things are forgotten?

13 Unto thee have I cried, O LORD; * and early shall
my prayer come before thee.

14 LORD, why abhorrest thou my soul, * and hidest thou
thy face from me?

15 I am in misery, and like unto him that is at the
point to die; * even from my youth up, thy terrors have I suffered with
a troubled mind.

16 Thy wrathful displeasure goeth over me, * and the
fear of thee hath undone me.

17 They came round about me daily like water, * and
compassed me together on every side.

18 My lovers and friends hast thou put away from me,
* and hid mine acquaintance out of my sight. 

Evening
Prayer.

Psalm
89. \emph{Misericordias Domini.}

MY
song shall be alway of the loving-kindness of the LORD; * with my mouth
will I ever be showing thy truth from one generation to another.

2 For I have said, Mercy shall be set up for ever;
* thy truth shalt thou stablish in the heavens.

3 I have made a covenant with my chosen; * I have sworn
unto David my servant:

4 Thy seed will I stablish for ever, * and set up thy
throne from one generation to another.

5 O LORD, the very heavens shall praise thy wondrous
works; * and thy truth in the congregation of the saints.

6 For who is he among the clouds, * that shall be compared
unto the LORD?

7 And what is he among the gods, * that shall be like
unto the LORD?

8 God is very greatly to be feared in the council of
the saints, * and to be had in reverence of all them that are round about
him.

9 O Lord God of hosts, who is like unto thee? * thy
truth, most mighty LORD, is on every side.

10 Thou rulest the raging of the sea; * thou stillest
the waves thereof when they arise.

11 Thou hast subdued Egypt, and destroyed it; * thou
hast scattered thine enemies abroad with thy mighty arm.

12 The heavens are thine, the earth also is thine;
* thou hast laid the foundation of the round world, and all that therein
is.

13 Thou hast made the north and the south; * Tabor
and Hermon shall rejoice in thy Name.

14 Thou hast a mighty arm; * strong is thy hand, and
high is thy right hand.

15 Righteousness and equity are the habitation of thy
seat; * mercy and truth shall go before thy face.

16 Blessed is the people, O LORD, that can rejoice
in thee; * they shall walk in the light of thy countenance.

17 Their delight shall be daily in thy Name; * and
in thy righteousness shall they make their boast.

18 For thou art the glory of their strength, * and
in thy loving-kindness thou shalt lift up our horns.

19 For the LORD is our defence; * the Holy One of Israel
is our King.

20 Thou spakest sometime in visions unto thy saints,
and saidst, * I have laid help upon one that is mighty, I have exalted
one chosen out of the people.

21 I have found David my servant; * with my holy oil
have I anointed him.

22 My hand shall hold him fast, * and my arm shall
strengthen him.

23 The enemy shall not be able to do him violence;
* the son of wickedness shall not hurt him.

24 I will smite down his foes before his face, * and
plague them that hate him.

25 My truth also and my mercy shall be with him; *
and in my Name shall his horn be exalted.

26 I will set his dominion also in the sea, * and his
right hand in the floods.

27 He shall call me, Thou art my Father, * my God,
and my strong salvation.

28 And I will make him my firstborn, * higher than
the kings of the earth.

29 My mercy will I keep for him for evermore, * and
my covenant shall stand fast with him.

30 His seed also will I make to endure for ever, *
and his throne as the days of heaven.31 But if his children forsake my
law, * and walk not in my judgments;

32 If they break my statutes, and keep not my commandments;
* I will visit their offences with the rod, and their sin with scourges.

33 Nevertheless, my loving-kindness will I not utterly
take from him, * nor suffer my truth to fail.

34 My covenant will I not break, nor alter the thing
that is gone out of my lips: * I have sworn once by my holiness, that
I will not fail David.

35 His seed shall endure for ever, * and his throne
is like as the sun before me.

36 He shall stand fast for evermore as the moon, *
and as the faithful witness in heaven.

37 But thou hast abhorred and forsaken thine anointed,
* and art displeased at him.

38 Thou hast broken the covenant of thy servant, *
and cast his crown to the ground.

39 Thou hast overthrown all his hedges, * and broken
down his strongholds.

40 All they that go by spoil him, * and he is become
a reproach to his neighbours.

41 Thou hast set up the right hand of his enemies,
* and made all his adversaries to rejoice.

42 Thou hast taken away the edge of his sword, * and
givest him not victory in the battle.

43 Thou hast put out his glory, * and cast his throne
down to the ground.

44 The days of his youth hast thou shortened, * and
covered him with dishonour.

45 LORD, how long wilt thou hide thyself? for ever?
* and shall thy wrath burn like fire?

46 O remember how short my time is; * wherefore hast
thou made all men for nought? 2. ... life draweth nigh unto
hell.

3. I am counted as one of them that go down into the
pit : and I have been even as a man that hath no strength.

4. Free among the dead, like unto them that are wounded, and lie in the
grave : ....
until 1928

47 What man is he that liveth, and shall not see death? * and shall he deliver
his soul from the power of the grave?

48 Lord, where are thy old loving-kindnesses, * which
thou swarest unto David in thy truth?

49 Remember, Lord, the rebuke that thy servants have,
* and how I do bear in my bosom the rebukes of many people; 47 ... his soul from the land
of hell.
until 1928

50 Wherewith thine enemies have blasphemed thee, * and slandered the footsteps
of thine anointed.

51
Praised be the LORD for evermore. * Amen, and Amen.

BOOK
IV.

The
Eighteenth Day.

Morning
Prayer.

Psalm
90. \emph{Domine, refugium.}

LORD,
thou hast been our refuge, * from one generation to another.

2 Before the mountains were brought forth, or ever
the earth and the world were made, * thou art God from everlasting, and
world without end.

3 Thou turnest man to destruction; * again thou sayest,
Come again, ye children of men.

50. Wherewith thine enemies
have blasphemed thee,

and slandered the footsteps of thine Anointed :

Praised be the LORD for evermore. Amen, and Amen.
until
1928

4 For a thousand years in thy sight are but as yesterday when it is past,
* and as a watch in the night.

5 As soon as thou scatterest them they are even as a
sleep; * and fade away suddenly like the grass.

6 In the morning it is green, and groweth up; * but in
the evening it is cut down, dried up, and withered.

7 For we consume away in thy displeasure, * and are afraid
at thy wrathful indignation.

8 Thou hast set our misdeeds before thee; * and our secret
sins in the light of thy countenance.

9 For when thou art angry all our days are gone: * we
bring our years to an end, as it were a tale that is told.

10 The days of our age are threescore years and ten;
and though men be so strong that they come to fourscore years, * yet is
their strength then but labour and sorrow; so soon passeth it away, and
we are gone.

4 ... but as yesterday :
seeing that is past as a watch in the night. until
1928

11 But who regardeth the power of thy wrath? * or feareth aright thy indignation?

12 So1 teach us to number our days, * that
we may apply our hearts unto wisdom.

13 Turn thee again, O LORD, at the last, * and be gracious
unto thy servants.

14 O satisfy us with thy mercy, and that soon: * so
shall we rejoice and be glad all the days of our life.

15 Comfort us again now after the time that thou hast
plagued us; * and for the years wherein we have suffered adversity.

16 Show thy servants thy work, * and their children
thy glory.

17 And the glorious majesty of the LORD our God be
upon us: * prosper thou the work of our hands upon us; O prosper thou
our handy-work.

Psalm
91. \emph{Qui habitat.}

WHOSO
dwelleth under the defence of the Most High, * shall abide under the shadow
of the Almighty.

2 I will say unto the LORD, Thou art my hope, and my
stronghold; * my God, in him will I trust.

3 For he shall deliver thee from the snare of the hunter,
* and from the noisome pestilence.

4 He shall defend thee under his wings, and thou shalt
be safe under his feathers; * his faithfulness and truth shall be thy
shield and buckler.

5 Thou shalt not be afraid for any terror by night,
* nor for the arrow that flieth by day;

6 For the pestilence that walketh in darkness, * nor
for the sickness that destroyeth in the noon-day.

7 A thousand shall fall beside thee, and ten thousand
at thy right hand; * but it shall not come nigh thee.

8 Yea, with thine eyes shalt thou behold, * and see
the reward of the ungodly.

9 For thou, LORD, art my hope; * thou hast set thine
house of defence very high.

10 There shall no evil happen unto thee, * neither
shall any plague come nigh thy dwelling.

11 For he shall give his angels charge over thee, *
to keep thee in all thy ways.

12 They shall bear thee in their hands, * that thou
hurt not thy foot against a stone.

13 Thou shalt go upon the lion and adder: * the young
lion and the dragon shalt thou tread under thy feet.

14 Because he hath set his love upon me, therefore
will I deliver him; * I will set him up, because he hath known my Name.

15 He shall call upon me, and I will hear him; * yea,
I am with him in trouble; I will deliver him, and bring him to honour.

16 With long life will I satisfy him, * and show him
my salvation.

Psalm
92. \emph{Bonum est confiteri.}

IT
is a good thing to give thanks unto the LORD, * and to sing praises unto
thy Name, O Most Highest; 2 To tell of thy loving-kindness
early in the morning, * and of thy truth in the night season;

3 Upon an instrument of ten strings, and upon the lute;
* upon a loud instrument, and upon the harp.

4 For thou, LORD, hast made me glad through thy works;
* and I will rejoice in giving praise for the operations of thy hands.

5 O LORD, how glorious are thy works! * thy thoughts
are very deep.

6 An unwise man doth not well consider this, * and
a fool doth not understand it.

7 When the ungodly are green as the grass, and when
all the workers of wickedness do flourish, * then shall they be destroyed
for ever; but thou, LORD, art the Most1 Highest for evermore.

11 ... of thy wrath : for
even thereafter as a man feareth, so is thy displeasure. until
1928 

1 O in
1892 Book

1 most
(i.e., not capitallized) in 1892 Book

8 For lo, thine enemies, O LORD, lo, thine enemies
shall perish; * and all the workers of wickedness shall be destroyed.

9 But my2 horn shall be exalted like the horn of an unicorn;
* for I am anointed with fresh oil.

10 Mine eye also shall see his lust of mine enemies,
* and mine ear shall hear his desire of the wicked that arise up against
me.

11 The righteous shall flourish like a palm-tree, *
and shall spread abroad like a cedar in Lebanon.

12 Such as are planted in the house of the LORD, *
shall flourish in the courts of the house of our God.

13 They also shall bring forth more fruit in their
age, * and shall be fat and well-liking;

14 That they may show how true the LORD my strength
is, * and that there is no unrighteousness in him. 

Evening
Prayer.

Psalm
93. \emph{Dominus regnavit.}

THE
LORD is King, and hath put on glorious apparel; * the LORD hath put on
his apparel, and girded himself with strength.

2 He hath made the round world so sure, * that it cannot
be moved.

3 Ever since the world began, hath thy seat been prepared:
* thou art from everlasting.

4 The floods are risen, O LORD, the floods have lift
up their voice; * the floods lift up their waves.

5 The waves of the sea are mighty, and rage horribly;
* but yet the LORD, who dwelleth on high, is mightier.

6 Thy testimonies, O LORD, are very sure: * holiness
becometh thine house for ever.

Psalm
94. \emph{Deus ultionum.}

O
LORD God, to whom vengeance belongeth, * thou God, to whom vengeance belongeth,
show thyself.

2 Arise, thou Judge of the world, * and reward the
proud after their deserving.

3 LORD, how long shall the ungodly, * how long shall
the ungodly triumph?

4 How long shall all wicked doers speak so disdainfully,
* and make such proud boasting?

5 They smite down thy people, O LORD, * and trouble
thine heritage.

6 They murder the widow and the stranger, * and put
the fatherless to death.

7 And yet they say, Tush, the LORD shall not see, *
neither shall the God of Jacob regard it.

8 Take heed, ye unwise among the people: * O ye fools,
when will ye understand?

9 He that planted the ear, shall he not hear? * or
he that made the eye, shall he not see? 2
mine in
1789 Book

10 Or he that instructeth1 the heathen, * it is he that teacheth
man knowledge; shall not he punish?

11 The LORD knoweth the thoughts of man, * that they
are but vain.

12 Blessed is the man whom thou chastenest, O LORD, *
and teachest him in thy law;

13 That thou mayest give him patience in time of adversity,
* until the pit be digged up for the ungodly.

14 For the LORD will not fail his people; * neither will
he forsake his inheritance;

15 Until righteousness turn again unto judgment: * all
such as are true in heart shall follow it.

16 Who will rise up with me against the wicked? * or
who will take my part against the evil doers?

17 If the LORD had not helped me, * it had not failed,
but my soul had been put to silence.

1 nurtureth until
1928

18 But when I said, My foot hath slipt2; * thy mercy, O LORD,
held me up.

19 In the multitude of the sorrows that I had in my heart,
* thy comforts have refreshed my soul. 2 slipped in
1789 Book

20 Wilt thou have any thing to do with the throne3 of wickedness,
* which imagineth mischief as a law?

21 They gather them together against the soul of the
righteous, * and condemn the innocent blood.

22 But the LORD is my refuge, * and my God is the strength
of my confidence.

23 He shall recompense them their wickedness, and destroy
them in their own malice; * yea, the LORD our God shall destroy them.

The Nineteenth
Day.

Morning
Prayer.

Psalm
95. \emph{Venite, exultemus.}

O
COME, let us sing unto the LORD; * let us heartily rejoice in the strength
of our salvation.

2 Let us come before his presence with thanksgiving;
* and show ourselves glad in him with psalms.

3 For the LORD is a great God; * and a great King above
all gods.

4 In his hand are all the corners of the earth; * and
the strength of the hills is his also.

5 The sea is his, and he made it; * and his hands prepared
the dry land.

6 O come, let us worship and fall down, * and kneel
before the LORD our Maker.

7 For he is the Lord our God; * and we are the people
of his pasture, and the sheep of his hand.

8 To-day if ye will hear his voice, harden not your
hearts * as in the provocation, and as in the day of temptation in the
wilderness;

9 When your fathers tempted me, * proved me, and saw
my works.

10 Forty years long was I grieved with this generation,
and said, * It is a people that do err in their hearts, for they have
not known my ways:

11 Unto whom I sware in my wrath, * that they should
not enter into my rest.

3
stool until
1928

Psalm
96. \emph{Cantate Domino.}

O
SING unto the LORD a new song; * sing unto the LORD, all the whole earth.

2 Sing unto the LORD, and praise his Name; * be telling
of his salvation from day to day.

3 Declare
his honour unto the heathen, * and his wonders unto all peoples1.

4 For the LORD is great, and cannot worthily be praised;*
he is more to be feared than all gods.

5 As for all the gods of the heathen, they are but idols;
* but it is the LORD that made the heavens.

6 Glory and worship are before him; * power and honour are
in his sanctuary.

7 Ascribe unto the LORD, O ye kindreds of the peoples, *
ascribe unto the LORD worship and power.

8 Ascribe unto the LORD the honour due unto his Name; *
bring presents, and come into his courts.

9 O worship the LORD in the beauty of holiness; * let the
whole earth stand in awe of him.

10 Tell it out among the heathen, that the LORD is King,
and that it is he who hath made the round world so fast that it cannot
be moved; * and how that he shall judge the peoples1 righteously.

11 Let the heavens rejoice, and let the earth be glad; *
let the sea make a noise, and all that therein is.

12 Let the field be joyful, and all that is in it; * then
shall all the trees of the wood rejoice before the LORD.

13 For he cometh, for he cometh to judge the earth; * and
with righteousness to judge the world, and the peoples1 with
his truth.
1 people until
1928

Psalm 97. \emph{Dominus regnavit.}

THE
LORD is King, the earth may be glad thereof; * yea, the multitude of the
isles may be glad thereof.

2 Clouds and darkness are round about him: * righteousness
and judgment are the habitation of his seat.

3 There shall go a fire before him, * and burn up his
enemies on every side.

4 His lightnings gave shine unto the world: * the earth
saw it, and was afraid.

5 The hills melted like wax at the presence of the
LORD; * at the presence of the Lord of the whole earth.

6 The heavens have declared his righteousness, * and
all the peoples have seen his glory.

7 Confounded be all they that worship carved images,
and that delight in vain gods: * worship him, all ye gods.

8 Sion heard of it, and rejoiced; and the daughters
of Judah were glad, * because of thy judgments, O LORD.

9 For thou, LORD, art higher than all that are in the
earth: * thou art exalted far above all gods.

10 O ye that love the LORD, see that ye hate the thing
which is evil: * the Lord preserveth the souls of his saints; he shall
deliver them from the hand of the ungodly.

11 There is sprung up a light for the righteous, *
and joyful gladness for such as are true-hearted.

12 Rejoice in the LORD, ye righteous; * and give thanks
for a remembrance of his holiness. 

Evening
Prayer.

Psalm
98. \emph{Cantate Domino.}

O
SING unto the LORD a new song; * for he hath done marvellous things.

2 With his own right hand, and with his holy arm, *
hath he gotten himself the victory.

3 The LORD declared his salvation; * his righteousness
hath he openly showed in the sight of the heathen.

4 He hath remembered his mercy and truth toward the
house of Israel; * and all the ends of the world have seen the salvation
of our God.

5 Show yourselves joyful unto the LORD, all ye lands;
* sing, rejoice, and give thanks.

6 Praise the LORD upon the harp; * sing to the harp
with a psalm of thanksgiving.

7 With trumpets also and shawms, * O show yourselves
joyful before the LORD, the King.

8 Let the sea make a noise, and all that therein is;
* the round world, and they that dwell therein.

9 Let the floods clap their hands, and let the hills
be joyful together before the LORD; * for he is come to judge the earth.

10 With righteousness shall he judge the world, * and
the peoples1 with equity.
1 people until
1928

Psalm 99. \emph{Dominus regnavit.}

THE
LORD is King, be the people never so impatient; * he sitteth between the
Cherubim1, be the earth never so unquiet.

2 The LORD is great in Sion, * and high above all people.

3 They shall give thanks unto thy Name, * which is
great, wonderful, and holy.

4 The King’s power loveth judgment; thou hast
prepared equity, * thou hast executed judgment and righteousness in
Jacob.

5 O magnify the LORD our God, and fall down before
his footstool; * for he is holy.

6 Moses and Aaron among his priests, and Samuel among
such as call upon his Name: * these called upon the LORD, and he heard
them.

7 He spake unto them out of the cloudy pillar; * for
they kept his testimonies, and the law that he gave them.

1 Cherubims
before 1793

8 Thou heardest them, O LORD our God; * thou forgavest
them, O God, though thou didst punish their wicked doings.

9 O magnify the LORD our God, and worship him upon his
holy hill; * for the LORD our God is holy.

8 ... O God, and punishedst their own inventions. until
1928

Psalm 100.\emph{Jubilate
Deo.}

O
BE1 joyful in the LORD, all ye lands: * serve the LORD with
gladness, and come before his presence with a song.

2 Be ye sure that the LORD he is God; it is he that
hath made us, and not we ourselves; * we are his people, and the sheep
of his pasture.

3 O go your way into his gates with thanksgiving,
and into his courts with praise; * be thankful unto him, and speak good
of his Name.

4 For the LORD is gracious, his mercy is everlasting;
* and his truth endureth from generation to generation.

1 O
BE ye 1822-1871

\begin{quote}

\emph{Psalms 1-50;
101-150 Return to the U. S. Book of Common Prayer: 1789;
1892; 1928}

\end{quote}

\xchapter{The Psalter, Book V}

\textbf{The Psalter from the 1789, 1892 and 1928 U. S. Books of Common
Prayer}

Psalm
101. \emph{Misericordiam et judicium.}

MY
song shall be of mercy and judgment; * unto thee, O LORD, will I sing.

2 O let me have understanding * in the way of godliness!

3 When wilt thou come unto me? * I will walk in my
house with a perfect heart.

4 I will take no wicked thing in hand; I hate the sins
of unfaithfulness; * there shall no such cleave unto me.

5 A froward heart shall depart from me; * I will not
know a wicked person.

6 Whoso privily slandereth his neighbour, * him will
I destroy. Go
to Psalms 1-50; 

51-100

Return to the U. S. Book of Common Prayer: 1789;
1892; 1928

7 Whoso hath also a haughty look and a proud heart, * I will not suffer
him.

8 Mine eyes look upon such as are faithful in the land,
* that they may dwell with me.

9 Whoso leadeth a godly life, * he shall be my servant.

10 There shall no deceitful person dwell in my house;
* he that telleth lies shall not tarry in my sight.

11 I shall soon destroy all the ungodly that are in
the land; * that I may root out all wicked doers from the city of the
LORD. 

The
Twentieth Day.

Morning
Prayer.

Psalm
102. \emph{Domine, exaudi.}

HEAR
my prayer, O LORD, * and let my crying come unto thee.

2 Hide not thy face from me in the time of my trouble;
* incline thine ear unto me when I call; O hear me, and that right soon.

3 For my days are consumed away like smoke, * and my
bones are burnt up as it were a firebrand.

4 My heart is smitten down, and withered like grass;
* so that I forget to eat my bread.

5 For the voice of my groaning, * my bones will scarce
cleave to my flesh.

6 I am become like a pelican in the wilderness, * and
like an owl that is in the desert.

7 I have watched, and am even as it were a sparrow,
* that sitteth alone upon the housetop.

8 Mine enemies revile me all the day long; * and they
that are mad upon me are sworn together against me.

9 For I have eaten ashes as it were bread, * and mingled
my drink with weeping;

10 And that, because of thine indignation and wrath;
* for thou hast taken me up, and cast me down.

11 My days are gone like a shadow, * and I am withered
like grass.

12 But thou, O LORD, shalt endure for ever, * and thy
remembrance throughout all generations.

13 Thou shalt arise, and have mercy upon Sion; * for
it is time that thou have mercy upon her, yea, the time is come.

14 And why? thy servants think upon her stones, * and
it pitieth them to see her in the dust.

7 Who hath also a proud look
and high stomach : ...
until 1928.

15 The nations1 shall fear thy Name, O LORD; * and all the kings
of the earth thy majesty;

16 When the LORD shall build up Sion, * and when his
glory shall appear;

17 When he turneth him unto the prayer of the poor destitute,
* and despiseth not their desire.

18 This shall be written for those that come after, *
and the people which shall be born shall praise the LORD.

19 For he hath looked down from his sanctuary; * out
of the heaven did the LORD behold the earth; 1 heathen
until 1928.

20 That he might hear the mournings2 of such as are in captivity,
* and deliver them that are appointed unto death;

21 That they may declare the Name of the LORD in Sion,
* and his worship at Jerusalem; 2 mourning 1793-1892

22 When the peoples3 are gathered together, * and the
kingdoms also, to serve the LORD.

23 He brought down my strength in my journey, * and
shortened my days.

24 But I said, O my God, take me not away in the midst
of mine age; * as for thy years, they endure throughout all generations.

25 Thou, Lord, in the beginning hast laid the foundation
of the earth, * and the heavens are the work of thy hands.

26 They shall perish, but thou shalt endure: * they
all shall wax old as doth a garment;

27 And as a vesture shalt thou change them, and they
shall be changed; * but thou art the same, and thy years shall not fail.

28 The children of thy servants shall continue, * and
their seed shall stand fast in thy sight.

Psalm
103. \emph{Benedic, anima mea.}

PRAISE
the LORD, O my soul; * and all that is within me, praise his holy Name.

2 Praise the LORD, O my soul, * and forget not all
his benefits:

3 Who forgiveth all thy sin, * and healeth all thine
infirmities;

4 Who saveth thy life from destruction, * and crowneth
thee with mercy and loving-kindness;

5 Who satisfieth thy mouth with good things, * making
thee young and lusty as an eagle.

6 The LORD executeth righteousness and judgment * for
all them that are oppressed with wrong.

7 He showed his ways unto Moses, * his works unto the
children of Israel.

8 The LORD is full of compassion and mercy, * long-suffering,
and of great goodness.

9 He will not alway be chiding; * neither keepeth he
his anger for ever.

10 He hath not dealt with us after our sins; * nor
rewarded us according to our wickednesses.

11 For look how high the heaven is in comparison of
the earth; * so great is his mercy also toward them that fear him.

12 Look how wide also the east is from the west; *
so far hath he set our sins from us.

13 Yea, like as a father pitieth his own children;
* even so is the LORD merciful unto them that fear him.

14 For he knoweth whereof we are made; * he remembereth
that we are but dust.

15 The days of man are but as grass; * for he flourisheth
as a flower of the field.

16 For as soon as the wind goeth over it, it is gone;
* and the place thereof shall know it no more.

17 But the merciful goodness of the LORD endureth for
ever and ever upon them that fear him; * and his righteousness upon children's
children;

18 Even upon such as keep his covenant, * and think
upon his commandments to do them.

19 The LORD hath prepared his seat in heaven, * and
his kingdom ruleth over all. 3 people
until 1928.

20 O praise the LORD, ye angels of his, ye that excel in strength;
* ye that fulfil his commandment, and hearken unto the voice of his word1.

21 O praise the LORD, all ye his hosts; * ye servants
of his that do his pleasure.

22 O speak good of the LORD, all ye works of his, in
all places of his dominion: * praise thou the LORD, O my soul.

Evening
Prayer.

Psalm
104. \emph{Benedic, anima mea.}

PRAISE
the LORD, O my soul: * O LORD my God, thou art become exceeding glorious;
thou art clothed with majesty and honour.

2 Thou deckest thyself with light as it were with a
garment, * and spreadest out the heavens like a curtain.

3 Who layeth the beams of his chambers in the waters,
* and maketh the clouds his chariot, and walketh upon the wings of the
wind.

4 He maketh his angels winds, * and his ministers a
flaming fire.

5 He laid the foundations of the earth, * that it never
should move at any time.

6 Thou coveredst it with the deep like as with a garment;
* the waters stand above the hills.

7 At thy rebuke they flee; * at the voice of thy thunder
they haste away.

8 They go up as high as the hills, and down to the
valleys beneath; * even unto the place which thou hast appointed for them.

9 Thou hast set them their bounds, which they shall
not pass, * neither turn again to cover the earth.

10 He sendeth the springs into the rivers, * which
run among the hills.

1 words
until
1845 and in 1892 Book

11 All beasts1 of the field drink thereof, * and the wild asses
quench their thirst.

12 Beside them shall the fowls of the air have their
habitation, * and sing among the branches.

13 He watereth the hills from above; * the earth is filled
with the fruit of thy works.

14 He bringeth forth grass for the cattle, * and green
herb for the service of men;

15 That he may bring food out of the earth, and wine
that maketh glad the heart of man; * and oil to make him a cheerful countenance,
and bread to strengthen man's heart. 1
All the beasts 1822-1892

16 The trees of the LORD also are full of sap; * even the cedars of Lebanon2
which he hath planted;

17 Wherein the birds make their nests; * and the fir-trees
are a dwelling for the stork.

18 The high hills are a refuge for the wild goats; *
and so are the stony rocks for the conies.

19 He appointed the moon for certain seasons, * and the
sun knoweth his going down. \\
2 Libanus
until 1928.

20 Thou makest darkness that it may be night; * wherein all the beasts
of the forest3 do move.

21 The lions, roaring after their prey, * do seek their
meat from God.

22 The sun ariseth, and they get them away together,
* and lay them down in their dens.

23 Man goeth forth to his work, and to his labour, *
until the evening.

24 O LORD, how manifold are thy works! * in wisdom hast
thou made them all; the earth is full of thy riches.

25 So is the great and wide sea also; * wherein are things
creeping innumerable, both small and great beasts.

26 There go the ships, and there is that leviathan, *
whom thou hast made to take his pastime therein.

27 These wait all upon thee, * that thou mayest give
them meat in due season.

28 When thou givest it them, they gather it; * and when
thou openest thy hand, they are filled with good.

29 When thou hidest thy face, they are troubled: * when
thou takest away their breath, they die, and are turned again to their dust.

30 When thou lettest thy breath go forth, they shall
be made; * and thou shalt renew the face of the earth.

31 The glorious majesty of the LORD shall endure for
ever; * the LORD shall rejoice in his works.

32 The earth shall tremble at the look of him; * if he
do but touch the hills, they shall smoke.

33 I will sing unto the LORD as long as I live; * I will
praise my God while I have my being.

34 And so shall my words please him: * my joy shall be
in the LORD. \\
3
forests 1845-1871

35 As for sinners, they shall be consumed out of the earth, * and
the ungodly shall come to an end.

36 Praise thou the LORD, O my soul. * Praise the LORD.

The
Twenty-first Day.

Morning
Prayer.

Psalm
105. \emph{Confitemini Domino.}

GIVE
thanks unto the LORD, and call upon his Name; * tell the people what things
he hath done.

2 O let your songs be of him, and praise him; * and
let your talking be of all his wondrous works.

3 Rejoice in his holy Name; * let the heart of them
rejoice that seek the LORD.

4 Seek the LORD and his strength; * seek his face evermore.

5 Remember the marvellous works that he hath done;
* his wonders, and the judgments of his mouth;

6 O ye seed of Abraham his servant, * ye children of
Jacob his chosen.

7 He is the LORD our God; * his judgments are in all
the world.

8 He hath been alway mindful of his covenant and promise,
* that he made to a thousand generations;

9 Even the covenant that he made with Abraham; * and
the oath that he sware unto Isaac;

10 And appointed the same unto Jacob for a law, * and
to Israel for an everlasting testament;

11 Saying, Unto thee will I give the land of Canaan,
* the lot of your inheritance:

12 When there were yet but a few of them, * and they
strangers in the land;

13 What time as they went from one nation to another,
* from one kingdom to another people;

14 He suffered no man to do them wrong, * but reproved
even kings for their sakes;

15 Touch not mine anointed, * and do my prophets no
harm.

16 Moreover, he called for a dearth upon the land,
* and destroyed all the provision of bread.

17 But he had sent a man before them, * even Joseph,
who was sold to be a bond-servant;

18 Whose feet they hurt in the stocks; * the iron entered
into his soul;

19 Until the time came that his cause was known: *
the word of the LORD tried him.

20 The king sent, and delivered him; * the prince of
the people let him go free.

21 He made him lord also of his house, * and ruler
of all his substance;

22 That he might inform his princes after his will,
* and teach his senators wisdom.

23 Israel also came into Egypt, * and Jacob was a stranger
in the land of Ham.

24 And he increased his people exceedingly, * and made
them stronger than their enemies;

25 Whose heart turned so, that they hated his people,
* and dealt untruly with his servants.

26 Then sent he Moses his servant, * and Aaron whom
he had chosen.

27 And these showed his tokens among them, * and wonders
in the land of Ham.

28 He sent darkness, and it was dark; * and they were
not obedient unto his word.

29 He turned their waters into blood, * and slew their
fish.

30 Their land brought forth frogs; * yea, even in their
kings' chambers.

31 He spake the word, and there came all manner of
flies, * and lice in all their quarters.

32 He gave them hailstones for rain; * and flames of
fire in their land.

33 He smote their vines also and fig-trees; * and destroyed
the trees that were in their coasts.

34 He spake the word, and the grasshoppers came, and
caterpillars innumerable, * and did eat up all the grass in their land,
and devoured the fruit of their ground.

35 He smote all the firstborn in their land; * even
the chief of all their strength.

36 He brought them forth also with silver and gold;
* there was not one feeble person among their tribes.

37 Egypt was glad at their departing; * for they were
afraid of them.

38 He spread out a cloud to be a covering, * and fire
to give light in the night season.

39 At their desire he brought quails; * and he filled
them with the bread of heaven.

40 He opened the rock of stone, and the waters flowed
out, * so that rivers ran in the dry places.

41 For why? he remembered his holy promise; * and Abraham
his servant.

42 And he brought forth his people with joy, * and
his chosen with gladness;

43 And gave them the lands of the heathen; * and they
took the labours of the people in possession;

44 That they might keep his statutes, * and observe
his laws. 

Evening
Prayer.

Psalm
106. \emph{Confitemini Domino.}

O
GIVE thanks unto the LORD; for he is gracious, * and his mercy endureth
for ever.

2 Who can express the noble acts of the LORD, * or show forth
all his praise?

3 Blessed are they that alway keep judgment, * and
do righteousness.

4 Remember me, O LORD, according to the favour that
thou bearest unto thy people; * O visit me with thy salvation;

5 That I may see the felicity of thy chosen, * and
rejoice in the gladness of thy people, and give thanks with thine inheritance.

6 We have sinned with our fathers; * we have done amiss,
and dealt wickedly.

7 Our fathers regarded not thy wonders in Egypt, neither
kept they thy great goodness in remembrance; * but were disobedient at
the sea, even at the Red Sea.

8 Nevertheless, he helped them for his Name's sake,
* that he might make his power to be known.

9 He rebuked the Red Sea also, and it was dried up;
* so he led them through the deep, as through a wilderness.

10 And he saved them from the adversary's hand, * and
delivered them from the hand of the enemy.

11 As for those that troubled them, the waters overwhelmed
them; * there was not one of them left.

12 Then believed they his words, * and sang praise
unto him.

13 But within a while they forgat his works, * and
would not abide his counsel.

14 But lust came upon them in the wilderness, * and
they tempted God in the desert.

15 And he gave them their desire, * and sent leanness
withal into their soul.

16 They angered Moses also in the tents, * and Aaron
the saint of the LORD.

17 So the earth opened, and swallowed up Dathan, *
and covered the congregation of Abiram.

18 And the fire was kindled in their company; * the
flame burnt up the ungodly.

19 They made a calf in Horeb, * and worshipped the
molten image.

20 Thus they turned their glory * into the similitude
of a calf that eateth hay.

21 And they forgat God their Saviour, * who had done
so great things in Egypt;

22 Wondrous works in the land of Ham; * and fearful
things by the Red Sea.

23 So he said he would have destroyed them, had not
Moses his chosen stood before him in the gap, * to turn away his wrathful
indignation, lest he should destroy them.

24 Yea, they thought scorn of that pleasant land, *
and gave no credence unto his word;

25 But murmured in their tents, * and hearkened not
unto the voice of the LORD.

26 Then lift he up his hand against them, * to overthrow
them in the wilderness;

27 To cast out their seed among the nations, * and
to scatter them in the lands.

28 They joined themselves unto Baal-peor, * and ate
the offerings of the dead.

29 Thus they provoked him to anger with their own inventions;
* and the plague was great among them. 35. As for sinners, they shall
be consumed out of the

earth, and the ungodly shall come to an end : praise

thou the LORD, O my soul, praise the LORD.
until
1928

30 Then stood up Phinehas, and interposed; * and so the plague ceased.

31 And that was counted unto him for righteousness,
* among all posterities for evermore.

32 They angered him also at the waters of strife, *
so that he punished Moses for their sakes;

33 Because they provoked his spirit, * so that he spake
unadvisedly with his lips.

34 Neither destroyed they the heathen, * as the LORD
commanded them;

35 But were mingled among the heathen, * and learned
their works.

36 Insomuch that they worshipped their idols, which
became a snare unto them; * yea, they offered their sons and their daughters
unto devils;

37 And shed innocent blood, even the blood of their
sons and of their daughters, * whom they offered unto the idols of Canaan;
and the land was defiled with blood.

38 Thus were they stained with their own works, * and
went a whoring with their own inventions.

39 Therefore was the wrath of the LORD
kindled against his people, * insomuch that he abhorred his own inheritance.

40 And he gave them over into the hand of the heathen;
* and they that hated them were lords over them.

41 Their enemies oppressed them, * and had them in
subjection.

42 Many a time did he deliver them; * but they rebelled
against him with their own inventions, and were brought down in their
wickedness.

43 Nevertheless, when he saw their adversity, * he
heard their complaint.

44 He thought upon his covenant, and pitied them, according
unto the multitude of his mercies; * yea, he made all those that led them
away captive to pity them.

45 Deliver us, O LORD our God,
and gather us from among the heathen; * that we may give thanks unto thy
holy Name, and make our boast of thy praise.

46
Blessed be the LORD God of Israel, from everlasting,
and world without end; * And let all the people say, Amen.

BOOK
V.

The
Twenty-second Day.

Morning
Prayer.

Psalm
107. \emph{Confitemini Domino.}

O
GIVE thanks unto the LORD, for he is gracious,
* and his mercy endureth for ever.

2 Let them give thanks whom the LORD
hath redeemed, * and delivered from the hand of the enemy;

3 And gathered them out of the lands, from the east,
and from the west; * from the north, and from the south.

4 They went astray in the wilderness out of the way,
* and found no city to dwell in.

5 Hungry and thirsty, * their soul fainted in them.

6 So they cried unto the LORD
in their trouble, * and he delivered them from their distress.

7 He led them forth by the right way, * that they might
go to the city where they dwelt.

8 O that men would therefore praise the LORD
for his goodness; * and declare the wonders that he doeth for the children
of men!

9 For he satisfieth the empty soul, * and filleth the
hungry soul with goodness.

10 Such as sit in darkness,
and in the shadow of death, * being fast bound in misery and iron;

11 Because they rebelled against the words of the Lord,
* and lightly regarded the counsel of the Most Highest;

12 He also brought down their heart through heaviness:
* they fell down, and there was none to help them.

13 So when they cried unto the LORD
in their trouble, * he delivered them out of their distress.

14 For he brought them out of darkness, and out of
the shadow of death, * and brake their bonds in sunder.

15 O that men would therefore praise the LORD
for his goodness; * and declare the wonders that he doeth for the children
of men !

16 For he hath broken the gates of brass, * and smitten
the bars of iron in sunder.

17 Foolish men are plagued
for their offence, * and because of their wickedness.

18 Their soul abhorred all manner of meat, * and they
were even hard at death's door.

19 So when they cried unto the LORD
in their trouble, * he delivered them out of their distress.

20 He sent his word, and healed them; * and they were
saved from their destruction.

21 O that men would therefore praise the LORD
for his goodness; * and declare the wonders that he doeth for the children
of men !

22 That they would offer unto him the sacrifice of
thanksgiving, * and tell out his works with gladness!

23 They that go down to the sea in ships, * and occupy
their business in great waters;

24 These men see the works of the LORD,
* and his wonders in the deep.

25 For at his word the stormy wind ariseth, * which
lifteth up the waves thereof.

26 They are carried up to the heaven, and down again
to the deep; * their soul melteth away because of the trouble. 30 Then stood up Phinees,
and prayed : ...
until 1928

27 They reel to and fro, and stagger like a drunken man, * and are at
their wit's1 end.

28 So when they cry unto the LORD
in their trouble, * he delivereth them out of their distress.

29 For he maketh the storm to cease, * so that the
waves thereof are still.

30 Then are they glad, because they are at rest; *
and so he bringeth them unto the haven where they would be.

31 O that men would therefore praise the LORD
for his goodness; * and declare the wonders that he doeth for the children
of men!

32 That they would exalt him also in the congregation
of the people, * and praise him in the seat of the elders!

33
He turneth the floods into a wilderness, * and drieth up the water-springs.

34 A fruitful land maketh he barren, * for the wickedness
of them that dwell therein.

35 Again, he maketh the wilderness a standing water,
* and water-springs of a dry ground.

36 And there he setteth the hungry, * that they may
build them a city to dwell in;

37 That they may sow their land, and plant vineyards,
* to yield them fruits of increase.

38 He blesseth them, so that they multiply exceedingly;
* and suffereth not their cattle to decrease.

39 And again, when they are minished and brought low
* through oppression, through any plague or trouble;

40 Though he suffer them to be evil entreated through
tyrants, * and let them wander out of the way in the wilderness;

41 Yet helpeth he the poor out of misery, * and maketh
him households like a flock of sheep.

42 The righteous will consider this, and rejoice; *
and the mouth of all wickedness shall be stopped.

43 Whoso is wise, will ponder these things; * and they
shall understand the loving-kindness of the LORD.

Evening
Prayer.

Psalm
108. \emph{Paratum cor meum.}

O
GOD, my heart is ready, my heart is ready; * I will sing, and give praise
with the best member that I have.

2 Awake, thou lute and harp; * I myself will awake
right early. \\
1
wits until
1793; and 1871-1892

3 I will give thanks unto thee, O LORD, among the
peoples1; * I will sing praises unto thee among the nations.

4 For thy mercy is greater than the heavens, * and thy
truth reacheth unto the clouds.

5 Set up thyself, O God, above the heavens, * and thy
glory above all the earth;

6 That thy beloved may be delivered: * let thy right
hand save them, and hear thou me.

1
people
until 1928

7 God hath spoken in his holiness; * I will rejoice therefore, and divide
Shechem2, and mete out the valley of Succoth.

8 Gilead is mine, and Manasseh is mine; * Ephraim also
is the strength of my head; Judah is my lawgiver;

9 Moab is my wash-pot; over Edom will I cast out my
shoe; * upon Philistia will I triumph.

10 Who will lead me into the strong city? * and who
will bring me into Edom?

11 Hast not thou forsaken us, O God? * and wilt not
thou, O God, go forth with our hosts?

12 O help us against the enemy: * for vain is the help
of man.

13 Through God we shall do great acts; * and it is
he that shall tread down our enemies.

Psalm
109. \emph{Deus, laudem.}

HOLD
not thy tongue, O God of my praise; * for the mouth of the ungodly, yea,
the mouth of the deceitful is opened upon me.

2 And they have spoken against me with false tongues;
* they compassed me about also with words of hatred, and fought against
me without a cause.

3 For the love that I had unto them, lo, they take
now my contrary part; * but I give myself unto prayer.

4 Thus have they rewarded me evil for good, * and hatred
for my good will.

2 Sichem
until
1928

8. Gilead is mine, and Manasses is mine : Ephraim

also is the strength of my head.

9. Judah is my law-giver, Moab is my wash-pot :

over Edom ...
until 1928

5 Set thou an ungodly man to be ruler over him, * and let an adversary stand
at his right hand.

6 When sentence is given upon him, let him be condemned;
* and let his prayer be turned into sin.

7 Let his days be few; * and let another take his office.

8 Let his children be fatherless, * and his wife a widow.

9 Let his children be vagabonds, and beg their bread;
* let them seek it also out of desolate places.

10 Let the extortioner consume all that he hath; * and
let the stranger spoil his labour.

11 Let there be no man to pity him, * nor to have compassion
upon his fatherless children.

12 Let his posterity be destroyed; * and in the next
generation let his name be clean put out.

13 Let the wickedness of his fathers be had in remembrance
in the sight of the LORD; * and let not the sin of his mother be done away.

14 Let them alway be before the LORD, * that he may root
out the memorial of them from off the earth;

15 And that, because his mind was not to do good; * but
persecuted the poor helpless man, that he might slay him that was vexed
at the heart.

16 His delight was in cursing, and it shall happen unto
him; * he loved not blessing, therefore shall it be far from him.

17 He clothed himself with cursing like as with a raiment,
* and it shall come into his bowels like water, and like oil into his bones.

18 Let it be unto him as the cloak that he hath upon
him, * and as the girdle that he is alway girded withal.

19 Let it thus happen from the LORD unto mine enemies,
* and to those that speak evil against my soul.

20 But deal thou with me, O LORD God, according unto thy Name; * for sweet
is thy mercy.

21 O deliver me, for I am helpless and poor, \textasciitilde{} and
my heart is wounded within me.

22 I go hence like the shadow that departeth, * and
am driven away as the grasshopper.

23 My knees are weak through fasting; * my flesh is
dried up for want of fatness.

24 I am become also a reproach unto them: * they that
look upon me shake their heads.

25 Help me, O LORD my God; * O save me according to
thy mercy;

26 And they shall know how that this is thy hand, *
and that thou, LORD, hast done it.

27 Though they curse, yet bless thou; * and let them
be confounded that rise up against me; but let thy servant rejoice.

28 Let mine adversaries be clothed with shame; * and
let them cover themselves with their own confusion, as with a cloak.

29 As for me, I will give great thanks unto the LORD
with my mouth, * and praise him among the multitude;

30 For he shall stand at the right hand of the poor,
* to save his soul from unrighteous judges. 

The
Twenty-third Day.

Morning
Prayer.

Psalm
110. \emph{Dixit Dominus.}

THE
LORD said unto my Lord, * Sit thou on my right hand, until I make thine
enemies thy footstool.

2 The LORD shall send the rod of thy power out of Sion:
* be thou ruler, even in the midst among thine enemies. \\
5 ... and let Satan stand at
...
until 1928

3 In the day of thy power shall thy1 people offer themselves
willingly with an holy worship: * thy young men come to thee as dew from
the womb of the morning.

4 The LORD sware, and will not repent, * Thou art a
Priest for ever after the order of Melchizedek.

5 The Lord upon thy right hand * shall wound even kings
in the day of his wrath.

6 He shall judge among the heathen; * he shall fill
the places with the dead bodies, and smite in sunder the heads over divers
countries.

7 He shall drink of the brook in the way; * therefore
shall he lift up his head.

Psalm
111. \emph{Confitebor tibi.}

I
WILL give thanks unto the LORD with my whole heart, * secretly among the
faithful, and in the congregation.

2 The works of the LORD are great, * sought out of
all them that have pleasure therein.

3 His work is worthy to be praised and had in honour,
* and his righteousness endureth for ever.

4 The merciful and gracious LORD hath so done his marvellous
works, * that they ought to be had in remembrance.

5 He hath given meat unto them that fear him; * he
shall ever be mindful of his covenant.

6 He hath showed his people the power of his works,
* that he may give them the heritage of the heathen.

7 The works of his hands are verity and judgment; *
all his commandments are true.

8 They stand fast for ever and ever, * and are done
in truth and equity.

9 He sent redemption unto his people; * he hath commanded
his covenant for ever; holy and reverend is his Name. 3. In the day of thy power
shall thy1 people offer thee free-will offerings with an holy
worship : the dew of thy birth is of the womb of the morning.
until 1928\\
1
the in
1789 Book

10 The fear of the LORD is the beginning of wisdom; * a good understanding
have all they that do thereafter; his praise endureth for ever.

Psalm
112. \emph{Beatus vir.}

BLESSED
is the man that feareth the LORD; * he hath great delight in his commandments.

2 His seed shall be mighty upon earth; * the generation
of the faithful shall be blessed.

3 Riches and plenteousness shall be in his house; *
and his righteousness endureth for ever.

4 Unto the godly there ariseth up light in the darkness;
* he is merciful, loving, and righteous.

5 A good man is merciful, and lendeth; * and will guide
his words with discretion.

6 For he shall never be moved: * and the righteous
shall be had in everlasting remembrance.

7 He will not be afraid of any evil tidings; * for
his heart standeth fast, and believeth in the LORD.

8 His heart is stablished, and will not shrink, * until
he see his desire upon his enemies.

9 He hath dispersed abroad, and given to the poor.
* and his righteousness remaineth for ever; his horn shall be exalted
with honour.

10 The ungodly shall see it, and it shall grieve him;
* he shall gnash with his teeth, and consume away; the desire of the ungodly
shall perish.

Psalm
113. \emph{Laudate, pueri.}

PRAISE
the LORD, ye servants; * O praise the Name of the LORD.

2 Blessed be the Name of the LORD * from this time
forth for evermore.

3 The LORD’S Name is praised * from the rising
up of the sun unto the going down of the same. 10 ... thereafter; the praise
of it endureth for ever. until
1928

4 The LORD is high above all nations1, * and his glory above
the heavens.

5 Who is like unto the LORD our God, that hath his
dwelling so high, * and yet humbleth himself to behold the things that
are in heaven and earth!

6 He taketh up the simple out of the dust, * and lifteth
the poor out of the mire;

7 That he may set him with the princes, * even with
the princes of his people.

8 He maketh the barren woman to keep house, * and to
be a joyful mother of children. 

Evening
Prayer.

Psalm
114. \emph{In exitu Israel.}

WHEN
Israel came out of Egypt, * and the house of Jacob from among the strange
people,

2 Judah was his sanctuary, * and Israel his dominion.

3 The sea saw that, and fled; * Jordan was driven back.

4 The mountains skipped like rams, * and the little
hills like young sheep.

5 What aileth thee, O thou sea, that thou fleddest?
* and thou Jordan, that thou wast driven back?

6 Ye mountains, that ye skipped like rams? * and ye
little hills, like young sheep?

7 Tremble, thou earth, at the presence of the Lord:
* at the presence of the God of Jacob;

8 Who turned the hard rock into a standing water, *
and the flint-stone into a springing well.

Psalm
115. \emph{Non nobis, Domine.}

NOT
unto us, O LORD, not unto us, but unto thy Name give the praise; * for
thy loving mercy, and for thy truth’s sake.

2 Wherefore shall the heathen say, * Where is now their
God?

3 As for our God, he is in heaven: * he hath done whatsoever
pleased him.

4 Their idols are silver and gold, * even the work
of men’s hands.

5 They have mouths, and speak not; * eyes have they,
and see not.

6 They have ears, and hear not; * noses have they,
and smell not.

7 They have hands, and handle not; feet have they,
and walk not; * neither speak they through their throat.

8 They that make them are like unto them; * and so
are all such as put their trust in them.

9 But thou, house of Israel, trust thou in the LORD;
* he is their helper and defender.

10 Ye house of Aaron, put your trust in the LORD; *
he is their helper and defender.

11 Ye that fear the LORD, put your trust in the LORD;
* he is their helper and defender.

12 The LORD hath been mindful of us, and he shall bless
us; * even he shall bless the house of Israel, he shall bless the house
of Aaron.

13 He shall bless them that fear the LORD, * both small
and great.

14 The LORD shall increase you more and more, * you
and your children.

15 Ye are the blessed of the LORD, * who made heaven
and earth.

16 All the whole heavens are the LORD’S; * the earth hath he given to the children of men.

17 The dead praise not thee, O LORD, * neither all
they that go down into silence.

18 But we will praise the LORD, * from this time forth
for evermore. Praise the LORD. 

The
Twenty-fourth Day.

Morning
Prayer.

Psalm
116. \emph{Dilexi, quoniam.} 1 heathen until
1928

MY
delight is in the LORD; * because he hath heard the voice of my prayer;

2 Because he hath inclined his ear unto me; * therefore
will I call upon him as long as I live.

3 The snares of death compassed me round about, * and
the pains of hell gat hold upon me. I AM well pleased
: that the LORD hath heard the voice of my prayer;

2. That he hath inclined ... until
1928

4 I found trouble and heaviness; then called I upon the Name of the LORD;
* O LORD, I beseech thee, deliver my soul.

5 Gracious is the LORD, and righteous; * yea, our God
is merciful.

6 The LORD preserveth the simple: * I was in misery,
and he helped me.

7 Turn again then unto thy rest, O my soul; * for the
LORD hath rewarded thee.

8 And why? thou hast delivered my soul from death, *
mine eyes from tears, and my feet from falling.

9 I will walk before the LORD * in the land of the living.

10 I believed, and therefore will I speak; but I was
sore troubled: * I said in my haste, All men are liars.

11 What reward shall I give unto the LORD * for all the
benefits that he hath done unto me?

12 I will receive the cup of salvation, * and call upon
the Name of the LORD.

13 I will pay my vows now in the presence of all his
people: * right dear in the sight of the LORD is the death of his saints.

4. I shall find trouble and
heaviness, and I will call

upon the Name of the LORD : ... until
1928

14 Behold, O LORD, how that I am thy servant; * I
am thy servant, and the son of thine1 handmaid; thou hast broken
my bonds in sunder.

15 I will offer to thee the sacrifice of thanksgiving,
* and will call upon the Name of the LORD.

16 I will pay my vows unto the LORD, in the sight of
all his people, * in the courts of the LORD’S house; even in the midst
of thee, O Jerusalem. Praise the LORD.

1 thy 1793-1892

Psalm
117. \emph{Laudate Dominum.}

O
PRAISE the LORD, all ye nations1; * praise him, all ye peoples.

2 For his merciful kindness is ever more and more toward2
us; * and the truth of the LORD endureth for ever. Praise the LORD.

Psalm
118. \emph{Confitemini Domino.}

O
GIVE thanks unto the LORD, for he is gracious; * because his mercy endureth
for ever.

2 Let Israel now confess that he is gracious, * and
that his mercy endureth for ever.

3 Let the house of Aaron now confess, * that his mercy
endureth for ever.

4 Yea, let them now that fear the LORD confess, * that
his mercy endureth for ever.

5 I called upon the LORD in trouble; * and the LORD
heard me at large.

6 The LORD is on my side; * I will not fear what man
doeth unto me.

7 The LORD taketh my part with them that help me; *
therefore shall I see my desire upon mine enemies.

8 It is better to trust in the LORD, * than to put
any confidence in man.

9 It is better to trust in the LORD, * than to put
any confidence in princes.

10 All nations compassed me round about; * but in the
Name of the LORD will I destroy them.

11 They kept me in on every side, they kept me in,
I say, on every side; * but in the Name of the LORD will I destroy them.

12 They came about me like bees, and are extinct even
as the fire among the thorns; * for in the Name of the LORD I will destroy
them.

13 Thou hast thrust sore at me, that I might fall;
* but the LORD was my help.

14 The LORD is my strength, and my song; * and is become
my salvation.

15 The voice of joy and health is in the dwellings
of the righteous; * the right hand of the LORD bringeth mighty things
to pass.

16 The right hand of the LORD hath the preeminence;
* the right hand of the LORD bringeth mighty things to pass.

17 I shall not die, but live, * and declare the works
of the LORD.

18 The LORD hath chastened and corrected me; * but
he hath not given me over unto death.

19 Open me the gates of righteousness, * that I may
go into them, and give thanks unto the LORD.

20 This is the gate of the LORD, * the righteous shall
enter into it.

21 I will thank thee; for thou hast heard me, * and
art become my salvation.

22 The same stone which the builders refused, * is
become the head-stone in the corner.

23 This is the LORD’S doing, * and it is marvellous
in our eyes.

24 This is the day which the LORD hath made; * we will
rejoice and be glad in it.

25 Help me now, O LORD: * O LORD, send us now prosperity.

1 heathen until
1928

2 towards
until
1892

26 Blessed be he that cometh in the Name of the LORD: * we have wished
you good luck, we1 that are of the house of the LORD.

27 God is the LORD, who hath showed us light: * bind
the sacrifice with cords, yea, even unto the horns of the altar.

28 Thou art my God, and I will thank thee; * thou art
my God, and I will praise thee.

29 O give thanks unto the LORD; for he is gracious,
* and his mercy endureth for ever. 

Evening
Prayer.

Psalm
119. I. \emph{Beati immaculati.}

BLESSED
are those that are undefiled in the way, and walk in the law of the LORD.

2 Blessed are they that keep his testimonies, * and
seek him with their whole heart; \\
1 ye until
1928

3 Even they who do no wickedness, * and walk in his ways.

4 Thou hast charged * that we shall diligently keep thy
commandments.

5 O that my ways were made so direct, * that I might
keep thy statutes!

6 So shall I not be confounded, * while I have respect
unto all thy commandments.

7 I will thank thee with an unfeigned heart, * when I
shall have learned the judgments of thy righteousness. 3. For they who do no wickedness
: walk in his ways. until
1928

8 I will
keep thy statutes1; * O forsake me not utterly.

II.
\emph{In quo corrigit?}

WHEREWITHAL
shall a young man cleanse his way? * even by ruling himself after thy
word.

10 With my whole heart have I sought thee; * O let
me not go wrong out of thy commandments. 1 ceremonies until
1928

11 Thy word2 have I hid within my heart, * that I should not
sin against thee.

12 Blessed art thou, O LORD; * O teach me thy statutes.

13 With my lips have I been telling * of all the judgments
of thy mouth.

14 I have had as great delight in the way of thy testimonies,
* as in all manner of riches.

15 I will talk of thy commandments, * and have respect
unto thy ways.

16 My delight shall be in thy statutes, * and I will
not forget thy word.

III.
\emph{Retribue servo tuo.}

O
DO well unto thy servant; * that I may live, and keep thy word.

18 Open thou mine eyes; * that I may see the wondrous
things of thy law.

19 I am a stranger upon earth; * O hide not thy commandments
from me.

20 My soul breaketh out for the very fervent desire
* that it hath alway unto thy judgments.

21 Thou hast rebuked the proud; * and cursed are they
that do err from thy commandments.

22 O turn from me shame and rebuke; * for I have kept
thy testimonies.

23 Princes also did sit and speak against me; * but
thy servant is occupied in thy statutes.

24 For thy testimonies are my delight, * and my counsellors.

IV. \emph{Adhaesit
pavimento.}

MY
soul cleaveth to the dust; * O quicken thou me, according to thy word.

26 I have acknowledged my ways, and thou heardest me:
* O teach me thy statutes.

27 Make me to understand the way of thy commandments;
* and so shall I talk of thy wondrous works.

28 My soul melteth away for very heaviness; * comfort
thou me according unto thy word.

29 Take from me the way of lying, * and cause thou
me to make much of thy law.

30 I have chosen the way of truth, * and thy judgments
have I laid before me.

31 I have stuck unto thy testimonies; * O LORD, confound
me not.

32 I will run the way of thy commandments, * when thou
hast set my heart at liberty. 

The
Twenty-fifth Day.

Morning
Prayer.

V. \emph{Legem
pone.}

TEACH
me, O LORD, the way of thy statutes, * and I shall keep it unto the end.

34 Give me understanding, and I shall keep thy law;
* yea, I shall keep it with my whole heart.

35 Make me to go in the path of thy commandments; *
for therein is my desire. 2 words until
1928

36 Incline my3 heart unto thy testimonies, * and not to covetousness.

37 O turn away mine eyes, lest they behold vanity; *
and quicken thou me in thy way.

38 O stablish thy word in thy servant, * that I may fear
thee.

39 Take away the rebuke that I am afraid of; * for thy
judgments are good.

40 Behold, my delight is in thy commandments; * O quicken
me in thy righteousness.

VI.
\emph{Et veniat super me.}

LET
thy loving mercy come also unto me, O LORD, * even thy salvation, according
unto thy word.

42 So shall I make answer unto my blasphemers; * for
my trust is in thy word.

43 O take not the word of thy truth utterly out of
my mouth; * for my hope is in thy judgments.

44 So shall I alway keep thy law; * yea, for ever and
ever.

45 And I will walk at liberty; * for I seek thy commandments.

46 I will speak of thy testimonies also, even before
kings, * and will not be ashamed.

47 And my delight shall be in thy commandments, * which
I have loved.

48 My hands also will I lift up unto thy commandments,
which I have loved; * and my study shall be in thy statutes.

VII.
\emph{Memor esto verbi tui.}

O
THINK upon thy servant, as concerning thy word, * wherein thou hast caused
me to put my trust.

50 The same is my comfort in my trouble; * for thy
word hath quickened me.

51 The proud have had me exceedingly in derision; *
yet have I not shrinked from thy law.

52 For I remembered thine everlasting judgments, O
LORD, * and received comfort.

53 I am horribly afraid, * for the ungodly that forsake
thy law.

54 Thy statutes have been my songs, * in the house
of my pilgrimage.

55 I have thought upon thy Name, O LORD, in the night
season, * and have kept thy law.

56 This I had, * because I kept thy commandments.

VIII.
\emph{Portio mea, Domine.}

THOU
art my portion, O LORD; * I have promised to keep thy law.

58 I made my humble petition in thy presence with my
whole heart; * O be merciful unto me, according to thy word.

59 I called mine own ways to remembrance, * and turned
my feet unto thy testimonies.

60 I made haste, and prolonged not the time, * to keep
thy commandments. 3
mine
until 1892

61 The snares of the ungodly have compassed me about; * but I have not forgotten
thy law.

62 At midnight I will rise to give thanks unto thee,
* because of thy righteous judgments.

63 I am a companion of all them that fear thee, * and
keep thy commandments.

64 The earth, O LORD, is full of thy mercy: * O teach
me thy statutes.

IX.
\emph{Bonitatem fecisti.}

O
LORD, thou hast dealt graciously with thy servant, * according unto thy
word. 61 The congregations of the
ungodly have robbed me : but I ... until
1928

66 O teach4 me true understanding and knowledge; * for I have
believed thy commandments.

67 Before I was troubled, I went wrong; * but now have
I kept thy word.

68 Thou art good and gracious; * O teach me thy statutes.

69 The proud have imagined a lie against me; * but I
will keep thy commandments with my whole heart.

70 Their heart is as fat as brawn; * but my delight hath
been in thy law.

71 It is good for me that I have been in trouble; * that
I may learn thy statutes.

72 The law of thy mouth is dearer unto me * than thousands
of gold and silver.

Evening Prayer.

X. \emph{Manus
tuae fecerunt me.}

THY
hands have made me and fashioned me: * O give me understanding, that I
may learn thy commandments.

74 They that fear thee will be glad when they see me;
* because I have put my trust in thy word.

75 I know, O LORD, that thy judgments are right, *
and that thou of very faithfulness hast caused me to be troubled.

76 O let thy merciful kindness be my comfort, * according
to thy word unto thy servant.

77 O let thy loving mercies come unto me, that I may
live; * for thy law is my delight.

78 Let the proud be confounded, for they go wickedly
about to destroy me; * but I will be occupied in thy commandments.

79 Let such as fear thee, and have known thy testimonies,
* be turned unto me.

80 O let my heart be sound in thy statutes, * that
I be not ashamed.

XI. \emph{Defecit
anima mea.}

MY
soul hath longed for thy salvation, * and I have a good hope because of
thy word.

82 Mine eyes long sore for thy word; * saying, O when
wilt thou comfort me?

83 For I am become like a bottle in the smoke; * yet
do I not forget thy statutes.

84 How many are the days of thy servant? * when wilt
thou be avenged of them that persecute me?

85 The proud have digged pits for me, * which are not
after thy law.

86 All thy commandments are true: * they persecute
me falsely; O be thou my help.

87 They had almost made an end of me upon earth; *
but I forsook not thy commandments.

88 O quicken me after thy loving-kindness; * and so
shall I keep the testimonies of thy mouth.

XII.
\emph{In aeternum, Domine.}

O
LORD, thy word * endureth for ever in heaven. 

90 Thy truth also remaineth from one generation to
another; * thou hast laid the foundation of the earth, and it abideth.

91 They continue this day according to thine ordinance;
* for all things serve thee.

92 If my delight had not been in thy law, * I should
have perished in my trouble.

93 I will never forget thy commandments; * for with
them thou hast quickened me.

94 I am thine: O save me, * for I have sought thy commandments.

95 The ungodly laid wait for me, to destroy me; * but
I will consider thy testimonies.

96 I see that all things come to an end; * but thy
commandment is exceeding broad.

XIII\emph{.
Quomodo dilexi!}

LORD,
what love have I unto thy law! * all the day long is my study in it.

98 Thou, through thy commandments, hast made me wiser
than mine enemies; * for they are ever with me.

99 I have more understanding than my teachers; * for
thy testimonies are my study.

100 I am wiser than the aged; * because I keep thy
commandments.

101 I have refrained my feet from every evil way, *
that I may keep thy word.

102 I have not shrunk from thy judgments; * for thou
teachest me.

103 O how sweet are thy words unto my throat; * yea,
sweeter than honey unto my mouth!

104 Through thy commandments I get understanding: *
therefore I hate all evil ways. 

The
Twenty-sixth Day.

Morning
Prayer.

XIV. \emph{Lucerna
pedibus meis.}

THY
word is a lantern unto my feet, * and a light unto my paths.

106 I have sworn, and am stedfastly purposed, * to
keep thy righteous judgments.

107 I am troubled above measure: * quicken me, O LORD,
according to thy word.

108 Let the free-will offerings of my mouth please
thee, O LORD; * and teach me thy judgments.

109 My soul is alway in my hand; * yet do I not forget
thy law.

110 The ungodly have laid a snare for me; * but yet
I swerved not from thy commandments.

111 Thy testimonies have I claimed as mine heritage
for ever; * and why? they are the very joy of my heart.

112 I have applied my heart to fulfil thy statutes
alway, * even unto the end.

XV. \emph{Iniquos
odio habui.}

I
HATE them that imagine evil things; * but thy law do I love.

114 Thou art my defence and shield; * and my trust
is in thy word.

115 Away from me, ye wicked; * I will keep the commandments
of my God.

116 O stablish me according to thy word, that I may
live; * and let me not be disappointed of my hope.

117 Hold thou me up, and I shall be safe; * yea, my
delight shall be ever in thy statutes.

118 Thou hast trodden down all them that depart from
thy statutes; * for they imagine but deceit.

119 Thou puttest away all the ungodly of the earth
like dross; * therefore I love thy testimonies.

120 My flesh trembleth for fear of thee; * and I am
afraid of thy judgments.

XVI.
\emph{Feci judicium.}

I
DEAL with the thing that is lawful and right; * O give me not over unto
mine oppressors.

122 Make thou thy servant to delight in that which
is good, * that the proud do me no wrong.

123 Mine eyes are wasted away with looking for thy
health, * and for the word of thy righteousness.

124 O deal with thy servant according unto thy loving
mercy, * and teach me thy statutes.

125 I am thy servant; O grant me understanding, * that
I may know thy testimonies.

126 It is time for thee, LORD, to lay to thine hand;
* for they have destroyed thy law. 4
learn until
1928

127 For I love thy commandments * above gold and precious stones5.

128 Therefore hold I straight all thy commandments;
* and all false ways I utterly abhor.

XVII.
\emph{Mirabilia.}

THY
testimonies are wonderful; * therefore doth my soul keep them.

130 When thy word goeth forth, * it giveth light and
understanding unto the simple.

131 I opened my mouth, and drew in my breath; * for
my delight was in thy commandments.

132 O look thou upon me, and be merciful unto me, *
as thou usest to do unto those that love thy Name.

133 Order my steps in thy word; * and so shall no wickedness
have dominion over me.

134 O deliver me from the wrongful dealings of men;
* and so shall I keep thy commandments.

135 Show the light of thy countenance upon thy servant,
* and teach me thy statutes.

136 Mine eyes gush out with water, * because men keep
not thy law.

XVIII.\emph{Justus es, Domine.}

RIGHTEOUS
art thou, O LORD; * and true are thy judgments.

138 The testimonies that thou hast commanded * are
exceeding righteous and true.

139 My zeal hath even consumed me; * because mine enemies
have forgotten thy words.

140 Thy word is tried to the uttermost, * and thy servant
loveth it.

141 I am small and of no reputation; * yet do I not
forget thy commandments.

142 Thy righteousness is an everlasting righteousness,
* and thy law is the truth.

143 Trouble and heaviness have taken hold upon me;
* yet is my delight in thy commandments.

144 The righteousness of thy testimonies is everlasting:
* O grant me understanding, and I shall live. 

Evening Prayer.

XIX. \emph{Clamavi
in toto corde meo.}

I
CALL with my whole heart; * hear me, O LORD; I will keep thy statutes.

146 Yea, even unto thee do I call; * help me, and I
shall keep thy testimonies.

147 Early in the morning do I cry unto thee; * for
in thy word is my trust. 

5 stone in
the Engl. Book

148 Mine eyes prevent the night watches; * that I might be occupied in thy
word6.

149 Hear my voice, O LORD, according unto thy lovingkindness;
* quicken me, according to thy judgments.

150 They draw nigh that of malice persecute me, * and
are far from thy law.

151 Be thou nigh at hand, O LORD; * for all thy commandments
are true.

152 As concerning thy testimonies, I have known long
since, * that thou hast grounded them for ever.

XX.
\emph{Vide humilitatem.}

O
CONSIDER mine adversity, and deliver me, * for I do not forget thy law.

154 Avenge thou my cause, and deliver me; * quicken
me according to thy word.

155 Health is far from the ungodly; * for they regard
not thy statutes.

156 Great is thy mercy, O LORD; * quicken me, as thou
art wont.

157 Many there are that trouble me, and persecute me;
* yet do I not swerve from thy testimonies.

158 It grieveth me when I see the transgressors; *
because they keep not thy law.

159 Consider, O LORD, how I love thy commandments;
* O quicken me, according to thy loving-kindness.

160 Thy word is true from everlasting; * all the judgments
of thy righteousness endure for evermore.

XXI. \emph{Principes
persecuti sunt.}

PRINCES
have persecuted me without a cause; * but my heart standeth in awe of
thy word.

162 I am as glad of thy word, * as one that findeth
great spoils.

163 As for lies, I hate and abhor them; * but thy law
do I love.

164 Seven times a day do I praise thee; * because of
thy righteous judgments.

6 words
until
1928

149 ... quicken me, according as thou art wont. until
1928

165 Great is the peace that they have who love thy law; * and they have
none occasion of stumbling.

166 LORD, I have looked for thy saving health, * and
done after thy commandments.

167 My soul hath kept thy testimonies, * and loved
them exceedingly.

168 I have kept thy commandments and testimonies; *
for all my ways are before thee.

XXII.\emph{Appropinquet deprecatio.}

LET
my complaint come before thee, O LORD; * give me understanding according
to thy word.

170 Let my supplication come before thee; * deliver
me according to thy word.

171 My lips shall speak of thy praise, * when thou
hast taught me thy statutes.

172 Yea, my tongue shall sing of thy word; * for all
thy commandments are righteous.

173 Let thine hand help me; * for I have chosen thy
commandments.

174 I have longed for thy saving health, O LORD; *
and in thy law is my delight.

175 O let my soul live, and it shall praise thee; *
and thy judgments shall help me.

176 I have gone astray like a sheep that is lost; *
O seek thy servant, for I do not forget thy commandments. 

The
Twenty-seventh Day.

Morning
Prayer.

Psalm
120. \emph{Ad Dominum.}

WHEN
I was in trouble, I called upon the LORD, * and he heard me.

2 Deliver my soul, O LORD, from lying lips, * and from
a deceitful tongue.

3 What reward shall be given or done unto thee, thou
false tongue? * even mighty and sharp arrows, with hot burning coals. 165 ... and they are not offended
at it. until
1928

4 Woe is me, that I am constrained to dwell with Meshech1, *
and to have my habitation among the tents of Kedar!

5 My soul hath long dwelt among them * that are enemies
unto peace.

6 I labour for peace; but when I speak unto them thereof,
* they make them ready to battle.

1
Mesech until
1928

Psalm
121. \emph{Levavi oculos.}

I
WILL lift up mine eyes unto the hills; * from whence cometh my help?1

2 My help cometh even from the LORD, * who hath made
heaven and earth.

3 He will not suffer thy foot to be moved; * and he
that keepeth thee will not sleep.

4 Behold, he that keepeth Israel * shall neither slumber
nor sleep.

5 The LORD himself is thy keeper; * the LORD is thy
defence upon thy right hand;6 So that the sun shall not burn thee by
day, * neither the moon by night.

7 The LORD shall preserve thee from all evil; * yea,
it is even he that shall keep thy soul.

8 The LORD shall preserve thy going out, and thy coming
in, * from this time forth for evermore.

Psalm
122. \emph{Lætatus sum.}

I
WAS glad when they said unto me, * We will go into the house of the LORD.

2 Our feet shall stand in thy gates, * O Jerusalem.

3 Jerusalem is built as a city * that is at unity in
itself.

4 For thither the tribes go up, even the tribes of
the LORD, * to testify unto Israel, to give thanks unto the Name of the
LORD.

5 For there is the seat of judgment, * even the seat
of the house of David.

6 O pray for the peace of Jerusalem; * they shall prosper
that love thee.

7 Peace be within thy walls, * and plenteousness within
thy palaces.

8 For my brethren and companions' sakes, * I will wish
thee prosperity.

9 Yea, because of the house of the LORD our God, *
I will seek to do thee good.

Psalm
123. \emph{Ad te levavi oculos meos.}

UNTO
thee lift I up mine eyes, * O thou that dwellest in the heavens.

2 Behold, even as the eyes of servants look unto the
hand of their masters, and as the eyes of a maiden unto the hand of her
mistress, * even so our eyes wait upon the LORD our God, until he have
mercy upon us.

3 Have mercy upon us, O LORD, have mercy upon us; *
for we are utterly despised.

4 Our soul is filled with the scornful reproof of the
wealthy, * and with the despitefulness of the proud.

Psalm
124. \emph{Nisi quia Dominus.}

IF
the LORD himself had not been on our side, now may Israel say; * if the
LORD himself had not been on our side, when men rose up against us;

\emph{Levavi oculos meos} until
1871

1 help. (i.
e., no question mark) until 1928

2 They
had swallowed us up alive1; * when they were so wrathfully
displeased at us.

3 Yea, the waters had drowned us, * and the stream
had gone over our soul.

4 The deep waters of the proud * had gone even over
our soul.

5 But praised be the LORD, * who hath not given us
over for a prey unto their teeth.

6 Our soul is escaped even as a bird out of the snare
of the fowler; * the snare is broken, and we are delivered.

7 Our help standeth in the Name of the LORD, * who
hath made heaven and earth.

Psalm
125. \emph{Qui confidunt.}

THEY
that put their trust in the LORD shall be even as the mount Sion, * which
may not be removed, but standeth fast for ever.

2 The hills stand about Jerusalem; * even so standeth
the LORD round about his people, from this time forth for evermore.

1 quick until
1928

3 For
the sceptre of the ungodly shall not abide upon the lot of the righteous;
* lest the righteous put their hand unto wickedness.

4 Do well, O LORD, * unto those that are good and true
of heart.

5 As for such as turn back unto their own wickedness,
* the LORD shall lead them forth with the evil doers; but peace shall
be upon Israel. 

Evening
Prayer.

Psalm
126. \emph{In convertendo.}

WHEN
the LORD turned again the captivity of Sion, * then were we like unto
them that dream.

2 Then was our mouth filled with laughter, * and our
tongue with joy.

3 Then said they among the heathen, * The LORD hath
done great things for them.

4 Yea, the LORD hath done great things for us already;
* whereof we rejoice.

5 Turn our captivity, O LORD, * as the rivers in the
south.

6 They that sow in tears * shall reap in joy.

7 He that now goeth on his way weeping, and beareth
forth good seed, * shall doubtless come again with joy, and bring his
sheaves with him.

Psalm
127.\emph{Nisi Dominus.}

EXCEPT
the LORD build the house, * their labour is but lost that build it.

2 Except the LORD keep the city, * the watchman waketh
but in vain.

3 It is but lost labour that ye haste to rise up early,
and so late take rest, and eat the bread of carefulness; * for so he giveth
his beloved sleep.

4 Lo, children, and the fruit of the womb, * are an
heritage and gift that cometh of the LORD.

5 Like as the arrows in the hand of the giant, * even
so are the young children.

6 Happy is the man that hath his quiver full of them;
* they shall not be ashamed when they speak with their enemies in the
gate.

Psalm
128. \emph{Beati omnes.}

BLESSED
are all they that fear the LORD, * and walk in his ways. 3. For the rod of the ungodly
cometh not into the lot of the righteous : ... until
1928

2 For thou shalt eat the labours1 of thine hands: * O well is
thee, and happy shalt thou be.

3 Thy wife shall be as the fruitful vine * upon the walls
of thine house;

4 Thy children like the olive-branches * round about
thy table.

5 Lo, thus shall the man be blessed * that feareth the
LORD.

6 The LORD from out of Sion shall so bless thee, * that
thou shalt see Jerusalem in prosperity all thy life long;

7 Yea, that thou shalt see thy children's children, *
and peace upon Israel.

Psalm
129. \emph{Sæpe expugnaverunt.}

MANY
a time have they fought against me from my youth up, * may Israel now
say;

2 Yea, many a time have they vexed me from my youth
up; * but they have not prevailed against me.

3 The plowers plowed upon my back, * and made long
furrows.

4 But the righteous LORD * hath hewn the snares of
the ungodly in pieces.

5 Let them be confounded and turned backward, * as
many as have evil will at Sion. 1 labour
in 1789 Book

6 Let
them be even as the grass upon the housetops, * which withereth afore
it be grown up;

7 Whereof the mower filleth not his hand, * neither
he that bindeth up the sheaves his bosom.

8 So that they who go by say not so much as, The LORD
prosper you; * we wish you good luck in the Name of the LORD.

Psalm
130. \emph{De profundis.}

OUT
of the deep have I called unto thee, O LORD; * Lord, hear my voice.

2 O let thine ears consider well * the voice of my
complaint.

3 If thou, LORD, wilt be extreme to mark what is done
amiss, * O Lord, who may abide it?

4 For there is mercy with thee; * therefore shalt thou
be feared.

5 I look for the LORD; my soul doth wait for him; *
in his word is my trust. 6. Let them be even as the
grass growing upon the house-tops : which withereth afore it be plucked
up; until
1928

6 My soul fleeth unto the Lord before the morning watch; * I say, before
the morning watch.

7 O Israel, trust in the LORD; for with the LORD there
is mercy, * and with him is plenteous redemption.

8 And he shall redeem Israel * from all his sins.

Psalm
131. \emph{Domine, non est.}

LORD,
I am not high-minded; * I have no proud looks.

2 I do not exercise myself in great matters * which
are too high for me.

3 But I refrain my soul, and keep it low, like as a
child that is weaned from his mother: * yea, my soul is even as a weaned
child.

4 O Israel, trust in the LORD * from this time forth
for evermore. 

The
Twenty-eighth Day.

Morning
Prayer.

Psalm
132. \emph{Memento, Domine.}

LORD,
remember David, * and all his trouble:

2 How he sware unto the LORD, * and vowed a vow unto
the Almighty God of Jacob:

3 I will not come within the tabernacle of mine house,
* nor climb up into my bed;

4 I will not suffer mine eyes to sleep, nor mine eyelids
to slumber; * neither the temples of my head to take any rest;

5 Until I find out a place for the temple of the LORD;
* an habitation for the Mighty God of Jacob.

6. My soul fleeth unto the
Lord : before the morning

watch, I say, ... until
1928

6 Lo, we heard of the same at Ephratah1, * and found it in the
wood.

7 We will go into his tabernacle, * and fall low on our
knees before his footstool.

8 Arise, O LORD, into thy resting-place; * thou, and
the ark of thy strength.

9 Let thy priests be clothed with righteousness; * and
let thy saints sing with joyfulness. 1
Ephrata until
1928

10 For thy servant David's sake, * turn not away the face of thine anointed.

11 The LORD hath made a faithful oath unto David, * and
he shall not shrink from it:

10 ... thurn not away the
presence of thine anointed. until
1928

12 Of the fruit of thy body * shall I set upon thy throne2.

13 If thy children will keep my covenant, and my testimonies
that I shall teach them; * their children also shall sit upon thy throne2
for evermore.

14 For the LORD hath chosen Sion to be an habitation
for himself; * he hath longed for her.

15 This shall be my rest for ever: * here will I dwell,
for I have a delight therein.

16 I will bless her victuals with increase, * and will
satisfy her poor with bread.

17 I will deck her priests with health, * and her saints
shall rejoice and sing.

18 There shall I make the horn of David to flourish:
* I have ordained a lantern for mine anointed.

19 As for his enemies, I shall clothe them with shame;
* but upon himself shall his crown flourish.

2 seat until
1928

Psalm
133. \emph{Ecce, quam bonum!}

BEHOLD,
how good and joyful a thing it is, * for brethren to dwell together in
unity!

2 It is like the precious oil upon the head, that ran
down unto the beard, * even unto Aaron's beard, and went down to the skirts
of his clothing.

3 Like as the dew of Hermon, * which fell upon the
hill of Sion.

4 For there the LORD promised his blessing, * and life
for evermore.

Psalm
134. \emph{Ecce nunc.}

BEHOLD
now, praise the LORD, * all ye servants of the LORD;

2 Ye that by night stand in the house of the LORD,
* even in the courts of the house of our God.

3 Lift up your hands in the sanctuary, * and praise
the LORD.

4 The LORD that made heaven and earth * give thee blessing
out of Sion.

Psalm
135. \emph{Laudate Nomen.}

O
PRAISE the LORD, laud ye the Name of the LORD; * praise it, O ye servants
of the LORD;

2 Ye that stand in the house of the LORD, * in the
courts of the house of our God.

3 O praise the LORD, for the LORD is gracious; * O
sing praises unto his Name, for it is lovely.

4 For why? the LORD hath chosen Jacob unto himself,
* and Israel for his own possession.

5 For I know that the LORD is great, * and that our Lord is above all gods.

... it is : brethren to dwell
... in 1892 Book

... it is, brethren, to dwell ... until
1892

6 Whatsoever the LORD pleased, that did he in heaven,
and in earth; * and1 in the sea, and in all deep places.

7 He bringeth forth the clouds from the ends of the world,
* and sendeth forth lightnings with the rain, bringing the winds out of
his treasuries.

8 He smote the firstborn of Egypt, * both of man and
beast.

9 He hath sent tokens and wonders into the midst of thee,
O thou land of Egypt; * upon Pharaoh, and all his servants.

10 He smote divers nations, * and slew mighty kings: \\
1 and added
in 1871

11 Sihon2, king of the Amorites; and Og, the king of Bashan3;
* and all the kingdoms of Canaan;

12 And gave their land to be an heritage, * even an
heritage unto Israel his people.

13 Thy Name, O LORD, endureth for ever; * so doth thy
memorial, O LORD, from one generation to another.

14 For the LORD will avenge his people, * and be gracious
unto his servants.

15 As for the images of the heathen, they are but silver
and gold; * the work of men's hands.

16 They have mouths, and speak not; * eyes have they,
but they see not.

17 They have ears, and yet they hear not; * neither
is there any breath in their mouths.

18 They that make them are like unto them; * and so
are all they that put their trust in them.

19 Praise the LORD, ye house of Israel; * praise the
LORD, ye house of Aaron.

20 Praise the LORD, ye house of Levi; * ye that fear
the LORD, praise the LORD.

21 Praised be the LORD out of Sion, * who dwelleth
at Jerusalem. 

Evening
Prayer.

Psalm
136.\emph{Confitemini.}

O
GIVE thanks unto the LORD, for he is gracious: * and his mercy endureth
for ever.

2 O give thanks unto the God of all gods: * for his
mercy endureth for ever.

3 O thank the LORD of all lords: * for his mercy endureth
for ever.

4 Who only doeth great wonders: * for his mercy endureth
for ever.

5 Who by his excellent wisdom made the heavens: * for
his mercy endureth for ever.

6 Who laid out the earth above the waters: * for his
mercy endureth for ever.

7 Who hath made great lights: * for his mercy endureth
for ever:

8 The sun to rule the day: * for his mercy endureth
for ever;

9 The moon and the stars to govern the night: * for
his mercy endureth for ever.

10 Who smote Egypt, with their firstborn: * for his
mercy endureth for ever;

11 And brought out Israel from among them: * for his
mercy endureth for ever;

12 With a mighty hand and stretched-out arm: * for
his mercy endureth for ever.

13 Who divided the Red Sea in two parts: * for his
mercy endureth for ever;

14 And made Israel to go through the midst of it: *
for his mercy endureth for ever.

15 But as for Pharaoh and his host, he overthrew them
in the Red Sea: * for his mercy endureth for ever.

16 Who led his people through the wilderness: * for
his mercy endureth for ever.

17 Who smote great kings: * for his mercy endureth
for ever;

18 Yea, and slew mighty kings: * for his mercy endureth
for ever: 2
Sehon until
1928 \\
3
Basan until
1928

19 Sihon1, king of the Amorites: * for his mercy endureth for
ever;

20 And Og, the king of Bashan2: * for his
mercy endureth for ever;

21 And gave away their land for an heritage: * for
his mercy endureth for ever;

22 Even for an heritage unto Israel his servant: *
for his mercy endureth for ever.

23 Who remembered us when we were in trouble: * for
his mercy endureth for ever;

24 And hath delivered us from our enemies: * for his
mercy endureth for ever.

25 Who giveth food to all flesh: * for his mercy endureth
for ever.

26 O give thanks unto the God of heaven: * for his
mercy endureth for ever.

27 O give thanks unto the Lord of lords: * for his
mercy endureth for ever.

Psalm
137. \emph{Super flumina.}

BY
the waters of Babylon we sat down and wept, * when we remembered thee,
O Sion.

2 As for our harps, we hanged them up * upon the trees
that are therein.

3 For they that led us away captive, required of us
then a song, and melody in our heaviness: * Sing us one of the songs of
Sion.

4 How shall we sing the LORD’S song * in a strange
land?

5 If I forget thee, O Jerusalem, * let my right hand
forget her cunning.

1 Sehon until
1928

2 Basan until
1928

6 If I do not remember thee, let my tongue cleave to the roof of my mouth;
* yea, if I prefer not Jerusalem above my chief joy.

7 Remember the children of Edom, O LORD, in the day of Jerusalem; * how
they said, Down with it, down with it, even to the ground.

8 O daughter of Babylon, wasted with misery; * yea,
happy shall he be that rewardeth thee as thou hast served us.

9 Blessed shall he be that taketh thy children, * and
throweth them against the stones.

Psalm
138. \emph{Confitebor tibi.}

I
WILL give thanks unto thee, O Lord, with my whole heart; * even before
the gods will I sing praise unto thee.

2 I will worship toward thy holy temple, and praise
thy Name, because of thy loving-kindness and truth; * for thou hast magnified
thy Name, and thy word, above all things.

3 When I called upon thee, thou heardest me; * and
enduedst my soul with much strength.

4 All the kings of the earth shall praise thee, O LORD;
* for they have heard the words of thy mouth. 6 ... Jerusalem in my mirth
until
1928

5 Yea, they shall sing of1 the ways of the LORD, * that great
is the glory of the LORD.

6 For though the LORD be high, yet hath he respect
unto the lowly; * as for the proud, he beholdeth them afar off.

7 Though I walk in the midst of trouble, yet shalt
thou refresh me; * thou shalt stretch forth thy hand upon the furiousness
of mine enemies, and thy right hand shall save me.

8 The LORD shall make good his loving-kindness toward
me; * yea, thy mercy, O LORD, endureth for ever; despise not then the
works of thine own hands. 

The
Twenty-ninth Day.

Morning
Prayer.

Psalm
139. \emph{Domine, probasti.}

O
LORD, thou hast searched me out, and known me. * Thou knowest my down-sitting,
and mine uprising; thou understandest my thoughts long before. 1 in in
the 1892 Book

2 Thou art about my path, and about my bed; * and art acquainted with all
my ways.

3 For lo, there is not a word in my tongue, * but thou,
O LORD, knowest it altogether. 2 ... : and spiest out all
my ways. until
1928

4 Thou hast beset1 me behind and before, * and laid thine hand
upon me.

5 Such knowledge is too wonderful and excellent for me;
* I cannot attain unto it.

6 Whither shall I go then from thy Spirit? * or whither
shall I go then from thy presence?

7 If I climb up into heaven, thou art there; * if I go
down to hell, thou art there also.

8 If I take the wings of the morning, * and remain in
the uttermost parts of the sea;

9 Even there also shall thy hand lead me, * and thy right
hand shall hold me.

10 If I say, Peradventure the darkness shall cover me;
* then shall my night be turned to day.

11 Yea, the darkness is no darkness with thee, but the
night is as clear as the day; * the darkness and light to thee are both
alike.

12 For my reins are thine; * thou hast covered me in
my mother's womb.

13 I will give thanks unto thee, for I am fearfully and
wonderfully made: * marvellous are thy works, and that my soul knoweth right
well.

14 My bones are not hid from thee, * though I be made
secretly, and fashioned beneath in the earth.

15 Thine eyes did see my substance, yet being imperfect;
* and in thy book were all my members written;

16 Which day by day were fashioned, * when as yet there
was none of them.

17 How dear are thy counsels unto me, O God; * O how great is the sum
of them!

18 If I tell them, they are more in number than the
sand: * when I wake up, I am present with thee.

19 Wilt thou not slay the wicked, O God? * Depart from
me, ye blood-thirsty men.

20 For they speak unrighteously against thee; * and
thine enemies take thy Name in vain.

21 Do not I hate them, O LORD, that hate thee? * and
am not I grieved with those that rise up against thee?

22 Yea, I hate them right sore; * even as though they
were mine enemies.

23 Try me, O God, and seek the ground of my heart;
* prove me, and examine my thoughts.

24 Look well if there be any way of wickedness in me;
* and lead me in the way everlasting.

Psalm
140. \emph{Eripe me, Domine.}

DELIVER
me, O LORD, from the evil man; * and preserve me from the wicked man;

2 Who imagine mischief in their hearts, * and stir
up strife all the day long.

3 They have sharpened their tongues like a serpent;
* adder's poison is under their lips.

4 Keep me, O LORD, from the hands of the ungodly; *
preserve me from the wicked men, who are purposed to overthrow my goings.

5 The proud have laid a snare for me, and spread a
net abroad with cords; * yea, and set traps in my way.

6 I said unto the LORD, Thou art my God, * hear the
voice of my prayers, O LORD.

7 O LORD God, thou strength of my health; * thou hast
covered my head in the day of battle.

8 Let not the ungodly have his desire, O LORD; * let
not his mischievous imagination prosper, lest they be too proud.

9 Let the mischief of their own lips fall upon the
head of them * that compass me about.

10 Let hot burning coals fall upon them; * let them
be cast into the fire, and into the pit, that they never rise up again.

11 A man full of words shall not prosper upon the earth:
* evil shall hunt the wicked person to overthrow him.

12 Sure I am that the LORD will avenge the poor, *
and maintain the cause of the helpless.

13 The righteous also shall give thanks unto thy Name;
* and the just shall continue in thy sight. 

Evening
Prayer.

Psalm
141. \emph{Domine, clamavi.}

LORD,
I call upon thee; haste thee unto me, * and consider my voice, when I
cry unto thee.

2 Let my prayer be set forth in thy sight as the incense;
* and let the lifting up of my hands be an evening sacrifice.

3 Set a watch, O LORD, before my mouth, * and keep
the door of my lips. 1 fashioned until
1928

4 O let not mine heart be inclined to any evil thing; * let me not be
occupied in ungodly works with the men that work wickedness, neither let
me eat of such things as please them.

5 Let the righteous rather smite me friendly, and reprove
me; * yea, let not my head refuse their precious balms.

6 As for the ungodly, * I will pray yet against their
wickedness.

7 Let their judges be overthrown in stony places, *
that they may hear my words; for they are sweet.

8 Our bones lie scattered before the pit, * like as
when one breaketh and heweth wood upon the earth.

9 But mine eyes look unto thee, O LORD God; * in thee
is my trust; O cast not out my soul.

10 Keep me from the snare that they have laid for me,
* and from the traps of the wicked doers.

11 Let the ungodly fall into their own nets together,
* and let me ever escape them.

Psalm
142. \emph{Voce mea ad Dominum.}

I
CRIED unto the LORD with my voice; * yea, even unto the LORD did I make
my supplication.

2 I poured out my complaints before him, * and showed
him of my trouble.

3 When my spirit was in heaviness, thou knewest my
path; * in the way wherein I walked, have they privily laid a snare for
me.

4 I looked also upon my right hand, * and saw there
was no man that would know me.

5 I had no place to flee unto, * and no man cared for
my soul.

6 I cried unto thee, O LORD,
and said, * Thou art my hope, and my portion in the land of the living.

7 Consider my complaint; * for I am brought very low.

8 O deliver me from my persecutors; * for they are
too strong for me.

9 Bring my soul out of prison, that I may give thanks
unto thy Name; * which thing if thou wilt grant me, then shall the righteous
resort unto my company.

Psalm
143. \emph{Domine, exaudi.}

HEAR
my prayer, O LORD, and consider my desire; * hearken unto me for thy truth
and righteousness' sake.

2 And enter not into judgment with thy servant; * for
in thy sight shall no man living be justified.

3 For the enemy hath persecuted my soul; he hath smitten
my life down to the ground; * he hath laid me in the darkness, as the
men that have been long dead.

4 Therefore is my spirit vexed within me, * and my
heart within me is desolate.

5 Yet do I remember the time past; I muse upon all
thy works; * yea, I exercise myself in the works of thy hands.

6 I stretch forth my hands unto thee; * my soul gaspeth
unto thee as a thirsty land.

7 Hear me, O LORD, and that soon; for my spirit waxeth
faint: * hide not thy face from me, lest I be like unto them that go down
into the pit.

8 O let me hear thy loving-kindness betimes in the
morning; for in thee is my trust: * show thou me the way that I should
walk in; for I lift up my soul unto thee.

9 Deliver me, O LORD, from mine enemies; * for I flee
unto thee to hide me.

10 Teach me to do the thing that pleaseth thee; for
thou art my God: * let thy loving Spirit lead me forth into the land of
righteousness.

11 Quicken me, O LORD, for thy Name's sake; * and for
thy righteousness' sake bring my soul out of trouble.

12 And of thy goodness slay mine enemies, * and destroy
all them that vex my soul; for I am thy servant. 

The
Thirtieth Day

Morning Prayer

Psalm
144. \emph{Benedictus Dominus.}

BLESSED
be the LORD my strength, * who teacheth my hands to war, and my fingers
to fight:

2 My hope and my fortress, my castle and deliverer,
my defender in whom I trust; * who subdueth my people that is under me.

3 LORD, what is man, that thou hast such respect unto
him? * or the son of man, that thou so regardest him?

4 Man is like a thing of nought; * his time passeth
away like a shadow.

5 Bow thy heavens, O LORD, and come down; * touch the
mountains, and they shall smoke.

6 Cast forth thy lightning, and tear them; * shoot
out thine arrows, and consume them. 4. ... that work wickedness,
lest I eat of such things as please them.\\
5. Let the righteous
rather smite me friendly : and reprove me.\\
6. But
let not their precious balms break my head : yea, I will pray yet against
their wickedness. until
1928

7 Send down thine hand from above; * deliver me, and take me out of the
great waters, from the hand of strangers1;

8 Whose mouth talketh of vanity, * and their right hand
is a right hand of wickedness.

9 I will sing a new song unto thee, O God; * and sing
praises unto thee upon a ten-stringed lute.

10 Thou hast given victory unto kings, * and hast delivered
David thy servant from the peril of the sword.

1 strange children until
1928

11 Save me, and deliver me from the hand of strangers1, * whose
mouth talketh of vanity, and their right hand is a right hand of iniquity:

12 That our sons may grow up as the young plants, *
and that our daughters may be as the polished corners of the temple;

13 That our garners may be full and plenteous with
all manner of store; * that our sheep may bring forth thousands, and
ten thousands in our fields;

14 That our oxen may be strong to labour; that there
be no decay, * no leading into captivity, and no complaining in our streets.

15 Happy are the people that are in such a case; *
yea, blessed are the people who have the LORD for their God.

Psalm
145. \emph{Exaltabo te, Deus.}

I
WILL magnify thee, O God, my King; * and I will praise thy Name for ever
and ever.

2 Every day will I give thanks unto thee; * and praise
thy Name for ever and ever.

3 Great is the LORD, and marvellous worthy to be praised;
* there is no end of his greatness.

4 One generation shall praise thy works unto another,
* and declare thy power.

5 As for me, I will be talking of thy worship, * thy
glory, thy praise, and wondrous works;

6 So that men shall speak of the might of thy marvellous
acts; * and I will also tell of thy greatness.

7 The memorial of thine abundant kindness shall be
showed; * and men shall sing of thy righteousness.

8 The LORD is gracious and merciful; * long-suffering,
and of great goodness.

9 The LORD is loving unto every man; * and his mercy
is over all his works.

10 All thy works praise thee, O LORD; * and thy saints
give thanks unto thee.

11 They show the glory of thy kingdom, * and talk of
thy power;

12 That thy power, thy glory, and mightiness of thy
kingdom, * might be known unto men.

13 Thy kingdom is an everlasting kingdom, * and thy
dominion endureth throughout all ages.

14 The LORD upholdeth all such as fall, * and lifteth
up all those that are down.

15 The eyes of all wait upon thee, O Lord; * and thou
givest them their meat in due season.

16 Thou openest thine hand, * and fillest all things
living with plenteousness.

17 The LORD is righteous in all his ways, * and holy
in all his works.

18 The LORD is nigh unto all them that call upon him;
* yea, all such as call upon him faithfully.

19 He will fulfil the desire of them that fear him;
* he also will hear their cry, and will help them.

20 The LORD preserveth all them that love him; * but
scattereth abroad all the ungodly.

21 My mouth shall speak the praise of the LORD; * and
let all flesh give thanks unto his holy Name for ever and ever.

Psalm
146. \emph{Lauda, anima mea.}

PRAISE
the LORD, O my soul: while I live, will I praise the LORD; * yea, as long
as I have any being, I will sing praises unto my God.

2 O put not your trust in princes, nor in any child
of man; * for there is no help in them.

3 For when the breath of man goeth forth, he shall
turn again to his earth, * and then all his thoughts perish.

4 Blessed is he that hath the God of Jacob for his
help, * and whose hope is in the LORD his God:

5 Who made heaven and earth, the sea, and all that
therein is; * who keepeth his promise for ever;

6 Who helpeth them to right that suffer wrong; * who
feedeth the hungry.

7 The LORD looseth men out of
prison; * the LORD giveth sight to the blind.

8 The LORD helpeth them that are fallen; * the LORD
careth for the righteous.

9 The LORD careth for the strangers; he defendeth the
fatherless and widow: * as for the way of the ungodly, he turneth it upside
down.

10 The LORD thy God, O Sion, shall be King for evermore,
* and throughout all generations. 

Evening
Prayer

Psalm
147. \emph{Laudate Dominum.}

PRAISE
the LORD, for it is a good thing to sing praises unto our God; * yea,
a joyful and pleasant thing it is to be thankful.

2 The LORD doth build up Jerusalem, * and gather together
the outcasts of Israel.

3 He healeth those that are broken in heart, * and
giveth medicine to heal their sickness.

4 He telleth the number of the stars, * and calleth
them all by their names.

5 Great is our Lord, and great is his power; * yea,
and his wisdom is infinite.

6 The LORD setteth up the meek, * and bringeth the
ungodly down to the ground.

7 O sing unto the LORD with thanksgiving; * sing praises
upon the harp unto our God:

8 Who covereth the heaven with clouds, and prepareth
rain for the earth; * and maketh the grass to grow upon the mountains,
and herb for the use of men;

9 Who giveth fodder unto the cattle, * and feedeth
the young ravens that call upon him.

10 He hath no pleasure in the strength of an horse;
* neither delighteth he in any man's legs.

11 But the LORD’S delight is in them that fear
him, * and put their trust in his mercy.

12 Praise the LORD, O Jerusalem; * praise thy God,
O Sion.

13 For he hath made fast the bars of thy gates, * and
hath blessed thy children within thee.

14 He maketh peace in thy borders, * and filleth thee
with the flour of wheat.

15 He sendeth forth his commandment upon earth, * and
his word runneth very swiftly.

16 He giveth snow like wool, * and scattereth the hoarfrost
like ashes.

17 He casteth forth his ice like morsels: * who is
able to abide his frost?

18 He sendeth out his word, and melteth them: * he
bloweth with his wind, and the waters flow.

19 He showeth his word unto Jacob, * his statutes and
ordinances unto Israel.

20 He hath not dealt so with any nation; * neither
have the heathen knowledge of his laws.

1 strange children until
1928

Psalm
148. \emph{Laudate Dominum.}

O
PRAISE the LORD from the heavens: * praise him in the heights.

2 Praise him, all ye angels of his: * praise him, all
his host1.

3 Praise him, sun and moon: * praise him, all ye stars
and light.

4 Praise him, all ye heavens, * and ye waters that
are above the heavens.

5 Let them praise the Name of the LORD: * for he spake
the word, and they were made; he commanded, and they were created.

6 He hath made them fast for ever and ever: * he hath
given them a law which shall not be broken.

O PRAISE the LORD
of heaven : praise him in the height. until
1928 

1 hosts 1822-1892

7 Praise the LORD from the earth, * ye dragons and
all deeps;

8 Fire and hail, snow and vapours, * wind and storm,
fulfilling his word;

9 Mountains and all hills; * fruitful trees and all cedars; 7 Praise the LORD
upon earth : ... until
1928

10 Beasts and all cattle; * creeping things and flying fowls;

11 Kings of the earth, and all peoples2;
* princes, and all judges of the world;

12 Young men and maidens, old men and children, praise
the Name of the LORD: * for his Name only is excellent, and his praise
above heaven and earth.

13 He shall exalt the horn of his people: all his saints
shall praise him; * even the children of Israel, even the people that
serveth him.

Psalm
149. \emph{Cantate Domino.}

O
SING unto the LORD a new song; * let the congregation of saints praise
him.

2 Let Israel rejoice in him that made him, * and let
the children of Sion be joyful in their King.

3 Let them praise his Name in the dance: * let them
sing praises unto him with tabret and harp.

4 For the LORD hath pleasure in his people, * and helpeth
the meek-hearted.

5 Let the saints be joyful with glory; * let them rejoice
in their beds.

6 Let the praises of God be in their mouth; * and a
two-edged sword in their hands; 10 Beasts and all cattle :
worms and feathered fowls; until
1928 \\
2 people until
1928

7 To be avenged of the nations1, * and to
rebuke the peoples;

8 To bind their kings in chains, * and their nobles with
links of iron;

9 To execute judgment upon them; as it is written, *
Such honour have all his saints. 1
heathen until
1928

9. That they may be avenged
of them, as it is ... until
1928

Psalm 150. \emph{Laudate Dominum.}

O
PRAISE God in his sanctuary1: * praise him in the firmament
of his power.

2 Praise him in his noble acts: * praise him according
to his excellent greatness.

3 Praise him in the sound of the trumpet: * praise
him upon the lute and harp.

1 holiness
until
1928

4 Praise him in the timbrels2 and dances: * praise him upon the
strings and pipe.

5 Praise him upon the well-tuned cymbals: * praise him
upon the loud cymbals.

6 Let every thing that hath breath * praise the LORD.

\emph{The
End of the Psalter.}

2
cymbals until
1928

\begin{quote}

\emph{Psalms 1-50;
51-100 Return to the U. S. Book of Common Prayer: 1789;
1892; 1928}

\end{quote}

\xchapter{The Form and Manner of Making, Ordaining, and Consecrating Bishops, Priests, and Deacons}

\textbf{The Book of Common Prayer (1928)}

The Ordinal

being the

Form of Making, Ordaining, and Consecrating

Bishops, Priests, and Deacons

\emph{together with}

The Form of Consecration of a Church

An Office of Institution of Ministers

The Form and Manner of

Making, Ordaining, and Consecrating

Bishops, Priests, and Deacons

according to the

Order of the Protestant Episcopal Church

in the United States of America,

as established by the Bishops, the Clergy, and Laity

of said Church, in General Convention,

September, A.D. 1792 

THE PREFACE.

\emph{I}\emph{T is evident unto all men,
diligently reading Holy Scripture and ancient Authors, that from the Apostles'
time there have been these Orders of Ministers in Christ's Church,—Bishops,
Priests, and Deacons. Which Offices were evermore had in such reverend
estimation, that no man might presume to execute any of them, except he
were first called, tried, examined, and known to have such qualities as
are requisite for the same; and also by public Prayer, with Imposition
of Hands, were approved and admitted thereunto by lawful Authority. And
therefore, to the intent that these Orders may be continued, and reverently
used and esteemed in this Church, no man shall be accounted or taken to
be a lawful Bishop, Priest, or Deacon, in this Church, or suffered to
execute any of the said Functions, except he be called, tried, examined,
and admitted thereunto, according to the Form hereafter following, or
hath had Episcopal Consecration or Ordination. And none shall be admitted a Deacon, Priest, or Bishop,
except he be of the age which the Canon in that case provided may require. And the Bishop, knowing either by himself, or by sufficient
testimony, any Person to be a man of virtuous conversation, and without
crime, and, after examination and trial, finding him sufficiently instructed
in the Holy Scripture, and otherwise learned as the Canons require, may,
at the times appointed, or else, on urgent occasion, upon some other day,
in the face of the Church, admit him a Deacon, in such manner and form
as followeth.}

The
Form and Manner of

Making
Deacons

¶ \emph{When the day appointed by the Bishop is come,
there shall be a Sermon, or Exhortation, declaring the Duty and Office
of such as come to be admitted Deacons; how necessary that Order is in
the Church of Christ, and also, how the People ought to esteem them in
their Office.}

¶ \emph{The Sermon being ended, a Priest shall present unto the Bishop,
sitting in his chair near to the Holy Table, such as desire to be ordained
Deacons, each of them being decently habited, saying these words,}

REVEREND
Father in God, I present unto you these persons present, to be admitted
Deacons.

¶
\emph{The Bishop.}

TAKE
heed that the persons, whom ye present unto us, be apt and meet, for their
learning and godly conversa-tion, to exercise their Ministry duly, to
the honour of God, and the edifying of his Church.

¶
\emph{The Priest shall answer,}

I
HAVE inquired concerning them, and also examined them, and think them
so to be.

¶
\emph{Then the Bishop shall say unto the People,}

BRETHREN,
if there be any of you who knoweth any Impediment, or notable Crime, in
any of these persons presented to be ordered Deacons, for the which he
ought not to be admitted to that Office, let him come forth in the Name
of God, and show what the Crime or Impediment is.

¶
\emph{And if any great Crime or Impediment be objected, the Bishop shall
cease from Ordering that person, until such time as the party accused
shall be found clear of that Crime.}

¶ \emph{Then the Bishop (commending such as shall be found meet to
be Ordered, to the Prayers of the congregation) shall, with the Clergy
and People present, say the Litany.}

¶
\emph{And Note, That after the Suffrage,}That it may please thee to
illuminate all Bishops, \emph{etc., shall be said the following Suffrage:}

THAT
it may please thee to bless these thy servants, now to be admitted to
the Order of Deacons, and to pour thy grace upon them; that they may duly
execute their Office, to the edifying of thy Church, and the glory of
thy holy Name; 

¶
\emph{And}NOTE\emph{further, That in the discretion
of the Bishop, instead of the Litany appointed, may be said the Litany
for Ordinations.}

¶
\emph{Then shall be said the Service for the Communion, with the Collect,
Epistle, and Gospel, as followeth.}

\emph{The
Collect.}

ALMIGHTY
God, who by thy divine providence hast appointed divers Orders of Ministers
in thy Church, and didst inspire thine Apostles to choose into the Order
of Deacons the first Martyr Saint Stephen, with others; Mercifully behold
these thy servants now called to the like Office and Administration: so
replenish them with the truth of thy Doctrine, and adorn them with innocency
of life, that, both by word and good example, they may faithfully serve
thee in this Office, to the glory of thy Name, and the edification of
thy Church; through the merits of our Saviour Jesus Christ, who liveth
and reigneth with thee and the Holy Ghost, now and for ever. \emph{Amen}.

\emph{The
Epistle}. 1 Timothy iii. 8.

LIKEWISE
must the deacons be grave, not double tongued, not given to much wine,
not greedy of filthy lucre; holding the mystery of the faith in a pure
conscience. And let these also first be proved; then let them use the
office of a deacon, being found blameless. Even so must their wives be
grave, not slanderers, sober, faithful in all things. Let the deacons
be the husbands of one wife, ruling their children and their own houses
well. For they that have used the office of a deacon well purchase to
them-selves a good degree, and great boldness in the faith which is in
Christ Jesus.

¶
\emph{Or else}.

\emph{The
Epistle}. Acts vi. 2.

THEN
the twelve called the multitude of the disciples unto them, and said,
It is not reason that we should leave the word of God, and serve tables.
Wherefore, breth-ren, look ye out among you seven men of honest report,
full of the Holy Ghost and wisdom, whom we may appoint over this business.
But we will give ourselves continually to prayer, and to the ministry
of the word. And the saying pleased the whole multitude: and they chose
Stephen, a man full of faith and of the Holy Ghost, and Philip, and Prochorus,
and Nicanor, and Timon, and Parmenas, and Nicolas a proselyte of Antioch:
whom they set before the apostles: and when they had prayed, they laid
their hands on them. And the word of God increased; and the number of
the disciples multiplied in Jerusalem greatly; and a great company of
the priests were obedient to the faith.

¶
\emph{Then, the People being seated, the Bishop shall examine every one
of those who are to be Ordered, in the presence of the People, after this
manner following.}

DO
you trust that you are inwardly moved by the Holy Ghost to take upon you
this Office and Ministration, to serve God for the promoting of his glory,
and the edifying of his people?

\emph{Answer}. I trust so.

\emph{Bishop}. Do you think that you are truly called,
according to the will of our Lord Jesus Christ, and according to the Canons
of this Church, to the Ministry of the same?

\emph{Answer}. I think so.

\emph{Bishop}. Are you persuaded that the Holy
Scriptures contain all Doctrine required as necessary for eternal salvation
through faith in Jesus Christ?

\emph{Answer}. I am so persuaded.

\emph{Bishop}. Will you diligently read the same
unto the people assembled in the Church where you shall be appointed to
serve?

\emph{Answer}. I will.

\emph{Bishop}. It appertaineth to the Office
of a Deacon, in the Church where he shall be appointed to serve, to assist
the Priest in Divine Service, and specially when he ministereth the Holy
Communion, and to help him in the distribution thereof; and to read Holy
Scriptures and Homilies in the Church; and to instruct the youth in the
Catechism; in the absence of the Priest to baptize infants; and to preach,
if he be admitted thereto by the Bishop. And furthermore, it is his Office,
where provision is so made, to search for the sick, poor, and impotent
people of the Parish, that they may be relieved with the alms of the Parishioners,
or others. Will you do this gladly and willingly?

\emph{Answer}. I will so do, by the help of God.

\emph{Bishop}. Will you apply all your diligence
to frame and fashion your own lives, and the lives of your families, according
to the Doctrine of Christ; and to make both yourselves and them, as much
as in you lieth, wholesome examples of the flock of Christ?

\emph{Answer}. I will so do, the Lord being my helper.

\emph{Bishop}. Will you reverently obey your Bishop,
and other chief Ministers, who, according to the Canons of the Church,
may have the charge and government over you; following with a glad mind
and will their godly admonitions?

\emph{Answer}. I will endeavour so to do, the Lord
being my helper.

¶
\emph{Then, the People standing, the Bishop shall lay his Hands severally
upon the Head of every one to be made Deacon, humbly kneeling before
him, and shall say,}

TAKE
thou Authority to execute the Office of a Deacon in the Church of God
committed unto thee; In the Name of the Father, and of the Son, and
of the Holy Ghost. Amen.

¶ \emph{Then shall the
Bishop deliver to every one of them the New Testament, saying,}

TAKE thou Authority to
read the Gospel in the Church of God, and to preach the same, if thou
be thereto licensed by the Bishop himself.

¶ \emph{Then one of them, appointed
by the Bishop, shall read the Gospel.}

\emph{The Gospel.} St. Luke xii. 35.

LET your loins be girded
about, and your lights burning; and ye yourselves like unto men that wait
for their lord, when he will return from the wedding; that when he cometh
and knocketh, they may open unto him immediately. Blessed are those servants,
whom the lord when he cometh shall find watching: verily I say unto you,
that he shall gird himself, and make them to sit down to meat, and will
come forth and serve them. And if he shall come in the second watch, or
come in the third watch, and find them so, blessed are those servants.

¶ \emph{Then shall the
Bishop proceed in the Communion; and all who are Ordered shall tarry,
and receive the Holy Communion the same day, with the Bishop.}

¶ \emph{The Communion
ended, after the last Collect, and immediately before the Benediction,
shall be said this Collect following.}

ALMIGHTY God, giver of
all good things, who of thy great goodness hast vouchsafed to accept and
take these thy servants unto the Office of Deacons in thy Church; Make
them, we beseech thee, O Lord, to be modest, humble, and constant in their
Ministration, to have a ready will to observe all spiritual Discipline;
that they, having always the testimony of a good conscience, and continuing
ever stable and strong in thy Son Christ, may so well behave themselves
in this inferior Office, that they may be found worthy to be called unto
the higher Ministries in thy Church; through the same thy Son our Saviour
Jesus Christ, to whom be glory and honour, world without end. \emph{Amen}.

THE Peace of God, which
passeth all understanding, keep your hearts and minds in the knowledge
and love of God, and of his Son Jesus Christ our Lord: And the Blessing
of God Almighty, the Father, the Son, and the Holy Ghost, be amongst you,
and remain with you always. \emph{Amen}.

¶ \emph{And here it
must be declared unto the Deacon, that he must continue in that Office
of a Deacon the space of a whole year, (except for reasonable causes it
shall otherwise seem good unto the Bishop,) to the intent he may be perfect
and well expert in the things appertaining to the Ecclesiastical Administration.
In executing whereof if he be found faithful and diligent, he may be admitted
by his Diocesan to the Order of Priesthood, at the times appointed in
the Canon; or else, on urgent occasion, upon some other day, in the face
of the Church, in such manner and form as hereafter.}

\emph{}

The
Form and Manner of

Ordering
Priests

¶
\emph{When the day appointed by the Bishop is come, there shall be a Sermon
or Exhortation, declaring the Duty and Office of such as come to be
admitted Priests; how necessary that Order is in the Church of Christ,
and also, how the People ought to esteem them in their Office.}

¶ \emph{A Priest shall present unto the Bishop, sitting in his chair
near to the Holy Table, all those who are to receive the Order of Priesthood
that day, each of them being decently habited, and shall say,}

REVEREND
Father in God, I present unto you these persons present, to be admitted
to the Order of Priesthood.

¶
\emph{The Bishop.}

TAKE
heed that the persons, whom ye present unto us, be apt and meet, for
their learning and godly conversation, to exercise their Ministry duly,
to the honour of God, and the edifying of his Church.

¶
\emph{The Priest shall answer,}

I
HAVE inquired concerning them, and also examined them, and think them
so to be.

¶
\emph{Then the Bishop shall say unto the People,}

GOOD
People, these are they whom we purpose, God willing, to receive this
day unto the holy Office of Priesthood; for, after due examination,
we find not to the contrary, but that they are lawfully called to their
Function and Ministry, and that they are persons meet for the same.
But yet, if there be any of you who knoweth any Impedi-ment, or notable
Crime, in any of them, for the which he ought not to be received into
this holy Ministry, let him come forth in the Name of God, and show
what the Crime or Impediment is.

¶
\emph{And if any great Crime or Impediment be objected, the Bishop shall
cease from Ordering that person, until such time as the party accused
shall be found clear of that Crime.}

¶ \emph{Then the Bishop (commending such as shall be found meet to
be Ordered, to the Prayers of the congregation) shall, with the Clergy
and People present, say the Litany.}

¶\emph{And}NOTE,\emph{That after the
Suffrage,}That it may please thee to illuminate all Bishops, \emph{etc.,
shall be said the following Suffrage:}

THAT
it may please thee to bless these thy servants, now to be admitted to
the Order of Priests, and to pour thy grace upon them; that they may
duly exercise their Office, to the edifying of thy Church, and the glory
of thy holy Name;

¶
\emph{And}NOTE\emph{further, That in the discretion
of the Bishop, instead of the Litany appointed, may be said the Litany
for Ordinations.}

¶ \emph{Then shall be said the Service for the Communion, with the
Collect, Epistle, and Gospel, as followeth.}

\emph{The
Collect.}

ALMIGHTY
God, giver of all good things, who by thy Holy Spirit hast appointed divers
Orders of Ministers in thy Church; Mercifully behold these thy servants
now called to the Office of Priesthood; and so replenish them with the
truth of thy Doctrine, and adorn them with inno-cency of life, that, both
by word and good example, they may faithfully serve thee in this Office,
to the glory of thy Name, and the edification of thy Church; through the
merits of our Saviour Jesus Christ, who liveth and reigneth with thee
and the same Holy Spirit, world without end. \emph{Amen}. 

\emph{The
Epist}\emph{le}.
Ephesians iv. 7.

UNTO
every one of us is given grace according to the measure of the gift
of Christ. Wherefore he saith, When he ascended up on high, he led captivity
captive, and gave gifts unto men. (Now that he ascended, what is it
but that he also descended first into the lower parts of the earth?
He that descended is the same also that ascended up far above all heavens,
that he might fill all things.) And he gave some, apostles; and some,
prophets; and some, evangelists; and some, pastors and teachers; for
the perfecting of the saints, for the work of the ministry, for the
edifying of the body of Christ: till we all come in the unity of the
faith, and of the knowledge of the Son of God, unto a perfect man, unto
the measure of the stature of the fulness of Christ.

\emph{The Gospel}. St. Matthew ix. 36.

WHEN Jesus saw the multitudes,
he was moved with compassion on them, because they fainted, and were scattered
abroad as sheep having no shepherd. Then saith he unto his disciples,
The harvest truly is plenteous, but the labourers are few; pray ye therefore
the Lord of the harvest, that he will send forth labourers into his harvest.

¶ \emph{Or else.}

\emph{The Gospel}. St. John x. 1.

VERILY, verily, I say unto
you, He that entereth not by the door into the sheepfold, but climbeth
up some other way, the same is a thief and a robber. But he that entereth
in by the door is the shepherd of the sheep. To him the porter openeth;
and the sheep hear his voice: and he calleth his own sheep by name, and
leadeth them out. And when he putteth forth his own sheep, he goeth before
them, and the sheep follow him: for they know his voice. And a stranger
will they not follow, but will flee from him: for they know not the voice
of strangers. This parable spake Jesus unto them: but they understood
not what things they were which he spake unto them. Then said Jesus unto
them again, Verily, verily, I say unto you, I am the door of the sheep.
All that ever came before me are thieves and robbers: but the sheep did
not hear them. I am the door: by me if any man enter in, he shall be saved,
and shall go in and out, and find pasture. The thief cometh not, but for
to steal, and to kill, and to destroy: I am come that they might have
life, and that they might have it more abundantly. I am the good shepherd:
the good shepherd giveth his life for the sheep. But he that is an hireling,
and not the shepherd, whose own the sheep are not, seeth the wolf coming,
and leaveth the sheep, and fleeth: and the wolf catcheth them, and scattereth
the sheep. The hireling fleeth, because he is an hireling, and careth
not for the sheep. I am the good shepherd; and know my sheep, and am known
of mine, even as the Father knoweth me, and I know the Father; and I lay
down my life for the sheep. And other sheep I have, which are not of this
fold: them also I must bring, and they shall hear my voice; and there
shall be one flock, and one shepherd. 

¶
\emph{Then, the People being seated, the Bishop shall say unto those who
are to be ordained Priests as followeth.}

YE
have heard, Brethren, as well in your private examination, as in the
exhortation which was now made to you, and in the holy Lessons taken
out of the Gospel, and the writings of the Apostles, of what dignity,
and of how great importance this Office is, whereunto ye are called.
And now again we exhort you, in the Name of our Lord Jesus Christ, that
ye have in remembrance, into how high a Dignity, and to how weighty
an Office and Charge ye are called: that is to say, to be Messengers,
Watchmen, and Stewards of the Lord; to teach, and to premonish, to feed
and provide for the Lord’s family; to seek for Christ’s
sheep that are dispersed abroad, and for his children who are in the
midst of this naughty world, that they may be saved through Christ for
ever.

Have always therefore printed in your remembrance,
how great a treasure is committed to your charge. For they are the sheep
of Christ, which he bought with his death, and for whom he shed his
blood. The Church and Congregation whom you must serve, is his Spouse,
and his Body. And if it shall happen that the same Church, or any Member
thereof, do take any hurt or hindrance by reason of your negligence,
ye know the greatness of the fault, and also the horrible punishment
that will ensue. Wherefore consider with yourselves the end of the Ministry
towards the children of God, towards the Spouse and Body of Christ;
and see that ye never cease your labour, your care and diligence, until
ye have done all that lieth in you, according to your bounden duty,
to bring all such as are or shall be committed to your charge, unto
that agreement in the faith and knowledge of God, and to that ripeness
and perfectness of age in Christ, that there be no place left among
you, either for error in religion, or for viciousness in life.

Forasmuch then as your Office is both of so
great excellency, and of so great difficulty, ye see with how great
care and study ye ought to apply yourselves, as well to show yourselves
dutiful and thankful unto that Lord, who hath placed you in so high
a dignity; as also to beware that neither you yourselves offend, nor
be occasion that others offend. Howbeit, ye cannot have a mind and will
thereto of yourselves; for that will and ability is given of God alone:
therefore ye ought, and have need, to pray earnestly for his Holy Spirit.
And seeing that ye cannot by any other means compass the doing of so
weighty a work, pertaining to the salvation of man, but with doctrine
and exhortation taken out of the Holy Scriptures, and with a life agreeable
to the same; consider how studious ye ought to be in reading and learning
the Scriptures, and in framing the manners both of yourselves, and of
them that specially pertain unto you, according to the rule of the same
Scriptures; and for this self-same cause, how ye ought to forsake and
set aside, as much as ye may, all worldly cares and studies.

We have good hope that ye have well weighed these
things with yourselves, long before this time; and that ye have clearly
determined, by God's grace, to give yourselves wholly to this Office,
whereunto it hath pleased God to call you: so that, as much as lieth
in you, ye will apply yourselves wholly to this one thing, and draw
all your cares and studies this way; and that ye will continually pray
to God the Father, by the mediation of our only Saviour Jesus Christ,
for the heavenly assistance of the Holy Ghost; that, by daily reading
and weighing the Scriptures, ye may wax riper and stronger in your Ministry;
and that ye may so endeavour yourselves, from time to time, to sanctify
the lives of you and yours, and to fashion them after the Rule and Doctrine
of Christ, that ye may be wholesome and godly examples and patterns
for the people to follow.

And now, that this present Congregation of Christ
may also understand your minds and wills in these things, and that this
your promise may the more move you to do your duties; ye shall answer
plainly to these things, which we, in the Name of God, and of his Church,
shall demand of you touching the same.

DO
you think in your heart, that you are truly called, according to the will
of our Lord Jesus Christ, and according to the Canons of this Church,
to the Order and Ministry of Priesthood?

\emph{Answer}. I think it.

\emph{Bishop}. Are you persuaded that the Holy
Scriptures contain all Doctrine required as necessary for eternal salvation
through faith in Jesus Christ? And are you determined, out of the said
Scriptures to instruct the people committed to your charge; and to teach
nothing, as necessary to eternal salvation, but that which you shall be
persuaded may be concluded and proved by the Scripture?

\emph{Answer}. I am so persuaded, and have so
determined, by God’s grace.

\emph{Bishop}. Will you then give your faithful
diligence always so to minister the Doctrine and Sacraments, and the Discipline
of Christ, as the Lord hath commanded, and as this Church hath received
the same, according to the Commandments of God; so that you may teach
the people committed to your Cure and Charge with all diligence to keep
and observe the same?

\emph{Answer}. I will so do, by the help of the Lord.

\emph{Bishop}. Will you be ready, with all faithful
diligence, to banish and drive away from the Church all erroneous and
strange doctrines contrary to God's Word; and to use both public and private
monitions and exhortations, as well to the sick as to the whole, within
your Cures, as need shall require, and occasion shall be given?

\emph{Answer}. I will, the Lord being my helper.

\emph{Bishop}. Will you be diligent in Prayers, and
in reading the Holy Scriptures, and in such studies as help to the knowledge
of the same, laying aside the study of the world and the flesh?

\emph{Answer}. I will endeavour so to do, the
Lord being my helper.

\emph{Bishop}. Will you be diligent to frame
and fashion your own selves, and your families, according to the Doctrine
of Christ; and to make both yourselves and them, as much as in you lieth,
wholesome examples and patterns to the flock of Christ?

\emph{Answer}. I will apply myself thereto,
the Lord being my helper.

\emph{Bishop}. Will you maintain and set forwards,
as much as lieth in you, quietness, peace, and love, among all Christian
people, and especially among them that are or shall be committed to your
charge?

\emph{Answer}. I will so do, the Lord being my helper.

\emph{Bishop.} Will you reverently obey your
Bishop, and other chief Ministers, who, according to the Canons of the
Church, may have the charge and government over you; following with a
glad mind and will their godly admonitions, and submitting yourselves
to their godly judgments?

\emph{Answer}. I will so do, the Lord being my helper.

¶ \emph{Then, all standing, shall the
Bishop say,}

ALMIGHTY God, who hath
given you this will to do all these things; Grant also unto you strength
and power to perform the same, that he may accomplish his work which he
hath begun in you; through Jesus Christ our Lord. \emph{Amen}.

¶ \emph{After this, the
Congregation shall be desired, secretly in their Prayers, to make their
humble supplications to God for all these things; for the which Prayers
there shall be silence kept for a space.}

¶ \emph{After which,
the Persons to be ordained Priests kneeling, and others standing, the
Bishop shall sing or say the Veni, Creator Spiritus; the Bishop beginning,
and the Priests, and others that are present, answering by verses, as
followeth.}

\emph{Veni,
Creator Spiritus.}

COME,
Holy Ghost, our souls inspire, 

\emph{And lighten with celestial fire.} 

Thou the anointing Spirit art, 

\emph{Who dost thy sevenfold gifts impart.}

Thy
blessed unction from above, 

\emph{Is comfort, life, and fire of love.}

Enable with perpetual light 

\emph{The dulness of our blinded sight.}

Anoint
and cheer our soiled face 

\emph{With the abundance of thy grace}. 

Keep far our foes, give peace at home; 

\emph{Where thou art guide, no ill can come.}

Teach us to know the Father, Son, 

\emph{And thee, of both, to be but One;}

That, through the ages all along, 

\emph{This may be our endless song:}

Praise to thy eternal merit, 

\emph{Father, Son, and Holy Spirit.}

¶ \emph{Or this.}

COME, Holy Ghost, Creator
blest, 

Vouchsafe within our souls to rest; 

\emph{Come with thy grace and heavenly aid, And fill the hearts which thou hast made.}

To thee, the Comforter, we cry;

To thee, the Gift of God most high; 

\emph{The Fount of life, the Fire of love, The soul's Anointing from above.}

The sevenfold gifts of grace are thine, 

O Finger of the Hand Divine; 

\emph{True Promise of the Father thou, Who dost the tongue with speech endow.}

Thy light to every sense impart, 

And shed thy love in every heart; 

\emph{Thine own unfailing might supply To strengthen our infirmity.}

Drive far away our ghostly foe, 

And thine abiding peace bestow; 

\emph{If thou be our preventing Guide, No evil can our steps betide.}

¶ \emph{That done, the Bishop shall
pray in this wise, and say,}

Let us pray. 

ALMIGHTY
God, and heavenly Father, who, of thine infinite love and goodness towards
us, hast given to us thy only and most dearly beloved Son Jesus Christ,
to be our Redeemer, and the Author of everlasting life; who, after he
had made perfect our redemption by his death, and was ascended into heaven,
sent abroad into the world his Apostles, Prophets, Evangelists, Doctors,
and Pastors; by whose labour and ministry he gathered together a great
flock in all the parts of the world, to set forth the eternal praise of
thy holy Name: For these so great benefits of thy eternal goodness, and
for that thou hast vouchsafed to call these thy servants here present
to the same Office and Ministry, appointed for the salvation of mankind,
we render unto thee most hearty thanks, we praise and worship thee; and
we humbly beseech thee, by the same thy blessed Son, to grant unto all,
which either here or elsewhere call upon thy holy Name, that we may continue
to show ourselves thankful unto thee for these and all thy other benefits;
and that we may daily increase and go forwards in the know-ledge and faith
of thee and thy Son, by the Holy Spirit. So that as well by these thy
Ministers, as by them over whom they shall be appointed thy Ministers,
thy holy Name may be for ever glorified, and thy blessed kingdom enlarged;
through the same thy Son Jesus Christ our Lord, who liveth and reigneth
with thee in the unity of the same Holy Spirit, world without end. \emph{Amen}.

¶ \emph{When this Prayer
is done, the Bishop with the Priests present, shall lay their Hands severally
upon the Head of every one that receiveth the Order of Priesthood; the
Receivers humbly kneeling, and the Bishop saying,}

RECEIVE the Holy Ghost
for the Office and Work of a Priest in the Church of God, now committed
unto thee by the Imposition of our hands. Whose sins thou dost forgive,
they are forgiven; and whose sins thou dost retain, they are retained.
And be thou a faithful Dispenser of the Word of God, and of his holy Sacraments;
In the Name of the Father, and of the Son, and of the Holy Ghost.\emph{Amen}.

¶ \emph{Or this.}

TAKE thou Authority to
execute the Office of a Priest in the Church of God, now committed to
thee by the Imposition of our hands. And be thou a faithful Dispenser
of the Word of God, and of his holy Sacraments; In the Name of the Father,
and of the Son, and of the Holy Ghost. \emph{Amen}.

¶ \emph{Then the Bishop
shall deliver to every one of them kneeling, the Bible into his hand,
saying,}

TAKE thou Authority to
preach the Word of God, and to minister the holy Sacraments in the Congregation,
where thou shalt be lawfully appointed thereunto.

¶ \emph{When this is
done, the Nicene Creed shall be said, and the Bishop shall go on in the
Service of the Communion, which all they who receive Orders shall take
together, and remain in the same place where Hands were laid upon them,
until such time as they have received the Communion.}

¶ \emph{The Communion
being done, after the last Collect, and immediately before the Benediction,
shall be said this Collect.}

MOST merciful Father,
we beseech thee to send upon these thy servants thy heavenly blessing;
that they may be clothed with righteousness, and that thy Word spoken
by their mouths may have such success, that it may never be spoken in
vain. Grant also, that we may have grace to hear and receive what they
shall deliver out of thy most holy Word, or agreeable to the same, as
the means of our salvation; that in all our words and deeds we may seek
thy glory, and the increase of thy kingdom; through Jesus Christ our Lord.
\emph{Amen}.

THE Peace
of God, which passeth all understanding, keep your hearts and minds in
the knowledge and love of God, and of his Son Jesus Christ our Lord: And
the Blessing of God Almighty, the Father, the Son, and the Holy Ghost,
be amongst you, and remain with you always. Amen.

¶ \emph{And if, on the same day, the Order of Deacons
be given to some, and the Order of Priesthood to others; the Deacons shall
be first presented, and then the Priests; and it shall suffice that the
Litany be once said for both. The Epistle shall be Ephesians iv. 7 to
13, as before in this Office. Immediately after which, they that are to
be made Deacons, shall be examined and Ordained, as is above prescribed.
Then one of them having read the Gospel, (which shall be either Saint
Matthew ix. 36 to 38, as before in this Office; or else Saint Luke xii.
35 to 38, as before in the Form for the Ordering of Deacons,) they that
are to be made Priests shall likewise be examined and Ordained, as is
in this Office before appointed. The Collect shall be as followeth.}

\emph{The Collect.}

ALMIGHTY God, giver of
all good things, who by thy Holy Spirit hast appointed divers Orders of
Ministers in thy Church; Mercifully behold these thy servants now called
to the Office of Deacon and these thy servants now called to the Office
of Priest; and so replenish them with the truth of thy Doctrine, and adorn
them with innocency of life, that, both by word and good example, they
may faithfully serve thee in their Ministry, to the glory of thy Name,
and the edification of thy Church; through the merits of our Saviour Jesus
Christ, who liveth and reigneth with thee and the same Holy Spirit, world
without end. \emph{Amen}.

The
Form of Ordaining or

Consecrating
a Bishop

¶
\emph{When all things are duly prepared in the Church, and set in order,
the Presiding Bishop, or some other Bishop appointed by the Bishops
present, shall begin the Communion Service, in which this shall be}

\emph{The
Collect.}

ALMIGHTY
God, who by thy Son Jesus Christ didst give to thy holy Apostles many
excellent gifts, and didst charge them to feed thy flock; Give grace,
we beseech thee, to all Bishops, the Pastors of thy Church, that they
may diligently preach thy Word, and duly administer the godly Discipline
thereof; and grant to the people, that they may obediently follow the
same; that all may receive the crown of everlasting glory; through the
same thy Son Jesus Christ our Lord. \emph{Amen}.

¶ \emph{And another Bishop shall read
the Epistle.}

\emph{The Epistle}. 1 Timothy iii. 1.

THIS is a true saying,
If a man desire the office of a bishop, he desireth a good work. A bishop
then must be blameless, the husband of one wife, vigilant, sober, of good
behaviour, given to hospitality, apt to teach; not given to wine, no striker,
not greedy of filthy lucre; but patient, not a brawler, not covetous;
one that ruleth well his own house, having his children in subjection
with all gravity; (for if a man know not how to rule his own house, how
shall he take care of the church of God?) not a novice, lest being lifted
up with pride he fall into the condemnation of the devil. Moreover he
must have a good report of them which are without; lest he fall into reproach
and the snare of the devil.

¶ \emph{Or this.}

\emph{For the Epistle}. Acts xx. 17.

FROM Miletus Paul sent
to Ephesus, and called the elders of the church. And when they were come
to him, he said unto them, Ye know, from the first day that I came into
Asia, after what manner I have been with you at all seasons, serving the
Lord with all humility of mind, and with many tears, and temptations,
which befell me by the lying in wait of the Jews: and how I kept back
nothing that was profitable unto you, but have shewed you, and have taught
you publickly, and from house to house, testifying both to the Jews, and
also to the Greeks, repentance toward God, and faith toward our Lord Jesus
Christ. And now, behold, I go bound in the spirit unto Jerusalem, not
knowing the things that shall befall me there: save that the Holy Ghost
witnesseth in every city, saying that bonds and afflictions abide me.
But none of these things move me, neither count I my life dear unto myself,
so that I might finish my course with joy, and the ministry, which I have
received of the Lord Jesus, to testify the gospel of the grace of God.
And now, behold, I know that ye all, among whom I have gone preaching
the kingdom of God, shall see my face no more. Wherefore I take you to
record this day, that I am pure from the blood of all men. For I have
not shunned to declare unto you all the counsel of God. Take heed therefore
unto yourselves, and to all the flock, over the which the Holy Ghost hath
made you overseers, to feed the church of God, which he hath purchased
with his own blood. For I know this, that after my departing shall grievous
wolves enter in among you, not sparing the flock. Also of your own selves
shall men arise, speaking perverse things, to draw away disciples after
them. Therefore watch, and remember, that by the space of three years
I ceased not to warn every one night and day with tears. And now, brethren,
I commend you to God, and to the word of his grace, which is able to build
you up, and to give you an inheritance among all them which are sanctified.
I have coveted no man's silver, or gold, or apparel. Yea, ye yourselves
know, that these hands have ministered unto my necessities, and to them
that were with me. I have shewed you all things, how that so labouring
ye ought to support the weak, and to remember the words of the Lord Jesus,
how he said, It is more blessed to give than to receive. 

¶
\emph{Then another Bishop shall read the Gospel.}

\emph{The
Gospel}. St. John xxi. 15.

JESUS
saith to Simon Peter, Simon, son of Jonas, lovest thou me more than these?
He saith unto him, Yea, Lord; thou knowest that I love thee. He saith
unto him, Feed my lambs. He saith to him again the second time, Simon,
son of Jonas, lovest thou me? He saith unto him, Yea, Lord; thou knowest
that I love thee. He saith unto him, Feed my sheep. He saith unto him
the third time, Simon, son of Jonas, lovest thou me? Peter was grieved
because he said unto him the third time, Lovest thou me? And he said unto
him, Lord, thou knowest all things; thou knowest that I love thee. Jesus
saith unto him, Feed my sheep.

¶ \emph{Or this.}

\emph{The Gospel.} St. John xx. 19.

THE same day at evening,
being the first day of the week, when the doors were shut where the disciples
were assembled for fear of the Jews, came Jesus, and stood in the midst,
and saith unto them, Peace be unto you. And when he had so said, he shewed
unto them his hands and his side. Then were the disciples glad, when they
saw the Lord. Then saith Jesus to them again, Peace be unto you: as my
Father hath sent me, even so send I you. And when he had said this, he
breathed on them, and saith unto them, Receive ye the Holy Ghost: whose
soever sins ye remit, they are remitted unto them; and whose soever sins
ye retain, they are retained.

¶ \emph{Or this.}

\emph{The Gospel}. St. Matthew xxviii. 18.

JESUS came and spake unto
them, saying, All power is given unto me in heaven and in earth. Go ye
therefore, and teach all nations, baptizing them in the name of the Father,
and of the Son, and of the Holy Ghost: teaching them to observe all things
whatsoever I have commanded you: and, lo, I am with you alway, even unto
the end of the world.

¶ \emph{Then shall follow
the Nicene Creed, and after that the Sermon; which being ended, the Elected
Bishop, vested with his Rochet, shall be presented by two Bishops of this
Church unto the Presiding Bishop, or to the Bishop appointed, sitting
in his chair, near the Holy Table; the Bishops who present him saying,}

REVEREND Father in God,
we present unto you this godly and well-learned man, to be Ordained and
Consecrated Bishop.

¶ \emph{Then shall the
Presiding Bishop demand Testimonials of the person presented for Consecration,
and shall cause them to be read.}

¶ \emph{He shall then
require of him the following Promise of Conformity to the Doctrine, Discipline,
and Worship of the Protestant Episcopal Church.}

IN the Name of God, Amen.
I, \emph{N}., chosen Bishop of the Protestant Episcopal Church in \emph{N.},
do promise conform-ity and obedience to the Doctrine, Discipline, and
Worship of the Protestant Episcopal Church in the United States of America.
So help me God, through Jesus Christ.

¶ \emph{Then the Presiding
Bishop shall move the Congregation present to pray, saying thus to them:}

BRETHREN, it is written
in the Gospel of Saint Luke, that our Saviour Christ continued the whole
night in prayer, before he chose and sent forth his twelve Apostles. It
is written also, that the holy Apostles prayed before they ordained Matthias
to be of the number of the Twelve. Let us, therefore, following the example
of our Saviour Christ, and his Apostles, offer up our prayers to Almighty
God, before we admit and send forth this person presented unto us, to
the work whereunto we trust the Holy Ghost hath called him.

¶
\emph{And then shall be said the Litany; save only, that after this place,}That it may please thee to illuminate all Bishops,\emph{etc., the
proper Suffrage shall be,}

THAT
it may please thee to bless this our Brother elected, and to send thy
grace upon him, that he may duly execute the Office whereunto he is
called, to the edifying of thy Church, and to the honour, praise, and
glory of thy Name;

\emph{Answer}. We beseech thee to hear us, good
Lord.

¶
\emph{And Note, That in the discretion of the Presiding Bishop, instead
of the Litany, may be said the Litany for Ordinations.}

¶ \emph{Then shall be said the Prayer following.}

ALMIGHTY
God, giver of all good things, who by thy Holy Spirit hast appointed divers
Orders of Ministers in thy Church; Mercifully behold this thy servant,
now called to the Work and Ministry of a Bishop; and so replenish him
with the truth of thy Doctrine, and adorn him with innocency of life,
that, both by word and deed, he may faithfully serve thee in this Office,
to the glory of thy Name, and the edifying and well-governing of thy Church;
through the merits of our Saviour Jesus Christ, who liveth and reigneth
with thee and the same Holy Spirit, world without end. \emph{Amen}.

¶
\emph{Then, the People being seated, the Presiding Bishop, sitting in his
chair shall say to him that is to be Consecrated,}

BROTHER,
forasmuch as the Holy Scripture and the ancient Canons command, that we
should not be hasty in laying on hands, and admitting any person to Government
in the Church of Christ, which he hath purchased with no less price than
the effusion of his own blood; before we admit you to this Administration,
we will examine you in certain Articles, to the end that the Congregation
present may have a trial, and bear witness, how you are minded to behave
yourself in the Church of God.

ARE
you persuaded that you are truly called to this Ministration, according
to the will of our Lord Jesus Christ, and the order of this Church?

\emph{Answer}. I am so persuaded.

\emph{Bishop}. Are you persuaded that the Holy
Scriptures contain all Doctrine required as necessary for eternal sal-vation
through faith in Jesus Christ? And are you deter-mined out of the same
Holy Scriptures to instruct the people committed to your charge; and to
teach or maintain nothing, as necessary to eternal salvation, but that
which you shall be persuaded may be concluded and proved by the same?

\emph{Answer}. I am so persuaded, and determined,
by God's grace.

\emph{Bishop}. Will you then faithfully exercise
yourself in the Holy Scriptures, and call upon God by prayer for the true
understanding of the same; so that you may be able by them to teach and
exhort with wholesome Doctrine, and to withstand and convince the gainsayers?

\emph{Answer}. I will so do, by the help of
God.

\emph{Bishop}. Are you ready, with all faithful diligence,
to banish and drive away from the Church all erroneous and strange doctrine
contrary to God's Word; and both privately and openly to call upon and
encourage others to the same?

\emph{Answer}. I am ready, the Lord being my helper.

\emph{Bishop}. Will you deny all ungodliness
and worldly lusts, and live soberly, righteously, and godly in this present
world; that you may show yourself in all things an example of good works
unto others, that the adversary may be ashamed, having nothing to say
against you?

\emph{Answer}. I will so do, the Lord being my helper.

\emph{Bishop}. Will you maintain and set forward,
as much as shall lie in you, quietness, love, and peace among all men;
and diligently exercise such discipline as by the authority of God's Word,
and by the order of this Church, is com-mitted to you?

\emph{Answer}. I will so do, by the help of God.

\emph{Bishop}. Will you be faithful in Ordaining,
sending, or laying hands upon others?

\emph{Answer}. I will so be, by the help of God.

\emph{Bishop}. Will you show yourself gentle,
and be merciful for Christ’s sake to poor and needy people, and
to all strangers destitute of help?

\emph{Answer}. I will so show myself, by God’s
help.

¶
\emph{Then, all standing, the Presiding Bishop shall say,}

ALMIGHTY
God, our heavenly Father, who hath given you a good will to do all these
things; Grant also unto you strength and power to perform the same; that,
he accomplishing in you the good work which he hath begun, you may be
found perfect and irreprehensible at the latter day; through Jesus Christ
our Lord. \emph{Amen}.

¶
\emph{Then shall the Bishop elect put on the rest of the Episcopal habit,
and shall kneel down; and the Veni, Creator Spiritus shall be sung or
said over him; the Presiding Bishop shall begin, and the Bishops, and
the others that are present, standing, shall answer by verses, as followeth.}

COME,
Holy Ghost, our souls inspire, 

\emph{And lighten with celestial fire.} 

Thou the anointing Spirit art, 

\emph{Who dost thy sevenfold gifts impart.}

Thy
blessed unction from above, 

\emph{Is comfort, life, and fire of love.}

Enable with perpetual light 

\emph{The dulness of our blinded sight.}

Anoint
and cheer our soiled face 

\emph{With the abundance of thy grace}. 

Keep far our foes, give peace at home; 

\emph{Where thou art guide, no ill can come.}

Teach
us to know the Father, Son, 

\emph{And thee, of both, to be but One;}

That, through the ages all along, 

\emph{This may be our endless song:}

Praise to thy eternal merit, 

\emph{Father, Son, and Holy Spirit.}

¶
\emph{Or this.}

COME,
Holy Ghost, Creator blest, 

Vouchsafe within our souls to rest; 

\emph{Come with thy grace and heavenly aid, And fill the hearts which thou hast made.}

To thee,
the Comforter, we cry;

To thee, the Gift of God most high; 

\emph{The Fount of life, the Fire of love, The soul's Anointing from above.}

The
sevenfold gifts of grace are thine, 

O Finger of the Hand Divine; 

\emph{True Promise of the Father thou, Who dost the tongue with speech endow.}

Thy
light to every sense impart, 

And shed thy love in every heart; 

\emph{Thine own unfailing might supply To strengthen our infirmity.}

Drive
far away our ghostly foe, 

And thine abiding peace bestow; 

\emph{If thou be our preventing Guide, No evil can our steps betide.}

¶
\emph{That ended, the Presiding Bishop shall say,}

\emph{Answer.} Lord,
hear our prayer.

And let our cry come unto thee.

Let us
pray.

ALMIGHTY
God, and most merciful Father, who, of thine infinite goodness, hast given
thy only and dearly beloved Son Jesus Christ, to be our Redeemer, and
the Author of everlasting life; who, after that he had made per-fect our
redemption by his death, and was ascended into heaven, poured down his
gifts abundantly upon men, making some Apostles, some Prophets, some Evangelists,
some Pastors and Doctors, to the edifying and making perfect his Church;
Grant, we beseech thee, to this thy servant, such grace, that he may evermore
be ready to spread abroad thy Gospel, the glad tidings of reconciliation
withthee; and use the authority given him, not to destruction, but to
salvation; not to hurt, but to help: so that, as a wise and faithful servant,
giving to thy family their portion in due season, he may at last be received
into everlasting joy; through the same Jesus Christ our Lord, who, with
thee and the Holy Ghost, liveth and reigneth, one God, world without end.
\emph{Amen}.

¶
\emph{Then the Presiding Bishop and Bishops present shall lay their hands
upon the head of the Elected Bishop, kneeling before them, the Presiding
Bishop saying,}

RECEIVE
the Holy Ghost for the Office and Work of a Bishop in the Church of God,
now committed unto thee by the Imposition of our hands; In the Name of
the Father, and of the Son, and of the Holy Ghost. Amen. And remem-ber
that thou stir up the grace of God, which is given thee by this Imposition
of our hands; for God hath not given us the spirit of fear, but of power,
and love, and soberness.

¶ \emph{Then the Presiding Bishop shall
deliver him the Bible, saying,}

GIVE heed unto reading,
exhortation, and doctrine. Think upon the things contained in this Book.
Be diligent in them, that the increase coming thereby may be manifest
unto all men; for by so doing thou shalt both save thyself and them that
hear thee. Be to the flock of Christ a shepherd, not a wolf; feed them,
devour them not. Hold up the weak, heal the sick, bind up the broken,
bring again the outcasts, seek the lost. Be so merciful, that you be not
too remiss; so minister discipline, that you forget not mercy; that when
the Chief Shepherd shall appear, you may receive the never-fading crown
of glory; through Jesus Christ our Lord. Amen.

¶ \emph{Then the Presiding
Bishop shall proceed in the Communion Service; with whom the newly consecrated
Bishop, with others, shall also communicate.}

¶ \emph{And immediately
before the Benediction, shall be said this Prayer.}

MOST merciful Father, send
down, we beseech thee, upon this thy servant thy heavenly blessing; and
so endue him with thy Holy Spirit, that he, preaching thy Word, may not
only be earnest to reprove, beseech, and rebuke, with all patience and
doctrine; but also may be, to such as believe, a wholesome example in
word, in conver-sation, in love, in faith, in chastity, and in purity;
that, faithfully fulfilling his course, at the latter day he may receive
the crown of righteousness, laid up by the Lord Jesus, the righteous Judge,
who liveth and reigneth with thee and the same Holy Spirit, one God, world
without end. \emph{Amen}.

THE Peace of God, which
passeth all understanding, keep your hearts and minds in the knowledge
and love of God, and of his Son Jesus Christ our Lord: And the Blessing
of God Almighty, the Father, the Son, and the Holy Ghost, be amongst you,
and remain with you always. \emph{Amen}.

The
Litany and Suffrages

for
Ordinations

O
GOD the Father, 

\emph{Have mercy upon us.}

O God the Son, 

\emph{Have mercy upon us.}

O God the Holy Ghost, 

\emph{Have mercy upon us.}

O holy Trinity, one God, 

\emph{Have mercy upon us.}

WE
beseech thee to hear us, good Lord; and that it may please thee to grant
peace to the whole world, and to thy Church;

\emph{We beseech thee to hear us, good Lord.}

That it may please thee to sanctify and bless thy
holy Church throughout the world;

\emph{We beseech thee to hear us, good Lord.}

That it may please thee to inspire all Bishops, Priests,
and Deacons, with love of thee and of thy truth;

\emph{We beseech thee to hear us, good Lord.}

That it may please thee to endue all Ministers of thy Church with devotion
to thy glory and to the salvation of souls;

\emph{We beseech thee to hear us, good Lord.}

¶ \emph{Here, at the Ordination of
Deacons or of Priests shall be said,}

That it may please thee to bless
these thy servants, now to be admitted to the Order of Deacons (or Priests),
and to pour thy grace upon them; that they may duly execute their Office
to the edifying of thy Church, and to the glory of thy holy Name;

\emph{We beseech thee to hear us, good Lord.}

\emph{¶ Here, at the Consecration
of a Bishop shall be said,}

That it may please thee to bless
this our Brother elected, and to send thy grace upon him, that he may
duly execute the Office whereunto he is called, to the edifying of thy
Church, and to the honour, praise, and glory of thy Name;

We beseech thee to hear us, good Lord.

That it may please thee to guide
by thy indwelling Spirit those whom thou dost call to the Ministry of
thy Church; that they may go forward with courage, and persevere unto
the end;

\emph{We beseech thee to hear us, good Lord.}

That it may please thee to increase the number
of the Ministers of thy Church, that the Gospel may be preached to all
people;

\emph{We beseech thee to hear us, good Lord.}

That it may please thee to hasten the fulfilment
of thy purpose, that thy Church may be one;

\emph{We beseech thee to hear us, good Lord.}

That it may please thee to grant that we, with
all thy saints, may be partakers of thy everlasting kingdom;

\emph{We beseech thee to hear us, good Lord.}

Lord, have mercy upon us.

\emph{Christ, have mercy upon us.}

Lord, have mercy upon us.

OUR Father, who art in
heaven, Hallowed be thy Name. Thy kingdom come. Thy will be done, On earth
as it is in heaven. Give us this day our daily bread. And forgive us our
trespasses, As we forgive those who trespass against us. And lead us not
into temptation, But deliver us from evil. Amen.

\emph{Minister}. Hearken unto
our voice, O Lord, when we cry unto thee; 

\emph{Answer}. Have mercy upon us and hear us.

\emph{Minister.} O Lord, arise, help us;

\emph{Answer}. And deliver us for thy Name's sake.

\emph{Minister}. Let thy priests be clothed
with righteousness;

\emph{Answer}. And let thy saints sing with joyfulness.

\emph{Minister}. Lord, hear our prayer;

\emph{Answer.} And let our cry come unto thee.

Let us pray.

GOD,
who dost ever hallow and protect thy Church; Raise up therein, through
thy Spirit, good and faithful stewards of the mysteries of Christ, that
by their ministry and example thy people may abide in thy favour and be
guided in the way of truth; through Jesus Christ our Lord, who liveth
and reigneth with thee in the unity of the same Spirit ever, one God,
world without end. \emph{Amen}. 

\begin{quote}

\emph{Return to the U.
S. Book of Common Prayer (1928)}

\end{quote}

\xchapter{The Form of Consecration of a Church or Chapel}

\textbf{The Book of Common Prayer (1928)}

The
Form of Consecration of a Church or Chapel

¶
\emph{The following Office may be used with the Order for the Holy Communion,
or at Morning Prayer or Evening Prayer, or separately.}

¶\emph{The Bishop is to be received at the entrance of the Church,
or Chapel, by the Church-wardens and Vestrymen, or some other persons
appointed for that purpose. The Bishop and the Clergy who are present
shall go up the aisle of the Church, or Chapel, to the Holy Table, repeating
the following Psalm alternately, the Bishop one verse and the Clergy another.}

\emph{Domini
est terra}. Psalm xxiv.

THE
earth is the Lord’s, and all that therein is; * the compass of the
world, and they that dwell therein.

2 For he hath founded it upon the seas, * and stablished
it upon the floods.

3 Who shall ascend into the hill of the LORD?
* or who shall rise up in his holy place?

4 Even he that hath clean hands, and a pure heart;
* and that hath not lift up his mind unto vanity, nor sworn to deceive
his neighbour.

5 He shall receive the blessing from the LORD,
* and righteousness from the God of his salvation.

6 This is the generation of them that seek him; * even
of them that seek thy face, O Jacob.

7 Lift up your heads, O ye gates; and be ye lift up,
ye everlasting doors; * and the King of glory shall come in.

8 Who is this King of glory? * It is the LORD
strong and mighty, even the LORD mighty in battle.

9 Lift up your heads, O ye gates; and be ye lift up,
ye everlasting doors; * and the King of glory shall come in.

10 Who is this King of glory? * Even the LORD
of hosts, he is the King of glory. 

Glory be to the Father, and to the Son, * and to the
Holy Ghost;

As it was in the beginning, is now, and ever shall
be, * world without end. Amen.

¶
\emph{The Bishop shall go within the rails, with such of the Clergy as can
be there accommodated. The Bishop, sitting in his chair, shall have the
instruments of Donation and Endowment, if there be any, presented to him;
and then standing up, and turning to the Congregation, he shall say,}

DEARLY
beloved in the Lord; forasmuch as devout and holy men, as well under the
Law as under the Gospel, moved either by the express command of God, or
by the secret inspiration of the blessed Spirit, and acting agreeably
to their own reason and sense of the natural decency of things, have erected
houses for the public worship of God, and separated them from all unhallowed,
worldly, and common uses, in order to fill men's minds with greater reverence
for his glorious Majesty, and affect their hearts with more devotion and
humility in his service; which pious works have been approved of and graciously
accepted by our heavenly Father; Let us not doubt but that he will also
favourably approve our godly purpose of setting apart this place in solemn
manner, for the several Offices of religious worship, and let us faithfully
and devoutly beg his blessing on this our undertaking. 

¶
\emph{Then the Bishop, kneeling, shall say the following Prayer.}

O
ETERNAL God, mighty in power, and of majesty incomprehensible, whom
the heaven of heavens cannot contain, much less the walls of temples
made with hands; and who yet hast been graciously pleased to promise
thy especial presence, wherever two or three of thy faithful servants
shall assemble in thy Name, to offer up their praises and supplications
unto thee; Vouchsafe, O Lord, to be present with us, who are here gathered
together with all humility and readiness of heart, to consecrate this
place to the honour of thy great Name; separating it henceforth from
all unhallowed, ordinary, and common uses; and dedi-cating it to thy
service, for reading thy holy Word, for cele-brating thy holy Sacraments,
for offering to thy glorious Majesty the sacrifices of prayer and thanksgiving,
for blessing thy people in thy Name, and for all other holy offices:
accept, O Lord, this service at our hands, and bless it with such success
as may tend most to thy glory, and the furtherance of our happiness
both temporal and spiritual; through Jesus Christ our blessed Lord and
Saviour. \emph{Amen}.

¶
\emph{After this the Bishop shall stand up, and turning his face towards
the Congregation, shall say,}

REGARD,
O Lord, the supplications of thy servants, and grant that whosoever
in this house shall be received by Baptism into the congregation of
Christ's flock, may be sanctified by the Holy Ghost, and may continue
Christ's faithful soldier and servant unto his life's end. \emph{Amen}.

GRANT,
O Lord, that they who at this place shall in their own persons renew
the promises and vows of their Baptism, and be Confirmed by the Bishop,
may receive such a measure of thy Holy Spirit, that they may grow in
grace unto their life's end. \emph{Amen}.

GRANT,
O Lord, that whosoever shall receive in this place the blessed Sacrament
of the Body and Blood of Christ, may come to that holy ordinance with
faith, charity, and true repentance; and being filled with thy grace
and heavenly benediction, may, to their great and endless comfort, obtain
remission of their sins, and all other benefits of his passion. \emph{Amen}.

GRANT,
O Lord, that by thy holy Word which shall be read and preached in this
place, and by thy Holy Spirit grafting it inwardly in the heart, the
hearers thereof may both perceive and know what things they ought to
do, and may have power and strength to fulfil the same. \emph{Amen}.

GRANT, O Lord, that whosoever shall be
joined together in this place in the holy estate of Matrimony, may faithfully
perform and keep the vow and covenant betwixt them made, and may remain
in perfect love to-gether unto their life's end. \emph{Amen}.

GRANT, we beseech thee, blessed Lord,
that whosoever shall draw near to thee in this place, to give thee thanks
for the benefits which they have received at thy hands, to set forth
thy most worthy praise, to confess their sins unto thee, and to ask
such things as are requisite and necessary, as well for the body as
for the soul, may do it with such steadiness of faith, and with such
seriousness, affection, and devotion of mind, that thou mayest accept
their bounden duty and service, and vouchsafe to give whatever in thy
infinite wisdom thou shalt see to be most expedient for them. All which
we beg for Jesus Christ's sake, our most blessed Lord and Saviour. \emph{Amen}.

¶
\emph{Then, the Bishop sitting in his chair, the Sentence of Consecration
is to be read by some person appointed by him, and then laid by him
upon the Communion Table; after which, the Bishop shall say,}

BLESSED
be thy Name, O Lord, that it hath pleased thee to put it into the hearts
of thy servants to appro-priate and devote this house to thy honour
and worship; and grant that all who shall enjoy the benefit of this
pious work, may show forth their thankfulness, by making a right use
of it, to the glory of thy blessed Name; through Jesus Christ our Lord.
Amen.

\emph{¶
When there is a Communion, the following shall be the Collect, Epistle,
and Gospel.}

\emph{The
Collect.}

O
MOST glorious God, whom the heaven of heavens cannot contain; Graciously
accept the Dedication of this place to thy service; and grant that all
who shall call upon thee here may worship thee in spirit and in truth,
and may in their lives show forth thy praise; through Jesus Christ our
Lord. \emph{Amen}.

\emph{For the Epistle.} Revelation xxi. 2.

AND I John saw the holy city, new Jerusalem,
coming down from God out of heaven, prepared as a bride adorned for
her husband. And I heard a great voice out of heaven saying, Behold,
the tabernacle of God is with men, and he will dwell with them, and
they shall be his people, and God himself shall be with them, and be
their God. And God shall wipe away all tears from their eyes; and there
shall be no more death, neither sorrow, nor crying, neither shall there
be any more pain: for the former things are passed away. And he that
sat upon the throne said, Behold, I make all things new. And he said
unto me, Write: for these words are true and faithful.

\emph{The Gospel}. St. John ii. 13.

AND the Jews’ passover was at hand,
and Jesus went up to Jerusalem, and found in the temple those that sold
oxen and sheep and doves, and the changers of money sitting: and when
he had made a scourge of small cords, he drove them all out of the temple,
and the sheep, and the oxen; and poured out the changers' money, and
overthrew the tables; and said unto them that sold doves, Take these
things hence; make not my Father's house an house of merchandise. And
his disciples remembered that it was written, The zeal of thine house
hath eaten me up.

¶ \emph{And immediately before the
final Blessing, the Bishop shall say this Prayer.}

BLESSED be thy Name, O Lord God, for
that it hath pleased thee to have thy habitation among the sons of men,
and to dwell in the midst of the assembly of the saints upon the earth;
Grant, we beseech thee, that in this place now set apart to thy service,
thy holy Name may be worshipped in truth and purity through all generations;
through Jesus Christ our Lord. \emph{Amen}.

THE Peace of God, which passeth all understanding,
keep your hearts and minds in the knowledge and love of God, and of
his Son Jesus Christ our Lord: And the Blessing of God Almighty, the
Father, the Son, and the Holy Ghost, be amongst you, and remain with
you always.\emph{Amen}.

\begin{quote}

\emph{Return to the U.
S. Book of Common Prayer (1928)}

\end{quote}

\xchapter{An Office of Institution of Ministers}

\textbf{The Book of Common Prayer (1928)}

An
Office of 

Institution of Ministers 

into Parishes or Churches

¶
\emph{The Bishop having received due Notice of the Election of a Minister
into a Parish or Church, as prescribed by Canon, and being satisfied
that the Person chosen is a qualified Minister of this Church, may proceed
to institute him into the Parish.}

¶ \emph{The following Office may be used with the Order for the Holy
Communion, or at Morning Prayer or Evening Prayer, or separately.}

¶ \emph{In any Diocese, the concluding Paragraph in the Letter of
Institution may be omitted, where it interferes with the Usages, Laws,
or Charters of the Church in the same.}

\emph{To
our well-beloved in Christ,}A. B.\emph{, Presbyter, Greeting.}

\emph{Sigillum. Signat.}

WE
do by these Presents give and grant unto you, in whose Learning Diligence,
sound Doctrine, and Prudence, we do fully con-fide, our Licence and Authority
to perform the Office of a Priest, in the Parish (\emph{or} Church)
of \emph{E}. And also hereby do institute you into said Parish, (\emph{or}
Church,) possessed of full power to perform every Act of sacerdotal Function
among the People of the same; you continuing in communion with us, and
complying with the rubrics and canons of the Church, and with such lawful
directions as you shall at any time receive from us.

And as a canonically instituted Priest into the Office
of Rector of —— Parish, (\emph{or} Church,) you are faithfully
to feed that portion of the flock of Christ which is now intrusted to
you; not as a man-pleaser, but as continually bearing in mind that you
are accountable to us here, and to the Chief Bishop and Sovereign Judge
of all, hereafter. 

And as the Lord hath ordained that they who serve at
the altar should live of the things belonging to the altar; so we authorize
you to claim and enjoy all the accustomed temporalities appertaining to
your cure, until some urgent reason or reasons occasion a wish in you,
or in the congregation committed to your charge, to bring about a separation,
and dissolution of all sacerdotal relation, between you and them; of all
which you will give us due notice; and in case of any difference between
you and your congregation, as to a separ-ation and dissolution of all
sacerdotal connection between you and them, we, your Bishop, with the
advice of our Presbyters, are to be the ultimate arbiter and judge.

In witness whereof, we have hereunto affixed our episcopal
seal and signature, at —— , this —— day of ——
, A. D. —— , and in the —— year of our consecration.

¶ \emph{At the time designated for the new Incumbent's Institution,
the Bishop, or the Institutor appointed by him, attended by the new
Incumbent, and by the other Clergy present, shall enter the Chancel.
Then all the Clergy present standing in the Chancel or Choir, except
the Bishop, or the Priest who acts as Institutor, who shall go within
the rails of the Altar; the Wardens (or, in case of their necessary
absence, two members of the Vestry) standing on the right and left of
the Altar, without the rails; the Senior Warden (or the member of the
Vestry supplying his place) holding the keys of the Church in his hand,
in open view, the Bishop, or the Priest who acts as the Institutor,
shall say,}

DEARLY
beloved in the Lord, we have assembled for the purpose of instituting
the Rev. \emph{A. B}. into this Parish, (\emph{or} Church,) as Priest
and Rector of the same; and we are possessed of your Vote that he has
been so elected; as also of the prescribed Letter of Institution. But
if any of you can show just cause why he may not be instituted, we proceed
no further, because we would not that an unworthy person should minister
among you.

¶ \emph{If any objection
be offered, the Bishop, or the Priest who acts as the Institutor, shall
judge whether it afford just cause to suspend the Service.}

¶ \emph{No objection
being offered, or the Institutor choosing to go on with the Service, then
shall be read the Letter of Institution.}

¶ \emph{And then shall
the Senior Warden (or the member of the Vestry supplying his place) present
the keys of the Church to the new Incumbent, saying,}

IN the name and behalf
of —— Parish (\emph{or} Church) I do receive and acknowledge
you, the Rev. \emph{A. B}., as Priest and Rector of the same; and in
token thereof, give into your hands the keys of this Church.

¶ \emph{Then the new Incumbent shall
say,}

I, \emph{A. B}., receive
these keys of the House of God at your hands, as the pledges of my Institution,
and of your parochial recognition, and promise to be a faithful shepherd
over you; In the Name of the Father, and of the Son, and of the Holy Ghost.

¶
\emph{Here the Institutor shall begin the Office.}

\emph{Minister}.
The Lord be with you.

\emph{Answer}. And with thy spirit.

Let us pray.

DIRECT us, O Lord, in all
our doings, with thy most gracious favour, and further us with thy continual
help; that in all our works begun, continued, and ended in thee, we may
glorify thy holy Name, and finally, by thy mercy, obtain everlasting life;
through Jesus Christ our Lord, who hath taught us to pray unto thee, O
Almighty Father, in his prevailing Name and words,

OUR Father, who art in
heaven, Hallowed be thy Name. Thy kingdom come. Thy will be done, On earth
as it is in heaven. Give us this day our daily bread. And forgive us our
trespasses, As we forgive those who trespass against us. And lead us not
into temptation, But deliver us from evil. For thine is the kingdom, and
the power, and the glory, for ever and ever. Amen.

¶ \emph{Then shall the
Institutor receive the Incumbent within the rails of the Altar, and present
him the Bible, Book of Common Prayer, and Books of Canons of the General
and Diocesan Convention, saying as follows.}

RECEIVE these Books; and
let them be the rule of thy conduct in dispensing the divine Word, in
leading the Devotions of the People, and in exercising the Discipline
of the Church; and be thou in all things a pattern to the flock committed
to thy care.

\emph{¶ Then shall be said or sung Exsurgat
Deus, Psalm lxviii., or Judica me, Domine, Psalm xxvi.}

\emph{Minister}. The Law was given
by Moses;

\emph{People}. But Grace and Truth came by Jesus
Christ:

\emph{Minister and People}. Who is God over all,
blessed for evermore. Amen.

Let us pray.

MOST gracious Father,
the giver of all good and perfect gifts, who of thy wise providence hast
appointed divers Orders in thy Church; Give thy grace, we beseech thee,
to thy servant, to whom the charge of this Congrega-tion is now committed;
and so replenish him with the truth of thy doctrine, and endue him with
innocency of life, that he may faithfully serve before thee, to the glory
of thy great Name, and the benefit of thy holy Church; through Jesus Christ,
our only Mediator and Redeemer. \emph{Amen}.

O HOLY Jesus, who hast
purchased to thyself an universal Church, and hast promised to be with
the Ministers of Apostolic Succession to the end of the world; Be graciously
pleased to bless the ministry and service of him who is now appointed
to offer the sacrifices of prayer and praise to thee in this house, which
is called by thy Name. May the words of his mouth, and the meditation
of his heart, be alway acceptable in thy sight, O Lord, our strength and
our Redeemer. \emph{Amen}.

O GOD, Holy Ghost, Sanctifier
of the faithful, visit, we pray thee, this Congregation with thy love
and favour; en-lighten their minds more and more with the light of the
everlasting Gospel; graft in their hearts a love of the truth; increase
in them true religion; nourish them with all good-ness; and of thy great
mercy keep them in the same, O blessed Spirit, whom, with the Father and
the Son together, we worship and glorify as one God, world without end.
\emph{Amen}.

\emph{Benediction.}

THE
God of peace, who brought again from the dead our Lord Jesus Christ, the
great Shepherd of the sheep, through the blood of the everlasting covenant;
Make you perfect in every good work to do his will, working in you that
which is well pleasing in his sight; through Jesus Christ, to whom be
glory for ever and ever. Amen.

¶ Then shall the Instituted
Minister kneel at the Altar, to present his supplication for himself,
in this form.

O LORD my God, I am not
worthy that thou shouldest come under my roof; yet thou hast honoured
thy servant with appointing him to stand in thy House, and to serve at
thy holy Altar. To thee and to thy service I devote myself, body, soul,
and spirit, with all their powers and faculties. Fill my memory with the
words of thy Law; enlighten my understanding with the illumination of
the Holy Ghost; and may all the wishes and desires of my will centre in
what thou hast commanded. And, to make me instrumental in promoting the
salvation of the people now committed to my charge, grant that I may faithfully
ad-minister thy holy Sacraments, and by my life and doctrine set forth
thy true and lively Word. Be ever with me in the performance of all the
duties of my ministry: in prayer, to quicken my devotion; in praises,
to heighten my love and gratitude; and in preaching, to give a readiness
of thought and expression suitable to the clearness and excellency of
thy holy Word. Grant this for the sake of Jesus Christ thy Son our Saviour.

¶ \emph{The Instituted Minister, standing
up, shall say,}

\emph{Answer.} The Lord be with
you.

And with thy spirit.

Let us pray.

O ALMIGHTY God, who hast
built thy Church upon the foundation of the Apostles and Prophets, Jesus
Christ himself being the chief corner-stone; Grant that, by the operation
of the Holy Ghost, all Christians may be so joined together in unity of
spirit, and in the bond of peace, that they may be an holy temple acceptable
unto thee. And especially to this Congregation present, give the abundance
of thy grace; that with one heart they may desire the prosperity of thy
holy Apostolic Church, and with one mouth may profess the faith once delivered
to the Saints. Defend them from the sins of heresy and schism; let not
the foot of pride come nigh to hurt them, nor the hand of the ungodly
to cast them down. And grant that the course of this world may be so peaceably
ordered by thy govern-ance, that thy Church may joyfully serve thee in
all godly quietness; that so they may walk in the ways of truth and peace,
and at last be numbered with thy Saints in glory everlasting; through
the merits of the same thy blessed Son Jesus Christ, the gracious Bishop
and Shepherd of our souls, who liveth and reigneth with thee and the same
Holy Ghost, one God, world without end. \emph{Amen.}

¶ \emph{Then shall follow
the Sermon. And after that, if there be a Communion, the Instituted Minister
shall proceed to that Service, and to administer the holy Eucharist to
his Congregation; and after the Benediction, (which he shall always pronounce,)
the Wardens, Vestry, and others, shall salute and welcome him, bidding
him God-speed.}

¶ \emph{When the Bishop
of the Diocese is present at the Institution of a Minister, he shall make
to him the address, as prescribed in the Office in the form of a letter.}

\begin{quote}

\emph{Return to the U.
S. Book of Common Prayer (1928)}

\end{quote}

\xchapter{Forms of Prayer to be Used in Families}

\textbf{The Book of Common Prayer (1928)}

Forms
of Prayer

to be used in Families

MORNING
PRAYER.

¶
\emph{The Master or Mistress having called together as many of the Family
as can conveniently be present, let one of them, or any other who may
be appointed, say as followeth, all kneeling, and repeating with him
the Lord’s Prayer.}

OUR
Father, who art in heaven, Hallowed be thy Name. Thy kingdom come. Thy
will be done, On earth as it is in heaven. Give us this day our daily
bread. And forgive us our trespasses, As we forgive those who trespass
against us. And lead us not into temptation, But deliver us from evil.
For thine is the kingdom, and the power, and the glory, for ever and
ever. Amen.

¶
\emph{Here may follow the Collect for the day.}

\emph{Acknowledgment
of God’s Mercy and Preservation, especially through the Night past.}

ALMIGHTY
and everlasting God, in whom we live and move and have our being; We,
thy needy creatures, render thee our humble praises, for thy preservation
of us from the beginning of our lives to this day, and especially for
having delivered us from the dangers of the past night. For these thy
mercies, we bless and magnify thy glorious Name; humbly beseeching thee
to accept this our morning sacrifice of praise and thanksgiving; for
his sake who lay down in the grave, and rose again for us, thy Son our
Saviour Jesus Christ. \emph{Amen}.

\emph{Dedication
of Soul and Body to God’s Service, with a Resolution to be growing daily in Goodness.}

AND
since it is of thy mercy, O gracious Father, that another day is added
to our lives; We here dedicate both our souls and our bodies to thee and
thy service, in a sober, righteous, and godly life: in which resolution,
do thou, O merciful God, confirm and strengthen us; that, as we grow in
age, we may grow in grace, and in the knowledge of our Lord and Saviour
Jesus Christ. \emph{Amen}.

\emph{Prayer for Grace to enable us to perform that Resolution}.

BUT, O God, who knowest
the weakness and corruption of our nature, and the manifold temptations
which we daily meet with; We humbly beseech thee to have compassion on
our infirmities, and to give us the constant assistance of thy Holy Spirit;
that we may be effectually restrained from sin, and incited to our duty.
Imprint upon our hearts such a dread of thy judgments, and such a grateful
sense of thy goodness to us, as may make us both afraid and ashamed to
offend thee. And, above all, keep in our minds a lively remembrance of
that great day, in which we must give a strict account of our thoughts,
words, and actions to him whom thou hast appointed the Judge of quick
and dead, thy Son Jesus Christ our Lord. \emph{Amen}.

\emph{For Grace to guide and keep us the following Day, and for God’s Blessing on the business of the Same.}

IN particular, we implore
thy grace and protection for the ensuing day. Keep us temperate in all
things, and diligent in our several callings. Grant us patience under
our afflictions. Give us grace to be just and upright in all our dealings;
quiet and peaceable; full of compassion; and ready to do good to all men,
according to our abilities and opportunities. Direct us in all our ways.
Defend us from all dangers and adversities; and be graciously pleased
to take us, and all who are dear to us, under thy fatherly care and protection.
These things, and whatever else thou shalt see to be necessary and convenient
to us, we humbly beg, through the merits and mediation of thy Son Jesus
Christ, our Lord and Saviour. \emph{Amen}.

THE grace of our Lord Jesus
Christ, and the love of God, and the fellowship of the Holy Ghost, be
with us all evermore. \emph{Amen}.

EVENING PRAYER.

¶
\emph{The Family being together, a little before bedtime, let the Master
or Mistress, or any other who may be appointed, say as followeth, all
kneeling, and repeating with him the Lord's Prayer.}

OUR
Father, who art in heaven, Hallowed be thy Name. Thy kingdom come. Thy
will be done, On earth as it is in heaven. Give us this day our daily
bread. And forgive us our trespasses, As we forgive those who trespass
against us. And lead us not into temptation, But deliver us from evil.
For thine is the kingdom, and the power, and the glory, for ever and ever.
Amen.

¶
\emph{Here may follow the Collect for the day.}

\emph{Confession
of Sins, with a Prayer for Contrition and Pardon.}

* \emph{Here
let him who reads make a short pause, that every one may secretly
confess the sins and failings of that day.}

MOST
merciful God, who art of purer eyes than to behold iniquity, and hast
promised forgiveness to all those who confess and forsake their sins;
We come before thee in an humble sense of our own unworthiness, acknowledging
our manifold transgressions of thy righteous laws.* But, O gracious Father,
who desirest not the death of a sinner, look upon us, we beseech thee,
in mercy, and forgive us all our transgressions. Make us deeply sensible
of the great evil of them; and work in us an hearty contrition; that we
may obtain forgiveness at thy hands, who art ever ready to receive humble
and penitent sinners; for the sake of thy Son Jesus Christ, our only Saviour
and Redeemer. \emph{Amen.}

\emph{Prayer
for Grace to reform and grow Better.}

AND
lest, through our own frailty, or the temptations which encompass us,
we be drawn again into sin, vouchsafe us, we beseech thee, the direction
and assistance of thy Holy Spirit. Reform whatever is amiss in the temper
and disposition of our souls; that no unclean thoughts, unlawful designs,
or inordinate desires, may rest there. Purge our hearts from envy, hatred,
and malice; that we may never suffer the sun to go down upon our wrath;
but may always go to our rest in peace, charity, and good-will, with a
conscience void of offence towards thee, and towards men; that so we may
be preserved pure and blameless, unto the coming of our Lord and Saviour
Jesus Christ. \emph{Amen}.

\emph{The
Intercession.}

AND
accept, O Lord, our intercessions for all mankind. Let the light of thy
Gospel shine upon all nations; and may as many as have received it, live
as becomes it. Be gracious unto thy Church; and grant that every member
of the same, in his vocation and ministry, may serve thee faithfully.
Bless all in authority over us; and so rule their hearts and strengthen
their hands, that they may punish wickedness and vice, and maintain thy
true religion and virtue. Send down thy blessings, temporal and spiritual,
upon all our relations, friends, and neighbours. Reward all who have done
us good, and pardon all those who have done or wish us evil, and give
them repentance and better minds. Be merciful to all who are in any trouble;
and do thou, the God of pity, administer to them according to their several
necessities; for his sake who went about doing good, thy Son our Saviour
Jesus Christ. \emph{Amen}.

\emph{The
Thanksgiving.}

TO
our prayers, O Lord, we join our unfeigned thanks for all thy mercies;
for our being, our reason, and all other endowments and faculties of soul
and body; for our health, friends, food, and raiment, and all the other
comforts and conveniences of life. Above all, we adore thy mercy in sending
thy only Son into the world, to redeem us from sin and eternal death,
and in giving us the knowledge and sense of our duty towards thee. We
bless thee for thy pa-tience with us, notwithstanding our many and great
provo-cations; for all the directions, assistances, and comforts of thy
Holy Spirit; for thy continual care and watchful providence over us through
the whole course of our lives; and particularly for the mercies and benefits
of the past day; beseeching thee to continue these thy blessings to us,
and to give us grace to show our thankfulness in a sincere obedience to
his laws, through whose merits and inter-cession we received them all,
thy Son our Saviour Jesus Christ. \emph{Amen}.

\emph{Prayer
for God’s Protection through the Night following.}

IN
particular, we beseech thee to continue thy gracious protection to us
this night. Defend us from all dangers and mischiefs, and from the fear
of them; that we may enjoy such refreshing sleep as may fit us for the
duties of the coming day. And grant us grace always to live in such a
state that we may never be afraid to die; so that, living and dying, we
may be thine, through the merits and satisfaction of thy Son Christ Jesus,
in whose Name we offer up these our imperfect prayers. \emph{Amen}.

THE
grace of our Lord Jesus Christ, and the love of God, and the fellowship
of the Holy Ghost, be with us all evermore. \emph{Amen}.

¶\emph{On Sundays, and on other days when it may be convenient, it will be proper
to begin with a Chapter, or part of a Chapter, from the New Testament.}

\emph{}

\emph{}

A SHORTER FORM.

MORNING.

¶ \emph{After the reading of a brief portion of Holy Scripture, let
the Head of the Household, or some other member of the family, say as
followeth, all kneeling, and repeating with him the Lord's Prayer.}

OUR
Father, who art in heaven, Hallowed be thy Name. Thy kingdom come. Thy
will be done, On earth as it is in heaven. Give us this day our daily
bread. And forgive us our trespasses, As we forgive those who trespass
against us. And lead us not into temptation, But deliver us from evil.
For thine is the kingdom, and the power, and the glory, for ever and ever.
Amen.

O
LORD, our heavenly Father, Almighty and everlasting God, who hast safely
brought us to the beginning of this day; Defend us in the same with thy
mighty power; and grant that this day we fall into no sin, neither run
into any kind of danger; but that all our doings, being ordered by thy
governance, may be righteous in thy sight; through Jesus Christ our Lord.
\emph{Amen}.

\emph{¶
Here may be added any special Prayers.}

THE
grace of our Lord Jesus Christ, and the love of God, and the fellowship
of the Holy Ghost, be with us all evermore. \emph{Amen}.

EVENING.

¶
\emph{After the reading of a brief portion of Holy Scripture, let the Head
of the Household, or some other member of the family, say as followeth,
all kneeling and repeating with him the Lord's Prayer.}

OUR
Father, who art in heaven, Hallowed be thy Name. Thy kingdom come. Thy
will be done, On earth as it is in heaven. Give us this day our daily
bread. And forgive us our trespasses, As we forgive those who trespass
against us. And lead us not into temptation, But deliver us from evil.
For thine is the kingdom, and the power, and the glory, for ever and ever.
Amen.

LIGHTEN
our darkness, we beseech thee, O Lord; and by thy great mercy defend us
from all perils and dangers of this night; for the love of thy only Son,
our Saviour, Jesus Christ. \emph{Amen}.

¶
\emph{Here may be added any special Prayers.}

THE
Lord bless us and keep us. The Lord make his face to shine upon us, and
be gracious unto us. The Lord lift up his countenance upon us, and give
us peace, this night and evermore. \emph{Amen.}

ADDITIONAL
PRAYERS.

\emph{For
the Spirit of Prayer.}

ALMIGHTY
God, who pourest out on all who desire it, the spirit of grace and of
supplication; Deliver us, when we draw nigh to thee, from coldness of
heart and wanderings of mind, that with stedfast thoughts and kindled
affections, we may worship thee in spirit and in truth; through Jesus
Christ our Lord.\emph{Amen}.

\emph{In
the Morning.}

O
GOD, the King eternal, who dividest the day from the darkness, and turnest
the shadow of death into the morning; Drive far off from us all wrong
desires, incline our hearts to keep thy law, and guide our feet into the
way of peace; that having done thy will with cheerfulness while it was
day, we may, when the night cometh, rejoice to give thee thanks; through
Jesus Christ our Lord. \emph{Amen}.

ALMIGHTY
God, who alone gavest us the breath of life, and alone canst keep alive
in us the holy desires thou dost impart; We beseech thee, for thy compassion's
sake, to sanctify all our thoughts and endeavours; that we may neither
begin an action without a pure intention nor continue it without thy blessing.
And grant that, having the eyes of the mind opened to behold things invisible
and unseen, we may in heart be inspired by thy wisdom, and in work be
upheld by thy strength, and in the end be accepted of thee as thy faithful
servants; through Jesus Christ our Saviour. \emph{Amen}.

\emph{At
Night.}

O
LORD, support us all the day long, until the shadows lengthen and the
evening comes, and the busy world is hushed, and the fever of life is
over, and our work is done. Then in thy mercy grant us a safe lodging,
and a holy rest, and peace at the last. \emph{Amen}.

O
GOD, who art the life of mortal men, the light of the faithful, the strength
of those who labour, and the repose of the dead; We thank thee for the
timely blessings of the day, and humbly supplicate thy merciful protection
all this night. Bring us, we beseech thee, in safety to the morning hours;
through him who died for us and rose again, thy Son, our Saviour Jesus
Christ. \emph{Amen}.

\emph{Sunday
Morning.}

O
GOD, who makest us glad with the weekly remembrance of the glorious resurrection
of thy Son our Lord; Vouchsafe us this day such blessing through our worship
of thee, that the days to come may be spent in thy service; through the
same Jesus Christ our Lord. \emph{Amen}.

\emph{For
Quiet Confidence.}

O
GOD of peace, who hast taught us that in returning and rest we shall be
saved, in quietness and in confidence shall be our strength; By the might
of thy Spirit lift us, we pray thee, to thy presence, where we may be
still and know that thou art God; through Jesus Christ our Lord. \emph{Amen}.

\emph{For
Guidance.}

O
GOD, by whom the meek are guided in judgment, and light riseth up in darkness
for the godly; Grant us, in all our doubts and uncertainties, the grace
to ask what thou wouldest have us to do, that the Spirit of Wisdom may
save us from all false choices, and that in thy light we may see light,
and in thy straight path may not stumble; through Jesus Christ our Lord.
\emph{Amen.}

\emph{For
Trustfulness.}

O
MOST loving Father, who willest us to give thanks for all things, to
dread nothing but the loss of thee, and to cast all our care on thee,
who carest for us; Preserve us from faithless fears and worldly anxieties,
and grant that no clouds of this mortal life may hide from us the light
of that love which is immortal, and which thou hast manifested unto
us in thy Son, Jesus Christ our Lord. \emph{Amen}.

O
HEAVENLY Father, thou understandest all thy children; through thy gift
of faith we bring our perplexities to the light of thy wisdom, and receive
the blessed encouragement of thy sympathy, and a clearer knowledge of
thy will. Glory be to thee for all thy gracious gifts. \emph{Amen}.

\emph{For
Joy in God's Creation.}

O
HEAVENLY Father, who hast filled the world with beauty; Open, we beseech
thee, our eyes to behold thy gracious hand in all thy works; that rejoicing
in thy whole creation, we may learn to serve thee with gladness; for the
sake of him by whom all things were made, thy Son, Jesus Christ our Lord.
\emph{Amen}.

\emph{For
the Children.}

ALMIGHTY
God, heavenly Father, who hast blessed us with the joy and care of children;
Give us light and strength so to train them, that they may love whatsoever
things are true and pure and lovely and of good report, following the
example of their Saviour Jesus Christ. \emph{Amen}.

\emph{For
the Absent.}

O
GOD, whose fatherly care reacheth to the uttermost parts of the earth;
We humbly beseech thee graciously to behold and bless those whom we love,
now absent from us. Defend them from all dangers of soul and body; and
grant that both they and we, drawing nearer to thee, may be bound together
by thy love in the communion of thy Holy Spirit, and in the fellowship
of thy saints; through Jesus Christ our Lord. \emph{Amen.}

\emph{For
Those We Love.}

ALMIGHTY
God, we entrust all who are dear to us to thy never-failing care and love,
for this life and the life to come; knowing that thou art doing for them
better things than we can desire or pray for; through Jesus Christ our
Lord. Amen.

\emph{For
the Recovery of a Sick Person.}

O
MERCIFUL God, giver of life and health; Bless, we pray thee, thy servant,
[\emph{N}.], and those who administer to \emph{him} of thy healing
gifts; that \emph{he} may be restored to health of body and of mind;
through Jesus Christ our Lord. \emph{Amen}.

\emph{For
One about to undergo an Operation.}

ALMIGHTY
God our heavenly Father, we beseech thee graciously to comfort thy servant
in \emph{his} suffering, and to bless the means made use of for \emph{his}
cure. Fill \emph{his}heart with confidence, that though \emph{he}
be sometime afraid, \emph{he} yet may put his trust in thee; through
Jesus Christ our Lord. \emph{Amen}.

\emph{For
a Birthday.}

WATCH
over thy child, O Lord, as \emph{his} days increase; bless and guide
him wherever \emph{he} may be, keeping \emph{him} unspotted from the
world. Strengthen \emph{him} when \emph{he} stands; comfort \emph{him}
when discouraged or sorrowful; raise \emph{him} up if \emph{he} fall;
and in \emph{his} heart may thy peace which passeth understanding abide
all the days of \emph{his} life; through Jesus Christ our Lord. \emph{Amen}.

\emph{For
an Anniversary of One Departed.}

ALMIGHTY
God, we remember this day before thee thy faithful servant [\emph{N}.],
and we pray thee that, having opened to \emph{him} the gates of larger
life, thou wilt receive \emph{him} more and more into thy joyful service;
that \emph{he} may win, with thee and thy servants everywhere, the
eternal victory; through Jesus Christ our Lord. \emph{Amen}.

\emph{For Those in Mental Darkness.}

O HEAVENLY Father, we beseech
thee to have mercy upon all thy children who are living in mental darkness.
Restore them to strength of mind and cheerfulness of spirit, and give
them health and peace; through Jesus Christ our Lord. \emph{Amen}.

\emph{For a Blessing on the Families of the Land.}

ALMIGHTY God, our heavenly
Father, who settest the solitary in families; We commend to thy continual
care the homes in which thy people dwell. Put far from them, we beseech
thee, every root of bitterness, the desire of vain-glory, and the pride
of life. Fill them with faith, virtue, knowledge, temperance, patience,
godliness. Knit together in constant affection those who, in holy wedlock,
have been made one flesh; turn the heart of the fathers to the children,
and the heart of the children to the fathers; and so enkindle fervent
charity among us all, that we be evermore kindly affectioned with brotherly
love; through Jesus Christ our Lord. Amen.

\emph{For all Poor, Homeless, and Neglected Folk.}

O GOD, Almighty and merciful,
who healest those that are broken in heart, and turnest the sadness of
the sorrowful to joy; Let thy fatherly goodness be upon all that thou
hast made. Remember in pity such as are this day destitute, homeless,
or forgotten of their fellow-men. Bless the congregation of thy poor.
Uplift those who are cast down. Mightily befriend innocent sufferers,
and sanctify to them the endurance of their wrongs. Cheer with hope all
discouraged and unhappy people, and by thy heavenly grace preserve from
falling those whose penury tempteth them to sin; though they be troubled
on every side, suffer them not to be distressed; though they be perplexed,
save them from despair. Grant this, O Lord, for the love of him, who for
our sakes became poor, thy Son, our Saviour Jesus Christ. \emph{Amen}.

\emph{For Faithfulness in the Use of this World’s
Goods.}

ALMIGHTY God, whose loving
hand hath given us all that we possess; Grant us grace that we may honour
thee with our substance, and remembering the account which we must one
day give, may be faithful stewards of thy bounty; through Jesus Christ
our Lord. \emph{Amen}.

\emph{A General Intercession.}

O GOD, at whose word man
goeth forth to his work and to his labour until the evening; Be merciful
to all whose duties are difficult or burdensome, and comfort them concerning
their toil. Shield from bodily accident and harm the workmen at their
work. Protect the efforts of sober and honest industry, and suffer not
the hire of the labourers to be kept back by fraud. Incline the heart
of employers and of those whom they employ to mutual forbearance, fairness,
and good-will. Give the spirit of governance and of a sound mind to all
in places of authority. Bless all those who labour in works of mercy or
in schools of good learning. Care for all aged persons, and all little
children, the sick and the afflicted, and those who travel by land or
by sea. Remember all who by reason of weakness are overtasked, or because
of poverty are forgotten. Let the sorrowful sighing of the prisoners come
before thee; and according to the greatness of thy power, preserve thou
those that are appointed to die. Give ear unto our prayer, O merciful
and gracious Father, for the love of thy dear Son, our Saviour Jesus Christ.
\emph{Amen}.

\emph{Grace before Meat.}

BLESS, O Father, thy gifts
to our use and us to thy service; for Christ’s sake. \emph{Amen}.

GIVE us grateful hearts,
our Father, for all thy mercies, and make us mindful of the needs of others;
through Jesus Christ our Lord. \emph{Amen}. 

\begin{quote}

\emph{Return to the U.
S. Book of Common Prayer (1928)}

\end{quote}

\xchapter{The Articles of Religion}

\textbf{The Articles of Religion as they appear in the 1789, 1892 and 1928 Books of
Common Prayer}

Articles
of Religion

\emph{As
established by the Bishops, the Clergy, and the Laity of the Protestant
Episcopal Church in the United States of America, in Convention, on the twelfth day of September, in the Year of our Lord, 1801.}

As
noted in the title, the Articles of Religion were not adopted in the
American Church until 1801, and do not appear in Prayer Books until
after about 1808. In the 1789 BCP they are placed between the Psalter
and the Ordinal; in the 1892 and 1928 Books they appear at the very
end of the Book.

ARTICLES OF RELIGION.

I.
\emph{Of Faith in the Holy Trinity.}

THERE
is but one living and true God, everlasting, without body, parts, or passions;
of infinite power, wisdom, and goodness; the Maker, and Preserver of all
things both visible and invisible. And in the unity of this Godhead there
be three Persons, of one substance, power, and eternity; the Father, the
Son, and the Holy Ghost.

II.
\emph{Of the Word or Son of God, which was made very Man.}

The
Son, which is the Word of the Father, begotten from everlasting of the
Father, the very and eternal God, and of one substance with the Father,
took Man's nature in the womb of the blessed Virgin, of her substance:
so that two whole and perfect Natures, that is to say, the Godhead and
Manhood, were joined together in one Person, never to be divided, whereof
is one Christ, very God, and very Man; who truly suffered, was crucified,
dead, and buried, to reconcile his Father to us, and to be a sacrifice,
not only for original guilt, but also for actual sins of men.

III.
\emph{Of the going down of Christ into Hell.}

As
Christ died for us, and was buried; so also it is to be believed, that
he went down into Hell.

IV.
\emph{Of the Resurrection of Christ.}

Christ
did truly rise again from death, and took again his body, with flesh,
bones, and all things appertaining to the perfection of Man’s nature;
wherewith he ascended into Heaven, and there sitteth, until he return
to judge all Men at the last day.

V.
\emph{Of the Holy Ghost.}

The
Holy Ghost, proceeding from the Father and the Son, is of one substance,
majesty, and glory, with the Father and the Son, very and eternal God.

VI.
\emph{Of the Sufficiency of the Holy Scriptures for Salvation.}

Holy
Scripture containeth all things necessary to salvation: so that whatsoever
is not read therein, nor may be proved thereby, is not to be required
of any man, that it should be believed as an article of the Faith, or
be thought requisite or necessary to salvation. In the name of the Holy
Scripture we do understand those canonical Books of the Old and New Testament,
of whose authority was never any doubt in the Church.

\emph{Of
the Names and Number of the Canonical Books}

The
1789 and 1892 Books add ``ART.''
to every article title.

The
1789 and 1892 Books also generally italicize personal and place names.

Genesis,

Exodus,

Leviticus,

Numbers,

Deuteronomy,

Joshua,

Judges, 

Ruth,

The
First Book of Samuel,

The Second Book of Samuel,

The First Book of Kings,

The Second Book of Kings,

The First Book of Chronicles,

The Second Book of Chronicles, 

The First Book of Esdras,

The Second Book of Esdras, The
Book of Esther ,

The Book of Job,

The Psalms,

The Proverbs,

Ecclesiastes or Preacher,

Cantica, or Songs of Solomon,

Four Prophets the greater,

Twelve Prophets the less.

And
the other Books (as Hierome saith) the Church doth read for example of
life and instruction of manners; but yet doth it not apply them to establish
any doctrine; such are these following:

The
Third Book of Esdras,

The Fourth Book of Esdras, 

The Book of Tobias,

The Book of Judith, 

The rest of the Book of Esther,

The Book of Wisdom,

Jesus the Son of Sirach,

Baruch
the Prophet,

The Song of the Three Children,

The Story of Susanna,

Of Bel and the Dragon,

The Prayer of Manasses,

The First Book of Maccabees,

The Second Book of Maccabees.

All
the Books of the New Testament, as they are commonly received, we do receive,
and account them Canonical.

VII.
\emph{Of the Old Testament.}

The
Old Testament is not contrary to the New: for both in the Old and New
Testament everlasting life is offered to Mankind by Christ, who is the
only Mediator between God and Man, being both God and Man. Wherefore they
are not to be heard, which feign that the old Fathers did look only for
transitory promises. Although the Law given from God by Moses, as touching
Ceremonies and Rites, do not bind Christian men, nor the Civil precepts
thereof ought of necessity to be received in any commonwealth; yet notwith-standing,
no Christian man whatsoever is free from the obedience of the Commandments
which are called Moral.

VIII.
\emph{Of the Creeds.}

The
Nicene Creed, and that which is commonly called the Apostles’ Creed,
ought thoroughly to be received and believed: for they may be proved by
most certain warrants of Holy Scripture.

Numeri,
Deuteronomium, Joshue, Hester
until 1845.

IX.
\emph{Of Original or Birth Sin.}

Original sin standeth not in the following of Adam, (as the Pelagians
do vainly talk;) but it is the fault and corruption of the Nature of every
man, that naturally is engendered of the offspring of Adam; whereby man
is very far gone from original righteousness, and is of his own nature
inclined to evil, so that the flesh lusteth always contrary to the Spirit;
and therefore in every person born into this world, it deserveth God’s
wrath and damnation. And this infection of nature doth remain, yea in
them that are regenerated; whereby the lust of the flesh, called in Greek,
φρονημα σαρκος,
(which some do expound the wisdom, some sensuality, some the affection,
some the desire, of the flesh), is not subject to the Law of God. And
although there is no condemnation for them that believe and are baptized;
yet the Apostle doth confess, that concupiscence and lust hath of itself
the nature of sin.

X. \emph{Of Free Will.}

The condition of Man after the fall of Adam is such, that he cannot turn
and prepare himself, by his own natural strength and good works, to faith,
and calling upon God. Wherefore we have no power to do good works pleasant
and acceptable to God, without the grace of God by Christ preventing us,
that we may have a good will, and working with us, when we have that good
will.

XI. \emph{Of the Justification of Man.}

We are accounted righteous before God, only for the merit of our Lord
and Saviour Jesus Christ by Faith, and not for our own works or deservings.
Wherefore, that we are justified by Faith only, is a most wholesome Doctrine,
and very full of comfort, as more largely is expressed in the Homily of
Justification.

XII. \emph{Of Good Works.}

Albeit that Good Works, which are the fruits of Faith, and follow after
Justification, cannot put away our sins, and endure the severity of God's
judgment; yet are they pleasing and acceptable to God in Christ, and do
spring out necessarily of a true and lively Faith; insomuch that by them
a lively Faith may be as evidently known as a tree discerned by the fruit.

XIII. \emph{Of Works before Justification.}

Works done before the grace of Christ, and the Inspiration of the Spirit,
are not pleasant to God, forasmuch as they spring not of faith in Jesus
Christ; neither do they make men meet to receive grace, or (as the School-authors
say) deserve grace of congruity: yea rather, for that they are not done
as God hath willed and commanded them to be done, we doubt not but they
have the nature of sin.

XIV. \emph{Of Works of Supererogation.}

Voluntary Works besides, over and above, God's Commandments, which they
call Works of Supererogation, cannot be taught without arrogancy and impiety:
for by them men do declare, that they do not only render unto God as much
as they are bound to, but that they do more for his sake, than of bounden
duty is required: whereas Christ saith plainly, When ye have done all
that are commanded to you, say, We are unprofitable servants.

XV. \emph{Of Christ alone without Sin.}

Christ in the truth of our nature was made like unto us in all things,
sin only except, from which he was clearly void, both in his flesh, and
in his spirit. He came to be the Lamb without spot, who, by sacrifice
of himself once made, should take away the sins of the world; and sin
(as Saint John saith) was not in him. But all we the rest, although baptized,
and born again in Christ, yet offend in many things; and if we say we
have no sin, we deceive ourselves, and the truth is not in us.

XVI.
\emph{Of Sin after Baptism.}

Not
every deadly sin willingly committed after Baptism is sin against the
Holy Ghost, and unpardonable. Wherefore the grant of repentance is not
to be denied to such as fall into sin after Baptism. After we have received
the Holy Ghost, we may depart from grace given, and fall into sin, and
by the grace of God we may arise again, and amend our lives. And therefore
they are to be condemned, which say, they can no more sin as long as they
live here, or deny the place of forgiveness to such as truly repent.

XVII. \emph{Of Predestination
and Election}

Predestination
to Life is the everlasting purpose of God, whereby (before the foundations
of the world were laid) he hath constantly decreed by his counsel secret
to us, to deliver from curse and damnation those whom he hath chosen in
Christ out of mankind, and to bring them by Christ to ever-lasting salvation,
as vessels made to honour. Wherefore, they which be endued with so excellent
a benefit of God, be called according to God's purpose by his Spirit working
in due season: they through Grace obey the calling: they be justified
freely: they be made sons of God by adoption: they be made like the image
of his only-begotten Son Jesus Christ: they walk religiously in good works,
and at length, by God's mercy, they attain to everlasting felicity.

As the godly consideration of Predestination, and our Election in Christ,
is full of sweet, pleasant, and unspeakable comfort to godly persons,
and such as feel in themselves the working of the Spirit of Christ, mortifying
the works of the flesh, and their earthly members, and drawing up their
mind to high and heavenly things, as well because it doth greatly establish
and confirm their faith of eternal Salvation to be enjoyed through Christ,
as because it doth fervently kindle their love towards God: So, for curious
and carnal persons, lacking the Spirit of Christ, to have continually
before their eyes the sentence of God’s Predestination, is a most
dangerous downfall, whereby the Devil doth thrust them either into desperation,
or into wrethchlessness of most unclean living, no less perilous than
desperation.

Furthermore, we must receive God's promises in such
wise, as they be generally set forth to us in Holy Scripture: and, in
our doings, that Will of God is to be followed, which we have expressly
declared unto us in the word of God.

XVIII.
\emph{Of obtaining eternal Salvation only by the Name of Christ.}

They
also are to be had accursed that presume to say, That every man shall
be saved by the Law or Sect which he professeth, so that he be diligent
to frame his life according to that Law, and the light of Nature. For
Holy Scripture doth set out unto us only the Name of Jesus Christ, whereby
men must be saved.

XIX. \emph{Of the Church.}

The
visible Church of Christ is a congregation of faithful men, in which the
pure Word of God is preached, and the Sacraments be duly ministered according
to Christ’s ordinance, in all those things that of necessity are
requisite to the same.

As the Church of Jerusalem, Alexandria, and Antioch,
have erred; so also the Church of Rome hath erred, not only in their living
and manner of Ceremonies, but also in matters of Faith.

XX. \emph{Of the Authority
of the Church.}

The
Church hath power to decree Rites or Ceremonies, and authority in Controversies
of Faith: and yet it is not lawful for the Church to ordain anything that
is contrary to God’s Word written, neither may it so expound one
place of Scripture, that it be repugnant to another. Wherefore, although
the Church be a witness and a keeper of Holy Writ, yet, as it ought not
to decree any thing against the same, so besides the same ought not to
enforce any thing to be believed for necessity of Salvation.

XXI.
\emph{Of the Authority of General Councils.}

[The
Twenty-first of the former Articles is omitted; because it is partly of
a local and civil nature, and is provided for, as to the remaining parts
of it, in other Articles.]

XXII. \emph{Of Purgatory.}

The
Romish Doctrine concerning Purgatory, Pardons, Worshipping and Adoration,
as well of Images as of Relics, and also Invocation of Saints, is a fond
thing, vainly invented, and grounded upon no warranty of Scripture, but
rather repugnant to the Word of God.

XXIII. \emph{Of Ministering
in the Congregation.}

It
is not lawful for any man to take upon him the office of public preaching,
or ministering the Sacraments in the Congregation, before he be lawfully
called, and sent to execute the same. And those we ought to judge lawfully
called and sent, which be chosen and called to this work by men who have
public authority given unto them in the Congregation, to call and send
Ministers into the Lord’s vineyard.

XXIV.
\emph{Of Speaking in the Congregation in such a Tongue as the people understandeth.}

It
is a thing plainly repugnant to the Word of God, and the custom of the
Primitive Church, to have public Prayer in the Church, or to minister
the Sacraments, in a tongue not understanded of the people.

XXV. \emph{Of the Sacraments.}

Sacraments
ordained of Christ be not only badges or tokens of Christian men’s
profession, but rather they be certain sure witnesses, and effectual signs
of grace, and God’s good will towards us, by the which he doth work
invisibly in us, and doth not only quicken, but also strengthen and confirm
our Faith in him.

There are two Sacraments ordained of Christ our Lord
in the Gospel, that is to say, Baptism, and the Supper of the Lord.

Those five commonly called Sacraments, that is to say,
Confirmation, Penance, Orders, Matrimony, and Extreme Unction, are not
to be counted for Sacraments of the Gospel, being such as have grown partly
of the corrupt following of the Apostles, partly are states of life allowed
in the Scriptures; but yet have not like nature of Sacraments with Baptism,
and the Lord’s Supper, for that they have not any visible sign or
ceremony ordained of God.

The Sacraments are not ordained of Christ to be gazed
upon, or to be carried about, but that we should duly use them. And in
such only as worthily receive the same, they have a wholesome effect or
operation: but they that receive them unworthily, purchase to themselves
damnation, as Saint Paul saith.

XXVI.
\emph{Of the Unworthiness of the Ministers, which hinders not the effect
of the Sacraments.}

Although
in the visible Church the evil be ever mingled with the good, and sometimes
the evil have chief authority in the Ministration of the Word and Sacraments,
yet forasmuch as they do not the same in their own name, but in Christ's,
and do minister by his commission and authority, we may use their Ministry,
both in hearing the Word of God, and in receiving the Sacraments. Neither
is the effect of Christ's ordinance taken away by their wickedness, nor
the grace of God’s gifts diminished from such as by faith, and rightly,
do receive the Sacraments ministered unto them; which be effectual, because
of Christ’s institution and promise, although they be ministered
by evil men.

Nevertheless, it appertaineth to the discipline of
the Church, that inquiry be made of evil Ministers, and that they be accused
by those that have knowledge of their offences; and finally, being found
guilty, by just judgment be deposed.


XXVII. \emph{Of Baptism.}

Baptism
is not only a sign of profession, and mark of difference, whereby Christian
men are discerned from others that be not christened, but it is also a
sign of Regeneration or New-Birth, whereby, as by an instrument, they
that receive Baptism rightly are grafted into the Church; the promises
of the forgiveness of sin, and of our adoption to be the sons of God by
the Holy Ghost, are visibly signed and sealed; Faith is confirmed, and
Grace increased by virtue of prayer unto God.

The Baptism of young Children is in any wise to be
retained in the Church, as most agreeable with the institution of Christ.

XXVIII.
\emph{Of the Lord’s Supper.}

The
Supper of the Lord is not only a sign of the love that Christians ought
to have among themselves one to another; but rather it is a Sacrament
of our Redemption by Christ’s death: insomuch that to such as rightly,
worthily, and with faith, receive the same, the Bread which we break is
a partaking of the Body of Christ; and likewise the Cup of Blessing is
a partaking of the Blood of Christ.

Transubstantiation (or the change of the substance
of Bread and Wine) in the Supper of the Lord, cannot be proved by Holy
Writ; but is repugnant to the plain words of Scripture, overthroweth the
nature of a Sacrament, and hath given occasion to many superstitions.

The Body of Christ is given, taken, and eaten, in the
Supper, only after an heavenly and spiritual manner. And the mean whereby
the Body of Christ is received and eaten in the Supper, is Faith.

The Sacrament of the Lord’s Supper was not by
Christ’s ordinance reserved, carried about, lifted up, or worshipped.

XXIX.
\emph{Of the Wicked, which eat not the Body of Christ in the use of the
Lord's Supper.}

The
Wicked, and such as be void of a lively faith, although they do carnally
and visibly press with their teeth (as Saint Augustine saith) the Sacrament
of the Body and Blood of Christ; yet in no wise are they partakers of
Christ: but rather, to their condemnation, do eat and drink the sign or
Sacrament of so great a thing.

XXX.\emph{Of both Kinds.}

The
Cup of the Lord is not to be denied to the Lay-people: for both the parts
of the Lord’s Sacrament, by Christ’s ordinance and commandment,
ought to be ministered to all Christian men alike.

This
note for Art. 21 appears as a footnote in the 1789 and 1892 Books.

XXXI. \emph{Of the one Oblation of Christ finished upon
the Cross.}

The Offering of Christ once made in that perfect redemption,
propitiation, and satisfaction, for all the sins of the whole world, both
original and actual; and there is none other satisfaction for sin, but
that alone. Wherefore the sacrifices* of Masses, in the which it was commonly
said, that the Priest did offer Christ for the quick and the dead, to
have remission of pain or guilt, were blasphemous fables, and dangerous
deceits.

XXXII. \emph{Of the Marriage of Priests.}

Bishops, Priests, and Deacons, are not commanded by God’s
Law, either to vow the estate of single life, or to abstain from marriage:
therefore it is lawful for them, as for all other Christian men, to marry
at their own discretion, as they shall judge the same to serve better
to godliness.

XXXIII. \emph{Of excommunicate Persons, how they are to
be avoided.}

That person which by open denunciation of the Church
is rightly cut off from the unity of the Church, and excommunicated, ought
to be taken of the whole multitude of the faithful, as an Heathen and
Publican, until he be openly reconciled by penance, and received into
the Church by a Judge that hath the authority thereunto.

XXXIV. \emph{Of the Traditions of the Church.}

It is not necessary that the Traditions and Ceremonies
be in all places one, or utterly like; for at all times they have been
divers, and may be changed according to the diversity of countries, times,
and men’s manners, so that nothing be ordained against God’s
Word. Whosoever, through his private judgment, willingly and purposely,
doth openly break the Traditions and Ceremonies of the Church, which be
not repugnant to the Word of God, and be ordained and approved by common
authority, ought to be rebuked openly, (that others may fear to do the
like,) as he that offendeth against the common order of the Church, and
hurteth the authority of the Magistrate, and woundeth the consciences
of the weak brethren.

Every particular or national Church hath authority
to ordain, change, and abolish, Ceremonies or Rites of the Church ordained
only by man’s authority, so that all things be done to edifying.

XXXV. \emph{Of the Homilies.}

The Second Book of Homilies, the several titles whereof
we have joined under this Article, doth contain a godly and wholesome
Doctrine, and necessary for these times, as doth the former Book of Homilies,
which were set forth in the time of Edward the Sixth; and therefore we
judge them to be read in Churches by the Ministers, diligently and distinctly,
that they may be understanded of the people.

\emph{Of the Names of the Homilies}

1
Of the right Use of the Church. 

2 Against Peril of Idolatry.

3 Of repairing and keeping clean of Churches.

4 Of good Works: first of Fasting. 

5 Against Gluttony and Drunkenness.

6 Against Excess of Apparel

7 Of Prayer.

8 Of the Place and Time of Prayer.

9 That Common Prayers and Sacraments ought to be ministered in a known
tongue.

10 Of the reverend Estimation of God’s Word. 11
Of Alms-doing.

12 Of the Nativity of Christ.

13 Of the Passion of Christ.

14 Of the Resurrection of Christ.

15 Of the worthy receiving of the Sacrament of the Body and Blood
of Christ.

16 Of the Gifts of the Holy Ghost.

17 For the Rogation-days.

18 Of the State of Matrimony.

19 Of Repentance.

20 Against Idleness.

21 Against Rebellion.

[This Article is received in this Church,
so far as it declares the Book of Homilies to be an explication of Christian
doctrine, and instructive in piety and morals. But all references to the
constitution and laws of England are considered as inapplicable to the
circum-stances of this Church; which also suspends the order for the reading
of said Homilies in churches, until a revision of them may be conveniently
made, for the clearing of them, as well from obsolete words and phrases,
as from the local references.]

XXXVI. \emph{Of Consecration of Bishops and Ministers.}

The Book of Consecration of Bishops, and Ordering of
Priests and Deacons, as set forth by the General Convention of this Church
in 1792, doth contain all things necessary to such Consecration and Ordering;
neither hath it any thing that, of itself, is superstitious and ungodly.
And, therefore, whosoever are consecrated or ordered according to said
Form, we decree all such to be rightly, orderly, and lawfully consecrated
and ordered.

XXXVII. \emph{Of the Power of the Civil Magistrates.}

The Power of the Civil Magistrate extendeth to all men,
as well Clergy as Laity, in all things temporal; but hath no authority
in things purely spiritual. And we hold it to be the duty of all men who
are professors of the Gospel, to pay respectful obedience to the Civil
Authority, regularly and legitimately constituted.

XXXVIII. \emph{Of Christian Men’s Goods, which are
not common.}

The Riches and Goods of Christians are not common, as
touching the right, title, and possession of the same; as certain Anabaptists
do falsely boast. Notwithstanding, every man ought, of such things as
he possesseth, liberally to give alms to the poor, according to his ability.

XXXIX. \emph{Of a Christian Man’s Oath.}

As we confess that vain and rash Swearing is forbidden
Christian men by our Lord Jesus Christ, and James his Apostle, so we judge,
that Christian Religion doth not prohibit, but that a man may swear when
the Magistrate requireth, in a cause of faith and charity, so it be done
according to the Prophet’s teaching in justice, judgment, and truth.

* sacrifice
until
1845.

\begin{quote}

\emph{Return to the}U.
S. Book of Common Prayer: 1789, 1892,
1928

\end{quote}
